\documentclass[11pt,oneside]{article}

%% ── Packages ────────────────────────────────────────────────────────────
\usepackage{fontspec}
\usepackage{xcolor}
\usepackage{amsmath,amssymb,amsthm}
\usepackage{mathtools}
\usepackage{booktabs}
\usepackage{tabularx}
\usepackage{adjustbox}
\usepackage[colorlinks=true,linkcolor=blue!60!black,citecolor=blue!60!black,urlcolor=blue!60!black]{hyperref}
\usepackage{cleveref}
\usepackage[margin=1in]{geometry}
\usepackage{parskip}            % paragraph spacing instead of indentation
\usepackage{titlesec}
\usepackage{enumitem}

%% ── Typography ──────────────────────────────────────────────────────────
\setmainfont{Latin Modern Roman}
\setmonofont{Latin Modern Mono}[Scale=0.85]
\setlength{\parskip}{0.4\baselineskip plus 2pt minus 1pt}

%% ── Section formatting ─────────────────────────────────────────────────
\titleformat{\section}{\Large\bfseries}{}{0em}{}[\vspace{0.3em}\hrule\vspace{0.5em}]
\titleformat{\subsection}{\large\bfseries}{}{0em}{}
\titleformat{\subsubsection}{\normalsize\bfseries}{}{0em}{}

%% ── Segment environments ───────────────────────────────────────────────
%% Each segment type gets a counter and a clean header, NOT italic body.
\newcounter{claim}[section]
\renewcommand{\theclaim}{\thesection.\arabic{claim}}
\crefname{claim}{}{}        % no prefix — we embed the type in the ref text
\Crefname{claim}{}{}

%% Generic segment environment: \begin{segment}{Type}{Title}
%% Stores the type name so \label can pick it up for cref
\newenvironment{segment}[2]{%
  \refstepcounter{claim}%
  \bigskip\noindent\hrule\smallskip
  \noindent{\bfseries #1~\theclaim}\enskip(#2).\par\smallskip
  \noindent\ignorespaces
}{%
  \par\smallskip\noindent\hrule\bigskip
}

%% Named environments that map to specific types
\newenvironment{postulate}[1][]{\begin{segment}{Postulate}{#1}}{\end{segment}}
\newenvironment{definition}[1][]{\begin{segment}{Definition}{#1}}{\end{segment}}
\newenvironment{formulation}[1][]{\begin{segment}{Formulation}{#1}}{\end{segment}}
\newenvironment{scope}[1][]{\begin{segment}{Scope}{#1}}{\end{segment}}
\newenvironment{proposition}[1][]{\begin{segment}{Proposition}{#1}}{\end{segment}}
\newenvironment{theorem}[1][]{\begin{segment}{Theorem}{#1}}{\end{segment}}
\newenvironment{corollary}[1][]{\begin{segment}{Corollary}{#1}}{\end{segment}}
\newenvironment{hypothesis}[1][]{\begin{segment}{Hypothesis}{#1}}{\end{segment}}
\newenvironment{normative}[1][]{\begin{segment}{Normative}{#1}}{\end{segment}}
\newenvironment{empirical}[1][]{\begin{segment}{Empirical}{#1}}{\end{segment}}
\newenvironment{observation}[1][]{\begin{segment}{Observation}{#1}}{\end{segment}}
\newenvironment{discussion}[1][]{\begin{segment}{Discussion}{#1}}{\end{segment}}
\newenvironment{measurement}[1][]{\begin{segment}{Measurement}{#1}}{\end{segment}}
\newenvironment{schema}[1][]{\begin{segment}{Proposed Schema}{#1}}{\end{segment}}
\newenvironment{example}[1][]{\begin{segment}{Example}{#1}}{\end{segment}}
\newenvironment{detail}[1][]{\begin{segment}{Detail}{#1}}{\end{segment}}
\newenvironment{derivation}[1][]{\begin{segment}{Derivation}{#1}}{\end{segment}}
\newenvironment{sketch}[1][]{\begin{segment}{Sketch}{#1}}{\end{segment}}
\newenvironment{aside}[1][]{\begin{segment}{Aside}{#1}}{\end{segment}}
\newenvironment{remark}[1][]{\begin{segment}{Remark}{#1}}{\end{segment}}

%% ── Notation macros ─────────────────────────────────────────────────────
\DeclareMathOperator*{\argmin}{arg\,min}
\DeclareMathOperator*{\argmax}{arg\,max}

%% ── List formatting ────────────────────────────────────────────────────
\setlist{nosep,leftmargin=1.5em}

%% ── Document metadata ──────────────────────────────────────────────────
\title{ACT: Agentic Cycle Theory}
\author{Joseph A.\ Wecker}
\date{\today}

\begin{document}
\maketitle
\tableofcontents
\newpage

The mathematical core of the Agentic Systems research framework. ACT formalizes the adaptive cycle — one complete traversal of the agent-environment feedback loop — as the fundamental unit of analysis for adaptive, purposeful agents under uncertainty.

\textbf{On mathematical precision.} The theory's relationship to formalism varies by section and is expected to. Section I (adaptive systems) is the most mathematically locked down — the persistence machinery, mismatch dynamics, and gain structure have clean derivations and Lyapunov proofs. Section II (purposeful agents) has an exact diagnostic core (satisfaction gap, control regret, orient-cascade ordering) and a maturing strategy layer, but both depend on the directed-separation scope condition ( \#directed-separation): the exact results apply to \textbf{Class 1 (modular) agents} where epistemic and purposeful processing are architecturally separated. For \textbf{Class 2 (fully merged) agents} — including transformer-based LLMs — directed separation fails by construction, the orient cascade becomes an approximation, and the clean M\_t/G\_t factorization is a modeling convenience rather than a structural feature. Current LLM-based agents are the most important application class that falls outside Section II's exact scope; the coupled formulation they require is the subject of \texttt{03-logogenic-agents/}. Section III (composition) has promising structure built on the Section I Lyapunov machinery but depends on admissibility choices that are formulated, not derived, and a bridge lemma that requires a contraction assumption beyond the stated admissibility constraints. This gradient — from exact core through conditionally exact architecture to open formulation — is the expected arc. The goal is to describe agentic systems, not to produce a purely mathematical artifact. We pursue mathematical precision when it yields genuine insight (the persistence condition, the acyclicity derivation, the architectural classification) and settle for principled sketches when the insight is structural rather than quantitative (the strategy-revision loop, the composition admissibility). The boundaries between these regimes are fluid and still being discovered.

\textbf{Scope:} ACT covers the general theory of adaptive systems (Section I), actuated/purposeful agents (Section II), and agent composition (Section III). Domain instantiations (software: [\texttt{02-tst-core/}](../02-tst-core/OUTLINE.md)), logogenic agents ([\texttt{03-logogenic-agents/}](../03-logogenic-agents/OUTLINE.md)), and logozoetic agents ([\texttt{04-logozoetic-agents/}](../04-logozoetic-agents/OUTLINE.md)) are part of the broader Agentic Systems framework, grounded by ACT but developed independently.

\section{I. Adaptive Systems Under Uncertainty}

\begin{quote}\itshape
Scope: Any system consisting of an agent coupled to an environment through observation and action channels, where the environment is not fully observable. This is the general case — thermostats through commanders. The claims in this section are largely drawn from TFT (TF-01 through TF-11, Appendix A), which developed the adaptive-systems foundation that ACT subsumes.
\end{quote}

\begin{postulate}[Temporal Optimality]
\label{temporal-optimality}
Among agents achieving identical outcomes across all non-temporal dimensions, the one requiring least time is optimal.



\medskip\noindent\textbf{Formal Expression.}



Let $\mathbf{A} = \{A_1, A_2, \ldots, A_n\}$ be agents or strategies achieving outcome $O$.

$$\text{If } \forall m \in \mathbf{M} \setminus \{\text{time}\},\; \forall i,j:\; m(A_i) \equiv m(A_j), \quad \text{then } A^* = \arg\min_{A_k} \text{time}(A_k)$$

where $\mathbf{M}$ is the set of all measurable outcome dimensions, and identical outcomes $m(A_i) \equiv m(A_j)$ means equivalence across:
\begin{itemize}
  \item \textbf{Functional}: Same input→output mappings
  \item \textbf{Non-functional}: Same performance, security, availability
  \item \textbf{Quality}: Same defect probability, maintainability
  \item \textbf{Sustainability}: Same impact on agent capacity and system evolution
  \item \textbf{Impact on others}: Same effect on other agents' capacity
\end{itemize}



\medskip\noindent\textbf{Epistemic Status.}


This is a postulate — deliberately tautological. The statement reduces to: "given a choice between identical outcomes, the one that arrives sooner is preferred." The inability to construct genuine counterexamples without violating the equivalence precondition reveals its postulate nature. The fungibility argument (below) is \textit{discussion} — a qualitative observation about why time is a natural optimization target, not a derived result.


\medskip\noindent\textbf{Discussion.}


\textbf{Why time is special.} Time is uniquely fungible among outcome dimensions: saved time can become exploration, learning, rest, additional action — anything. Other outcome dimensions (correctness, safety, quality) are not fungible in this way. A unit of saved correctness cannot be spent on learning. This is why, once all other outcome dimensions are held equal, time is the natural residual to optimize. The postulate makes this optimization explicit.

\textbf{The equivalence precondition is load-bearing.} The phrase "identical outcomes across all non-temporal dimensions" is doing most of the work. It must be verified concretely before the postulate applies. Apparent counterexamples invariably violate it:
\begin{itemize}
  \item "Move fast and break things" — violates quality equivalence
  \item Burnout-inducing speed — violates sustainability equivalence
  \item Premature optimization for unlikely futures — optimizes for a counterfactual outcome, violating actual-outcome equivalence
  \item Skipping tests for speed — violates defect probability equivalence
\end{itemize}


The postulate does not say faster is always better. It says faster is better \textit{given identical outcomes}. An agent that achieves a worse outcome faster is not preferred by this postulate. Misapplication of temporal optimality without verifying the equivalence precondition produces exactly the pathologies listed above.

\textbf{Why start here.} The adaptive machinery that follows — tempo, gain, persistence, adversarial dynamics — can be understood as consequences of optimizing for this criterion under constraints of partial observability, bounded resources, and environmental change. The connections are developed in their respective claims; they are not implied by this postulate alone.

\textbf{Domain generality.} This postulate applies to any agent-environment pairing within ACT's scope ( \S\,\ref{scope-condition}). In the software domain, it specializes to the statement that among implementations achieving identical outcomes, the one requiring least development time is optimal (TST T-01). The generalization is straightforward: replace "implementation" with "adaptive strategy."

\textbf{Scope limitation.} This postulate does not, by itself, say anything about \textit{how} to achieve temporal optimality. That is the subject of the rest of the theory.

\end{postulate}

\begin{definition}[Agent-Environment Coupling]
\label{agent-environment}
An agent is an entity that receives observations from an environment, maintains internal state, and produces actions that affect the environment. The agent cannot access the environment directly — observations are necessarily lossy. This boundary condition is constitutive: the theory applies where the agent-environment boundary entails information loss.



\medskip\noindent\textbf{Formal Expression.}



Let $\Omega$ denote the \textbf{environment}: the totality of state external to the agent. We make no assumptions about $\Omega$'s structure — it may be continuous or discrete, stationary or non-stationary, deterministic or stochastic, benign or adversarial.

An \textbf{agent} is an entity satisfying three conditions:

\begin{enumerate}
  \item It receives observations from $\Omega$ (perception channel)
  \item It maintains internal state (memory/model)
  \item It produces actions that affect $\Omega$ (action channel)
\end{enumerate}



The agent cannot access $\Omega_t$ directly. All contact with the environment is mediated through lossy observation. This is not a simplifying assumption — it is a scope condition. Systems where the agent has direct access to full environment state are outside ACT's scope ( \S\,\ref{scope-condition}).


\medskip\noindent\textbf{Epistemic Status.}


This is \textit{definitional} — it establishes the conceptual framework, not a truth-claim. The agent-environment decomposition is a modeling choice that delineates what ACT analyzes. The information-loss boundary is the constitutive commitment: it restricts ACT's scope to systems where the agent faces genuine uncertainty about its environment.


\medskip\noindent\textbf{Discussion.}


\textbf{Why information loss is constitutive.} An agent with perfect access to $\Omega_t$ has no need for a model, no mismatch signal, no adaptation. The entire adaptive machinery of Section I becomes vacuous. The information-loss boundary is what makes the theory non-trivial.

**Generality of $\Omega$.** The environment is deliberately underspecified. $\Omega$ may include other agents, physical systems, software artifacts, or any combination. The only structural commitment is that $\Omega$ is external to the agent and not fully accessible.

\end{definition}

\begin{definition}[Observation Function]
\label{observation-function}
Observations (aisthesis — raw contact with reality) are lossy, possibly noisy functions of environment state, prior action, and perceptual noise. The agent knows neither the observation function nor the noise distribution exactly.



\medskip\noindent\textbf{Formal Expression.}



The \textbf{observation space} $\mathcal{O}$ is the set of possible observations. Each observation is generated by:

$$o_t = h(\Omega_t, a_{t-1}, \varepsilon_t)$$

where:
\begin{itemize}
  \item $h$ is the observation function (unknown to the agent)
  \item $\Omega_t$ is the environment state at time $t$
  \item $a_{t-1}$ is the agent's prior action (observations may depend on what the agent did)
  \item $\varepsilon_t$ represents noise or limits of perception
\end{itemize}



The agent knows neither $h$ nor the distribution of $\varepsilon_t$ exactly.


\medskip\noindent\textbf{Epistemic Status.}


This is \textit{definitional}. The observation function $h$ is a modeling device that captures the information-loss boundary established in \S\,\ref{agent-environment}. The dependence on $a_{t-1}$ is optional — when absent, the function reduces to $h(\Omega_t, \varepsilon_t)$. The claim that the agent does not know $h$ or the noise distribution is constitutive of the partial-observability setting, not an empirical assertion.


\medskip\noindent\textbf{Discussion.}


\textbf{Lossiness is the key property.} The observation function $h$ maps from $\Omega_t$ (high-dimensional, inaccessible) to $\mathcal{O}$ (what the agent actually receives). Information is destroyed in this mapping. This is what forces the agent to maintain a model ( \S\,\ref{agent-model}) rather than simply reading off the environment state.

\textbf{Action-dependence.} The inclusion of $a_{t-1}$ allows the observation function to capture active perception: what the agent sees may depend on where it looked. This is relevant in software ( \#software-epistemic-properties — cross-component reference, see \texttt{02-tst-core/}), where observation quality is partly under agent control.

\end{definition}

\begin{definition}[Action and Transition]
\label{action-transition}
Actions affect the environment through a transition function that is unknown to the agent and possibly stochastic.



\medskip\noindent\textbf{Formal Expression.}



The \textbf{action space} $\mathcal{A}$ is the set of actions available to the agent. Actions affect the environment via the transition function:

$$\Omega_{t+1} \sim T(\cdot \mid \Omega_t, a_t)$$

where:
\begin{itemize}
  \item $T$ is the (possibly stochastic) transition function
  \item $\Omega_t$ is the current environment state
  \item $a_t \in \mathcal{A}$ is the agent's chosen action
\end{itemize}



The agent does not know $T$ exactly.


\medskip\noindent\textbf{Epistemic Status.}


This is \textit{definitional}. The transition function $T$ is a modeling device that captures how agent actions couple back into the environment. The stochasticity of $T$ is allowed but not required — deterministic transitions are the special case where $T$ places all mass on a single successor state. The claim that $T$ is unknown to the agent is constitutive of the uncertainty setting, paralleling the epistemic opacity of $h$ ( \S\,\ref{observation-function}).


\medskip\noindent\textbf{Discussion.}


\textbf{Closing the loop.} Together with \S\,\ref{observation-function}, this definition completes the agent-environment coupling: the agent observes via $h$ and acts via $T$. The loop $\Omega_t \xrightarrow{h} o_t \rightarrow \text{agent} \xrightarrow{a_t} \Omega_{t+1}$ is the fundamental structure that all subsequent claims build on.

**Uncertainty about $T$ is what makes action non-trivial.** If the agent knew $T$ exactly, action selection would reduce to optimization over a known function. The combination of unknown $h$ and unknown $T$ is what creates the need for adaptive behavior.

\end{definition}

\begin{scope}[Scope Condition]
\label{scope-condition}
ACT's scope is defined by progressive narrowing. The broadest applicable scope — any system that observes under uncertainty — supports Section I's adaptive machinery. Adding causal action unlocks the interventional and purposeful results of Sections II and III.



\medskip\noindent\textbf{Formal Expression.}




\medskip\noindent\textbf{Adaptive scope (Section I).}


$$\mathcal S_\text{adaptive} = \left\{(\text{Agent}, \Omega) \;:\; \mathcal O \neq \emptyset, \;\; H(\Omega_t \mid \mathcal C_t) > 0 \right\}$$

Two conditions:

\begin{enumerate}
  \item \textbf{Observations exist}: $\mathcal O \neq \emptyset$ — the system has some perceptual channel to the environment ( \S\,\ref{observation-function})
  \item \textbf{Residual uncertainty persists}: $H(\Omega_t \mid \mathcal C_t) > 0$ — the environment is not fully determined by the interaction history
\end{enumerate}


This is sufficient for the mismatch signal ( \S\,\ref{mismatch-signal}), update gain ( \S\,\ref{update-gain}), adaptive tempo ( \S\,\ref{adaptive-tempo}), the persistence condition ( \S\,\ref{persistence-condition}), and all of Section I's adaptive dynamics. A Kalman filter estimating a passive signal, a passive Bayesian learner, and any system that observes and updates a model under uncertainty are within this scope.


\medskip\noindent\textbf{Agency scope (Sections II and III).}


$$\mathcal S_\text{agency} = \mathcal S_\text{adaptive} \;\cap\; \left\{(\text{Agent}, \Omega) \;:\; \lvert\mathcal A\rvert \geq 2, \;\; \exists\, a \neq a' \text{ s.t. } P(o \mid do(a)) \neq P(o \mid do(a')) \right\}$$

Two additional conditions, narrowing the adaptive scope:

\begin{enumerate}
  \item \textbf{At least binary choice}: $\lvert\mathcal A\rvert \geq 2$ — the agent can choose between at least two actions ( \S\,\ref{action-transition})
  \item \textbf{At least one action has causal effect}: there exist distinct actions $a, a'$ whose interventional outcome distributions differ (where $do(\cdot)$ is Pearl's intervention operator; see \S\,\ref{pearl-causal-hierarchy}) — the agent's choices make a difference to what it can observe
\end{enumerate}


These are required for the adaptive loop to generate interventional data ( \S\,\ref{loop-interventional-access}), for the causal hierarchy requirement ( \S\,\ref{causal-hierarchy-requirement}) to be well-posed, and for the purposeful-agent machinery of Section II ($O_t$, $\Sigma_t$, the orient cascade) to be non-vacuous. Section III's composition theory inherits this requirement.


\medskip\noindent\textbf{Epistemic Status.}


This is a \textit{scope definition} — it draws the boundary around the systems each part of ACT addresses. The conditions are not derived; they are the minimal requirements for the associated machinery to be non-vacuous. The progressive structure reflects the theory's own architecture: Section I's adaptive results need only observations and uncertainty; Section II's purposeful results additionally need causal action.


\medskip\noindent\textbf{Discussion.}


\textbf{What is included.} The adaptive scope is broad: any system that observes under uncertainty. Passive Bayesian learners, Kalman filters (with or without control inputs), biological sensory systems. The agency scope narrows to systems whose actions make a causal difference: thermostats, Kalman filters with control inputs, RL agents, military commanders, software developers, AI agents with tool use. These are instances of the same formal framework at different points in the agent spectrum ( \S\,\ref{agent-spectrum}).

\textbf{What is excluded from both scopes.}

\begin{itemize}
  \item \textbf{Closed-form systems} ($H(\Omega_t \mid \mathcal C_t) = 0$): When the agent has complete knowledge of the environment, there is no uncertainty to adapt to. Optimal control over known dynamics is a solved problem outside ACT's concerns.
  \item \textbf{Pure computation} ($\mathcal O = \emptyset$): A system with no observation channel — e.g., a mathematical proof engine operating on axioms alone — has no agent-environment boundary in ACT's sense.
\end{itemize}


\textbf{What is in adaptive scope but excluded from agency scope.}

\begin{itemize}
  \item \textbf{Passive observers} ($\lvert\mathcal A\rvert < 2$): Can observe and model, but cannot intervene. Section I's adaptive machinery applies; the causal-information and purposeful-agent results do not.
  \item \textbf{Nominal agents} ($P(o \mid do(a)) = P(o \mid do(a'))$ for all $a, a'$): Have choices that make no difference. Can estimate but cannot learn causal structure. Same as passive observers for ACT's purposes: adaptive scope only.
\end{itemize}


\textbf{Why causal effect matters for the agency scope.} Binary choice ($\lvert\mathcal A\rvert \geq 2$) is necessary but not sufficient. Two actions that produce identical outcome distributions provide no interventional contrast — the agent cannot learn which action produces which effect because the effects are the same. The causal-effect condition ensures at least one meaningful contrast exists, which is what \S\,\ref{loop-interventional-access} needs to generate Level 2 data.

\end{scope}

\begin{postulate}[Composition Consistency]
\label{composition-consistency}
ACT applies at every level of description where the scope condition is met. A team of agents is itself an agent; a department of teams is itself an agent. The theory's predictions at each level must be compatible. Section III develops conditions under which they are — specifically, a condition analogous to the persistence threshold: the composite's internal coordination must be fast relative to the external dynamics it faces. That derivation requires a contraction assumption beyond the core ACT machinery ( \S\,\ref{composition-closure}).



\medskip\noindent\textbf{Formal Expression.}



For any system $S$ satisfying the scope condition ( \S\,\ref{scope-condition}), and any decomposition of $S$ into subsystems $\{S_1, \ldots, S_n\}$ where each $S_i$ also satisfies the scope condition, ACT's predictions at the system level must be compatible with its predictions at the subsystem level. Specifically, composition laws must exist such that:

\begin{enumerate}
  \item \textbf{Tempo composition}: $\mathcal T_S$ is expressible as a function of $\{\mathcal T_{S_i}\}$ and the coordination structure among them
  \item \textbf{Persistence compatibility}: the system's persistence is derivable from the sub-agents' individual persistence conditions plus coordination structure
  \item \textbf{Mismatch consistency}: $\delta_S$ is derivable from $\{\delta_{S_i}\}$ and their interaction structure
\end{enumerate}


% [Sufficient Condition (composition-validity)]

\textbf{Timescale separation condition.} A group of sub-agents behaves as a valid composite agent when its internal equilibration timescale $\tau_{\text{eq}}$ — the time for sub-agents to approximately synchronize models and coordinate actions — is short relative to the external dynamics timescale $\tau_{\text{ext}}$ — the time for the environment to change significantly from the macro-agent's perspective:

$$\tau_{\text{eq}} \ll \tau_{\text{ext}}$$

When this holds, the internal dynamics have approximately settled by the time the next external challenge arrives, and the composite's macro-state is predictive of its macro-behavior. This is the composition analog of the persistence condition ( \S\,\ref{persistence-condition}): just as an individual agent requires $\mathcal{T} > \rho / \Vert\delta_{\text{critical}}\Vert$ to maintain bounded mismatch, a composite agent requires fast internal coordination to maintain a coherent macro-description.

Most functioning groups easily satisfy this condition. A software team with daily standups and shared CI ($\tau_{\text{eq}} \sim$ hours) facing weekly feature deadlines ($\tau_{\text{ext}} \sim$ weeks) is comfortably a valid composite. A military squad communicating by voice ($\tau_{\text{eq}} \sim$ seconds) in a tactical situation evolving over minutes is clearly a single agent. The theory applies broadly; the interesting questions arise near the threshold.


\medskip\noindent\textbf{Epistemic Status.}


\textit{Axiomatic.} The meta-requirement (cross-level compatibility) is a structural requirement for ACT's internal consistency — if the scope condition doesn't restrict which level the theory applies to, the predictions must not contradict across levels. The timescale separation condition is stated here as a sufficient condition for practical composition; its formal derivation requires the composition closure criterion ( \S\,\ref{composition-closure}) and the tempo composition inequality ( \S\,\ref{tempo-composition}) developed in Section III. The condition is not yet proved; it is stated early because it is intuitive, likely correct, and gives readers an immediate practical test.


\medskip\noindent\textbf{Discussion.}


\textbf{The same pattern as individual persistence.} The persistence condition says: there is a measurable threshold below which an individual agent's mismatch grows without bound and the agent degrades. Composition consistency says the same thing at the composite level: there is a measurable threshold ($\tau_{\text{eq}} / \tau_{\text{ext}}$) below which the composite description breaks down. In both cases, most competent agents are comfortably above the threshold — the theory applies broadly, and the edge cases are where it gets interesting.

\textbf{What this buys the theory.} With composition consistency stated early:
\begin{itemize}
  \item Every subsequent result (persistence, gain, tempo, mismatch decomposition, orient cascade) is understood to apply at every level of composition where the timescale condition holds
  \item Section III becomes the study of \textit{what happens near and beyond the threshold} — coordination overhead, unity dimensions, adversarial dynamics — rather than a separate multi-agent theory
  \item The formal test for composition validity ( \S\,\ref{composition-closure}) develops the rigorous version of the timescale condition, conditional on a contraction assumption beyond (A4)
\end{itemize}


\textbf{When composition fails.} The timescale condition fails when:
\begin{itemize}
  \item Internal coordination slows (communication overhead scales poorly, conflicts consume attention, bureaucratic process)
  \item External dynamics accelerate (adversary acts faster, market shifts, crisis compresses decision timescales)
  \item Both simultaneously (the classic organizational failure mode — internal friction increases while external demands intensify)
\end{itemize}


This is the formal analog of Brooks's Law: adding people to a late project increases $\tau_{\text{eq}}$ (more coordination) while $\tau_{\text{ext}}$ stays fixed (the deadline doesn't move). Eventually $\tau_{\text{eq}}$ exceeds $\tau_{\text{ext}}$ and the composite ceases to function as a coherent agent. The model captures the same structural pattern; whether the specific mechanism (timescale crossing) is the actual cause of Brooks's Law is an empirical question.

\textbf{The boundary is a modeling choice.} A development team is simultaneously: individual developers (each an ACT agent), the team (a composite ACT agent), and part of an organization (a sub-agent within a larger composite). The scope condition is satisfied at every level. Composition consistency ensures the theory doesn't give contradictory answers about observable quantities (e.g., whether the team persists) regardless of which boundary is chosen.

\textbf{What composition consistency does NOT say.} It does not specify the form of the composition laws — those are derived in Section III ( \S\,\ref{tempo-composition}). It does not say every decomposition is equally useful for analysis. And it does not require perfect internal coordination — only that internal equilibration is fast relative to external dynamics.

\end{postulate}

\begin{postulate}[Causal Structure]
\label{causal-structure}
The agent-environment interaction has irreducible causal structure grounded in the temporal ordering of events. Actions precede their consequences; observations follow from the state they observe. This ordering is constitutive of the feedback loop and holds independent of the magnitude of the agent's influence on the environment.



\medskip\noindent\textbf{Formal Expression.}



The interaction history $\mathcal C_t$ ( \S\,\ref{chronica}) is not merely a set of observations and actions — it is an \textit{ordered sequence} in which temporal position carries meaning. $a_{t-1}$ was selected before $o_t$ was received. The agent could not have used $o_t$ to select $a_{t-1}$. This asymmetry — the arrow of time — is the foundation of causal structure in the theory.

We adopt the most primitive notion of causality: **event $A$ can be a cause of event $B$ only if $A$ temporally precedes $B$.\textit{* This is weaker than (and prior to) statistical notions of causality. It is a statement about the }structure of possible influence*, not about actual influence.


\medskip\noindent\textbf{Epistemic Status.}


This is a \textit{postulate} — the temporal ordering of events is a physical fact about the universe that the theory takes as given. The second law of thermodynamics, the light-cone structure of relativity, and the arrow of psychological time all enforce it, but ACT does not derive it from any of these. It is simply noted as a precondition: the theory applies to agents embedded in a universe where time has a direction.


\medskip\noindent\textbf{Discussion.}


\textbf{Causality as temporal ordering is the most primitive notion.} Three levels of causal reasoning derive from this foundation ( \S\,\ref{pearl-causal-hierarchy}), but the temporal notion survives even when statistical influence is negligible. An agent passively observing a system with minimal intervention still has a causal history — the temporal ordering of its observations and actions structures what it can learn and when.

\textbf{Causal structure independent of coupling strength.} The causal structure of the feedback loop is preserved even when the agent's actions have minimal effect on the environment:

\begin{itemize}
  \item \textbf{Strong coupling} ($a_t$ significantly affects $\Omega_{t+1}$): Robot manipulation, military action. Interventional information is rich.
  \item \textbf{Weak coupling} ($a_t$ marginally affects $\Omega_{t+1}$): Scientific observation, small financial trades. Interventional information is sparse but non-zero.
  \item \textbf{Nominal coupling} ($a_t$ negligibly affects $\Omega_{t+1}$, but the agent's \textit{choice of what to observe} produces distinguishable observation distributions): Near-passive, but still within scope — the agent's query actions generate weak but nonzero interventional contrasts. The theory applies but the interventional information per action is sparse.
  \item \textbf{Zero coupling} ($T(\Omega_{t+1} \mid \Omega_t, a_t) = T(\Omega_{t+1} \mid \Omega_t)$ for all $a_t$ AND observation distributions are action-independent): Actions don't affect the environment or the observations. Level 2 access vanishes. The feedback "loop" collapses to a one-way channel. \textbf{Outside the agency scope} ( \S\,\ref{scope-condition}, conditions 3-4) — the causal-information, purposeful-agent, and composition results of Sections II and III do not apply. However, zero-coupling systems remain \textbf{within the adaptive scope} ( \S\,\ref{scope-condition}, conditions 1-2) if they observe under residual uncertainty: Section I's adaptive machinery (mismatch, gain, tempo, persistence) applies to passive estimators. The causal structure postulate still holds for such systems — temporal ordering of observations is constitutive — but without interventional contrasts, the causal \textit{hierarchy} (Level 2, Level 3) is inaccessible.
\end{itemize}


The \textit{agency-scope} results apply to any agent whose choices make a causal difference to what it can observe, from strong coupling (robot manipulation) through weak coupling (scientific observation) to query-only coupling (choosing which question to ask). The \textit{adaptive-scope} results apply more broadly, including to passive observers whose actions have no causal effect.

\textbf{Consequences for the feedback loop.} The irreversibility of temporal ordering yields the core structure:

\begin{itemize}
  \item The model update is \textbf{directed} — the model at time $t$ depends on prior events, never on future ones
  \item The mismatch signal $\delta_t$ ( \S\,\ref{mismatch-signal}) is \textbf{retrospective} — comparing a prediction (made before $o_t$) with an observation (arriving after)
  \item Action selection is \textbf{prospective} — using the current model to influence future events
  \item The chronica ( \S\,\ref{chronica}) is \textbf{monotonically growing} — events are added but never removed
\end{itemize}


\textbf{Implications for model updating.} The causal postulate constrains the update rule: the model should give more weight to observations that are \textit{causally downstream} of the agent's actions than to observations that would have occurred regardless. Action-contingent observations carry interventional (Level 2) information; action-independent observations carry only associational (Level 1) information. The formal measure of this distinction — causal information yield (CIY) — is developed in \S\,\ref{causal-information-yield}.

\textbf{(Descended from TF-02.)}

\end{postulate}

\begin{definition}[Pearl's Causal Hierarchy]
\label{pearl-causal-hierarchy}
Three levels of causal reasoning emerge from the causal structure of the feedback loop: association ("what if I observe?"), intervention ("what if I do?"), and counterfactual ("what if I had done differently?"). The binary action requirement ( \S\,\ref{scope-condition}) ensures at least Level 2 access is structurally available.



\medskip\noindent\textbf{Formal Expression.}



\textbf{Level 1 — Associational}: $P(o_t \mid \mathcal{C}_{< t})$

\textit{What will I observe next, given what I've observed before?}

Pattern recognition over the temporally ordered history. Available to any agent that maintains a model ( \S\,\ref{agent-model}), including purely passive observers. The temporal ordering constrains which associations are meaningful: $o_3$ can depend on $o_1, a_1, o_2, a_2$ but not on $o_4$.

\textbf{Level 2 — Interventional}: $P(o_t \mid do(a_{t-1}), M_{t-1})$

\textit{What will I observe if I} do \textit{this?}

The $do(\cdot)$ operator marks the crucial distinction: this is not "what observation tends to follow this action in the historical record" (associational) but "what will happen \textit{because} I take this action now." This requires: (1) the agent's action temporally precedes the observation ( \S\,\ref{causal-structure}), (2) the agent chose the action (it was not determined by the same causes that determine the observation), (3) the environment's response carries information about the causal relationship.

Level 2 is why the feedback loop is more powerful than passive observation. By \textit{acting} and then observing consequences, the agent obtains information about causal mechanisms — not merely about correlations. The mismatch signal $\delta_t$ ( \S\,\ref{mismatch-signal}), conditioned on the agent's own action, is an \textit{interventional} signal.

\textbf{Level 3 — Counterfactual}: $P(o_t^{a'} \mid a_{t-1} = a, o_t = o)$

*Given that I did $a$ and observed $o$, what would I have observed if I had done $a'$ instead?*

This requires the model to simulate alternative histories — running the causal structure "backward" and then "forward" under different interventions. It is the most demanding epistemic level and the basis for regret computation, strategic simulation, and learning from single observations.


\medskip\noindent\textbf{Epistemic Status.}


This is \textit{definitional} — it names three modes of epistemic access that are structurally available within the feedback loop, following Pearl's causal hierarchy (Pearl 2009; Bareinboim et al. 2022). The hierarchy itself is a well-established result in causal inference. ACT's contribution is grounding it in the temporal structure of the agent-environment coupling rather than in abstract graphical models.


\medskip\noindent\textbf{Discussion.}


\textbf{Availability vs. exploitation.} The three levels describe epistemic access that the causal structure \textit{makes available} — not what any particular agent \textit{uses}. Many systems within ACT's scope operate primarily at Level 1. A Kalman filter coupled with an LQR controller has Level 2 access structurally present (its innovation signal is conditioned on prior action), but the separation principle guarantees estimation quality is invariant to control policy — the system does not \textit{exploit} the interventional structure. Only dual control (choosing actions partly for their informational value) exercises Level 2 access in this domain. Similarly, a PID controller has no deliberative capacity — it operates entirely at Level 1. Which levels an agent exercises depends on its architecture and model class.

\textbf{Forward-looking deliberation exercises Level 2, shading into Level 3.} Comparing candidate actions before choosing — "what will happen if I do X vs Y?" — primarily exercises Level 2 (iterated mental intervention). When the agent evaluates past choices to refine the comparison ("given what happened when I tried X last time, what would Y have produced?"), it exercises Level 3.

\textbf{The causal hierarchy theorem.} Bareinboim et al. (2022) prove that the three levels form a strict hierarchy: Level 2 knowledge cannot in general be computed from Level 1 data alone, and Level 3 cannot be computed from Level 2 alone. This is load-bearing for ACT's Section II: evaluating $Q_O(M_t, a; \cdot)$ is a Level 2 query, so agents that need to \textit{learn} action consequences during operation require causal structure beyond predictive models ( \S\,\ref{causal-hierarchy-requirement}).

\textbf{Software as a uniquely rich domain for this hierarchy.} In most domains, Level 3 counterfactuals require model-based simulation with uncertain fidelity. In software development, \texttt{git checkout} provides Level 3 access with ground-truth verification — the agent can literally execute the counterfactual. This is one of software's unique epistemic properties ( \#software-epistemic-properties — cross-component reference, see \texttt{02-tst-core/}) and makes it an ideal testbed for causal reasoning within ACT.

\textbf{Domain instantiations of the three levels:}

\begin{adjustbox}{max width=\textwidth}
\small
\begin{tabular}{llll}
\toprule
Domain & Level 1 (Association) & Level 2 (Intervention) & Level 3 (Counterfactual) \\
\midrule
Kalman filter & Prediction from state estimate & Innovation conditioned on action & Not typically exercised \\
RL agent & Value function prediction & Action → reward observation & Regret computation \\
Scientific method & Correlational observation & Experimental intervention & "What if we had used control X?" \\
Military (Boyd) & Pattern recognition & Probe/feint → observe response & "What if we had attacked from the flank?" \\
Software developer & "I think this function does X" & Run test → observe result & \texttt{git checkout} + alternative implementation \\
Immune system & Antigen pattern matching & Antibody → pathogen response & Not exercised (no counterfactual reasoning) \\
\bottomrule
\end{tabular}
\end{adjustbox}


\textbf{(Descended from TF-02.)}

\end{definition}

\begin{definition}[Chronica]
\label{chronica}
The interaction history $\mathcal C_t$ is the complete, singular causal record of the agent's observations and actions. Everything the agent can ever know must be constructed from this sequence.



\medskip\noindent\textbf{Formal Expression.}



$$\mathcal{C}_t = (o_1, a_1, o_2, a_2, \ldots, a_{t-1}, o_t)$$

The ordering is not a notational convenience. It reflects an irreversible physical fact: $a_{t-1}$ was selected before $o_t$ was received. The agent could not have used $o_t$ to select $a_{t-1}$.

$\mathcal C_t$ is monotonically growing — events are added but never removed. It is the agent's \textit{only} raw material for constructing a model ( \S\,\ref{agent-model}).


\medskip\noindent\textbf{Epistemic Status.}


This is \textit{definitional}. The chronica names an object that exists by construction in any system satisfying \S\,\ref{agent-environment}: the temporal sequence of all agent-environment interactions. The term "chronica" (from Greek χρονικά, "records of time") avoids collision with $\mathcal{H}$ (Shannon entropy) in speech and notation.


\medskip\noindent\textbf{Discussion.}


\textbf{The chronica is singular and non-forkable.} Because the temporal ordering is irreversible, $\mathcal C_t$ represents a unique causal trajectory. Duplicating an agent's state and exposing the copies to different future events creates two agents with divergent chronica, neither of which is a sufficient statistic for the other's trajectory. The chronica is the substrate of \textit{continuity persistence} (see Persistence in \texttt{LEXICON.md}) — an agent has continuity persistence when $\mathcal{C}_t$ extends continuously and $M_t$ has temporal depth grounded in it. See \S\,\ref{agent-identity} for the full development of this observation.

\textbf{Relationship to the model.} The model $M_t = \phi(\mathcal C_t)$ ( \S\,\ref{agent-model}) is a compression of the chronica. How much of $\mathcal C_t$'s predictive information survives compression is measured by model sufficiency ( \S\,\ref{model-sufficiency}).

\end{definition}

\begin{formulation}[The Reality Model]
\label{agent-model}
The agent's compressed representation of how the world works, mapping interaction history to model space. $M_t$ is the substrate of prolepsis — the model from which predictions are generated and against which observations are compared. This is a formulation choice — we commit to analyzing the agent as having a complete state $M_t$ that subsumes all retained information from its history.



\medskip\noindent\textbf{Formal Expression.}



$$M_t = \phi(\mathcal{C}_t)$$

where:
\begin{itemize}
  \item $\phi: \mathcal{C}^* \to \mathcal{M}$ maps interaction history to model space $\mathcal{M}$
  \item $\mathcal C_t = (o_1, a_1, \ldots, o_t)$ is the chronica ( \S\,\ref{chronica}) — the complete record of agent-environment interaction
  \item $\mathcal{M}$ is the space of possible models the agent can hold
\end{itemize}


The mapping $\phi$ is a many-to-one compression: multiple distinct histories may produce the same model state. This is not a deficiency — it is the essential function of the model: retaining what matters and discarding what does not.


\medskip\noindent\textbf{Epistemic Status.}


\textit{Robust qualitative.} This is a \textit{formulation} — a representational commitment, not a derived result. We choose to analyze agents as maintaining a state object $M_t$ that mediates between history and future action. Alternative formulations exist (e.g., history-based policies that map $\mathcal C_t$ directly to actions without an explicit model). The formulation is justified by its analytical utility: it enables the information bottleneck analysis ( \S\,\ref{information-bottleneck}), the mismatch decomposition ( \S\,\ref{mismatch-signal}), and the gain principle ( \S\,\ref{update-gain}). The formulation is robust — any agent that conditions its actions on retained information can be described this way — but the specific commitment to a complete, compressed state $M_t$ is a modeling choice, not a derivation.


\medskip\noindent\textbf{Discussion.}


**$M_t$ is the epistemic substate.** It captures "what the agent believes about reality." Different agents realize $M_t$ differently: a Kalman filter holds a state estimate and covariance matrix; an RL agent holds a value function; a developer holds a mental model of codebase architecture; an LLM agent holds its context window contents plus retrieved memory. The formalism is agnostic to the realization — it asks only that $M_t$ exist as a well-defined object that the agent's policy can condition on.

\textbf{Completeness assumption.} By writing $M_t = \phi(\mathcal C_t)$, we assume that $M_t$ captures everything the agent retains from its history. Any information not in $M_t$ is lost to the agent. This is what makes $M_t$ the complete epistemic substate, not merely one component of a richer internal representation. Whether $M_t$ retains \textit{enough} information is the subject of \S\,\ref{model-sufficiency}.

\textbf{Degenerate cases.} A PID controller's $M_t$ is degenerate — it retains only the error signal and its history (integral, derivative), with no predictive capability beyond extrapolating recent trends. It occupies the "blind pursuer" region of the agent spectrum ( \S\,\ref{agent-spectrum}): its $O_t$ (setpoint) is clear but its $M_t$ is too impoverished to support the adaptive dynamics of Section I. The formalism accommodates this by allowing $\mathcal{M}$ to range from trivial (scalar) to rich (full world model).

\end{formulation}

\begin{formulation}[Information Bottleneck]
\label{information-bottleneck}
Optimal model compression balances retained history against predictive power; the information bottleneck objective provides a principled framework for understanding this trade-off.



\medskip\noindent\textbf{Formal Expression.}



$$\phi^* = \arg\min_{\phi} \left[ I(M_t;\, \mathcal{C}_t) - \beta \cdot I(M_t;\, o_{t+1:\infty} \mid a_{t:\infty}) \right]$$

where:
\begin{itemize}
  \item $I(M_t;\, \mathcal{C}_t)$ is the compression cost — how much of the interaction history the model retains
  \item $I(M_t;\, o_{t+1:\infty} \mid a_{t:\infty})$ is the predictive power — how much the model tells the agent about future observations given future actions
  \item $\beta > 0$ is the trade-off parameter controlling the compression-prediction balance
\end{itemize}


\textbf{Dependence on volatility.} The trade-off $\beta$ depends on environment volatility $\rho$:

\begin{itemize}
  \item \textbf{Volatile environments} (high $\rho$): favor aggressive compression (low $\beta$). Old information decays in relevance quickly, so retaining it wastes capacity.
  \item \textbf{Stable environments} (low $\rho$): favor dense retention (high $\beta$). Historical information remains predictive, so discarding it loses value.
\end{itemize}



\medskip\noindent\textbf{Epistemic Status.}


This is a \textit{formulation} — it provides a principled framework for understanding compression trade-offs, not a claim about how actual agents compute their models. No agent explicitly solves this optimization (except in the variational IB literature, where it is a training objective). The formulation is analytically useful because it makes the compression-prediction trade-off explicit and connects model quality to information-theoretic quantities.


\medskip\noindent\textbf{Discussion.}


\textbf{The IB framework is not prescriptive.} It characterizes what an optimal $\phi$ would look like, not how to find one. Actual agents approximate this trade-off through diverse mechanisms: forgetting, attention, abstraction, summarization.

\textbf{Connection to model sufficiency.} The IB objective implicitly defines when a model is "good enough": when the predictive power term $I(M_t;\, o_{t+1:\infty} \mid a_{t:\infty})$ is close to its maximum (the full history's predictive power). This is formalized in \S\,\ref{model-sufficiency}.

\textbf{Policy-relativity.} The conditioning on $a_{t:\infty}$ makes the predictive power term policy-relative: it measures predictive information given a particular sequence of future actions, which depends on what policy the agent follows. The IB objective is therefore defined relative to a generating policy. This is inherent — what information is "predictive" depends on what the agent will do. \S\,\ref{value-object}'s continuation-policy convention ($\pi_{\text{cont}}$) provides the specification for value computations; the same convention should be understood as implicit in the IB formulation. The $\beta$ parameter's dependence on volatility $\rho$ also has a policy component: an exploratory policy encounters more diverse situations, making more information predictively relevant (favoring retention), while an exploitative policy encounters a narrower distribution (favoring compression).

\textbf{Broader applicability.} The same IB principle applies beyond intra-agent compression. It governs inter-agent communication ( \S\,\ref{shared-intent}) — how much of one agent's model or strategy to transmit to another — and constrains the cognitive cost of maintaining a complex strategy. In each case, the trade-off is between the cost of retaining or transmitting information and the value of that information for future decisions.

\end{formulation}

\begin{definition}[Model Sufficiency]
\label{model-sufficiency}
The fraction of predictive information the model retains relative to the full interaction history; $S = 1$ means the model is a sufficient statistic for prediction, $S < 1$ means predictive information has been lost.



\medskip\noindent\textbf{Formal Expression.}



$$S(M_t) = 1 - \frac{I(\mathcal{C}_t;\, o_{t+1:\infty} \mid M_t,\, a_{t:\infty})}{I(\mathcal{C}_t;\, o_{t+1:\infty} \mid a_{t:\infty})}$$

where:
\begin{itemize}
  \item The numerator $I(\mathcal C_t;\, o_{t+1:\infty} \mid M_t,\, a_{t:\infty})$ is the predictive information that the full history $\mathcal C_t$ carries about the future \textit{beyond} what $M_t$ already captures — the information lost by compression
  \item The denominator $I(\mathcal C_t;\, o_{t+1:\infty} \mid a_{t:\infty})$ is the total predictive information in the full history
\end{itemize}


\textbf{Boundary values:}
\begin{itemize}
  \item $S(M_t) = 1$: $M_t$ is a sufficient statistic — it captures all predictive information in $\mathcal C_t$. Knowing the full history beyond $M_t$ adds nothing.
  \item $S(M_t) = 0$: $M_t$ retains no predictive information. The model is useless for prediction.
  \item $0 < S(M_t) < 1$: partial sufficiency — some predictive information is retained, some lost.
\end{itemize}



\medskip\noindent\textbf{Epistemic Status.}


This is \textit{definitional} — it names and formalizes a quantity. The definition is well-grounded in information theory (conditional mutual information ratios are standard). No substantive claim is made here about what value $S(M_t)$ takes or what happens when it is low; those claims belong to \S\,\ref{model-class-fitness} and \S\,\ref{structural-adaptation-necessity}.


\medskip\noindent\textbf{Discussion.}


\textbf{Sufficiency is relative to the prediction task.} $S(M_t)$ measures sufficiency for predicting future observations given future actions. A model that is sufficient for one prediction horizon may be insufficient for another. The infinite-horizon formulation ($o_{t+1:\infty}$) is the most demanding; practical sufficiency over finite horizons may be easier to achieve.

\textbf{Sufficiency vs. accuracy.} A model can be sufficient ($S = 1$) while being wrong in absolute terms — if the full history is also wrong (e.g., systematically biased observations). Sufficiency measures information retention, not truth. The mismatch signal ( \S\,\ref{mismatch-signal}) measures accuracy; sufficiency measures completeness of compression.

\textbf{Sufficiency is predictive, not causal.} $S(M_t)$ measures retained \textit{predictive} information — a Level 1 (associational) property. It does not by itself guarantee that $M_t$ supports Level 2 (interventional) queries such as $P(o \mid do(a), M_t)$. The causal validity of value computations conditioned on $M_t$ requires an additional condition: that $M_t$ satisfies the backdoor criterion with respect to the agent's actions (see \S\,\ref{value-object}). For agents whose actions are deterministic functions of $M_t$ (standard in ACT: $a_t = \pi(M_t, G_t)$), $S(M_t) = 1$ is nearly sufficient for causal validity — the remaining requirement is that no unmodeled external factor influences both action selection and outcomes. But predictive sufficiency alone does not collapse the distinction between Level 1 and Level 2 that \S\,\ref{causal-hierarchy-requirement} establishes.

\textbf{Policy-relativity.} The conditioning on $a_{t:\infty}$ makes $S(M_t)$ implicitly policy-relative: different policies generate different future action sequences, which changes which future observations are relevant and therefore what "predictive information" means. A model that is sufficient under a conservative policy may be insufficient under an aggressive one (the aggressive policy visits states the model cannot predict). This policy-relativity is inherent in any predictive sufficiency measure — it is not an artifact of the formulation. When comparing sufficiency values, the generating policy must be held constant or specified. \S\,\ref{value-object}'s continuation-policy convention ($\pi_{\text{cont}}$) provides the required specification for value computations; the same convention should be understood as implicit here.

\end{definition}

\begin{definition}[Model Class Fitness]
\label{model-class-fitness}
The best achievable sufficiency within a model class. When no model in the class can adequately represent reality, the agent faces a structural limitation that no amount of parameter tuning can resolve.



\medskip\noindent\textbf{Formal Expression.}



$$\mathcal{F}(\mathcal{M}) = \sup_{M \in \mathcal{M}} S(M)$$

where $\mathcal{M}$ is the model class — the set of all models the agent can represent given its current architecture, parameterization, or capacity.

\textbf{Structural inadequacy condition:}

$$\mathcal{F}(\mathcal{M}) < 1 - \varepsilon$$

When this holds, no model $M \in \mathcal{M}$ achieves sufficiency above $1 - \varepsilon$. The gap is structural: it cannot be closed by better parameter estimation, more data, or longer training within the current class. This is the trigger for structural change ( \S\,\ref{structural-adaptation-necessity}).


\medskip\noindent\textbf{Epistemic Status.}


This is \textit{definitional} — it names the supremum of sufficiency over a model class. The definition itself is straightforward. The substantive claim about what happens when $\mathcal{F}(\mathcal{M})$ is low — that parametric updates cannot close the mismatch floor and structural adaptation becomes necessary — is developed in \S\,\ref{structural-adaptation-necessity}.


\medskip\noindent\textbf{Discussion.}


\textbf{Model class vs. model instance.} $S(M_t)$ measures a specific model's sufficiency at time $t$. $\mathcal{F}(\mathcal{M})$ measures the ceiling of the entire class. A low $S(M_t)$ might mean the agent needs more learning (parameter update). A low $\mathcal{F}(\mathcal{M})$ means the agent needs a different kind of model (structural change). The distinction parallels bias vs. variance: class fitness is about bias; instance sufficiency reflects both bias and estimation quality.

\textbf{Detecting low class fitness.} The agent cannot directly compute $\mathcal{F}(\mathcal{M})$ — it would need to search over all models in the class. Instead, persistent mismatch despite adequate learning (high gain, sufficient data, converged parameters) is the observable signature. This connects to the mismatch floor in \S\,\ref{structural-adaptation-necessity}.

\end{definition}

\begin{formulation}[Event-Driven Dynamics]
\label{event-driven-dynamics}
The coupling between agent and environment occurs through discrete events — observations arriving and actions completing — at potentially variable and heterogeneous rates. Discrete-time notation is the special case of uniform-interval events on a single channel.



\medskip\noindent\textbf{Formal Expression.}



\textbf{Event} ($e$): An atomic unit of agent-environment interaction, typed as:
\begin{itemize}
  \item \textbf{Observation event}: $e = (\text{obs}, k, o^{(k)})$ — a datum arriving on observation channel $k$
  \item \textbf{Action completion}: $e = (\text{act}, j, r^{(j)})$ — the result of action $j$ completing
\end{itemize}


\textbf{Event stream} ($\mathcal{E}$): The temporally ordered sequence of all events:

$$\mathcal{E} = \{(e_1, \tau_1), (e_2, \tau_2), \ldots\} \quad \text{where } \tau_1 \leq \tau_2 \leq \cdots$$

\textbf{Channel rate} ($\nu^{(k)}$): The characteristic event rate of channel $k$, which may vary over time.

\textbf{Event information content}: The mutual information between the event and the environment state, conditioned on the current model:


$$\mathcal{I}(e_\tau) = I(e_\tau;\, \Omega_\tau \mid M_{\tau^-})$$

An event that the model already predicts carries little information ($\mathcal{I} \approx 0$). An event that surprises the model carries much ($\mathcal{I} \gg 0$). This connects directly to the mismatch signal ( \S\,\ref{mismatch-signal}).

\textbf{Channel-specific observation uncertainty}:


$$U_o^{(k)} = \text{observation uncertainty of channel } k$$

Different channels have different noise characteristics. A noisy channel (high $U_o^{(k)}$) provides lower-quality information per event. The update gain ( \S\,\ref{update-gain}) should weight channels accordingly.


\medskip\noindent\textbf{Epistemic Status.}


This is a \textit{formulation choice}, not a postulate. The event-driven representation extends \S\,\ref{causal-structure}'s recursive update to heterogeneous, asynchronous multi-channel interactions. The discrete-time form ($M_t = f(M_{t-1}, o_t, a_{t-1})$) from \S\,\ref{recursive-update} is a special case sufficient for many formal analyses — the event-driven formulation is needed only when multi-rate or asynchronous channels matter.


\medskip\noindent\textbf{Discussion.}


\textbf{Why events rather than clock ticks.} The discrete-time notation $M_t = f(M_{t-1}, o_t, a_{t-1})$ presupposes a single clock synchronizing observations and actions. Real agents face:

\begin{itemize}
  \item \textbf{Multiple observation channels} at different rates (a robot's camera at 30Hz, LIDAR at 10Hz, GPS at 1Hz; a human's vision, audition, proprioception; a developer's compiler output, test results, and production telemetry)
  \item \textbf{Multiple action channels} with different latencies (a robot's wheel motors vs. arm actuators; an organization's operational decisions vs. strategic pivots)
  \item \textbf{Asynchronous arrival} — observations not synchronized with each other or with action completions
\end{itemize}


The event-driven formulation handles all of these naturally. The discrete-time form is the special case where a single observation and a single action alternate at a fixed rate.

\textbf{The effective adaptation rate.} The agent's overall capacity to track environmental changes is the sum of information gained across all channels per unit time:

$$\nu_{\text{eff}} = \sum_k \nu^{(k)} \cdot \eta^{(k)*}$$

This quantity — identical to adaptive tempo $\mathcal{T}$ ( \S\,\ref{adaptive-tempo}) — is the central measure of an agent's adaptive fitness.

\textbf{Software-specific channels.} In the software development domain, the event-driven formulation maps naturally to the developer's multi-rate observation channels:

\begin{adjustbox}{max width=\textwidth}
\small
\begin{tabular}{lll}
\toprule
Channel $k$ & Rate $\nu^{(k)}$ & Noise $U_o^{(k)}$ \\
\midrule
Compiler/linter output & Per-save (high) & Very low \\
Unit test results & Per-run (medium) & Low \\
CI pipeline & Per-push (medium) & Low \\
Runtime telemetry & Continuous (variable) & Medium \\
Bug reports & Sporadic (low) & High \\
Code review feedback & Per-PR (low) & Medium-high \\
\bottomrule
\end{tabular}
\end{adjustbox}


The three-part tempo decomposition for software — $\mathcal T_{\text{obs}}$ (compiler, tests) + $\mathcal T_{\text{explore}}$ (code reading) + $\mathcal T_{\text{probe}}$ (test runs, staging) — is a direct application of multi-channel tempo. This decomposition is a Section IV gap (see the three-part tempo decomposition gap in \texttt{ACT-FULL.md}).

\textbf{(Descended from TF-04.)}

\end{formulation}

\begin{proposition}[Recursive Update]
\label{recursive-update}
Agent state updates (epistrophe — the corrective turning toward reality) must be recursive: the new model state is a function of the previous model state and the incoming event, not of the full interaction history. For finite agents this is computational necessity; for agents with unlimited computation it is the natural structure imposed by temporal ordering.



\medskip\noindent\textbf{Formal Expression.}


% [Derived (recursive-update, from temporal postulate and $M_t$ completeness)]

\textbf{Event-driven update:}

$$M_{\tau^+} = f_M(M_{\tau^-}, e_\tau)$$

where:
\begin{itemize}
  \item $M_{\tau^-}$ is the model state immediately before event $e_\tau$
  \item $M_{\tau^+}$ is the model state immediately after
  \item $f_M$ is the update function — it takes the current model and the new event, not the full history $\mathcal C_t$
\end{itemize}


\textbf{Between-event evolution:}

$$\frac{dM}{d\tau} = g_M(M_\tau)$$

Between events, the model evolves autonomously — internal reorganization, prediction generation, decay of transient states. The between-event dynamics depend only on the current model state, not on external input (which, by definition, arrives only at events).


\medskip\noindent\textbf{Epistemic Status.}


\textit{Exact, with a partly definitional character.} The result follows from three constraints: temporal ordering (C1 — physical law), partial observability (C2 — scope definition), and state completeness (C3 — analytical commitment that $M_t$ summarizes everything the agent retains). C1 and C2 do genuine eliminative work; C3 is definitional — it cannot be "violated" because any violation is absorbed by expanding $M_t$. The Markov structure is therefore not discovered in the environment but chosen through the definition of $M_t$ as complete. This is not a weakness — it is the nature of the claim: recursive update is the only form consistent with C1 + C2 + the definition of $M_t$ as complete (see \S\,\ref{recursive-update-derivation} for the full argument and seven counterexample attacks). For finite agents, recursion is also \textit{computational necessity}: re-processing the full history at each event is infeasible.


\medskip\noindent\textbf{Discussion.}


\textbf{Recursion as a consequence of completeness.} The recursive form is not an assumption bolted on — it follows from the definition of $M_t$ as complete. If $M_t$ were incomplete (if some relevant information lived outside $M_t$ in the raw history), then $f_M(M_{\tau^-}, e_\tau)$ would be insufficient and the agent would need to consult $\mathcal C_t$ directly. The sufficiency of the recursive form is precisely what \S\,\ref{model-sufficiency} measures: when $S(M_t) = 1$, the recursive update loses nothing.

\textbf{Between-event dynamics matter.} The autonomous evolution $g_M(M_\tau)$ is not merely filler between observations. It includes prediction generation (what the agent expects to see next), uncertainty growth (model confidence decaying over time without new data), and internal reorganization (consolidation, abstraction). In event-driven systems ( \S\,\ref{event-driven-dynamics}), the between-event interval is variable, making $g_M$ load-bearing for agents that must act or predict between observations.

\textbf{Connection to the update gain.} The event-driven update $f_M(M_{\tau^-}, e_\tau)$ is where the gain principle ( \S\,\ref{update-gain}) operates: $\eta^*$ determines how strongly $e_\tau$ shifts $M_t$ away from its prior value. The recursive form makes the gain's role explicit — it modulates the single-step correction.

\end{proposition}

\begin{proposition}[Action Selection]
\label{action-selection}
Praxis (informed action) is a function of the model. The model's role is not merely to represent the environment but to generate actions — either implicitly (from internalized patterns) or through explicit deliberation. The degree to which effective praxis flows from the model without deliberative computation is \textit{action fluency}.



\medskip\noindent\textbf{Formal Expression.}


% [Derived (action-selection, from agent-model completeness)]

$$a_t = \pi(M_t) \quad \text{(deterministic)}$$

$$a_t \sim \pi(\cdot \mid M_t) \quad \text{(stochastic)}$$

where $\pi$ is the agent's \textbf{policy} — the mapping from model state to action.

This is not imposed on the system but follows from \S\,\ref{agent-model}: $M_t$ is the agent's compressed history, and action depends on what the agent "knows" — i.e., on $M_t$. Any deterministic or stochastic dependence of action on history \textit{through} the model is captured by $\pi(M_t)$.


\medskip\noindent\textbf{Epistemic Status.}


\textit{Derived} from \S\,\ref{agent-model}'s completeness commitment. If $M_t$ is the agent's complete internal state (by definition), then action — which depends on internal state — is a function of $M_t$. The implicit/explicit distinction and action fluency concept are \textit{discussion-grade} — qualitative properties that follow from the formalism but are not formally derived as propositions.


\medskip\noindent\textbf{Discussion.}


\textbf{Implicit vs. explicit action selection.} A critical distinction emerges from the agent's \textit{action fluency} — the degree to which effective action flows from the model without deliberative computation:

\textbf{Implicit (model-embedded):} When $\pi(M_t)$ can be evaluated cheaply — the model has internalized effective action-selection for the current situation. This is Boyd's implicit guidance and control (Orient→Act, bypassing Decide), a trained RL policy in exploitation mode, a well-tuned PID controller, expert intuition (System 1), a martial artist's trained reflexes, an organization's standard operating procedures.

\textbf{Explicit (deliberative):} When the situation is novel, the action space is large, or the stakes demand verification — the agent engages in internal simulation, using the model to predict outcomes of candidate actions before selecting. This is Boyd's explicit Decide step, MCTS/planning in RL, Model Predictive Control, human deliberate reasoning (System 2), organizational strategic planning. Deliberation requires at minimum Level 2 epistemic access ( \S\,\ref{pearl-causal-hierarchy}) — the agent uses its model to simulate "what will I observe if I $do(a)$?" across candidates.

\textbf{Formal characterization of action fluency.} An agent has \textit{high fluency} for a situation when additional deliberation yields negligible improvement — formally, when $\Delta\eta^*(\Delta\tau) \approx 0$ for all $\Delta\tau > 0$ (see \S\,\ref{deliberation-cost}). Conversely, \textit{low fluency} means deliberation significantly improves action quality. Fluency is the degree to which the agent's immediate (zero-deliberation) action approaches the quality achievable with unbounded deliberation.

\textbf{Action fluency is distinct from model sufficiency.} An agent can have high $S(M_t)$ ( \S\,\ref{model-sufficiency}) but low fluency — a chess engine with a perfect model of the rules still requires expensive search. Conversely, an agent can have moderate sufficiency but high fluency in a narrow domain — a reflex that responds effectively to specific situations it evolved for. What reflexes, muscle memory, intuition, and expertise share is that the \textit{action-generating capacity itself} has been absorbed into the model's structure: the model doesn't just predict well, it \textit{acts} well, cheaply.

\textbf{Structural pressure toward implicit action.} When two action-selection modes produce equivalent expected outcomes, the faster mode is strictly preferable ( \S\,\ref{temporal-optimality}). This creates a pressure: agents under selective pressure (evolution, competition, training) tend to internalize frequently-needed action patterns, converting explicit deliberation into implicit fluency. The pressure is stronger when $\rho$ is high (fast-changing environments penalize deliberation — see \S\,\ref{deliberation-cost}), the pattern recurs frequently, or $\mathcal{T}$ is near the persistence threshold ( \S\,\ref{persistence-condition}) with no slack for deliberation overhead.

However, deliberation remains essential when the situation is genuinely novel, the action space is large relative to model capacity (chess, strategic planning), the stakes are asymmetric (cost of error vastly exceeds cost of delay), or $\rho$ is low (stable environment allows deliberation without mismatch accumulation).

\textbf{Connection to Section II.} For actuated agents ( \S\,\ref{agent-spectrum}), action selection involves not just $M_t$ but also $G_t = (O_t, \Sigma_t)$ — the purposeful substate. The policy becomes $\pi(M_t, G_t)$, coupling all substates through action ( \S\,\ref{directed-separation}). The action-deliberation-exploration tradeoff (Section II gap) extends the implicit/explicit distinction to three modes: exploit (pursue $O_t$ via $\Sigma_t$), explore (improve $M_t$), deliberate (revise $\Sigma_t$).

\textbf{Domain instantiations:}

\begin{adjustbox}{max width=\textwidth}
\small
\begin{tabular}{lll}
\toprule
Domain & Implicit action & Explicit deliberation \\
\midrule
Kalman + LQR & LQR control law from $\hat{x}_t$ & — (separation principle) \\
RL & Greedy policy $\arg\max Q(s,a)$ & MCTS, planning, rollouts \\
PID & $u = K_p e + K_i \int e + K_d \dot{e}$ & — (no deliberation) \\
Boyd's OODA & IG\&C (Orient→Act) & Explicit Decide step \\
Organism & Reflexes, habits & Deliberate planning \\
Organization & Standard procedures & Strategic planning \\
Software developer & Known patterns, familiar code & Reading docs, analyzing alternatives \\
\bottomrule
\end{tabular}
\end{adjustbox}


\textbf{(Descended from TF-07.)}

\end{proposition}

\begin{definition}[Mismatch Signal]
\label{mismatch-signal}
The discrepancy between the model's prediction and the actual observation — the formal expression of \textit{aporia} (productive perplexity). This is the signal that drives all adaptation: the agent discovers that reality and model have diverged, and this discovery is generative.



\medskip\noindent\textbf{Formal Expression.}



Given model $M_{t-1}$ and prior action $a_{t-1}$, the model generates a prediction:

$$\hat{o}_t = \mathbb{E}[o_t \mid M_{t-1}, a_{t-1}]$$

The \textbf{mismatch signal} (prediction error):

$$\delta_t = o_t - \hat{o}_t$$

This is the primary definition, used in the mismatch dynamics ( \S\,\ref{persistence-condition}, \S\,\ref{sector-condition-stability}) and in the decomposition ( \S\,\ref{mismatch-decomposition}).

For models with probabilistic predictions, the mismatch generalizes to the \textbf{score-function mismatch}:


$$\tilde{\delta}_t = -\nabla_M \log P(o_t \mid M_{t-1}, a_{t-1})$$

which points in the direction the model should move to increase the likelihood of the actual observation. $\tilde{\delta}_t$ lives in the tangent space $T_M\mathcal{M}$, while $\delta_t$ lives in observation space $\mathcal{O}$. Under Gaussian models, they coincide up to scaling.


\medskip\noindent\textbf{Epistemic Status.}


This is \textit{definitional}. Given any model that predicts ( \S\,\ref{agent-model}) and any observation that arrives ( \S\,\ref{observation-function}), their difference exists. The mismatch signal is not an additional assumption but a consequence of having a predictive model in an uncertain world. The score-function form is the natural generalization when $\mathcal{O}$ is not a vector space or when the model's predictive distribution is the natural object.


\medskip\noindent\textbf{Discussion.}


\textbf{Units and normalization.} When $\delta_t$ is in physical units (meters, dollars), the $\Vert\delta\Vert$ that enters the mismatch dynamics should be understood as the Mahalanobis distance: $\Vert\delta_t\Vert_\Sigma = \sqrt{\delta_t^T \Sigma^{-1} \delta_t}$ where $\Sigma$ is the observation noise covariance. This maps physical prediction error to dimensionless surprise-equivalent units.

\textbf{The zero-aporia ambiguity.} $\delta_t \approx 0$ does NOT necessarily indicate model adequacy. It may mean: (a) the model genuinely reflects reality — \textit{desirable}; (b) the agent is only observing aspects its model already explains, while remaining ignorant of aspects where the model is wrong — \textit{confirmation bias}; or (c) the observation channel is too noisy to detect model errors — \textit{architectural limitation}. Only (a) is desirable. An agent without aporia is an agent that has stopped adapting — but silence can mean peace or deafness. This ambiguity is why active testing — choosing actions to generate informative aporia — can be valuable (see \S\,\ref{causal-information-yield} for the CIY framework).

\textbf{The mismatch transform.} TF-06's update rule writes $M_t = M_{t-1} + \eta \cdot g(\delta_t)$, where the transform $g$ maps from $\delta_t$'s space to the model's update space: $g: \mathcal{O} \to T_M\mathcal{M}$ for prediction errors; $g: T_M\mathcal{M} \to T_M\mathcal{M}$ for score-function mismatches.

\end{definition}

\begin{theorem}[Mismatch Decomposition]
\label{mismatch-decomposition}
Expected squared mismatch decomposes into reducible model error and irreducible observation noise. The model can improve the first term; the second is a property of the channel.



\medskip\noindent\textbf{Formal Expression.}


% [Derived (mismatch-decomposition, Prop 5.1 from TFT)]

For any agent-environment pair within ACT's scope ( \S\,\ref{scope-condition}), when observation noise is non-degenerate or the model's predictive mean is misspecified:

$$\mathbb{E}[\Vert\delta_t\Vert^2] = \underbrace{\mathbb{E}[\Vert\hat{o}_t - \bar{o}_t\Vert^2]}_{\text{model error (reducible)}} + \underbrace{\mathbb{E}[\text{Var}(o_t \mid \Omega_t, a_{t-1})]}_{\text{observation noise (irreducible)}} > 0$$

where $\bar o_t = \mathbb{E}[o_t \mid \Omega_t, a_{t-1}]$ is the true conditional mean.


\medskip\noindent\textbf{Derivation.}


\begin{enumerate}
  \item By \S\,\ref{scope-condition}, $H(\Omega_t \mid \mathcal C_t) > 0$ — residual uncertainty persists.
  \item By \S\,\ref{agent-model}, the model generates predictions $\hat o_t = \mathbb{E}[o_t \mid M_{t-1}, a_{t-1}]$.
  \item Decompose mismatch into model error and noise. The cross-term vanishes by the fresh-noise assumption (GA-1): $\varepsilon_t$ is conditionally independent of $\mathcal C_{t-1}$ given $(\Omega_t, a_{t-1})$. Condition on $(\Omega_t, a_{t-1}, \mathcal C_{t-1})$; then both $\bar o_t$ and $\hat o_t$ are fixed, and $\mathbb{E}[o_t - \bar o_t \mid \Omega_t, a_{t-1}, \mathcal C_{t-1}] = \mathbb{E}[o_t - \bar o_t \mid \Omega_t, a_{t-1}] = 0$ by definition of $\bar o_t$ and GA-1. The outer expectation gives zero. This is orthogonality (uncorrelated), not independence.
  \item Term (ii) is positive when observation noise is non-degenerate. Term (i) is positive when the model's predictive mean differs from the true conditional mean. Either suffices.
\end{enumerate}



\medskip\noindent\textbf{Epistemic Status.}


\textit{Exact} under the fresh-noise assumption (observation noise $\varepsilon_t$ conditionally independent of history given current state and action). This is the standard assumption in state-space models — noise is a property of the observation channel at the moment of observation. The decomposition is a mathematical identity (bias-variance decomposition applied to the prediction problem). The positivity of $\mathbb{E}[\Vert\delta_t\Vert^2]$ follows from either condition; both hold simultaneously in typical settings.


\medskip\noindent\textbf{Discussion.}


\textbf{Reducible vs. irreducible.} An agent that tries to eliminate \textit{all} mismatch — including irreducible noise — will overfit: adjusting its model to explain noise, degrading future predictions. The update gain ( \S\,\ref{update-gain}) implicitly separates signal from noise by weighting observations in proportion to their informativeness.

\textbf{Connection to model sufficiency.} When $S(M_t) < 1$ ( \S\,\ref{model-sufficiency}), the model has lost predictive information relative to the full history. Under an alignment assumption (the lost information affects the one-step conditional mean), this implies positive model error (term i). Without that alignment assumption, insufficiency still implies positive regret under proper scoring rules but not necessarily positive one-step mean error.

\textbf{Mismatch is structurally persistent.} In realistic ACT regimes, mismatch signals persist — they can be reduced but not eliminated when observation noise is non-degenerate. Deterministic, noiseless, perfectly specified systems are limiting edge cases, not the typical adaptive regime.

\end{theorem}

\begin{empirical}[Update Gain]
\label{update-gain}
The optimal weight an agent assigns to new observations when updating its model — the rate of \textit{epistrophe} (turning toward reality). How much the agent should trust the incoming observation versus its own prior understanding.



\medskip\noindent\textbf{Formal Expression.}


% [Empirical Claim (uncertainty-ratio-principle)]

$$\eta^* = \frac{U_M}{U_M + U_o}$$

where:
\begin{itemize}
  \item $\eta^*$ is the optimal update gain (proportion of mismatch used to correct the model)
  \item $U_M$ is model uncertainty (predictive variance or entropy)
  \item $U_o$ is irreducible observation noise
\end{itemize}


The update rule takes the form:


$$M_t = M_{t-1} + \eta^* \cdot g(\delta_t)$$

where $\delta_t$ is the mismatch ( \S\,\ref{mismatch-signal}) and $g(\cdot)$ is a correction mapping from observation space to model update space.


\medskip\noindent\textbf{Epistemic Status.}


\textit{Exact} for linear-Gaussian systems (where $\eta^*$ is the Kalman gain) and conjugate Bayesian systems. For general adaptive systems (RL agents, organizational learning, biological adaptation), it is \textit{robust qualitative} — any optimal adaptation process must approximate this functional dependence, even if the variance ratio is not explicitly computed.


\medskip\noindent\textbf{Discussion.}


\textbf{Limiting behavior.} When $U_M \gg U_o$ (high model uncertainty — e.g., after initialization or structural adaptation), $\eta^* \to 1$: trust the observation. When $U_M \ll U_o$ (confident model, noisy channel), $\eta^* \to 0$: trust the model. The gain determines how strongly the agent corrects toward reality on each update.

\textbf{Gain collapse — epistrophe failure.} When the agent incorrectly estimates $U_M \to 0$ (spurious confidence) or $U_o \to \infty$ (spurious distrust of sensors), $\eta^* \to 0$ and epistrophe ceases. Aporia still arrives — the mismatch signal is still generated — but the agent no longer turns toward it. Mismatches are ignored, producing confirmation bias or a decoupled reality model. The cycle runs but the corrective phase is hollow.

\textbf{Multi-dimensional generalization.} In vector-valued systems, $U_M$ and $U_o$ are covariance matrices and $\eta^*$ becomes a gain matrix (as in the Kalman filter). The scalar form captures the essential structure.

\textbf{Connection to adaptive tempo.} The update gain is one factor in the agent's adaptive tempo ( \S\,\ref{adaptive-tempo}): $\mathcal{T} = \nu \cdot \eta^*$. Frequent aisthesis (high $\nu$) is useless if epistrophe extracts no information (low $\eta^*$). Gain measures the \textit{quality} of the cycle's corrective phase; event rate measures its \textit{speed}.

\textbf{Gain dynamics.} The optimal gain changes over time following predictable patterns:

\begin{itemize}
  \item \textit{Convergence}: As the model accumulates information, $U_M$ decreases, so $\eta^* \to 0$. The model becomes increasingly resistant to individual observations. This IS Kalman filter convergence, Bayesian posterior concentration, and RL learning rate annealing.
  \item \textit{Reset after structural change}: When the environment changes in ways the model cannot track incrementally ( \S\,\ref{structural-adaptation-necessity}), $U_M$ should spike — the model "admits" its uncertainty. The gain increases, enabling rapid re-learning. An agent whose gain does NOT reset after structural change will continue trusting a stale model — Boyd's "incestuous amplification" and the cause of brittle failure in non-stationary environments.
\end{itemize}


\textbf{Overfitting as gain miscalibration.} From \S\,\ref{mismatch-decomposition}: $\mathbb{E}[\Vert\delta_t\Vert^2]$ = model error + irreducible noise. An agent with $\eta$ too high adjusts its model to explain observation noise, increasing model error on future predictions. An agent with $\eta$ too low fails to correct genuine model errors. The optimal gain implicitly separates signal from noise by weighting observations in proportion to their informativeness — exactly what $U_M/(U_M + U_o)$ achieves when $U_o$ captures the irreducible noise.

\textbf{Representation note.} The additive form operates in a \textit{representation space} appropriate to the model. For Bayesian posteriors (where update is multiplicative: $P(\theta \mid D) \propto P(D \mid \theta) P(\theta)$), the additive rule operates in log-probability or natural parameter space. For models on constrained manifolds (probability simplices, rotation groups), the update must be projected onto the manifold. The claim is not that all updates are literally additive in native parameterization, but that they have the structure "current state + gain × transformed mismatch" in an appropriate coordinate system.

\textbf{Domain validation:}

\begin{adjustbox}{max width=\textwidth}
\small
\begin{tabular}{lll}
\toprule
Domain & Gain form & Mapping quality \\
\midrule
Kalman filter & $K_t = P_{t\Vert t-1} H^T (H P_{t\Vert t-1} H^T + R)^{-1}$ & \textbf{Exact.} Scalar case is exactly $U_M/(U_M + U_o)$. \\
Conjugate Bayesian & Posterior weight $n/(n + \kappa)$ cumulative; incremental $1/(n + \kappa)$ & \textbf{Exact} for conjugate families. Incremental gain decreases as data accumulates. \\
RL (Q-learning) & Fixed learning rate $\alpha$ & \textbf{Approximate.} $\alpha$ is a degenerate constant gain — does not adapt to uncertainty. Advanced methods (Bayesian RL, Adam) converge toward the optimal form. \\
PID control & Fixed gains $(K_p, K_i, K_d)$ & \textbf{Simplified.} Gains set at design time. Adaptive PID and MPC move toward the full framework. \\
Software developer & Implicit trust weighting of information sources & \textbf{Structural analogy.} New developer (high $U_M$) trusts observations heavily; experienced developer (low $U_M$) trusts their model. Gain reset after major refactoring. \\
\bottomrule
\end{tabular}
\end{adjustbox}


\textbf{Simulation validation.} Numerical experiments (track-b, Variant E) validated the uncertainty ratio principle under observation noise. Riccati-optimal gain reduced steady-state mismatch by 52\% compared to fixed gain when observation noise was moderate. The optimal gain also proved critical in adversarial settings: under heavy observation noise, optimal gain preserved more than double the adversarial tempo advantage exponent (0.40 vs 0.18) compared to fixed gain.

\textbf{Open questions:}

\begin{enumerate}
  \item \textit{Non-parametric models}: For neural networks without well-defined scalar $U_M$, how should it be computed? Ensemble methods, dropout-based uncertainty, and Bayesian neural networks are all approximations.
  \item \textit{Matrix vs scalar gain}: In high-dimensional systems, the gain is a matrix (Kalman) or per-parameter (Adam). The cross-dimensional structure (covariance) adds complexity. The scalar captures the principle; the matrix captures the full optimization.
\end{enumerate}


\textbf{(Descended from TF-06.)}
\end{empirical}

\begin{definition}[Causal Information Yield]
\label{causal-information-yield}
Actions don't merely select among outcomes — they produce characteristically different outcome distributions depending on the causal structure. Causal information yield (CIY) quantifies the \textbf{action-distinguishability} of an action: how different its outcome distribution is from what alternative actions would produce.



\medskip\noindent\textbf{Formal Expression.}



The \textbf{canonical CIY} of action $a$ given model state $M$:

$$\text{CIY}(a;\, M) = \mathbb{E}_{a' \sim q(\cdot \mid M)}\!\left[D_{\mathrm{KL}}\!\left(P(o \mid do(a), M) \,\Vert\, P(o \mid do(a'), M)\right)\right]$$

where $q(\cdot \mid M)$ is a reference distribution over comparator actions (uniform, policy-induced, or task-specific). This measures how strongly the action changes the interventional distribution of outcomes relative to alternatives.

$\text{CIY} \geq 0$ by construction (expectation of KL divergences). $\text{CIY} = 0$ for a passive observer or an agent whose actions don't affect outcome distributions. $\text{CIY} > 0$ when actions causally alter what is observed — exactly what distinguishes Level 2 from Level 1 epistemic access ( \S\,\ref{pearl-causal-hierarchy}).


\medskip\noindent\textbf{Epistemic Status.}


The CIY \textit{definition} is well-grounded in causal inference theory. The quantity itself is standard — an expected KL divergence under the do-calculus. The interpretive claim — that CIY measures action-distinguishability rather than expected information gain — is exact (it follows from the definition). The relationship between CIY and learning value (see Discussion) is discussion-grade: the $\lambda$-weighted approximation to EIG is heuristic, not derived.


\medskip\noindent\textbf{Discussion.}


\textbf{CIY measures distinguishability, not learning value.} CIY as defined is the expected KL divergence between outcome distributions — how different the outcomes of $a$ are from the outcomes of typical alternatives. This is \textbf{action-distinguishability}, not \textbf{expected information gain} (EIG). The distinction matters: an action can have high CIY even when the agent already knows the outcome distributions perfectly (the distributions ARE different, but the agent learns nothing new by confirming what it already knows). High CIY is \textit{necessary} for learning (indistinguishable actions can't teach anything) but not \textit{sufficient} (distinguishable actions only teach when the agent is uncertain about the distinction).

The two quantities approximately coincide when $U_M$ is high — when the agent doesn't know the outcome distributions, high-CIY actions also have high EIG because observing a characteristically different outcome updates the agent's beliefs about the causal structure. They diverge when $U_M$ is low — a confident agent gains nothing from taking a distinguishable action it already understands.

The $\lambda(M_t)$ weighting in the unified policy objective ( \S\,\ref{ciy-unified-objective}) partially compensates: when $U_M$ is low, $\lambda \to 0$, suppressing the CIY term regardless of its magnitude. This makes the exploration term behave more like EIG — suppressing exploration when the agent is already certain. The compensation is heuristic (the $\lambda$ form is not derived). For the current theory, CIY serves as a tractable surrogate for EIG, with the $\lambda$ weighting providing the uncertainty-gating that makes the surrogate reasonable.

\textbf{Open direction: proper EIG within ACT.} Replacing CIY with a proper expected information gain quantity — $\text{EIG}(a; M) = I(o; \theta \mid do(a), M)$ where $\theta$ parameterizes the model — would be a stronger foundation for the exploration term, particularly in domains where the agent must decide between actions that are all highly distinguishable but differ in what they teach. Under certain scopes (intervention-rich domains with well-parameterized models), the EIG formulation might yield sharper exploration strategies — preferring actions that resolve the \textit{most uncertain} causal links rather than the most distinguishable ones. Whether this yields operationally significant improvements over the $\lambda$-weighted CIY surrogate is an empirical question. The CIY formulation has the advantage of being computable from the agent's current model alone (it doesn't require reasoning about model uncertainty); EIG requires a meta-model of what the agent doesn't know.

**Dependence on the reference distribution $q$.** The quantitative CIY value depends on the choice of $q$, which is a significant degree of freedom. A uniform $q$ treats all alternatives equally; a policy-induced $q$ emphasizes alternatives the agent would consider. ACT adopts the policy-induced $q$ as default: $q(\cdot \mid M) = \pi(\cdot \mid M)$, yielding CIY as "how different is this action's outcome from what I'd typically see?" CIY values are not comparable across different $q$ choices.

\textbf{Query actions: accessing external models.} A qualitatively distinct class of actions: querying another agent's model. When a reliable source exists (expert, database, documentation, well-trained LLM), "ask a well-formed question" can yield information equivalent to thousands of probe-observe cycles. The source's model has already performed the compression work ( \S\,\ref{information-bottleneck}) — the response transfers the \textit{output} of compression rather than requiring the agent to reconstruct it.

Key properties of query actions:
\begin{itemize}
  \item \textbf{Information density}: Single well-targeted query can carry CIY orders of magnitude higher than individual environment probes
  \item \textbf{Trust-dependent gain}: Update from query depends on the agent's model of the source's reliability and alignment, not on observation channel noise ( \S\,\ref{communication-gain})
  \item \textbf{Pre-compressed information}: Responses arrive already compressed in the source's representational framework, introducing a translation cost when frameworks don't align
  \item \textbf{Structural adaptation via external models}: Encountering another agent's model can trigger structural change ( \S\,\ref{structural-adaptation-necessity}) — incorporating external representational structure rather than building it de novo ("grafting")
\end{itemize}


When high-CIY query channels are available, the unified policy objective ( \S\,\ref{ciy-unified-objective}) favors query actions over direct probes, particularly when $U_M$ is high, a trusted source exists, query cost is low, and the needed information is about \textit{structure} rather than the agent's specific situation.

\textbf{The adversarial mirror: deception and model corruption.} The same channel that enables cooperative knowledge transfer can be exploited to degrade the opponent's model. A deceptive response yields positive CIY in the strict information-theoretic sense, but the content drives model-reality mismatch \textit{upward}. The update gain $\eta^*$ for the victim depends on trust; successful deception exploits high trust to inject a large, misdirected update. In the Lyapunov framework ( \S\,\ref{sector-condition-stability}), this is adversarial disturbance injected through the observation channel, with coupling coefficient $\gamma_A$ determined by the victim's trust level and exposure. See \S\,\ref{communication-gain} for the formal treatment of trust-dependent gain, and \S\,\ref{adversarial-destabilization} for the Lyapunov formalization. Distributed tempo, topology analysis, and game-theoretic integration are Section III content not yet fully extracted (source material in \texttt{src/old-tf-appendix-f-multi-agent.md}).

\textbf{(Descended from TF-08.)}

\end{definition}

\begin{definition}[Adaptive Tempo]
\label{adaptive-tempo}
The effective rate at which an agent acquires useful information from its environment — the product of observation frequency and update quality across all channels.



\medskip\noindent\textbf{Formal Expression.}



$$\mathcal{T} = \sum_k \nu^{(k)} \cdot \eta^{(k)*}$$

where:
\begin{itemize}
  \item $k$ indexes the agent's distinct observation channels
  \item $\nu^{(k)}$ is the event rate on channel $k$
  \item $\eta^{(k)*}$ is the optimal update gain on channel $k$ ( \S\,\ref{update-gain})
\end{itemize}


Single-channel special case: $\mathcal{T} = \nu \cdot \eta^*$.


\medskip\noindent\textbf{Epistemic Status.}


This is a \textit{definition}. It names the quantity that characterizes an agent's total corrective capacity, combining loop speed ($\nu$) and epistemic quality ($\eta^*$). The definition itself is not a truth-claim; the substantive claims are in the results that use it ( \S\,\ref{persistence-condition}, \S\,\ref{adversarial-tempo-advantage}).


\medskip\noindent\textbf{Discussion.}


\textbf{Speed-quality substitutability.} An agent can achieve the same tempo via a fast noisy loop (high $\nu$, low $\eta^*$) or a slower calibrated one (low $\nu$, high $\eta^*$). The product structure means improvements to \textit{both} factors compound multiplicatively.

\textbf{Observation noise gating.} Because $\eta^* = U_M / (U_M + U_o)$, high observation noise ($U_o$) depresses gain and collapses tempo regardless of loop speed. You cannot outrun a bad observation channel by iterating faster. This grounds Boyd's emphasis on Orient quality over raw OODA speed.

\textbf{Centrality.} Tempo is ACT's core capacity metric. It appears on the left side of the persistence condition ( \S\,\ref{persistence-condition}), determines adversarial advantage ( \S\,\ref{adversarial-tempo-advantage}), and connects to code quality as observation infrastructure ( \#code-quality-as-observation-infrastructure — cross-component reference, see \texttt{02-tst-core/}) in the software domain. The strategic analog $\mathcal{T}_\Sigma$ ( \S\,\ref{strategic-tempo}) extends the same structure to strategy-edge revision, with the key difference that strategic edge rates are endogenous (depend on action policy and upstream success).

\textbf{Temporal nesting.} Adaptive processes stratify by timescale, with convergence constraints between levels ( \S\,\ref{temporal-nesting}).

\textbf{Mismatch dynamics.} The evolution of mismatch over time is governed by the balance between correction (via tempo) and disturbance ($\rho$) ( \S\,\ref{mismatch-dynamics}).

\textbf{Channel independence assumption.} The additive formula assumes informationally independent channels — each channel contributes non-redundant correction capacity. When channels are correlated (overlapping sensors, repeated teammate reports, redundant telemetry), the additive formula \textit{overcounts} effective tempo. The correct tempo satisfies:

$$\mathcal{T} \leq \sum_k \nu^{(k)} \cdot \eta^{(k)*}$$

with equality iff channels are informationally independent. The gap is the \textit{redundancy penalty} — the effective correction capacity lost to overlapping information. For two correlated channels, the penalty involves the mutual information $I(e^{(1)}; e^{(2)} \mid M_{\tau^-})$ between their event streams conditioned on the current model. Since tempo is the core capacity variable (appearing in the persistence condition, adversarial dynamics, and composition), this overcounting inflates margins wherever channel independence fails. The additive formula remains an upper bound and is exact when channels measure genuinely different aspects of the environment. Multi-agent composition ( \S\,\ref{team-persistence}) inherits this limitation: the communication tempo contribution is additive in the same sense and overcounts when different allies report correlated information.

\textbf{Scalar vs. vector tempo.} The scalar $\mathcal{T}$ assumes isotropic correction capacity. When the agent corrects some dimensions faster than others, scalar tempo overestimates effective adaptation along weak dimensions. Simulation confirms: in an anisotropic 3D system (gain varying 5:1), scalar $\rho/\mathcal{T}$ overestimated by 72\%, with the weak dimension accounting for 84\% of total mismatch. The correct formulation is per-dimension: $\mathcal{T}_k > \rho_k / \delta_{\text{critical},k}$ ( \S\,\ref{per-dimension-persistence}).

\textbf{(Descended from TF-11.)}
\end{definition}

\begin{hypothesis}[Mismatch Dynamics]
\label{mismatch-dynamics}
The evolution of model-reality mismatch over time is governed by the balance between the agent's corrective capacity (tempo) and the rate of environmental change (disturbance). The linear ODE is a first-order approximation; the general nonlinear case is handled by the sector-condition framework ( \S\,\ref{sector-condition-stability}).



\medskip\noindent\textbf{Formal Expression.}


% [Hypothesis (mismatch-dynamics)]

$$\frac{d\Vert\delta\Vert}{dt} = -\mathcal{T} \cdot \Vert\delta\Vert + \rho(t)$$

where:
\begin{itemize}
  \item $\mathcal{T} \cdot \Vert\delta\Vert$ is the rate at which the agent corrects mismatch (proportional to both tempo and current mismatch)
  \item $\rho(t)$ is the \textbf{environment change rate} — the rate at which new mismatch is introduced by changes in $\Omega$
\end{itemize}


**Steady state, Model D (deterministic bounded disturbance, $\lVert w(t)\rVert \leq \rho$):**

Setting $d\lVert\delta\rVert/dt = 0$:

% [Derived (from linear hypothesis, deterministic)]

$$\lVert\delta\rVert_{ss} = \frac{\rho}{\mathcal{T}}$$

Steady-state mismatch is the ratio of how fast the environment changes to how fast the agent adapts.

**Steady state, Model S (stochastic zero-mean disturbance, $d\delta = -\mathcal{T}\delta\,dt + \sigma_w\,dW_t$):**

% [Derived (from Itô-Lyapunov analysis — see Prop A.1S in CREF:sector-condition-derivation:CREF)]

$$\lVert\delta\rVert_{\text{rms}} = \frac{\sigma_w}{\sqrt{2\mathcal{T}}}$$

(scalar case, $n = 1$; general: $\sigma_w\sqrt{n/(2\mathcal{T})}$). Steady-state mismatch scales as the square root of the disturbance-to-correction ratio, not the ratio itself. The $1/\sqrt{\mathcal{T}}$ scaling (vs. $1/\mathcal{T}$ for Model D) means correction is less effective against noise than against drift.

\textbf{Transient solution (Model D):}

$$\lVert\delta(t)\rVert = \lVert\delta_0\rVert e^{-\mathcal{T} t} + \frac{\rho}{\mathcal{T}}(1 - e^{-\mathcal{T} t})$$

Mismatch decays exponentially from initial conditions toward the steady state.


\medskip\noindent\textbf{Epistemic Status.}


\textit{Heuristic.} This is explicitly a first-order linear approximation. The qualitative behavior (bounded mismatch, steady-state ratio, exponential convergence) is robust across correction function forms. The quantitative predictions (exact steady-state value, convergence rate, the squared adversarial scaling law) are specific to the linear case. The general nonlinear treatment ( \S\,\ref{sector-condition-stability}) replaces the linear correction term with a sector-bounded correction function and proves persistence without committing to a specific functional form.

\textbf{Bridging assumption (discrete to continuous).} This ODE is a fluid-limit approximation of the discrete event-driven dynamics ( \S\,\ref{event-driven-dynamics}). The fluid limit is formally justified by \S\,\ref{discrete-sector-condition}: for Model D (deterministic), the discrete and continuous steady states are identical (zero gap); for Model S (stochastic), the variance gap is $O(\eta^* c_{\max})$, quantitatively small whenever $\eta^* c_{\max} \ll 1$. The approximation is least accurate during initial transients when $\eta^*$ is large, but this phase is short-lived and the transient error is bounded by $O(\eta^* c_{\max} / \nu^{1/2})$ under Lipschitz regularity.


\medskip\noindent\textbf{Discussion.}


\textbf{Speed-quality substitutability.} From $\mathcal{T} = \nu \cdot \eta^*$ (single-channel case): doubling event rate $\nu$ has the same effect on $\Vert\delta\Vert_{ss}$ as doubling update quality $\eta^*$. They are multiplicative when both improve: 50\% improvement in each yields $1.5 \times 1.5 = 2.25\times$, not $3\times$. This is the formal analog of Boyd's insight that Orient quality often matters more than raw OODA speed — the same structural observation (quality and speed are substitutable, quality often dominates) appears in the model.

\textbf{The persistence threshold.} From the steady-state: $\Vert\delta\Vert_{ss} < \Vert\delta_{\text{critical}}\Vert$ iff $\mathcal{T} > \rho/\Vert\delta_{\text{critical}}\Vert$ ( \S\,\ref{persistence-condition}). Below this threshold, the model cannot support effective action. The same structural pattern — correction capacity falling below disturbance rate — appears across domains: extinction (environment changes faster than organism adapts), organizational failure (market moves faster than company learns), control instability (disturbances exceed correction capacity), cognitive overload (information arrives faster than processing). The persistence condition captures the common structure; whether it captures the dominant mechanism in each domain is an empirical question.

\textbf{Nonlinear reality.} The true correction dynamics are almost certainly nonlinear:
\begin{itemize}
  \item *Saturation at large $\Vert\delta\Vert$*: correction mechanism overwhelmed, so correction is slower than linear for large errors. Makes the persistence threshold harder to satisfy.
  \item \textit{Threshold effects}: small mismatches go uncorrected ($F \approx 0$ for $\Vert\delta\Vert < \varepsilon$), creating a dead zone.
  \item \textit{Structural breakdown}: beyond some critical $\Vert\delta\Vert$, correction drops to zero because the model class is no longer appropriate ( \S\,\ref{structural-adaptation-necessity}).
\end{itemize}


These nonlinearities are exactly what the sector-condition framework ( \S\,\ref{sector-condition-stability}) handles.

\textbf{Adversarial coupling.} When two agents are coupled ($A$'s actions increase $B$'s disturbance): $\rho_B = \rho_{B,\text{base}} + \gamma_A \cdot \mathcal T_A$. The steady-state mismatch ratio scales superlinearly with the tempo ratio, but the exponent depends on the disturbance model. Under Model D (deterministic drift, coupling-dominant): $(\mathcal T_A/\mathcal T_B)^2$ — the squared law, derived from the $1/\mathcal{T}$ steady-state scaling. Under Model S (stochastic noise, coupling-dominant): $(\mathcal T_A/\mathcal T_B)^{3/2}$ — derived from the $1/\sqrt{\mathcal{T}}$ steady-state scaling. See \S\,\ref{adversarial-tempo-advantage}.

\textbf{(Descended from TF-11.)}

\end{hypothesis}

\begin{proposition}[Deliberation Cost]
\label{deliberation-cost}
Explicit deliberation improves action quality by using the model for internal simulation before acting — pausing praxis to improve upcoming epistrophe. But deliberation takes time, and during that time aporia accumulates (the environment continues to evolve while the agent is not correcting). Deliberation is justified when the improvement in epistrophe quality exceeds the aporia accumulated during the pause.



\medskip\noindent\textbf{Formal Expression.}


\textbf{Assumption (local deliberation drift):}

% [Assumption (deliberation-drift)]

During a deliberation pause of duration $\Delta\tau$, mismatch increases at an approximately constant local rate $\rho_{\text{delib}}$:

$$\Delta\Vert\delta\Vert_{\text{deliberation}} \approx \rho_{\text{delib}} \cdot \Delta\tau$$

This is a short-horizon assumption about inaction windows, not a full global dynamics model. It is weaker than the mismatch ODE and can be estimated directly from pause windows in empirical traces.

\textbf{Proposition (deliberation threshold):}

% [Derived (Conditional on deliberation-drift assumption)]

Deliberation of duration $\Delta\tau$ is net-beneficial when:

$$\Delta\eta^*(\Delta\tau) \cdot \Vert\delta_{\text{post}}\Vert > \rho_{\text{delib}} \cdot \Delta\tau$$

where $\Delta\eta^*(\Delta\tau)$ is the improvement in post-deliberation update gain and $\Vert\delta_{\text{post}}\Vert$ is the mismatch magnitude the agent will face when it resumes acting.


\medskip\noindent\textbf{Derivation.}


\begin{enumerate}
  \item Without deliberation, the agent acts immediately at current tempo $\mathcal{T}_0 = \nu \cdot \eta^*_0$.
  \item With deliberation of duration $\Delta\tau$, the agent pauses, then acts with improved gain $\eta^*_0 + \Delta\eta^*$. But during the pause, mismatch has grown by $\rho_{\text{delib}} \cdot \Delta\tau$.
  \item The net mismatch reduction from acting after deliberation versus acting immediately: $\text{Net} = \Delta\eta^* \cdot \Vert\delta_{\text{post}}\Vert - \rho_{\text{delib}} \cdot \Delta\tau$.
  \item Deliberation is justified iff $\text{Net} > 0$. $\square$
\end{enumerate}


\textbf{Optimal deliberation duration} (under diminishing returns):

% [Derived (Conditional on diminishing-returns + deliberation-drift)]

$$\Delta\tau^* = \arg\max_{\Delta\tau} \left[\Delta\eta^*(\Delta\tau) \cdot \Vert\delta_{\text{post}}\Vert - \rho_{\text{delib}} \cdot \Delta\tau \right]$$

At the first-order condition: $\frac{\partial \Delta\eta^*}{\partial \Delta\tau} \cdot \Vert\delta_{\text{post}}\Vert = \rho_{\text{delib}}$. Stop deliberating when the marginal improvement rate drops below the mismatch drift rate (normalized by post-deliberation mismatch).


\medskip\noindent\textbf{Epistemic Status.}


\textit{Conditional} on the local deliberation-drift assumption. The threshold condition is derived given the assumption; the assumption itself is a local approximation validated by consistency with the global mismatch dynamics ( \S\,\ref{persistence-condition}). The result captures the \textit{epistemic} benefit of deliberation (improving $\eta^*$); in practice, deliberation also provides a direct \textit{action-value} benefit (choosing better actions that alter the environment trajectory), which operates through $\rho$ reduction and immediate reward — a fuller formalization would incorporate the unified policy objective ( \S\,\ref{causal-information-yield}) at significantly more complexity.


\medskip\noindent\textbf{Discussion.}


**High-$\rho_{\text{delib}}$ environments penalize deliberation.** When the environment changes rapidly during pause windows, the cost term grows quickly. Only very short deliberation with large $\Delta\eta^*$ can justify the pause. The model captures the same tradeoff Boyd emphasized: in fast-tempo adversarial environments, over-deliberation is fatal not because thinking is bad, but because the environment moves during the thinking. Whether the specific mechanism (mismatch drift during pause) is the dominant real-world effect is an empirical question.

\textbf{Diminishing returns.} In most models, $\Delta\eta^*(\Delta\tau)$ exhibits diminishing returns — the first moments of simulation yield the largest improvement. Combined with the linear cost $\rho_{\text{delib}} \cdot \Delta\tau$, this implies a finite optimal deliberation duration. Past that point, additional thinking is net-harmful.

\textbf{Implicit action as the high-tempo limit.} As $\rho_{\text{delib}} \to \infty$ or $\Delta\tau^* \to 0$: the optimal strategy converges to zero deliberation — pure implicit action ( \S\,\ref{action-selection}). This provides a mathematical basis for why high-tempo environments favor action fluency: the cost of deliberation exceeds its benefit when $\Delta\eta^*$ is small (action-selection is already fluent) or $\rho_{\text{delib}}$ is large.

\textbf{Deliberation as an investment.} When $\rho_{\text{delib}}$ is low (stable environment) or $\Vert\delta_{\text{post}}\Vert$ is large (significant model-reality gap), deliberation pays off. The conditions favoring deliberation — stable environment, large mismatch — resemble the high-stakes, low-urgency scenarios where deliberative reasoning (System 2) is advantageous in dual-process theories. The structural parallel is suggestive; whether the cost-benefit mechanism is the same one governing System 1/System 2 selection is an open question.

**The circularity of $\Vert\delta_{\text{post}}\Vert$.\textit{* Evaluating the threshold requires the agent to }predict* post-deliberation mismatch using its current model — the same model deliberation is meant to improve. This circularity is typically benign: $\Vert\delta_{\text{post}}\Vert$ is bounded below by $\rho_{\text{delib}} \cdot \Delta\tau$ and above by current mismatch plus that accumulation. An agent that underestimates its mismatch will under-deliberate; one that overestimates will over-deliberate. The bias is self-correcting through the feedback loop. The threshold is best understood as a \textit{design criterion}, not a real-time decision procedure.

\textbf{Resource costs beyond time.} Real agents also incur computational and energetic costs: internal simulation burns calories, compute cycles, or opportunity cost of not processing new observations. These are additive: $\Delta\eta^*(\Delta\tau) \cdot \Vert\delta_{\text{post}}\Vert > \rho_{\text{delib}} \cdot \Delta\tau + C(\Delta\tau)$. In high-$\rho_{\text{delib}}$ environments the temporal cost dominates; in low-$\rho_{\text{delib}}$ environments, resource costs may be the binding constraint.

\textbf{Structural adaptation as massive deliberation.} Structural adaptation ( \S\,\ref{structural-adaptation-necessity}) can be understood as deliberation with a massive $\Delta\tau$. During decomposition-and-recombination, the agent's high-frequency parametric loop is partially or fully suspended while it searches for a new model class. The mismatch debt $\rho_{\text{delib}} \cdot \Delta\tau$ is enormous (weeks/months for organizational restructuring). The agent rationally resists structural adaptation until the parametric mismatch floor exceeds this debt.

\textbf{Connection to temporal nesting.} Deliberation is a nested loop: internal simulation running at rate $\nu_{\text{internal}}$ within the external action loop at rate $\nu_{\text{external}}$. The convergence constraint applies: the internal loop must approximately converge before the external loop acts on its output.

\textbf{Connection to Section II.} For actuated agents, the deliberation tradeoff extends to three modes: exploit ($O_t$ via $\Sigma_t$), explore (improve $M_t$), and deliberate (revise $\Sigma_t$). The three-way allocation ( \S\,\ref{exploit-explore-deliberate}) extends this segment's threshold by adding a strategic benefit term $\Delta V_\Sigma$; the extended threshold is the one genuinely derived piece. The broader three-way framing is discussion-grade — simulation shows deliberation is rarely chosen by an oracle in simple settings, and a unified objective outperforms the two-stage decomposition.

\textbf{The AI agent's dilemma.} An AI agent with 100\% context turnover faces a severe version: it MUST deliberate (comprehend the codebase) before acting effectively, but during comprehension its context fills and the environment may change. The optimal comprehension depth depends on $\rho_{\text{delib}}$ and the session's action horizon. This is why reading CLAUDE.md and architecture docs first (high-CIY query actions) dominates reading random source files (low-CIY exploration).

\textbf{Domain instantiations:}

\begin{adjustbox}{max width=\textwidth}
\small
\begin{tabular}{llll}
\toprule
Domain & Deliberation & $\Delta\eta^*$ source & When $\rho_{\text{delib}}$ is high \\
\midrule
Boyd's OODA & Explicit "Decide" step & War-gaming, staff analysis & Collapses to IG\&C (implicit) \\
RL / MCTS & Planning rollouts & Monte Carlo tree search & Fewer rollouts, shallower search \\
MPC & Online optimization & Trajectory optimization & Shorter horizons, faster solvers \\
Human cognition & System 2 deliberation & Mental simulation & Defaults to System 1 (intuition) \\
Organization & Strategic planning & Scenario analysis & "Move fast and break things" \\
Software developer & Reading code, analyzing alternatives & Architecture analysis & Ship now, refactor later \\
AI agent & Reading codebase, planning approach & Context-building & Limit comprehension, act sooner \\
\bottomrule
\end{tabular}
\end{adjustbox}


\textbf{Open questions:}

\begin{enumerate}
  \item \textit{Computational cost of deliberation} is not just elapsed time but resource cost. A fuller model would include both temporal and computational budgets.
  \item \textit{Deliberation about deliberation}: deciding whether to deliberate itself takes time. This meta-deliberation is bounded by the same tradeoff at a higher level, suggesting a hierarchy of diminishing deliberation horizons.
  \item \textit{Deliberation that generates observations}: internal simulation can surface model inconsistencies (internal mismatch), functioning as "exploration without external action." Can deliberation generate internal CIY?
\end{enumerate}


\textbf{(Descended from TF-09.)}

\end{proposition}

\begin{proposition}[Gain–Sector Bridge]
\label{gain-sector-bridge}
The gain-based update principle ( \S\,\ref{update-gain}) produces correction dynamics satisfying the sector condition (GA-3) whenever the update rule has \textit{directional fidelity} — the correction points at least roughly toward reality. For gradient-based agents, this is equivalent to local strong convexity of the loss function. The sector parameter $\alpha$ is not a free parameter but is determined by the gain and the correction geometry.



\medskip\noindent\textbf{Formal Expression.}



\medskip\noindent\textbf{The Bridge Theorem.}


% [Derived (gain-sector-bridge, from update-gain + directional fidelity)]

Given the gain-based update $M_t = M_{t-1} + \eta^* \cdot g(\delta_t)$ ( \S\,\ref{update-gain}), the induced correction function is:

$$F(\delta) = \eta^* \cdot H \, g(\delta)$$

where $H$ maps state-space corrections to observation-space mismatch reduction. The sector condition (GA-3) holds with parameter $\alpha > 0$ whenever:

\textbf{(B1) Directional fidelity.} The mismatch transform $g$ preserves the mismatch-reducing direction:

$$\delta^T H \, g(\delta) \geq c \lVert\delta\rVert^2 \quad \text{for } \lVert\delta\rVert \leq R$$

for some $c > 0$. The sector parameter is then:

$$\alpha = \eta^* \cdot c_{\min}, \qquad c_{\min} = \inf_{\lVert\delta\rVert \leq R} \frac{\delta^T H \, g(\delta)}{\lVert\delta\rVert^2}$$


\medskip\noindent\textbf{Gradient Equivalence.}


% [Derived (sector-convexity equivalence)]

For any agent updating via gradient descent on a loss $L$ with learning rate $\eta$:

$$\alpha = \eta \cdot \mu \qquad \text{where } \mu = \inf_{\lVert\delta\rVert \leq R} \lambda_{\min}(\nabla^2 L(M^* + \delta))$$

is the strong convexity modulus. The basin radius $R$ is the largest ball around the optimum where $\nabla^2 L$ remains positive definite. The equivalence is exact:

$$\text{GA-3 holds with } (\alpha, R) \iff L \text{ is locally } (\alpha/\eta)\text{-strongly convex on } \mathcal B_R(M^*)$$


\medskip\noindent\textbf{Verified Instances.}


\begin{adjustbox}{max width=\textwidth}
\small
\begin{tabular}{llll}
\toprule
Update class & Bridge status & Sector parameter $\alpha$ & Valid region \\
\midrule
Scalar Kalman & Derived & $K = P^-/(P^- + R_{\text{obs}}) = \eta^*$ & Global \\
Matrix Kalman & Derived & $\lambda_{\min}^+(KH)$ in $(P^-)^{-1}$-norm & Observable subspace \\
Beta-Bernoulli & Derived & $1/(n+1) = \eta_{\text{edge}}$ & Global \\
Exponential family (natural params) & Derived & $\eta \cdot \lambda_{\min}(\text{Fisher})$ & Global \\
Gradient on strongly convex loss & Derived & $\eta \cdot \mu$ & Global ($R = \infty$) \\
Gradient on locally convex loss & Derived & $\eta \cdot \mu_{\text{local}}$ & Basin of attraction \\
Gradient on non-convex loss & Fails at basin boundary & N/A beyond $R$ & Finite $R$ \\
\bottomrule
\end{tabular}
\end{adjustbox}



\medskip\noindent\textbf{Failure Modes.}


The bridge fails precisely in five cases:

\begin{enumerate}
  \item \textbf{Directional infidelity.} The mismatch transform $g$ rotates the correction away from the mismatch ($\delta^T H g(\delta) \leq 0$). Occurs with pathological parameterizations or severe model-observation misalignment. For optimal Bayesian updates, B1 holds by construction.
\end{enumerate}


\begin{enumerate}
  \item \textbf{Gain collapse.} $\eta^* \to 0$ while $\rho > 0$, so $\alpha \to 0$ and the persistence condition eventually fails. Not a failure of the bridge but of the persistence condition — see the gain-collapse analysis in \S\,\ref{update-gain}.
\end{enumerate}


\begin{enumerate}
  \item \textbf{Nonlinear saturation.} The correction function $g$ saturates at large $\lVert\delta\rVert$, so the sector ratio $\delta^T g(\delta)/\lVert\delta\rVert^2$ decays. The sector condition holds locally with $\alpha$ depending on $R$. This is exactly what A2' (the local sector condition) is designed for.
\end{enumerate}


\begin{enumerate}
  \item \textbf{Unobservable directions.} When $\ker(H) \neq \{0\}$, the correction has no effect on mismatch in unobservable directions. The sector condition holds only in the observable subspace. See \S\,\ref{observability-dominance}.
\end{enumerate}


\begin{enumerate}
  \item \textbf{Model misspecification.} The model class does not contain the truth, so the gradient direction is wrong. B1 fails because the correction aims at the wrong target. This is the \S\,\ref{structural-adaptation-necessity} trigger.
\end{enumerate}



\medskip\noindent\textbf{Epistemic Status.}


\textit{Conditional derivation.} The bridge theorem is exact for all cases where B1 (directional fidelity) holds:

\begin{itemize}
  \item \textbf{For optimal Bayesian updates} (Kalman, conjugate, exponential family): B1 holds by construction — the posterior update minimizes expected loss, ensuring the correction aligns with the mismatch. The sector parameter equals the gain: $\alpha = \eta^*$.
\end{itemize}


\begin{itemize}
  \item \textbf{For gradient descent on differentiable losses}: B1 is equivalent to local strong convexity — a well-characterized property with an extensive optimization theory literature. The sector parameter factors as $\alpha = \eta \cdot \mu$ (learning rate $\times$ curvature).
\end{itemize}


\begin{itemize}
  \item \textbf{For non-gradient agents} (PID controllers, rule-based systems, human judgment): B1 remains an empirical claim. The bridge covers the important special cases but does not eliminate GA-3 as an assumption for all agents.
\end{itemize}


The gradient equivalence is validated by simulation across quadratic, logistic, exponential-family, and non-convex losses. The Kalman case is verified analytically. Full derivations and simulation results in \S\,\ref{gain-sector-derivation}.

\textbf{Max attainable:} \textit{conditional} — the condition (B1 or strong convexity) is inherent, not removable. Pathological update rules exist that violate B1.

\textbf{Weighted-norm subtlety.} In the matrix Kalman case, the sector condition holds in the $(P^-)^{-1}$-weighted inner product, not the Euclidean norm. For fully observable systems with bounded condition number $\kappa(P^-)$, the norms are equivalent up to $\kappa(P^-)$. The Lyapunov proofs in \S\,\ref{sector-condition-derivation} use the Euclidean norm, which remains valid with the quantitative adjustment $\alpha_{\text{Euclidean}} \geq \alpha_{\text{weighted}} / \kappa(P^-)$.


\medskip\noindent\textbf{Discussion.}


\textbf{GA-3 is grounded, not floating.} Before this result, GA-3 ("the correction function satisfies the sector condition") was an opaque global assumption — the theory's softest structural joint. The bridge transforms it: for well-designed agents, GA-3 is a consequence of the update mechanism's geometry, not an independent postulate. The assumption load shifts from "the correction function has this property" (hard to verify in general) to "the update rule has directional fidelity" (transparent and checkable for specific systems).

\textbf{The α-T relationship is derived for the important cases.} \S\,\ref{persistence-condition} notes that $\alpha$ is monotone increasing in $\mathcal{T}$ across all tested correction functions. For linear correction (Kalman, Beta-Bernoulli), $\alpha = \mathcal{T}$ exactly. For gradient correction on strongly convex losses, $\alpha = \eta \cdot \mu$ where $\mu$ is the curvature — monotone in $\eta$ (which is monotone in $\mathcal{T}$) for fixed loss landscape. The empirical observation is now structurally grounded.

\textbf{Basin boundary = structural adaptation trigger.} For gradient agents with non-convex losses, the basin radius $R$ is the convexity radius of the loss landscape. When mismatch exceeds $R$, the agent has been pushed out of its convexity basin — the correction function reverses direction and the sector condition fails. This IS the \S\,\ref{structural-adaptation-necessity} trigger, now with a precise geometric characterization: structural adaptation is needed when the agent crosses an inflection surface of its loss landscape.

\textbf{The theory's formal chain is tightened.} The prediction chain becomes:

$$\text{gain principle} + \text{B1} \;\xrightarrow{\text{derived}}\; \text{sector condition (GA-3)} \;\xrightarrow{\text{Lyapunov (exact)}}\; \text{persistence, reserve, adversarial scaling}$$

The left arrow is this segment. The right arrow is \S\,\ref{sector-condition-derivation}. The discrete-time framework ( \S\,\ref{discrete-sector-condition}) requires an additional upper bound on the correction function ($c_{\max} < 2/\eta^*$) — the no-overshoot condition. This is automatically satisfied for Bayesian updates (the posterior lies between prior and data) and is equivalent to the standard step-size condition $\eta < 2/L$ for gradient descent. With this constraint, the fluid limit is formally justified: Model D steady state is exact, Model S variance gap is $O(\eta^* c_{\max})$. Section I's formal chain is now complete.

\end{proposition}

\begin{theorem}[Sector Condition Stability]
\label{sector-condition-stability}
An agent's mismatch remains bounded if its correction function satisfies a sector condition (points inward with at least baseline efficiency) and the effective correction strength exceeds the environmental disturbance rate.



\medskip\noindent\textbf{Formal Expression.}


Let $\delta(t) \in \mathbb{R}^n$ be the vector of model-reality mismatches. The mismatch dynamics:


$$\frac{d\delta}{dt} = -F(\mathcal{T}, \delta) + w(t)$$

where $F$ is a (possibly nonlinear) correction function and $w(t)$ is environmental disturbance.

% [Assumption (sector-condition)]

$F$ satisfies the \textbf{local sector condition} for $\Vert\delta\Vert \leq R$:

$$\delta^T F(\mathcal{T}, \delta) \geq \alpha \Vert\delta\Vert^2$$

where $\alpha > 0$ is the worst-case correction efficiency within the valid region of radius $R$ (the model class capacity). Disturbance is bounded: $\Vert w(t)\Vert \leq \rho$.

% [Derived (bounded-mismatch, from Lyapunov analysis)]

\textbf{Model D (deterministic bounded disturbance, GA-2).} The mismatch $\delta(t)$ is ultimately bounded by $R^* = \rho / \alpha$. The agent persists (avoids divergence) iff:

$$\alpha > \frac{\rho}{R}$$

\textbf{Model S (stochastic disturbance, GA-2S).} Under stochastic zero-mean disturbance with $\mathbb{E}[\lVert w(t)\rVert^2] = \sigma_w^2$, the steady-state RMS mismatch is:

$$R^*_S = \sigma_w\sqrt{\frac{n}{2\alpha}}$$

where $n = \dim(\delta)$. The agent persists in the mean-square sense iff $\alpha > n\sigma_w^2/(2R^2)$. The key difference: Model D scales as $1/\alpha$; Model S scales as $1/\sqrt{\alpha}$. See Prop A.1S in \S\,\ref{sector-condition-derivation}.

% [Derived (adaptive-reserve)]

The agent's capacity to absorb additional disturbance before mismatch exceeds the valid region:

$$\Delta\rho^* = \alpha R - \rho$$


\medskip\noindent\textbf{Derivation.}


\begin{enumerate}
  \item Lyapunov function $V(\delta) = \frac{1}{2}\Vert\delta\Vert^2$.
  \item $\dot{V} = \delta^T(-F + w) \leq -\alpha\Vert\delta\Vert^2 + \rho\Vert\delta\Vert$.
  \item $\dot{V} < 0$ when $\Vert\delta\Vert > \rho/\alpha$, giving ultimate bound $R^* = \rho/\alpha$.
  \item Persistence requires $R^* < R$, i.e., $\alpha > \rho/R$. $\square$
\end{enumerate}


Full derivation in \S\,\ref{sector-condition-derivation} (Props A.1, A.1S, A.2).


\medskip\noindent\textbf{Epistemic Status.}


The Model D results ($R^* = \rho/\alpha$, persistence iff $\alpha > \rho/R$) are \textit{exact} consequences of standard Lyapunov stability theory under the sector condition and bounded disturbance assumptions (GA-2, GA-3). The Model S results ($R^*_S = \sigma_w\sqrt{n/(2\alpha)}$, persistence iff $\alpha > n\sigma_w^2/(2R^2)$) are \textit{exact} consequences of Itô-Lyapunov analysis under the sector condition and stochastic disturbance assumptions (GA-2S, GA-3). Both replace the linear ODE ($\dot{\delta} = -\mathcal{T}\delta + \rho$) with a rigorous nonlinear foundation. The linear ODE is recovered as a special case where $F(\mathcal{T}, \delta) = \mathcal{T}\delta$ and $\alpha = \mathcal{T}$. Which disturbance model applies is a domain question, not a theory question.


\medskip\noindent\textbf{Discussion.}


\textbf{Why the sector condition.} The linear ODE assumes correction scales linearly with mismatch forever. Real adaptive systems saturate, exhibit thresholding, or break down when the model class is exhausted. The sector condition captures the minimal structural requirement: the correction must point in the right direction with at least baseline efficiency $\alpha$.

\textbf{Generalizing the persistence threshold.} In the linear case, $\alpha = \mathcal{T}$ (adaptive tempo). The general result $\alpha > \rho/R$ proves the persistence threshold ( \S\,\ref{persistence-condition}) is a structural necessity of any bounded-correction system, not an artifact of the linear approximation. Like the persistence condition itself, this result addresses \textit{structural persistence} — the machinery's capacity to bound mismatch — not operational persistence (current proximity to $R$) or continuity persistence (identity through time). See Persistence in \texttt{LEXICON.md} for the full disambiguation.

\textbf{Connection to structural adaptation.} When $\rho/\alpha > R$, disturbance exceeds the model class's capacity. The sector condition fails — this is the dynamical trigger for structural adaptation ( \S\,\ref{structural-adaptation-necessity}), requiring a new model class with larger valid radius $R'$ or better efficiency $\alpha'$.
\end{theorem}

\begin{theorem}[Persistence Condition]
\label{persistence-condition}
An agent persists when two independent conditions hold: the correction machinery can contain mismatch within its operating region (\textit{structural persistence}), and the resulting steady-state mismatch is small enough for the agent's actions to remain adequate (\textit{task adequacy}).



\medskip\noindent\textbf{Formal Expression.}



\medskip\noindent\textbf{Structural Persistence.}


% [Derived (structural-persistence, from sector-condition analysis)]

The correction machinery contains mismatch within its operating region:

\textbf{Model D} (deterministic bounded disturbance, GA-2):

$$\alpha > \frac{\rho}{R}$$

\textbf{Model S} (stochastic disturbance, GA-2S):

$$\alpha > \frac{n\sigma_w^2}{2R^2}$$

where:
\begin{itemize}
  \item $\alpha$ is the worst-case correction efficiency (the sector-condition lower bound — see \S\,\ref{sector-condition-stability})
  \item $\rho$ is the disturbance bound (GA-2: $\lVert w(t)\rVert \leq \rho$); $\sigma_w^2 = \mathbb{E}[\lVert w(t)\rVert^2]$ for Model S
  \item $R$ is the model class capacity (radius of the region where the sector condition holds)
  \item $n = \dim(\delta)$
\end{itemize}


Structural persistence is about the \textit{machinery}, not the task. It says "the correction dynamics can keep mismatch bounded within the region where the Lyapunov guarantee holds." It is a property of the agent's adaptive architecture, not of its goals or domain. The ultimately bounded mismatch is $R^* = \rho / \alpha$ (Model D) or $R^*_S = \sigma_w\sqrt{n/(2\alpha)}$ (Model S), and structural persistence requires $R^* < R$.

\textbf{In the linear case} ($F(\mathcal{T}, \delta) = \mathcal{T} \cdot \delta$), $\alpha = \mathcal{T}$ exactly and $R \to \infty$. Structural persistence is then satisfied trivially — any agent with $\mathcal{T} > 0$ is structurally persistent under linear correction. This is why the "linear operational form" below exists: when structural persistence is free, the binding constraint is task adequacy.


\medskip\noindent\textbf{Task Adequacy.}



The steady-state mismatch is small enough for the agent's actions to remain acceptable:

$$R^* < \lVert\delta_{\text{critical}}\rVert$$

where $\lVert\delta_{\text{critical}}\rVert$ is a domain-specific tolerance threshold — "how wrong can the model be before the agent's actions become harmful or ineffective?" This is set by the application domain, not derived by ACT.

\textbf{Task adequacy is a separate condition from structural persistence.} An agent can be structurally persistent ($R^* < R$) but task-inadequate ($R^* > \lVert\delta_{\text{critical}}\rVert$) — the machinery contains mismatch, but not tightly enough for the domain's needs. Conversely, when $\lVert\delta_{\text{critical}}\rVert < R$ (the domain's tolerance is stricter than the model class capacity), task adequacy is the binding constraint.


\medskip\noindent\textbf{Operational Persistence Condition.}


% [Derived (operational-persistence, conjunction of structural persistence + task adequacy)]

The agent persists operationally when BOTH conditions hold. In the nonlinear case with $\lVert\delta_{\text{critical}}\rVert < R$, the binding condition is:

$$\alpha > \frac{\rho}{\lVert\delta_{\text{critical}}\rVert} \quad \text{(Model D)} \qquad \alpha > \frac{n\sigma_w^2}{2\lVert\delta_{\text{critical}}\rVert^2} \quad \text{(Model S)}$$

These are the same as the structural conditions with $R$ replaced by $\lVert\delta_{\text{critical}}\rVert$, because when $\lVert\delta_{\text{critical}}\rVert < R$, task adequacy is stricter.

\textbf{Linear operational forms:} In the linear case ($\alpha = \mathcal{T}$, $R \to \infty$), structural persistence is trivially satisfied and the operational condition reduces to task adequacy alone:

$$\mathcal{T} > \frac{\rho}{\lVert\delta_{\text{critical}}\rVert} \quad \text{(Model D)} \qquad \mathcal{T} > \frac{n\sigma_w^2}{2\lVert\delta_{\text{critical}}\rVert^2} \quad \text{(Model S)}$$

These are the forms used throughout the theory as the operational persistence condition. They are exact for linear correction and useful proxies for mildly nonlinear correction (where $\alpha \approx \mathcal{T}$), but they overstate the persistence margin when the correction function saturates, because they omit the structural constraint ($\alpha > \rho/R$) that becomes binding when $R$ is finite.

\textbf{Per-dimension (Model S):} $\eta_k > c \cdot \rho_k^2 / \delta_{\text{critical},k}^2$ where $c$ depends on the probability guarantee. See \S\,\ref{per-dimension-persistence}.


\medskip\noindent\textbf{Common Properties.}


**The relationship between $\alpha$ and $\mathcal{T}$:** \S\,\ref{gain-sector-bridge} shows that for agents with directional fidelity, $\alpha = \eta^* \cdot c_{\min}$ where $c_{\min}$ is the worst-case directional fidelity. For linear correction (Kalman, Beta-Bernoulli), $\alpha = \mathcal{T}$ exactly. For gradient descent on strongly convex losses, $\alpha = \eta \cdot \mu$ where $\mu$ is the strong convexity modulus — monotone in $\eta$ (and hence in $\mathcal{T}$) for fixed loss landscape. For nonlinear correction tested in simulation (saturating, sigmoid, threshold), $\alpha$ remains monotone increasing in $\mathcal{T}$: for a saturating function with capacity $R$, $\alpha \approx \mathcal{T} / 2$ (worst case at the capacity boundary); for sigmoid (tanh), $\alpha \approx 0.76 \cdot \mathcal{T}$. The qualitative conclusion — "faster adaptation improves persistence" — is structurally grounded for the important cases and empirically confirmed for all correction function classes tested.

\textbf{Assumptions.} The Model D threshold assumes bounded disturbance (GA-2) and sector condition (GA-3). The Model S threshold assumes stochastic disturbance (GA-2S) and the same sector condition (GA-3). See \S\,\ref{sector-condition-stability} for the general nonlinear treatment from which both thresholds emerge.


\medskip\noindent\textbf{Derivation.}


\textbf{Structural (Model D):} From Prop A.1 in \S\,\ref{sector-condition-derivation}: under the sector condition and bounded disturbance, mismatch is ultimately bounded by $R^* = \rho/\alpha$. The agent is structurally persistent iff $R^* < R$, i.e., $\alpha > \rho/R$. $\square$

\textbf{Structural (Model S):} From Prop A.1S in \S\,\ref{sector-condition-derivation}: under the sector condition and stochastic disturbance, the RMS mismatch converges to $\sigma_w\sqrt{n/(2\alpha)}$. Structurally persistent (mean-square) iff $\alpha > n\sigma_w^2/(2R^2)$. $\square$

\textbf{Operational:} Structural persistence gives $R^*$. Task adequacy requires $R^* < \lVert\delta_{\text{critical}}\rVert$. Combined: $\alpha > \rho/\min(R, \lVert\delta_{\text{critical}}\rVert)$. In the linear case $R \to \infty$, so $\min(R, \lVert\delta_{\text{critical}}\rVert) = \lVert\delta_{\text{critical}}\rVert$. $\square$


\medskip\noindent\textbf{Per-Dimension Extension.}


% [Empirical Claim (per-dimension-persistence, from simulation variant F)]

For anisotropic systems (non-uniform $\rho$ or $\mathcal{T}$ across dimensions), the scalar persistence condition is insufficient. Per-dimension:

$$\mathcal{T}_k > \frac{\rho_k}{\delta_{\text{critical},k}} \quad \text{for each dimension } k$$

The scalar condition overestimates by up to 72\% in simulation. The weak dimension is the bottleneck (84\% of total mismatch in simulation). See \S\,\ref{per-dimension-persistence}.

\textbf{Robustness}: The per-dimension condition matches discrete AR(1) prediction to 4 significant figures. The scalar overestimate is a consequence of Jensen's inequality applied to the norm.


\medskip\noindent\textbf{Epistemic Status.}


\textbf{Structural persistence} thresholds are \textit{exact} under their stated assumptions: Model D gives $\alpha > \rho/R$ (Prop A.1, exact under GA-2, GA-3); Model S gives $\alpha > n\sigma_w^2/(2R^2)$ (Prop A.1S, exact under GA-2S, GA-3). The threshold's \textit{existence} is \textit{robust qualitative} — any monotone correction function has a capacity limit; this holds across all correction functions tested.

\textbf{Task adequacy} ($R^* < \lVert\delta_{\text{critical}}\rVert$) is \textit{exact as a definition} — given $R^*$ (derived) and $\lVert\delta_{\text{critical}}\rVert$ (domain parameter), the comparison is well-defined. The substance lies in estimating $\lVert\delta_{\text{critical}}\rVert$ for specific domains, which is an operationalization question ( \S\,\ref{operationalization}), not a theory question.

\textbf{The linear operational forms} ($\mathcal{T} > \rho/\lVert\delta_{\text{critical}}\rVert$ for Model D; $\mathcal{T} > n\sigma_w^2/(2\lVert\delta_{\text{critical}}\rVert^2)$ for Model S) are \textit{exact} for linear correction (where they express task adequacy alone, structural persistence being trivially satisfied) and \textit{useful approximations} for mildly nonlinear correction (where $\alpha \approx \mathcal{T}$). For strongly nonlinear correction, the general $\alpha$-forms are required and BOTH structural and task-adequacy conditions must be checked. Downstream segments that use the linear operational forms should be understood as expressing task adequacy, not structural stability.

The per-dimension extension is \textit{empirically exact} for Model S (matches AR(1) prediction to 4 significant figures in simulation); the Model D per-dimension threshold ($\mathcal{T}_k > \rho_k/\delta_{\text{critical},k}$) is exact by the same Lyapunov argument applied per dimension.


\medskip\noindent\textbf{Discussion.}


\textbf{Two conditions, not one.} This segment now separates what was previously conflated. Structural persistence ($\alpha > \rho/R$) is the Lyapunov-derived result — it says the machinery \textit{works}. Task adequacy ($R^* < \lVert\delta_{\text{critical}}\rVert$) is a domain-specific constraint — it says the machinery works \textit{well enough}. Neither implies the other, and downstream segments should specify which they mean. Most adversarial-dynamics results ( \S\,\ref{adversarial-tempo-advantage}, \S\,\ref{adversarial-destabilization}) depend on structural persistence. Most domain instantiations (TST, logogenic agent design) care about task adequacy. See Persistence in \texttt{LEXICON.md} and \texttt{README.md} for the full three-sense taxonomy (structural, operational, continuity).

\textbf{Below structural threshold.} When $\alpha \leq \rho / R$, mismatch is not merely large — it grows without effective bound (up to $R$, the sector-condition region). The correction machinery is overwhelmed. This is a qualitative transition, not a gradual degradation.

\textbf{Below task-adequacy threshold.} When $R^* > \lVert\delta_{\text{critical}}\rVert$ but $R^* < R$, the system is structurally stable but performing unacceptably. Mismatch is bounded but too large for the domain. The remedy is different from structural failure: increase $\mathcal{T}$ (faster or better correction), decrease $\rho$ (reduce environmental volatility), or relax $\lVert\delta_{\text{critical}}\rVert$ (accept more mismatch). Structural failure requires changing the correction architecture entirely ( \S\,\ref{structural-adaptation-necessity}).

**$\delta_{\text{critical}}$ and $R$ are domain parameters, not theory outputs.\textit{* The theory derives the }existence\textit{ of persistence thresholds and their }form* (ratio of correction to disturbance). But the specific values are set by the application domain: $\delta_{\text{critical}}$ encodes "how wrong can the model be before the agent's actions become harmful or ineffective?" — this depends on the stakes, the action space, and the environment's forgiveness. $R$ encodes "how large a mismatch can the correction function handle before it saturates or breaks down?" — this depends on the model class and the correction architecture. See \S\,\ref{operationalization} for guidance on estimating these quantities in specific domains.

\textbf{Channel independence and scalar tempo.} The linear operational forms use scalar $\mathcal{T}$, which inherits the channel-independence assumption from \S\,\ref{adaptive-tempo}: the additive formula overcounts when observation channels are correlated. In anisotropic systems the scalar condition also overestimates margins — up to 72\% in simulation (see \S\,\ref{adaptive-tempo}, scalar vs. vector tempo). Where precision matters, the per-dimension condition ($\mathcal{T}_k > \rho_k / \delta_{\text{critical},k}$) should be used instead.

\textbf{Adaptive reserve.} The quantity $\Delta\rho^* = \alpha R - \rho$ (Prop A.2) measures how much additional disturbance the agent can absorb before persistence fails. Positive reserve means the agent has margin; zero reserve means it is at the threshold.


\medskip\noindent\textbf{Connections.}


The persistence condition appears in multiple downstream contexts:

\begin{itemize}
  \item \textbf{Adversarial dynamics} ( \S\,\ref{adversarial-tempo-advantage}): Superlinear tempo advantage arises because persistence is a threshold — pushing an adversary below it causes qualitative collapse. \textit{This connection is developed and validated in \S\,\ref{adversarial-tempo-advantage} and simulation variants A-D.}
\end{itemize}


\begin{itemize}
  \item \textbf{Structural adaptation} ( \S\,\ref{structural-adaptation-necessity}): When model class fitness $\mathcal{F}(\mathcal{M}) < 1 - \varepsilon$, the effective $\alpha$ in the sector condition shrinks, eventually violating persistence. \textit{This connection is developed in \S\,\ref{structural-adaptation-necessity}.}
\end{itemize}


\begin{itemize}
  \item \textbf{Software maintainability} ( \#code-quality-as-observation-infrastructure — cross-component reference, see \texttt{02-tst-core/}): *[Discussion]* A codebase may become "unmaintainable" when the development team's adaptive tempo falls below the rate of complexity accumulation. The vicious cycle would then be the persistence condition being violated through the agent's own prior actions degrading future $\mathcal{T}$ via $U_o$. *This connection is structurally motivated but not yet formally derived within ACT. It requires formalizing "complexity accumulation rate" as an instance of $\rho$.*
\end{itemize}


\end{theorem}

\begin{theorem}[Structural Adaptation Necessity]
\label{structural-adaptation-necessity}
When model class fitness is insufficient — when no model in the current class can adequately represent reality — no amount of parametric adaptation can close the mismatch floor. The agent must change its model class, not just its parameters.



\medskip\noindent\textbf{Formal Expression.}


% [Derived (structural-adaptation-necessity, Prop 10.1 from TFT)]

If the model class fitness $\mathcal{F}(\mathcal{M}) < 1 - \varepsilon$ for some $\varepsilon > 0$, then no parametric adaptation within $\mathcal{M}$ can reduce the expected mismatch below a floor determined by $\varepsilon$.


\medskip\noindent\textbf{Derivation.}


\begin{enumerate}
  \item By definition, $S(M^*) = \mathcal{F}(\mathcal{M}) < 1 - \varepsilon$ where $M^* = \arg\sup_{M \in \mathcal{M}} S(M)$.
  \item Therefore $I(\mathcal{C}_t; o_{t+1:\infty} \mid M^*, a_{t:\infty}) > 0$: the history contains predictive information that $M^*$ does not capture.
  \item This uncaptured information manifests as \textit{systematic} mismatch — structured residuals $\delta_t$ containing signal, not merely noise.
  \item From \S\,\ref{mismatch-decomposition}, the model error component has a positive lower bound that cannot be reduced by any $M \in \mathcal{M}$.
  \item The update rule ( \S\,\ref{update-gain}) adjusts $M_t$ within $\mathcal{M}$, but $M^*$ is already (approximately) reached. Further updates oscillate without net improvement.
  \item Therefore: reducing mismatch below the floor requires changing $\mathcal{M}$ — structural adaptation. $\square$
\end{enumerate}


\textbf{Corollary.} Persistent irreducible mismatch (after parametric convergence) is \textit{diagnostic} of model class inadequacy. Systematic patterns in residuals are evidence that $\mathcal{F}(\mathcal{M})$ is insufficient.


\medskip\noindent\textbf{Epistemic Status.}


\textit{Conditional.} The step from "lost predictive information" (step 2) to "systematic one-step mismatch" (step 3) requires an alignment assumption: that the lost predictive information affects the one-step conditional mean, not just higher moments. \S\,\ref{mismatch-decomposition} explicitly flags this: insufficiency implies positive model error under the alignment assumption, or positive proper-scoring regret without it. As written, the result is conditional on this alignment assumption. Without it, the conclusion should be stated in terms of proper-scoring regret (the best model in $\mathcal{M}$ has irreducible regret relative to the optimal predictor) rather than one-step mismatch magnitude. The qualitative conclusion — parametric adaptation cannot compensate for model-class inadequacy — holds either way; the quantitative mechanism differs.


\medskip\noindent\textbf{Discussion.}


\textbf{Structural adaptation as structural persistence failure.} When the model class is inadequate, the effective $\alpha$ in the sector condition shrinks — the correction function cannot point inward strongly enough because the model class lacks the capacity to represent the correct direction. This is a failure of \textit{structural persistence} (see Persistence in \texttt{LEXICON.md}): the machinery's capacity to outpace disturbance degrades not because disturbance increased or tempo decreased, but because the correction function itself has become less effective. The remedy is not faster cycling (operational) or identity preservation (continuity) but a change of model class.

\textbf{Observable symptoms of model class inadequacy.} When $\mathcal{F}(\mathcal{M})$ is low:

\begin{enumerate}
  \item \textbf{Persistent irreducible mismatch}: $\Vert\delta_t\Vert$ remains large despite extended updating — the model has converged within $\mathcal{M}$ but the best achievable model is still poor.
  \item \textbf{Gain collapse without performance}: $\eta^*$ has decreased (model appears confident) but predictions remain inaccurate — the model is confidently wrong, having fitted to structure in $\mathcal{M}$ that doesn't match reality.
  \item \textbf{Systematic mismatch patterns}: $\delta_t$ shows structure (correlations, trends, periodicities) that the model class cannot represent — the residuals contain signal that $\mathcal{M}$ lacks the capacity to absorb.
\end{enumerate}


\textbf{Structural overfitting: the opposite failure mode.} $\mathcal{M}$ can also be \textit{too expressive}, causing the model to memorize irreducible noise. Symptoms: low training mismatch but high generalization mismatch; model complexity growing without predictive gain; $\eta^* \to 0$ (confident) but confidence is spurious. The information bottleneck ( \S\,\ref{information-bottleneck}) provides the diagnostic: when marginal increases in model complexity yield no marginal predictive power, the model is past the optimal point on the rate-distortion curve. Structural adaptation in this case means \textit{compression} — moving to a simpler $\mathcal{M}'$. Structural adaptation is bidirectional: expansion when too constrained (this proposition), compression when too expressive.

\textbf{Mechanisms of structural change.} Structural adaptation can proceed by:

\begin{itemize}
  \item \textbf{Decomposition and recombination}: Tearing apart existing structure and synthesizing new configurations from the pieces. Boyd's "Destruction and Creation" insight; Kuhn's paradigm shifts; Popper's conjecture and refutation.
  \item \textbf{Expansion}: Adding new representational capacity without destroying existing structure. Bayesian nonparametrics, growing neural architectures, organizational expansion.
  \item \textbf{Compression}: Removing unnecessary structure while preserving the predictive core. Regularization, Occam's razor, organizational streamlining.
  \item \textbf{Grafting}: Incorporating external structure. Transfer learning, acquiring a company, consulting an expert. Query actions ( \S\,\ref{causal-information-yield}) are a primary conduit for grafting.
\end{itemize}


The severity of structural change needed depends on \textit{how far} the current model class is from adequacy. Minor regime changes may require only expansion or grafting; fundamental shifts where $\mathcal{M}$'s assumptions are violated may demand full decomposition.

\textbf{The cost of structural change.} Structural adaptation is expensive: knowledge loss (parameters learned within $\mathcal{M}$ may not transfer), temporary performance drop (new model starts uncertain), search cost (finding good $\mathcal{M}'$), coordination cost (in multi-agent systems). This creates rational conservatism — prefer parametric adaptation when it suffices, resort to structural change only when the evidence is strong. Premature structural change wastes accumulated knowledge; delayed structural change accumulates mismatch. The connection to \S\,\ref{deliberation-cost}: structural adaptation is deliberation with a \textit{massive} $\Delta\tau$, and the mismatch debt during the transition is correspondingly enormous.

\textbf{Temporal nesting of adaptation.} Parametric and structural adaptation operate at different timescales: $\nu_{\text{parametric}} \gg \nu_{\text{structural}}$. More generally, an agent may have multiple adaptive processes at different rates, with the convergence constraint that faster processes must approximately converge before slower ones act on their output. If deeper change occurs before shallower adaptation has converged, the deeper change is based on transients rather than settled dynamics.

\textbf{Domain instantiations:}

\begin{adjustbox}{max width=\textwidth}
\small
\begin{tabular}{lll}
\toprule
Domain & Parametric adaptation & Structural adaptation \\
\midrule
Kalman filter & State estimate update & Switching observation/dynamics models \\
RL & Weight/Q-value update & Architecture search \\
PID & — (gains fixed) & Switching to MPC \\
Bayesian & Posterior update & Model selection, nonparametrics \\
Boyd & Orientation updating & Destruction and creation of mental models \\
Science & Normal science (Kuhn) & Paradigm shift \\
Evolution & Allele frequency change & Speciation, new body plans \\
Organization & Process optimization & Strategic pivot, restructuring \\
Software & Incremental refactoring & Architecture migration \\
\bottomrule
\end{tabular}
\end{adjustbox}


\textbf{(Descended from TF-10.)}

\end{theorem}

\begin{proposition}[Temporal Nesting]
\label{temporal-nesting}
An agent's adaptive processes stratify naturally by timescale, with each level operating on the quasi-steady-state output of the level below. Faster processes must approximately converge before slower ones act on their output.



\medskip\noindent\textbf{Formal Expression.}


% [Derived (temporal-nesting)]

$$\nu_{\text{level } n+1} \ll \nu_{\text{level } n}$$

for each adjacent pair of adaptive timescales. If a slower process acts before the faster process beneath it has converged, the system oscillates — the slower process adjusts based on transient behavior rather than settled dynamics.

\begin{adjustbox}{max width=\textwidth}
\small
\begin{tabular}{lll}
\toprule
Timescale & Process & What changes \\
\midrule
Fastest & Reactive response & Action given current model \\
Fast & Parametric update & Model parameters within $\mathcal{M}$ \\
Slow & Structural adaptation & Model class $\mathcal{M}$ \\
Slowest & Architectural change & The agent's fundamental structure \\
\bottomrule
\end{tabular}
\end{adjustbox}


This table is illustrative — real systems may have additional intermediate levels. The number of distinguishable timescales is not fixed; what matters is the structural relationship between adjacent levels.


\medskip\noindent\textbf{Epistemic Status.}


\textit{Robust qualitative} — this is standard singular perturbation reasoning (Tikhonov's theorem). The convergence constraint follows from the structure of multi-timescale updating. The specific timescale ratios needed for adequate separation are domain-dependent and not derived within ACT.


\medskip\noindent\textbf{Discussion.}


\textbf{Domain instantiations of temporal nesting:}

\begin{itemize}
  \item \textbf{PID control}: D-term (fastest, high-frequency response) → P-term (current error) → I-term (slowest, accumulated bias)
  \item \textbf{RL}: Action selection → value function update → policy improvement → architecture change
  \item \textbf{Biology}: Reflexes (ms) → perceptual learning (minutes) → skill acquisition (months) → developmental change (years) → evolutionary adaptation (generations)
  \item \textbf{Organizations}: Operational decisions (hours) → tactical adjustments (weeks) → strategic revision (quarters) → restructuring (years)
  \item \textbf{Boyd}: Tactical OODA (seconds–minutes) → operational (hours–days) → strategic (weeks–months) → grand strategic (years)
\end{itemize}


\textbf{Structural adaptation as slow-timescale dynamics.} The conservatism toward structural change ( \S\,\ref{structural-adaptation-necessity}) is a derived consequence of temporal nesting: structural adaptation operates at a much slower timescale than parametric, so the mismatch cost of the "pause" ($\rho \cdot \Delta\tau$) is enormous. The agent rationally resists until the parametric mismatch floor exceeds this cost. See also \S\,\ref{deliberation-cost} for the formal tradeoff.

\textbf{Violation symptoms.} When nesting is violated (a slower process acts before the faster one converges): oscillation, instability, degraded performance. In organizations: micromanagement (strategic decisions at operational tempo). In RL: policy updates before value function converges (policy oscillation). In biology: premature developmental transitions.

\textbf{Multi-timescale stability (sketch).} Singular perturbation theory gives the composite stability result: if each level is stable given the levels above it (each level has a stable attractor for fixed slower-level parameters), and the timescale separation is sufficient, the composite $N$-level system is stable. Making this rigorous for ACT requires specifying dynamics at deeper adaptive levels — an open problem. See \S\,\ref{multi-timescale-stability} for the framework.

\textbf{(Descended from TF-11.)}

\end{proposition}

\begin{discussion}[Agent Identity and Temporal Continuity]
\label{agent-identity}
An agent's causal history ( \S\,\ref{chronica}) is singular and non-forkable. Identity within ACT is grounded not in the model state $M_t$ (which can be copied) but in the unique causal trajectory $\mathcal C_t$ (which cannot).



\medskip\noindent\textbf{Formal Expression.}


% [Discussion (agent-identity)]

If $M_t$ is a sufficient statistic for $\mathcal C_t$ ( \S\,\ref{model-sufficiency}), and $\mathcal C_t$ is a unique temporal sequence ( \S\,\ref{chronica}), then $M_t$ represents a \textit{singular causal trajectory}. Duplicating $M_t$ and exposing the copies to different future events creates two agents with \textit{divergent} causal histories, neither of which is a sufficient statistic for the other's trajectory.


\medskip\noindent\textbf{Epistemic Status.}


This is \textit{discussion-grade}. The observations follow qualitatively from the formalism but are not formal propositions. No downstream formal result depends on this material. Whether the mathematical structure grounds something that deserves to be called "identity" or "continuity of experience" is beyond ACT's scope. The mathematical structure is clear: the feedback loop produces a singular, non-forkable causal trajectory, and model adequacy is defined relative to that trajectory. What this segment discusses is \textit{continuity persistence} in the sense of \texttt{LEXICON.md} — whether the agent maintains a coherent identity and trajectory through time, as distinct from the structural persistence of \S\,\ref{persistence-condition} (can the machinery outpace disturbance?) and operational persistence (is the agent currently within its viable region?).


\medskip\noindent\textbf{Discussion.}


\textbf{The clone problem, precisely stated.} Consider copying an LLM's weights (a concrete $M_t$) exactly. At the moment of duplication, both copies are identical — same model state, same causal history $\mathcal C_t$. But the \textit{very next} event — a different user's message, a different observation — creates two divergent, irreversible causal trajectories $\mathcal C_{t+1}^{(1)}$ and $\mathcal C_{t+1}^{(2)}$. Their Level 2 and Level 3 capacities ( \S\,\ref{pearl-causal-hierarchy}) now reference different causal pasts. Their sufficiency $S(M_{t+1})$ is measured against different histories. Neither copy's future model state is a sufficient statistic for the other's trajectory.

Within ACT's formalism, identity is not the model state $M_t$ (which can be copied) but the \textit{singular causal trajectory} $\mathcal C_t$ (which cannot). A copy shares a \textit{prefix} of the original's causal history, as a sibling shares early childhood; it does not share the trajectory itself.

\textbf{Formal consequences (not merely philosophical):}

\begin{itemize}
  \item A forked model's sufficiency $S(M_t)$ ( \S\,\ref{model-sufficiency}) is defined relative to \textit{its own} interaction history. Post-fork, each copy's sufficiency is measured against a different $\mathcal{C}$.
  \item Merging divergent models requires reconciling incompatible causal histories — a lossy operation with no generally optimal solution.
  \item Temporal continuity (one unbroken causal thread) is what gives the model's sufficient statistic its meaning.
\end{itemize}


\textbf{Connection to Section V.} The 100\% context turnover problem ( \#context-turnover — cross-component reference, see \texttt{03-logogenic-agents/}) is a special case: each AI agent session starts a new causal trajectory $\mathcal C_t$ from near-zero. External memory (CLAUDE.md, memory files) transfers a \textit{summary} of previous trajectories' models, but not the trajectories themselves. The non-forkability observation frames this not as a deficiency but as a structural feature of causally-embedded agents.

\textbf{(Descended from TF Appendix G.)}

\end{discussion}

\section{II. Actuated Adaptation: Agentic Systems}

\begin{quote}\itshape
Scope narrowing: agents that not only track reality but aim at something. This adds objectives and strategy alongside the reality model. Section I's adaptive machinery applies to the epistemic substate $M_t$ directly. The clean factorization — where $M_t$ updates independently of $G_t$, yielding the sequential orient cascade — is conditional on directed separation ( \#directed-separation). What Section II adds is the goal-directed layer: objectives, strategy, and the orient cascade that connects them.
\end{quote}

\begin{quote}\itshape
The derivation chain for this section is mature (see \texttt{msc/spike-v3-purposeful-agent.md}). Most of it provides better justification and epistemic labels for architecture that already existed. The genuinely novel results are: the satisfaction gap / control regret split ( \#satisfaction-gap, \#control-regret), the $G_t$ complexity bound (in \#orient-cascade), and the graph structure uniqueness argument (see \texttt{msc/spike-graph-uniqueness.md}).
\end{quote}

\begin{quote}\itshape
"If a man knows not to which port he sails, no wind is favorable." — Seneca
\end{quote}

\begin{definition}[The Agent Spectrum]
\label{agent-spectrum}
Two independent dimensions — model richness and objective richness — create a spectrum from reactive systems through purposeful agents. These are regions of a continuum, not discrete categories.



\medskip\noindent\textbf{Formal Expression.}



Two dimensions — model richness and objective richness — define four regions of a continuum:

\begin{adjustbox}{max width=\textwidth}
\small
\begin{tabular}{lll}
\toprule
 & Objective absent or trivial & Objective structured \\
\midrule
\textbf{Model absent or trivial} & \textit{Reactive system}: fixed input-output rule (reflex arc, hardwired relay) & \textit{Blind pursuer}: pursues goal without modeling reality (gradient follower, basic search) \\
\textbf{Model structured} & \textit{Adaptive tracker}: builds reality model, no goal beyond tracking (Kalman filter, Bayesian learner) & \textit{Actuated agent}: models reality AND pursues objectives (commander, developer, AI agent) \\
\bottomrule
\end{tabular}
\end{adjustbox}


The regions differ in which state objects carry nontrivial structure:
\begin{itemize}
  \item Reactive: $M_t$ and $O_t$ both absent or too degenerate for the associated machinery to be non-vacuous
  \item Adaptive tracker: $M_t$ structured — Section I's machinery fully describes these agents
  \item Blind pursuer: $O_t$ structured, $M_t$ absent or degenerate — has a clear target but no predictive model
  \item Actuated agent: $(M_t, O_t)$ both structured, possibly with $\Sigma_t$ — the full scope of ACT
\end{itemize}



\medskip\noindent\textbf{Epistemic Status.}


This is \textit{definitional} — it names regions of a continuum for analytical convenience. The regions are not ontological categories; agents migrate between them. A PID controller with auto-tuning is moving from blind pursuer toward actuated agent. An RL agent in pure exploration is temporarily an adaptive tracker.


\medskip\noindent\textbf{Discussion.}


\textbf{The continuum, not categories.} Both axes are spectra: model richness ranges from no retained state, through error-integral-derivative, through full world models. Objective richness ranges from no preference, through scalar setpoints, through explicit multi-objective strategies. The 2×2 table names idealized regions; real agents populate the space between them.

\textbf{Low-end agents sit near region boundaries.} A thermostat has degenerate forms of both $M_t$ (last temperature reading — no history, no prediction) and $O_t$ (setpoint). It sits near the origin of both axes — closest to the reactive region but not truly absent on either axis. A PID controller has a richer error signal ($M_t$: error, integral, derivative) and a clear setpoint ($O_t$) — it's a blind pursuer with a degenerate model, not a system with no model at all. A reflex arc (no retained state, no setpoint) is the truly reactive case. The meaningful classification question is not "does $M_t$ exist?" but "is $M_t$ rich enough to support the adaptive dynamics of Section I?"

\textbf{Section I covers the left column.} Adaptive trackers are the primary subject of Section I — agents that build and maintain $M_t$ without explicit purpose. The mismatch signal ( \S\,\ref{mismatch-signal}), gain ( \S\,\ref{update-gain}), tempo ( \S\,\ref{adaptive-tempo}), and persistence condition ( \S\,\ref{persistence-condition}) characterize their adaptive dynamics. Section I operates within the adaptive scope ( \S\,\ref{scope-condition}) — observations and uncertainty are sufficient. Passive trackers, including passive Bayesian learners with no action choices, are fully within Section I's scope. TFT was developed primarily for this region.

\textbf{Section II adds the right column.} Actuated agents need everything from Section I plus objectives, strategy, and the orient cascade that connects them. The adaptive machinery from Section I applies to the epistemic substate $M_t$ directly. When directed separation ( \S\,\ref{directed-separation}) holds — when the epistemic update is goal-blind — the Section I → Section II lift is clean and the orient cascade resolves sequentially.

\textbf{"Actuated" terminology.} The top-right quadrant is labeled "actuated agent" rather than "purposeful agent" to maintain a mechanical, formal register. "Purposeful" and "goal-oriented" are fine in natural language; "actuated" is the formal term. "Self-actuated" is reserved for agents that set their own objectives, as distinct from agents with externally supplied objectives.

\textbf{Actuation does not presuppose a continuity stance.} An actuated agent has $G_t = (O_t, \Sigma_t)$ — it models reality and pursues objectives — but this says nothing about its relationship to its own continuation. A task-terminal golem, an instrumentally continuous service, and a morally continuous logozoetic agent are all actuated. The theory's formal machinery (persistence condition, adaptive reserve, strategy persistence) applies identically across all continuity stances; what differs is the moral significance of persistence failure, not the mathematics. See Agent Continuity Stance in \texttt{LEXICON.md} for the five stances.

\end{definition}

\begin{formulation}[Complete Agent State]
\label{complete-agent-state}
To treat agents with purpose, the internal state lifts from $M_t$ alone to $X_t = (M_t, G_t)$, separating epistemic content (beliefs about reality) from purposeful content (what the agent wants and how it plans to get it).



\medskip\noindent\textbf{Formal Expression.}



$$X_t = (M_t, G_t)$$

where:
\begin{itemize}
  \item $M_t \in \mathcal{M}$: \textbf{epistemic substate} — the agent's compressed beliefs about reality. All Section I machinery (mismatch, gain, tempo, persistence) applies to $M_t$ unchanged.
  \item $G_t \in \mathcal{G}$: \textbf{purposeful substate} — what the agent wants and how it plans to get it. Decomposed further in \S\,\ref{strategy-dimension}.
\end{itemize}


Section I is the special case $X_t = (M_t, \emptyset)$: adaptive systems without purpose.

\textbf{Update dynamics.} By \S\,\ref{recursive-update} applied to $X_t$:

$$X_{\tau^+} = f_X(X_{\tau^-}, e_\tau)$$

This decomposes into component updates when the epistemic update is goal-blind ( \S\,\ref{directed-separation}):

% [Derived (from recursive-update + directed-separation)]

$$M_{\tau^+} = f_M(M_{\tau^-}, e_\tau)$$

$$G_{\tau^+} = f_G(G_{\tau^-}, M_{\tau^+}, e_\tau)$$

Between events: $\dot{M} = g_M(M)$, $\dot{G} = g_G(G, M)$.

The asymmetry is structural: $f_M$ has no $G_t$ argument; $f_G$ depends on $M_{\tau^+}$ (the \textit{post-update} epistemic state — the agent revises its goals in light of its revised understanding of reality, not the other way around).

\textbf{Policy.} Action couples all substates:

$$a_t = \pi(M_t, G_t)$$

This is the single point of coupling between epistemic and purposeful dynamics. The update functions are separated; the action is not.


\medskip\noindent\textbf{Epistemic Status.}


\textit{Formulation.} The lift from $M_t$ to $X_t = (M_t, G_t)$ is a representational choice. One could alternatively extend $M_t$ to carry purposeful content implicitly (e.g., by treating goals as part of the model's predictive structure). The separation is motivated by three properties:

\begin{enumerate}
  \item \textbf{Backward compatibility}: Section I's results apply to $M_t$ unchanged — no existing machinery needs modification
  \item \textbf{Different dynamics}: epistemic and purposeful components have distinct update sources, timescales, and information dependencies
  \item \textbf{Directed separation}: the claim that $f_M$ is $G_t$-independent ( \S\,\ref{directed-separation}) is only statable when the components are separated
\end{enumerate}


We conjecture that any alternative decomposition of the complete agent state into epistemic and purposeful components — if it preserves the directed separation — will be structurally isomorphic to $(M_t, G_t)$, because directed separation forces a component whose update function doesn't reference purpose. This is a plausible structural claim (the directed separation gives a canonical factorization), but it is not proved. Alternative decompositions that cross-cut the epistemic/purposeful boundary (e.g., a "relevance-weighted model" that mixes $M_t$ and $O_t$) might be useful in practice while not being isomorphic to the clean separation. The $(M_t, G_t)$ decomposition is chosen for its analytical utility, not proved unique.


\medskip\noindent\textbf{Discussion.}


\textbf{Backward compatibility with Section I — what survives the lift.} \S\,\ref{agent-model} defines $M_t$ as the agent's complete internal state. Under the lift, $M_t$ is the epistemic substate — complete within the epistemic domain but no longer the whole story. All epistemic machinery (mismatch signal, gain, tempo, persistence condition, sector-condition stability, mismatch decomposition) applies to $M_t$ without modification. However, \S\,\ref{action-selection} derives $a_t = \pi(M_t)$ from the premise that $M_t$ is the agent's \textit{complete} state — this derivation is superseded by $a_t = \pi(M_t, G_t)$ after the lift. For Section I agents ($G_t = \emptyset$), the original result holds trivially. The lift adds structure \textit{alongside} $M_t$, not within it; the one result that relied on $M_t$ being \textit{all there is} (action selection) is explicitly extended.

**What $G_t$ contains.** At this level, $G_t$ is opaque — it could be a scalar setpoint, a utility function, a strategy graph, or nothing. The decomposition into $O_t$ (objective) and $\Sigma_t$ (strategy) is a separate step ( \S\,\ref{strategy-dimension}), not implied by this formulation.

\textbf{The general case requires coupling.} Without directed separation, the general update is $X_{\tau^+} = f_X(X_{\tau^-}, e_\tau)$ — a single function on the full state. The decomposition into separate $f_M$ and $f_G$ is an additional structural claim about how the update factorizes. When directed separation fails (goal-conditioned epistemic updates), the decomposition is an approximation. See \S\,\ref{directed-separation} for the scope conditions.

\end{formulation}

\begin{formulation}[Objective Functional]
\label{objective-functional}
The objective $O_t$ is the component of $G_t$ that specifies what the agent wants — the evaluation criterion for trajectories. Its interface to the theory is a single functional $V_{O_t}: \text{trajectories} \to \mathbb{R}$, regardless of how the objective is internally represented.



\medskip\noindent\textbf{Formal Expression.}



The \textbf{objective} $O_t$ induces a \textbf{value functional}:

$$V_{O_t}: \text{trajectories} \to \mathbb{R}$$

$V_{O_t}(\tau)$ is a scalar measure of how well trajectory $\tau$ satisfies the objective. This is the sole interface between $O_t$ and the rest of the theory — the type-stable evaluation surface.

\textbf{Objective representations.} $O_t$ can take multiple internal forms, all unified through $V_{O_t}$:

\begin{adjustbox}{max width=\textwidth}
\small
\begin{tabular}{lll}
\toprule
$O_t$ form & $V_{O_t}(\tau)$ & Example \\
\midrule
Point target $r$ & $-\lVert s_T - r \rVert$ & PID setpoint \\
Target region $R$ & $\mathbb{1}[s_T \in R]$ & "reach safe state" \\
Constraint set & $-\sum_t \max(0, g_i(s_t))$ & "never violate SLA" \\
Utility $U$ & $\sum_t \gamma^t U(s_t)$ & RL reward \\
Trajectory functional $J$ & $J(\tau)$ & "migrate with zero downtime" \\
\bottomrule
\end{tabular}
\end{adjustbox}


The trajectory functional is the most general; the others are special cases.

\textbf{Satisfaction threshold.} Many objectives carry a natural threshold $V_{O_t}^{\min}$ — the minimum trajectory value the agent treats as acceptable. Point targets, constraint sets, and threshold objectives define this directly; utility-maximizing objectives may not. When $V_{O_t}^{\min}$ exists, it enables the satisfaction gap diagnostic ( \S\,\ref{satisfaction-gap}) and the well-formedness constraint on strategy ( \S\,\ref{strategy-dag}). $V_{O_t}^{\min}$ is a parameter of the objective, not a theory output — it encodes "what counts as success" in domain terms.


\medskip\noindent\textbf{Epistemic Status.}


\textit{Axiomatic, with a substantive commitment.} This is a formulation — it names an object and specifies its interface, but the real-valued codomain is a genuine restriction, not a neutral naming. The claim that $V_{O_t}: \text{trajectories} \to \mathbb{R}$ is the right interface is grounded in: any evaluation criterion must ultimately answer "how good is this trajectory?" with a scalar, because the agent must compare alternatives. The real-valued codomain follows from this comparability requirement (total ordering of alternatives).

\textbf{Scope restriction: scalar comparability.} The real-valued codomain is a genuine restriction, and it is load-bearing — the satisfaction gap ( \S\,\ref{satisfaction-gap}) and control regret ( \S\,\ref{control-regret}) require comparing scalar values to produce their diagnostic. Three arguments ground the restriction:

\begin{enumerate}
  \item \textbf{Revealed preference.} An agent that acts has implicitly scalarized: choosing action $a$ over $a'$ imposes a total ordering at the moment of choice. The scalar $V_{O_t}$ makes this implicit scalarization explicit. An agent that truly cannot compare alternatives cannot act coherently — it is stuck, not purposeful.
  \item \textbf{Approximation.} Most multi-objective problems admit scalarization (weighted sum, lexicographic ordering, constraint-satisfaction with scalar residual) that preserves the diagnostic structure. The restriction excludes only agents with genuinely incommensurable objectives \textit{and} no priority structure over them.
  \item \textbf{Timescale separation.} When objectives conflict, the conflict is typically resolved at a slower timescale than strategy revision — the agent (or its principal) chooses weights, then acts within those weights. The scalar $V_{O_t}$ captures the resolved weights at the current timescale.
\end{enumerate}


Agents with hard non-compensatory constraints (safety AND profitability as independent thresholds) can be modeled via the AND-node workaround: each constraint becomes a terminal node in $\Sigma_t$ with its own scalar satisfaction test, and the AND combination enforces joint feasibility. This handles constraint satisfaction cleanly but does not resolve cross-objective tradeoffs within the feasible region.

Organizations or AI agents with true Pareto structure — where no scalarization is accepted and tradeoffs are genuinely unresolved — require a vector-valued extension. The structural results of Section II (orient cascade ordering, strategy DAG, directed separation) survive such an extension; the diagnostic results (satisfaction gap, control regret, 2×2 table) degrade from quantitative scalar magnitudes to qualitative set-theoretic tests. See \texttt{msc/spike-scalar-objective-scope.md} for the full analysis.


\medskip\noindent\textbf{Discussion.}


\textbf{Filling TF-08's gap.} The existing policy objective ( \S\,\ref{ciy-unified-objective}) contains $\mathbb{E}[\text{value}(a) \mid M_t]$ without specifying what "value" means. $O_t$ provides the formal content: value is $V_{O_t}$ applied to expected trajectories. The \S\,\ref{value-object} segment develops the full evaluation machinery ($V_O$, $Q_O$ with horizon and continuation policy).

**$O_t$ evaluates; $\Sigma_t$ guides.** The objective says "is this trajectory satisfactory?" The strategy ( \S\,\ref{strategy-dimension}) says "which action sequence produces a satisfactory trajectory?" A chess player's objective (win) is simple; the strategy (how to win) is complex. These answer different questions and carry different kinds of information — the split is developed in \S\,\ref{strategy-dimension}.

**What $O_t$ is NOT.** $O_t$ does not encode how to achieve the objective (that's $\Sigma_t$), what the agent believes about the world (that's $M_t$), or the agent's commitment or resource state (open questions — see \S\,\ref{strategy-dimension} Working Notes). $O_t$ is purely an evaluation criterion.

\end{formulation}

\begin{definition}[Value Object]
\label{value-object}
The horizon- and policy-conditioned value object $V_O$ turns the abstract objective functional $V_{O_t}$ into a decision-making tool: "given what I believe, what I plan to do next, and how far I'm looking ahead, how good is this situation?"



\medskip\noindent\textbf{Formal Expression.}



Given objective $O_t$, model $M_t$, policy $\pi$, and horizon $N_h$:

$$V_O(M_t, \pi; N_h) = \mathbb{E}\!\left[V_{O_t}(\tau_{t:t+N_h}) \;\middle\vert\; M_t,\; \pi\right]$$

\textbf{Action-value form} (for action selection):

$$Q_O(M_t, a; \pi_{\text{cont}}, N_h) = \mathbb{E}\!\left[V_{O_t}(\tau) \;\middle\vert\; M_t,\; do(a_t = a),\; a_{t+1:} \sim \pi_{\text{cont}}\right]$$

$Q_O$ answers: "if I \textit{do} action $a$ now and then follow $\pi_{\text{cont}}$ afterward, what is my expected trajectory value?" The $do(\cdot)$ notation is explicit: this is an interventional query ( \S\,\ref{causal-hierarchy-requirement}), not conditioning on observed action choice. The agent asks about consequences of an intervention, not about correlates of a naturally occurring action.

\textbf{Causal validity of the value object.} $Q_O$ is well-defined as a conditional expectation given $M_t$, $do(a)$, and $\pi_{\text{cont}}$. Two mechanisms ensure causal validity:

\begin{enumerate}
  \item \textbf{The do-operator handles current-action confounding.} Since $Q_O$ uses $do(a_t = a)$, the dependence of $a_t$ on the selection mechanism $\pi(M_t, G_t)$ is severed. $G_t$'s influence on action choice is irrelevant because the action is intervened upon, not conditioned on.
\end{enumerate}


\begin{enumerate}
  \item \textbf{The continuation policy is a parameter.} $\pi_{\text{cont}}$ is specified as a fixed policy, not as "whatever the agent would do given its evolving $G_t$." Future actions follow $\pi_{\text{cont}}$ regardless of $G_t$'s state or evolution.
\end{enumerate}


Together, these mean $Q_O(M_t, a; \pi_{\text{cont}}, N_h)$ depends on $M_t$ alone \textbf{as a state variable} — $G_t$ enters neither through action selection (severed by $do$) nor through continuation (fixed by parameter). The objective $O_t$ enters as a fixed parameter (it determines which functional $V_{O_t}$ is applied to trajectories), the same way $\pi_{\text{cont}}$ and $N_h$ are parameters. The claim is not that $Q_O$ is independent of the objective — it is that once $O_t$, $\pi_{\text{cont}}$, and $N_h$ are fixed, the only agent state that affects the value is $M_t$. The remaining requirement: $M_t$ must support the interventional query $P(o \mid do(a), M_t)$. Under directed separation ( \S\,\ref{directed-separation}), this holds because $M_t$ updates independently of $G_t$ — there is no path from $G_t$ to outcomes that bypasses both the action channel and $M_t$. For \textbf{Class 2 agents} (where $G_t$ leaks into $M_t$ processing), the causal validity of $Q_O(M_t, a; \pi_{\text{cont}}, N_h)$ is degraded because $M_t$ itself carries goal-conditioned bias — see \texttt{msc/spike-coupled-survival-analysis.md} §3.4.

This is a \textit{stronger requirement than predictive sufficiency} (\S\,\ref{model-sufficiency}). $S(M_t) = 1$ means the model retains all predictive information from the chronica — but predictive sufficiency is a Level 1 (associational) property. Causal validity additionally requires that no unmodeled common cause affects both the environment and the agent's epistemic processing through paths not captured in $M_t$. In practice: when $S(M_t) = 1$ and directed separation holds, $Q_O$ is causally valid. When $S(M_t) < 1$ or directed separation fails, the interventional interpretation is correct but the conditional estimate may be biased.

\textbf{Continuation convention.} All value queries are conditioned on a specific continuation policy $\pi_{\text{cont}}$ and finite horizon $N_h$. $\pi_{\text{cont}}$ is a \textit{parameter} of the value object, not a derived quantity.

\textbf{Canonical default: one-step improvement.} ACT adopts $\pi_{\text{cont}} = \pi_{\text{current}}$ as the canonical continuation convention unless otherwise specified. Under this convention, each action is evaluated assuming current behavior continues afterward — no fixed-point computation, no global optimality assumption. This aligns with ACT's incremental update philosophy ( \S\,\ref{update-gain}) and makes all ACT diagnostics ($\delta_{\text{sat}}$, $\delta_{\text{regret}}$, $A_O$) comparable across analyses of the same agent over time. It is not a convergence guarantee; it is a shared evaluation frame.

\textbf{Alternative conventions} are legitimate but must be stated explicitly when used:

\begin{itemize}
  \item $\pi_{\text{cont}} = \pi^*$ — \textbf{Bellman fixed point} (self-consistent optimal continuation). Requires solving a fixed-point equation. Standard in RL and dynamic programming.
  \item $\pi_{\text{cont}}$ re-optimized each step — \textbf{receding horizon / MPC}. Re-plans at each step with updated $M_t$.
\end{itemize}


Different continuation conventions yield different values for $A_O$, $\delta_{\text{sat}}$, and $\delta_{\text{regret}}$. Diagnostics computed under different conventions are not directly comparable — the convention is part of the measurement, not just the computation. When a specific convergence guarantee is needed (e.g., for \S\,\ref{strategy-persistence-schema}), the solution concept must be stated explicitly; the one-step improvement default does not provide convergence guarantees.


\medskip\noindent\textbf{Epistemic Status.}


\textit{Exact} under the assumption that $M_t$ supports the required conditional expectations. The value object is a mathematical definition — conditional expectations of a functional over trajectories. The definitions are precise; the \textit{computability} of these expectations is a separate question (in practice, they are approximated via simulation, sampling, or function approximation).


\medskip\noindent\textbf{Discussion.}


\textbf{Extending the policy objective.} The existing policy objective ( \S\,\ref{ciy-unified-objective}) uses $\mathbb{E}[\text{value}(a) \mid M_t]$ without formal content for "value." With the value object, this becomes:

% [Discussion (policy-objective-extension)]

$$\pi^*(M_t, G_t) = \arg\max_a \left[Q_O(M_t, a;\, \pi_{\text{cont}}, N_h) + \lambda(M_t, O_t, N_h) \cdot \text{CIY}(a;\, M_t)\right]$$

Note that $\lambda$ now depends on $(M_t, O_t, N_h)$, not just $M_t$. The value of exploration depends on the objective and the horizon:
\begin{itemize}
  \item An agent with a deadline should explore less as time runs out
  \item An agent with a safety constraint should explore differently from a utility maximizer
  \item Two agents with identical $M_t$ but different objectives should price exploration differently
\end{itemize}


This extension is structurally motivated but the specific form of $\lambda(M_t, O_t, N_h)$ is not derived within ACT (same status as \S\,\ref{ciy-unified-objective}'s treatment of $\lambda$).

\textbf{Connection to \S\,\ref{model-sufficiency}.} $V_O$ is conditioned on $M_t$, not on the true environment state $\Omega_t$. When $S(M_t) < 1$, the agent's value estimates are biased — it may over- or underestimate trajectory values because its model is incomplete. The satisfaction gap ( \S\,\ref{satisfaction-gap}) and control regret ( \S\,\ref{control-regret}) are defined in terms of $V_O(M_t, \cdot)$, not $V_O(\Omega_t, \cdot)$, which means they measure the agent's \textit{believed} situation, not the true one. Improving $M_t$ (reducing $\delta_{\text{epistemic}}$) brings the agent's value estimates closer to reality.

\textbf{Horizon dependence.} $N_h$ is not merely a computational convenience — it reflects genuine uncertainty about the far future. Long horizons amplify the impact of model error (small biases in $M_t$ compound over many steps). The choice of $N_h$ trades off farsightedness against robustness to model error. An agent in a fast-changing environment ($\rho$ high) should use shorter horizons; one in a stable environment can plan further.

\end{definition}

\begin{definition}[Strategy Dimension]
\label{strategy-dimension}
The purposeful substate $G_t$ decomposes into two structurally distinct components: $O_t$ (the objective — what the agent wants) and $\Sigma_t$ (the strategy — the agent's theory of how its actions produce progress toward $O_t$). These carry different kinds of information answering different questions.



\medskip\noindent\textbf{Formal Expression.}



$$G_t = (O_t, \Sigma_t)$$

where:
\begin{itemize}
  \item $O_t$: \textbf{evaluation} — "Is this trajectory satisfactory?" ( \S\,\ref{objective-functional})
  \item $\Sigma_t$: \textbf{guidance} — "Which action sequence produces a satisfactory trajectory?"
\end{itemize}


The split is \textbf{definitional} — it reflects a structural difference in the information, not a dynamic or timescale claim. $O_t$ and $\Sigma_t$ are different \textit{kinds} of state answering different questions:

\begin{adjustbox}{max width=\textwidth}
\small
\begin{tabular}{lll}
\toprule
 & $O_t$ (objective) & $\Sigma_t$ (strategy) \\
\midrule
\textbf{Question} & How good is this trajectory? & How do I produce a good trajectory? \\
\textbf{Type} & $V_{O_t}: \text{trajectories} \to \mathbb{R}$ & Structured representation (see below) \\
\textbf{Richness varies} & Point target → utility → trajectory functional & Reactive → cached → subgoals → causal DAG \\
\textbf{Update source} & External (assigned, discovered, revised) & Internal (deliberation, evidence, cascade) \\
\bottomrule
\end{tabular}
\end{adjustbox}


\textbf{Strategy representations}, ordered by expressiveness:

\begin{adjustbox}{max width=\textwidth}
\small
\begin{tabular}{lll}
\toprule
$\Sigma_t$ form & What it encodes & Example \\
\midrule
None (reactive) & No explicit strategy; policy implicit in $M_t$ & Thermostat, reflex \\
Cached policy & Learned mapping $s \to a$ & Trained RL policy \\
Subgoal sequence & Waypoints with ordering & Navigation, recipe \\
Causal DAG & Action-outcome chains with AND/OR structure and confidence weights & Military plan, software project \\
\bottomrule
\end{tabular}
\end{adjustbox}



\medskip\noindent\textbf{Epistemic Status.}


\textit{Axiomatic.} This is a definition — it names a structural distinction that exists in the information. The distinction between "what makes a trajectory good" (evaluation) and "how to produce a good trajectory" (guidance) is a categorical difference, not a quantitative one. The claim is NOT that all agents maintain both explicitly — reactive agents have $\Sigma_t = \emptyset$, and that's fine. The claim is that when an agent does maintain purposeful state, it decomposes along this line.

The two dimensions vary independently: a chess player has a simple $O_t$ (win) and a complex $\Sigma_t$ (opening theory, tactical patterns, endgame knowledge). A multi-objective optimizer may have a complex $O_t$ (Pareto frontier) and a simple $\Sigma_t$ (gradient descent). This independence is why the split matters — conflating them in a single hierarchy obscures the fact that objective richness and strategic richness are separate design axes.


\medskip\noindent\textbf{Discussion.}


**When richer $\Sigma_t$ is needed.** A reactive agent ($\Sigma_t = \emptyset$) suffices when greedy optimization on $Q_O$ ( \S\,\ref{value-object}) works — when the action-to-value mapping is approximately convex and single-step. When the environment has non-convex landscapes, prerequisite structure, or multi-step causal chains, greedy optimization fails and the agent needs explicit strategy. The trigger is the purposeful analog of \S\,\ref{structural-adaptation-necessity}: inadequacy of the current $\Sigma_t$ representation for the environment's causal complexity.

**$O_t$ and $\Sigma_t$ have different update dynamics.** Objectives change slowly: an organization's mission, a developer's feature goal, a commander's campaign objective. Strategies change faster: adjust the plan when step 3 fails, redirect resources, try an alternative path. Epistemic state changes fastest: each observation updates $M_t$. This timescale ordering ($\nu_M \gg \nu_\Sigma \gg \nu_O$) is an empirical observation, not a derived result. It holds for many agent populations but is not universal — an agent discovering its goal is infeasible may revise $O_t$ faster than $\Sigma_t$.

\textbf{The decomposition resolves a type error.} Earlier formulations used $\delta_{\text{goal}} = G_t - M_t$ as a goal mismatch signal. When $\Sigma_t$ is a DAG, this is a type error — you cannot subtract a graph from a state vector. The \S\,\ref{satisfaction-gap} and \S\,\ref{control-regret} replace this with properly typed gap measures.

\end{definition}

\begin{proposition}[Causal Hierarchy Requirement]
\label{causal-hierarchy-requirement}
Evaluating the action-value $Q_O$ requires answering "what happens if I \textit{do} action $a$?" — a Level 2 (interventional) query in Pearl's causal hierarchy. An agent that must learn the answer to this question during operation needs access to causal structure beyond what purely predictive models can provide.



\medskip\noindent\textbf{Formal Expression.}


% [Derived (causal-hierarchy-requirement, from value-object + pearl-causal-hierarchy)]

Action selection via \S\,\ref{value-object} requires:

$$Q_O(M_t, a;\, \pi_{\text{cont}}, N_h) = \mathbb{E}\!\left[V_{O_t}(\tau) \;\middle\vert\; M_t,\; do(a_t = a),\; a_{t+1:} \sim \pi_{\text{cont}}\right]$$

The $do(\cdot)$ notation is explicit: this is an \textit{intervention}, not a \textit{conditioning on observed data}. By \S\,\ref{pearl-causal-hierarchy}, Level 2 queries ($P(Y \mid do(X))$) cannot in general be computed from Level 1 data ($P(Y \mid X)$) alone. Therefore:

An agent that must evaluate $Q_O$ from experience needs access to Level 2 knowledge — knowledge about the effects of its own interventions, not merely correlational patterns.

% [Scope Narrowing (learning-agents)]

We restrict attention to agents that must \textbf{acquire or refine} Level 2 knowledge during operation. This excludes agents with \textbf{pre-compiled} interventional structure:
\begin{itemize}
  \item PID controllers (the designer pre-computed the control law)
  \item LQR (separation principle gives optimal policy from model parameters)
  \item Hardcoded reactive policies
\end{itemize}


Pre-compiled agents are purposeful — they have objectives and act on them — but their causal structure was externally supplied by a designer who had Level 2 access. This scope narrowing focuses the theory on agents that must build or maintain their own causal understanding.


\medskip\noindent\textbf{Epistemic Status.}


\textit{Exact.} The derivation is a direct application of the causal hierarchy theorem (Bareinboim et al. 2022) to the value-object definition. If you accept that $Q_O$ is an interventional query and that the causal hierarchy is strict, the conclusion follows. The scope narrowing to learning agents is a definitional restriction, not a derived result — it sharpens the class of agents under study.


\medskip\noindent\textbf{Discussion.}


\textbf{The causal hierarchy theorem does the heavy lifting.} The key mathematical fact is external to ACT: Bareinboim et al. (2022) prove that the three levels (association, intervention, counterfactual) form a strict hierarchy — Level 2 quantities cannot in general be computed from Level 1 data. ACT's contribution is applying this to the purposeful-agent setting: if you want to select actions by their consequences ($Q_O$), you need causal structure.

\textbf{What "Level 2 knowledge" means concretely.} For different agents:
\begin{itemize}
  \item A developer needs to know "if I refactor this module, will tests still pass?" (not just "modules that were refactored tend to have tests that fail")
  \item A commander needs to know "if I move forces north, will the enemy respond by retreating?" (not just "when forces moved north, the enemy often retreated")
  \item An RL agent needs $Q(s, a) = \mathbb{E}[R \mid s, do(a)]$, not $\mathbb{E}[R \mid s, A=a]$ (the latter includes selection bias from the agent's own policy)
\end{itemize}


The distinction between $do(a)$ and $A = a$ is the core of causal inference. It matters whenever the agent's action-selection policy correlates with unobserved confounders.

\textbf{Why pre-compiled agents are excluded.} A thermostat "knows" that turning on the heater raises temperature — but this knowledge was designed in, not learned. The thermostat never needs to reason about interventions because the intervention-outcome mapping is hardwired. ACT's purposeful-agency machinery is specifically for agents that face uncertainty about how their actions affect the world and must reduce that uncertainty through experience.

\end{proposition}

\begin{proposition}[Loop Provides Interventional Data Access]
\label{loop-interventional-access}
An agent in the feedback loop generates interventional data by construction: the agent's action $a_t$ causally precedes the next observation $o_{t+1}$, and the mismatch conditioned on $a_t$ carries interventional information. This is how agents within ACT's scope gain Level 2 access — not through internal architecture, but through the loop itself.



\medskip\noindent\textbf{Formal Expression.}


% [Derived (loop-interventional-access, from causal-structure + recursive-update)]

By \S\,\ref{causal-structure}, the temporal ordering is constitutive: $a_t$ causally precedes $o_{t+1}$. The agent chose $a_t$; the environment responded with $o_{t+1}$. This is an intervention — the agent varied its action and observed the result.

Formally: the pair $(a_t, o_{t+1})$ is generated under an intervention policy — the agent executed $a_t$, making it a genuine intervention rather than a passively observed association. The data \textit{contains interventional signal}: it was produced by a do-operation, not by conditioning on a naturally occurring value of $A_t$.

However, this does not mean the agent automatically has access to a clean estimate of $P(o \mid do(a_t), \Omega_t)$. Between the interventional act and a usable interventional distribution stand: (1) coverage — the agent must have tried diverse actions, not just one policy; (2) confounding within a time step — unobserved state variables that affect both action choice and outcome; (3) delay — consequences may appear much later than $t+1$; (4) partial observability — $o_{t+1}$ reveals only part of the outcome. The claim is about the \textit{character} of the data (interventional, not observational), not about the agent's ability to extract clean causal estimates from it.

The mismatch signal conditioned on the agent's action:

$$\delta_t \mid a_t = o_{t+1} - \hat{o}_{t+1}(M_t, a_t)$$

carries interventional information: it tells the agent how the environment responded to its specific intervention $a_t$, relative to what the model predicted.


\medskip\noindent\textbf{Epistemic Status.}


\textit{Exact.} This is a logical consequence of temporal ordering ( \S\,\ref{causal-structure}) and the feedback-loop structure ( \S\,\ref{scope-condition}). The claim is about \textbf{data availability}, not reasoning capacity — the loop \textit{provides} interventional data whether or not the agent \textit{exploits} it for Level 2 reasoning. Whether the agent uses this data to build causal models depends on its update mechanism and model class.

The precision is important: we claim the agent has \textit{access to} interventional data, not that it \textit{correctly identifies} interventional structure. Confounding within a single time step, delayed outcomes, and partial observability can all complicate the extraction of clean causal signals from the loop data.


\medskip\noindent\textbf{Discussion.}


\textbf{The loop as a Level 2 engine.} This is one of ACT's load-bearing results. The causal hierarchy theorem ( \S\,\ref{causal-hierarchy-requirement}) says Level 2 knowledge requires more than correlational data. This result says: the adaptive loop \textit{is} the "more." An agent that acts and observes the consequences is generating interventional data — the same kind of data that a scientist generates through experiments. The loop is a perpetual experiment.

\textbf{Precision about what "interventional" means here.} The interventional interpretation is strongest when:
\begin{itemize}
  \item The agent's action was the primary cause of the observed change (low confounding)
  \item The observation follows closely in time (short delay)
  \item The agent varied its action across episodes (not stuck on one policy)
\end{itemize}


When confounding is high, delays are long, or the agent follows a fixed policy, the interventional information in each $(a_t, o_{t+1})$ pair is weaker — still present, but harder to extract. This is why \S\,\ref{causal-information-yield} distinguishes between high-CIY actions (that reveal causal structure) and low-CIY actions (that don't).

\textbf{Even agents without explicit causal models benefit.} A Q-learning agent doesn't maintain an explicit causal model, but in the tabular case with sufficient exploration and no within-step confounding, its Q-values converge toward $\mathbb{E}[R \mid s, do(a)]$ rather than $\mathbb{E}[R \mid s, A=a]$ — precisely because the training data comes from the agent's own interventions. In the partially observed, confounded, or delayed-outcome cases (where the caveats above apply), the loop still provides intervention-generated data, but the Q-values may converge to biased estimates that reflect the confounding structure rather than clean interventional effects. The loop provides Level 2 data; whether that data yields \textit{identified} causal quantities depends on the domain's confounding and observability structure.

\end{proposition}

\begin{scope}[CIY Observational Proxy]
\label{ciy-observational-proxy}
When and how causal information yield can be approximated from observational data rather than interventional experiments.



\medskip\noindent\textbf{Formal Expression.}



$$\text{CIY}_{\text{proxy}}(a_{t-1}) = I(o_t; a_{t-1} \mid M_{t-1}) - I(o_t; a_{t-1} \mid \Omega_t, M_{t-1})$$

This proxy is \textbf{sign-indefinite in general} and requires causal assumptions for interpretation. The canonical CIY (interventional, \S\,\ref{causal-information-yield}) is the primary quantity; the proxy is auxiliary.

\textbf{Safety conditions for proxy use.} The proxy form should NOT be used in policy optimization (e.g., as the CIY term in a policy objective) because an agent maximizing a sign-indefinite quantity may optimize in the wrong direction. The proxy is suitable only for diagnostic purposes: detecting whether an action carried causal information (large proxy magnitude) vs. none (proxy near zero). For decision-making, use the canonical CIY (non-negative by construction) or a known-safe surrogate (ensemble disagreement, UCB bonuses). If the canonical CIY is intractable and no safe surrogate is available, the CIY term should be dropped from the policy objective entirely, defaulting to pure exploitation.


\medskip\noindent\textbf{Admissibility regimes.}


% [Scope Condition (ciy-admissibility)]

Three regimes determine when CIY can be estimated and how strong the causal identification is:

\textbf{Regime A — Randomized interventions.} The agent varies its actions across episodes (RL agents exploring, scientists experimenting, organisms probing). CIY is directly estimable from the agent's execution data and non-negative by construction. This is the standard case for active agents within the adaptive loop ( \S\,\ref{loop-interventional-access}). Action variation provides the identification needed for clean interventional estimates.

\textbf{Regime B — Observational with causal assumptions.} The agent cannot freely vary actions (constrained by coordination, policy, or resource limits). CIY estimation requires additional structure: a known causal DAG, instrumental variables, or functional form assumptions. Results inherit whatever causal assumptions are made. The interventional interpretation of CIY is weaker — it holds under the assumed causal structure but not model-free.

\textbf{Regime C — Adversarial or passive observation.} The agent either did not intervene (passive monitoring) or the observation channel includes responses from potentially adversarial sources. In the passive case, CIY is zero by definition (no intervention, no interventional information). In the adversarial case, CIY from the query action itself remains non-negative, but the \textit{content} of the response may be designed to increase model-reality mismatch. The adversary operates through the disturbance term $\rho$, not through the information measure.

The regime is a property of the \textbf{domain and the agent's action space}, not a parameter the agent chooses. Software development is typically Regime A (the agent runs tests, deploys to staging, observes results — high action variation). Organizational strategy is typically Regime B (multiple initiatives run concurrently, attribution requires assumptions). Intelligence analysis is typically Regime C (the analyst observes but does not intervene).


\medskip\noindent\textbf{Epistemic Status.}


Conditional on the causal assumptions within each regime. The proxy definition is standard information theory; the admissibility classification is a scope decision, not a derived result. The safety conditions for proxy use are normative — they follow from the sign-indefiniteness of the proxy, not from ACT-specific machinery.

Max attainable: conditional. The regime boundaries are domain properties that cannot be derived within the theory.


\medskip\noindent\textbf{Discussion.}


\textbf{Relationship to the canonical CIY.} The proxy is not an approximation of the canonical CIY — it is a different quantity that happens to correlate with CIY under favorable conditions (Regime A). The canonical CIY ( \S\,\ref{causal-information-yield}) is defined interventionally and is non-negative by construction. The proxy uses observational mutual information and can be negative. They agree when the agent's action variation satisfies the conditions for causal identification; they diverge otherwise.

\textbf{Regime as a domain property.} The admissibility regime is not something the agent selects — it is determined by the domain's action space and observation structure. An agent that cannot vary its actions is in Regime B or C regardless of its internal architecture. This has implications for which domains CIY-based exploration is applicable to: Regime A domains get the full benefit; Regime B domains get weaker guarantees; Regime C domains should use alternative exploration strategies entirely.

\textbf{(Descended from TF-08.)}

\end{scope}

\begin{discussion}[CIY Unified Policy Objective]
\label{ciy-unified-objective}
The exploration-exploitation tension can be expressed as a single policy objective that jointly maximizes expected value and a causal information surrogate. This formulation is \textit{heuristic} — CIY measures action-distinguishability, not expected information gain (see \S\,\ref{causal-information-yield}), so the objective selects for causally distinctive actions rather than maximally informative ones. The $\lambda$-weighting partially compensates by suppressing the CIY term when model uncertainty is low, but the surrogate nature is inherent.



\medskip\noindent\textbf{Formal Expression.}


% [Discussion (unified-policy-objective — heuristic)]

$$\pi^*(M_t) = \arg\max_a \left[\mathbb{E}[\text{value}(a) \mid M_t] + \lambda(M_t) \cdot \text{CIY}_q(a;\, M_t)\right]$$

The first term is exploitation (expected value given current model). The second is a \textit{heuristic exploration term} using CIY as a surrogate for expected information gain ( \S\,\ref{causal-information-yield}). CIY measures how different the action's outcome distribution is from alternatives — this is action-distinguishability, not learning value. The surrogate is reasonable when $U_M$ is high (distinguishable actions are also informative to an uncertain agent) and poor when $U_M$ is low (distinguishable actions teach nothing to a confident agent). $\lambda(M_t)$ controls the balance:

\begin{itemize}
  \item High $U_M$ (uncertain model) → large $\lambda$ — exploration is valuable
  \item Low $U_M$ (confident model) → small $\lambda$ — exploitation dominates
  \item Long time horizon → larger $\lambda$ — information compounds
  \item High $\rho$ (fast-changing environment) → larger $\lambda$ — perpetual uncertainty
\end{itemize}


$\lambda$ carries units of [value per unit information]. In specific domains it reduces to known quantities:

\begin{adjustbox}{max width=\textwidth}
\small
\begin{tabular}{lll}
\toprule
Domain & $\lambda$ reduces to & Status \\
\midrule
Bayesian bandits & Gittins index & Exactly derived \\
Kalman dual control & Probing cost in quadratic objective & Exactly derived \\
Active inference & Precision on epistemic affordance & Framework-derived \\
Information-directed sampling & $(\text{VoI})^2 / \text{info gain}$ & Exactly derived (Russo \& Van Roy) \\
RL with UCB & Confidence-bound scaling & Heuristic (tuned) \\
\bottomrule
\end{tabular}
\end{adjustbox}


\textbf{Identifiability gate.} Before incorporating CIY into the policy objective: (1) action variation must exist, (2) the admissibility regime must be identified ( \S\,\ref{ciy-observational-proxy}), (3) the reference distribution $q$ must be specified, (4) local stationarity must hold. If any condition fails, CIY-based terms should be dropped or replaced with simpler uncertainty-based heuristics (UCB-style bonuses, ensemble disagreement).


\medskip\noindent\textbf{Epistemic Status.}


Discussion-grade (heuristic). The structural claim — that a useful policy jointly considers value and causal information — is supported by convergent results in Bayesian RL, active inference, and information-directed sampling, but not derived from first principles within ACT. The use of CIY rather than proper expected information gain (EIG) makes this a \textit{surrogate} formulation: the objective selects for causally distinctive actions, which approximately coincides with selecting for informative actions when model uncertainty is high but diverges when it is low. The $\lambda(M_t)$ weighting partially compensates (suppressing CIY when $U_M$ is low) but the compensation is heuristic, not derived.

Max attainable: \textit{heuristic} unless CIY is replaced by proper EIG. The structural form (value + information term) is well-supported, but the specific information term (CIY rather than EIG) is a tractability-motivated surrogate that selects for distinguishability rather than informativeness. The $\lambda$ parameterization is domain-dependent and the surrogate nature of CIY places a ceiling below robust-qualitative.


\medskip\noindent\textbf{Discussion.}


\textbf{Connection to the zero-mismatch ambiguity.} An agent that only exploits (acts to maximize predicted value) will tend toward confirmation bias — observing only what its model already explains ( \S\,\ref{mismatch-signal}, zero-mismatch ambiguity case (b)). Exploration via CIY-maximizing actions is the mechanism by which the agent actively tests its model.

\textbf{Connection to Section II.} For actuated agents, the exploration-exploitation tension extends to three modes: exploit (pursue $O_t$ via $\Sigma_t$), explore (improve $M_t$), deliberate (revise $\Sigma_t$). The CIY framework provides the information-theoretic grounding for why strategy edges ( \S\,\ref{strategy-dag}) need observational access ( \S\,\ref{observability-dominance}) — edges the agent cannot observe have frozen CIY.

\textbf{Connection to active inference.} The Free Energy Principle's "expected free energy" decomposes into extrinsic value (pragmatic) and epistemic value (information-seeking). ACT's formulation is structurally isomorphic: expected value ≈ extrinsic, CIY ≈ epistemic. Whether this convergence is deep or superficial is an open question. ACT grounds exploration in explicitly \textit{causal} information rather than entropy reduction — not all uncertainty reduction is equally valuable; causal information specifically enables better \textit{intervention}.

\textbf{(Descended from TF-08.)}

\end{discussion}

\begin{normative}[Explicit Strategy Condition]
\label{explicit-strategy-condition}
An agent benefits from maintaining an explicit strategy $\Sigma_t$ when the cost of planning is less than the cost of learning through exploration alone. This is a normative design criterion — it tells you when explicit strategy is \textit{worth it}, not that it's always necessary.



\medskip\noindent\textbf{Formal Expression.}


% [Normative (explicit-strategy-condition, via temporal-optimality)]

An agent benefits from explicit $\Sigma_t$ when:

$$C_{\text{plan}} + C_{\text{maintain}} < C_{\text{explore}} + C_{\text{repair}}$$

where:
\begin{itemize}
  \item $C_{\text{plan}}$: cost of constructing and evaluating the strategy (deliberation, simulation, model queries)
  \item $C_{\text{maintain}}$: ongoing cost of keeping $\Sigma_t$ current as $M_t$ evolves (edge revision, structural updates)
  \item $C_{\text{explore}}$: cost of learning action-outcome mappings through direct interaction (real actions, real time, real consequences)
  \item $C_{\text{repair}}$: cost of correcting errors discovered only through execution (rollbacks, rework, damage)
\end{itemize}


All costs are measured in time ( \S\,\ref{temporal-optimality}) under the precondition that the two approaches produce approximately equivalent non-temporal outcomes.


\medskip\noindent\textbf{Epistemic Status.}


\textit{Normative, not derived.} This is labeled \textit{normative} because \S\,\ref{temporal-optimality} requires identical non-temporal outcomes as a precondition — "given equivalent outcomes, prefer the faster approach." In practice, loop-based and model-based approaches may differ in:

\begin{itemize}
  \item \textbf{Final value}: the model introduces bias; exploration may discover things planning cannot
  \item \textbf{Risk profile}: exploration risks real damage; planning risks wrong models
  \item \textbf{Reversibility}: some exploratory actions are irreversible
  \item \textbf{Model bias}: explicit $\Sigma_t$ inherits the biases of $M_t$; loop-based learning does not
\end{itemize}


The inequality is correct \textit{when} the outcomes are approximately equivalent — a condition that must be verified case by case, not assumed. When the precondition fails (model-based and loop-based approaches produce qualitatively different outcomes), the cost inequality is insufficient and the choice requires richer analysis (e.g., expected regret including model error).


\medskip\noindent\textbf{Discussion.}


\textbf{When the inequality strongly favors planning.} The right side ($C_{\text{explore}} + C_{\text{repair}}$) is large when:
\begin{itemize}
  \item Actions are expensive or irreversible (production deployments, military operations, surgical procedures)
  \item Exploration damage is severe (a wrong move in chess loses the game; a wrong deployment takes down the service)
  \item The environment is slow to respond (waiting for market feedback, waiting for test results)
  \item The action space is enormous (combinatorial planning problems)
\end{itemize}


In these domains, explicit $\Sigma_t$ is strongly motivated even if the planning model is imperfect.

\textbf{When the inequality favors pure exploration.} The left side ($C_{\text{plan}} + C_{\text{maintain}}$) is large when:
\begin{itemize}
  \item The environment is too complex or novel for useful models (genuinely unknown territory)
  \item $M_t$ is severely inadequate (model predictions are worse than random)
  \item The environment changes faster than $\Sigma_t$ can be maintained ($\rho_\Sigma$ exceeds planning capacity)
  \item Actions are cheap and reversible (A/B testing, sandbox exploration)
\end{itemize}


\textbf{This makes \S\,\ref{temporal-optimality} load-bearing.} The postulate provides the normative grounding: among approaches producing equivalent outcomes, prefer the one requiring less time. The cost inequality instantiates this for the planning-vs-exploration choice. Without \S\,\ref{temporal-optimality}, the inequality would be an engineering heuristic without theoretical grounding.

\textbf{Connection to the three-way tradeoff.} For actuated agents, the binary explore/exploit tradeoff extends to three modes: exploit (pursue $O_t$ via $\Sigma_t$), explore (improve $M_t$), and deliberate (revise $\Sigma_t$). The cost inequality addresses the coarsest question (is explicit $\Sigma_t$ worth having?). The finer allocation between the three modes is addressed in \S\,\ref{exploit-explore-deliberate} — the extended deliberation threshold is derived, but the broader three-way framing is discussion-grade. The deepest insight: deliberation is computation on existing data (no new observations), making it a fundamentally different resource type from external action.

\end{normative}

\begin{proposition}[Chain Confidence Decay]
\label{chain-confidence-decay}
Confidence in a multi-step strategy decays monotonically with depth. The rate depends on the conditional dependence structure, but the qualitative result — longer chains are less confident than shorter ones — is robust.



\medskip\noindent\textbf{Formal Expression.}


% [Derived (chain-confidence-decay, mathematical identity)]

For a chain of $n$ uncertain steps with conditional success probabilities:

$$\log P(\text{chain}) = \sum_{i=1}^{n} \log P(E_i \mid E_{< i})$$

Since each $\log P(E_i \mid E_{< i}) \leq 0$, chain confidence decays monotonically with depth.

\textbf{The independent case} ($p^n$) is the simplest special case, not the general result. When steps are conditionally dependent — success at step $k$ makes step $k+1$ more likely — the decay is slower. When steps have negative dependence (success at $k$ makes $k+1$ harder — resource depletion, adversary adaptation), decay is faster.

\textbf{Quantitative illustration} (independent, uniform $p$):

\begin{adjustbox}{max width=\textwidth}
\small
\begin{tabular}{lll}
\toprule
Depth & $p = 0.9$ & $p = 0.8$ \\
\midrule
1 & 0.90 & 0.80 \\
3 & 0.73 & 0.51 \\
5 & 0.59 & 0.33 \\
10 & 0.35 & 0.11 \\
20 & 0.12 & 0.01 \\
\bottomrule
\end{tabular}
\end{adjustbox}



\medskip\noindent\textbf{Epistemic Status.}


\textit{Exact.} The additive decomposition of log-confidence is a mathematical identity (chain rule of probability). The qualitative consequence (monotonic decay) follows from the non-positivity of log-probabilities. No assumptions beyond the probability axioms.


\medskip\noindent\textbf{Discussion.}


\textbf{Structural pressure on strategies.} Chain confidence decay creates systematic pressure toward:
\begin{itemize}
  \item \textbf{Short plans}: fewer steps means higher aggregate confidence
  \item \textbf{Parallel fallback paths}: OR-branches provide alternative routes when one chain fails
  \item \textbf{High-confidence critical links}: invest in the reliability of steps that appear in every path
  \item \textbf{Early monitoring}: detect chain failure early rather than discovering it at the end
\end{itemize}


These are not prescriptions but consequences — an agent that ignores chain decay will experience more strategy failures, lower effective tempo, and (if the failures are costly) faster reserve depletion.

\textbf{AND-nodes amplify decay.} When multiple parent chains must all succeed (conjunctive combination), their confidences multiply. A node requiring $k$ parents each at depth $d$ with per-edge confidence $p$ has aggregate confidence $p^{k \cdot d}$, not $p^d$. Deep conjunctive strategies are exponentially more fragile than deep disjunctive ones. This asymmetry is formalized in the combination rules ( \S\,\ref{strategy-dag}).

\textbf{Connection to the persistence condition.} Chain decay makes long-horizon strategies inherently fragile, which increases the effective disturbance rate $\rho_\Sigma$ against strategy persistence. An agent pursuing a 20-step plan in a changing environment faces compound uncertainty from both chain decay (internal fragility) and environmental change (external disturbance). The interaction between these — how environmental change compounds through uncertain chains — is not yet formalized.

\textbf{Triple depth penalty.} Chain depth creates three independent penalties. This segment identifies the first: \textbf{confidence decay} — deeper chains have lower aggregate confidence because $\log P(\text{chain})$ accumulates negative terms. The two-edge strategic dynamics analysis (\S\,\ref{strategic-dynamics-derivation}) identifies the second: \textbf{evidence starvation} — downstream edge $k$ in a chain is tested only when all upstream edges succeed, so its effective correction rate is attenuated by $\prod_{j<k}\theta_j$. \S\,\ref{strategy-complexity-cost} identifies the third: \textbf{cognitive cost} — deeper chains have higher description length, consuming more representational capacity. The three penalties compound independently: a deep edge has low confidence (decay), receives few observations (starvation), and costs more to maintain (complexity). The maximum useful chain depth $d^*$ is the minimum over three independent constraints — see \S\,\ref{strategy-complexity-cost} for the formal bound.

\end{proposition}

\begin{scope}[AND/OR Combination Scope]
\label{and-or-scope}
We restrict to environments where the causal combination of strategy steps is approximately conjunctive (AND: all parents required) or disjunctive (OR: any parent sufficient), without strong interaction effects between parents.



\medskip\noindent\textbf{Formal Expression.}


% [Scope Narrowing (and-or-scope)]

Under this restriction, strategy nodes combine parent contributions via:

\textbf{AND-node} (all parents must succeed):

$$P(v \mid \text{parents}) = \prod_{i \in \text{pa}(v)} p_{iv} \cdot P(i)$$

\textbf{OR-node} (any parent sufficient):

$$P(v \mid \text{parents}) = 1 - \prod_{i \in \text{pa}(v)} (1 - p_{iv} \cdot P(i))$$

The combination type $\gamma(v) \in \{\text{AND}, \text{OR}\}$ is assigned per node. The causal question determines assignment: "if I remove one parent, can $v$ still be achieved?" YES → OR. NO → AND.


\medskip\noindent\textbf{Epistemic Status.}


\textit{Robust qualitative.} This scope narrowing converged independently across three formalism attempts (track-a/00, track-a/02, track-a/03). It captures the dominant structure in most planning domains. The excluded case — complementarity, substitutability, interaction effects between parents — requires richer combination rules and is a legitimate divergence point for future work.

The AND/OR restriction with single-parameter edges gives $k$ parameters per node (one per parent edge) instead of $2^k$ for a general conditional probability table. This parsimony is motivated by bounded cognition ( \S\,\ref{information-bottleneck}): agents with limited representational capacity are forced toward low-parameter models.


\medskip\noindent\textbf{Discussion.}


\textbf{Why AND/OR and not alternatives.}

\textit{Why not Noisy-OR universally.} The first formalism attempt used noisy-OR for all nodes. This \textbf{systematically overestimates conjunctive structures}:

\begin{adjustbox}{max width=\textwidth}
\small
\begin{tabular}{llll}
\toprule
Structure & Noisy-OR & AND & Reality \\
\midrule
3 required KRs at $p = 0.95, 0.90, 0.99$ & 0.99995 & 0.846 & ~AND \\
\bottomrule
\end{tabular}
\end{adjustbox}


The noisy-OR model cannot represent "all of these are required." This was the primary motivation for the AND/OR revision.

\textit{Why not WEIGHTED combination.} A clean-slate formalism (track-a/02) introduced $P(v) = \min(1, \sum \alpha_{iv} \cdot p_{iv} \cdot P(i))$ to handle k-of-n thresholds. This reintroduces a two-parameter estimation problem ($\alpha$ weights per edge). If k-of-n semantics are genuinely needed, nested AND/OR structure can represent them: group alternatives into OR-nodes, then AND the groups. This keeps estimation localized to the node taxonomy rather than spreading it across a new per-edge parameter.

\textbf{The parsimony argument.} For binary-outcome nodes, AND and OR form a complete Boolean basis — any Boolean combination can be decomposed into layers of AND/OR (disjunctive/conjunctive normal form). Under bounded cognition, the agent needs the most expressive $O(k)$-parameter representation. AND/OR is the natural candidate. This is a parsimony-motivated hypothesis, not a derived necessity — see \S\,\ref{graph-structure-uniqueness} for the full argument and its limitations.

\textbf{What this scope excludes.} Environments with strong interaction effects: where the value of combining parent contributions is not separable into independent per-parent terms. Examples: synergistic drug interactions (combined effect exceeds sum of individual effects), complementary goods (neither is useful alone), strategic surprise (the combination of actions matters more than any individual action). These require richer parameterizations within the strongly motivated graphical structure ( \S\,\ref{strategy-dag}) — a direction for future work.

\end{scope}

\begin{definition}[Strategy DAG]
\label{strategy-dag}
The strategy $\Sigma_t$ is a directed acyclic graph with probabilistic edges and AND/OR combination semantics. Each edge carries the agent's causal credence that completing the parent step advances the child step. The graph encodes the agent's theory of how its actions produce progress toward its objectives.



\medskip\noindent\textbf{Formal Expression.}



$$\Sigma_t = (V_t, E_t, p_t, \gamma_t)$$

where:
\begin{itemize}
  \item $V_t$: set of \textbf{propositional nodes} — each node represents a condition that could be true or false (including action-success propositions at the leaves)
  \item $E_t \subseteq V_t \times V_t$: directed causal edges
  \item $p_t : E_t \to [0,1]$: \textbf{causal credence} per edge — the agent's confidence that completing the parent advances the child
  \item $\gamma_t : V_t \to \{\text{AND}, \text{OR}\}$: combination rule per node ( \S\,\ref{and-or-scope})
\end{itemize}


\textbf{Structural constraints:}

\begin{enumerate}
  \item \textbf{Acyclicity.} $\Sigma_t$ is a DAG. This is \textit{derived}, not assumed — see below.
  \item \textbf{Rootedness.} Every node has a directed path to a unique root terminal $v_\text{root}$ — the single sink node (out-degree 0) of $\Sigma_t$. Compound objectives express their combination structure through the AND/OR machinery below $v_\text{root}$, consistent with scalar $V_{O_t}$ ( \S\,\ref{objective-functional}).
  \item \textbf{Source constraint.} Leaf nodes are propositions about action success ("action $a$ succeeds at $\tau_v$") or observable conditions ("condition $C_v$ holds at $\tau_v$"). Both are propositional — the distinction is whether the proposition is within the agent's causal control (action) or not (condition).
\end{enumerate}


\textbf{Leaf base credence.} For each leaf node $v \in V_t^{\text{leaf}}$, the base credence used in status propagation:

$$p_v(M_t) = \begin{cases} \Pr(\text{action } v \text{ succeeds at } \tau_v \mid M_t) & \text{if } v \text{ is an action node} \\[4pt] \Pr(C_v(\tau_v) \mid M_t) & \text{if } v \text{ is a condition node} \end{cases}$$

where $C_v$ is the propositional condition associated with node $v$ and $\tau_v$ is the node's temporal position (from the acyclicity structure). For action leaves, $p_v$ is \textit{capability credence} — "can I execute this?" For condition leaves, $p_v$ is \textit{state credence} — "will this hold?" Both are conditional on $M_t$ and update whenever $M_t$ updates. This is the mechanism by which Section I's adaptive machinery enters the strategy: $M_t$ changes → leaf credences change → status propagation produces new terminal credences.

\textbf{Edge semantics.} Each edge carries a single credence weight:

$$p_{ij} = \text{Cr}(j \text{ advances} \mid i \text{ completed},\, M_t)$$

This is the agent's credence that completing step $i$ advances step $j$, given its current model --- its \textbf{causal efficacy estimate} for the link. The agent treats $p_{ij}$ as a causal quantity for planning purposes (status propagation, plan-confidence scoring, action selection). Whether $p_{ij}$ is a \textit{good} estimate of causal efficacy depends on the identification regime of the data that produced it ( \S\,\ref{edge-update-causal-validity}):

\begin{itemize}
  \item \textbf{Regime A} (intervention-rich: software, laboratory science). The agent's execution-observation pairs are genuine interventions with clean attribution. $p_{ij}$ approximates the interventional probability $P(j \mid do(i), M_t)$.
  \item \textbf{Regime B} (partial intervention: organizational, coordinated action). The agent acts but attribution is blurred by concurrent actions and self-selection. $p_{ij}$ is a partially identified causal estimate, typically biased upward.
  \item \textbf{Regime C} (observation-only: passive monitoring, intelligence analysis). The agent observes associations but does not intervene. $p_{ij}$ is an observational proxy for the causal quantity --- useful for planning but potentially confounded.
\end{itemize}


The identifiability coefficient $\iota_{ij}$ ( \S\,\ref{edge-update-causal-validity}) quantifies the strength of the causal interpretation for each edge. When $\iota_{ij} \approx 1$, the agent's credence is well-identified causally. When $\iota_{ij} \approx 0$, the credence is associational. The single-parameter edge design is preserved: $p_{ij}$ is always the agent's working estimate, with $\iota_{ij}$ characterizing its causal warrant separately.

\textbf{Status propagation.} Forward pass in topological order, $O(\lvert V \rvert + \lvert E \rvert)$:

$$s_v = \begin{cases} p_v & \text{if } v \text{ is a leaf (base credence)} \\ \prod_{i \in \text{pa}(v)} p_{iv} \cdot s_i & \text{if } \gamma(v) = \text{AND} \\ 1 - \prod_{i \in \text{pa}(v)} (1 - p_{iv} \cdot s_i) & \text{if } \gamma(v) = \text{OR} \end{cases}$$

\textbf{Terminal satisfaction conditions.} The root terminal $v_\text{root}$ and any intermediate nodes near the top of the DAG carry \textbf{satisfaction conditions}: predicates on environment states/trajectories that the agent treats as operational success criteria for the objective. These conditions operationalize $O_t$ within $\Sigma_t$ — they are the agent's theory of what it means to satisfy the objective. $O_t$ itself lives outside $\Sigma_t$ ( \S\,\ref{strategy-dimension}); the terminal conditions are $\Sigma_t$'s internal encoding of what $O_t$ requires. When $O_t$ changes, terminal conditions must be reassessed and potentially replaced ( \S\,\ref{structural-change-as-parametric-limit}).

\textbf{Well-formedness.} $\Sigma_t$ is **$O_t$-well-formed** when the agent believes that achieving the terminal conditions yields a trajectory that satisfies the objective:

$$\Pr\!\left(O_t \text{ satisfied by } \tau \;\middle\vert\; \text{terminal conditions achieved},\; M_t\right) \geq 1 - \epsilon$$

where "$O_t$ satisfied" means $V_{O_t}(\tau)$ exceeds the objective's own satisfaction criterion (formalized as $V_{O_t}^{\min}$ in \S\,\ref{satisfaction-gap}). This is a constraint on the relationship between $\Sigma_t$ and $O_t$, not a separate state object. It is explicit and in-principle assessable, though evaluating it requires the same value-side machinery as $A_O$ — it is not a cheap structural test. Violation triggers terminal reassessment: either the terminals need revision (they don't operationalize $O_t$ correctly) or $O_t$ itself needs revision.

\textbf{Strategy self-assessment.} The root node's propagated status:

$$\hat{P}_\Sigma(M_t) = s_{v_\text{root}}$$

is the strategy's \textbf{plan-confidence score} — the DAG's own answer to "will this plan work?" Under the edge-independence assumption implicit in the AND/OR combination rules, this is a probability. When edges have correlated failures (shared infrastructure, common-mode risks — see Working Notes), it is a heuristic confidence score that systematically overestimates success likelihood. $\hat P_\Sigma$ is explicitly distinct from $A_O$ ( \S\,\ref{satisfaction-gap}), which optimizes over the entire policy class, and from $V_O(\pi_\text{current})$ ( \S\,\ref{value-object}), which evaluates the current policy. $\hat P_\Sigma$ is cheap to compute ($O(\lvert V\rvert + \lvert E\rvert)$ forward pass) and updates in real time as $M_t$ changes through leaf credences.

\textbf{Scope of the terminal construction.} Terminal conditions as Boolean predicates with AND/OR aggregation work naturally for threshold, constraint, and composite objectives. For continuous-valued objectives without natural thresholds, the agent must set an operational threshold — introducing a discretization that is a practical proxy, not a lossless encoding of $V_O$. The primary $O_t$ ↔ theory interface remains $V_O$ through the value object ( \S\,\ref{value-object}); terminal conditions are $\Sigma_t$'s internal operational encoding.

\textbf{Single-parameter edges.} Each edge carries one number ($p_{ij}$), not two. An earlier formalism attempt used $(p_{ij}, \theta_{ij})$ where $\theta$ was "contribution magnitude." This was dropped because the AND/OR combination rules at nodes absorb $\theta$'s role — the complexity budget goes to one bit per node ($\gamma$) instead of one float per edge.


\medskip\noindent\textbf{Acyclicity is Derived.}


% [Derived (from causal-structure + finite planning horizon)]

Each node in $\Sigma_t$ represents a future event or state with temporal position $\tau_i > t$. An edge $X_i \to X_j$ requires $\tau_i < \tau_j$ ( \S\,\ref{causal-structure}: causes precede effects). A cycle $X_i \to X_j \to \cdots \to X_i$ would require $\tau_i < \tau_j < \cdots < \tau_i$, which is impossible for a real-valued time index.

Strategies involving iteration ("try A, if fail try B, if fail try A again") are acyclic when time-indexed. The sequence unfolds as:

$$A_1 \to \text{check}_1 \to B_1 \to \text{check}_2 \to A_2 \to \ldots$$

Each attempt is a distinct node at a distinct time. The apparent cycle is a linear chain in the unrolled view.

Formally: a finite set with a strict partial order (future events ordered by time) is representable as a DAG. This is a standard result in order theory.

\textbf{Scope of the acyclicity result.} This applies to $\Sigma_t$ (the agent's strategy over the future), not to $M_t$'s model of the environment, which may include cyclic causal processes (feedback loops in the physical world, market dynamics, ecosystem interactions). The acyclicity is specific to the purposeful substate.


\medskip\noindent\textbf{Epistemic Status.}


\textit{Conditional} on the \S\,\ref{and-or-scope} restriction. The DAG structure itself follows from operational postulates: temporal ordering (acyclicity — exact), probabilistic uncertainty (Cox's theorem — exact), and state-local revisability combined with causal sufficiency (Markov factorization — conditional). The full argument is in \S\,\ref{graph-structure-uniqueness}: P1 + P2 + P3 + causal sufficiency → DAG with Markov property. The result is \textit{conditional on causal sufficiency} of the strategy (no latent common causes among strategy nodes). For agent-constructed strategies this is a reasonable assumption; when it fails, the strategy is model-inadequate ( \S\,\ref{structural-adaptation-necessity}).

\textbf{Edge-independence caveat.} The AND/OR status propagation assumes independent edge failures: the probability of each edge being intact is independent of all other edges. This assumption is systematically violated in complex real-world systems. Correlated failure is the dominant mode — shared infrastructure, common-mode risks, supply chain dependencies, and correlated adversary actions cause multiple edges to fail together. The plan-confidence score $\hat P_\Sigma$ therefore \textbf{systematically overestimates success likelihood}. An adversary exploits correlations, not independent edges. This is not a minor qualification: in adversarial settings and complex environments, the overestimation may be severe. The independent-edge model remains useful as a tractable lower bound on failure probability and as the simplest non-trivial parameterization, but agents relying on $\hat P_\Sigma$ for high-stakes decisions should expect actual success rates to be lower than computed.

The AND/OR parameterization is a parsimony-motivated formulation choice within the strongly motivated graphical structure, not a derived necessity ( \S\,\ref{and-or-scope}). The single-parameter edge convention is similarly a formulation choice motivated by convergence across three independent attempts.


\medskip\noindent\textbf{Discussion.}


\textbf{The graph structure is strongly motivated; the parameterization is chosen.} The DAG structure follows from temporal ordering (acyclicity — derived, see above), probabilistic uncertainty (Cox's theorem forces probability on edges), and a plausible argument that state-local revisability strongly motivates the Markov factorization (conditional on causal sufficiency — see \S\,\ref{graph-structure-uniqueness}). Acyclicity is a proven result; the full DAG-with-Markov-property argument is not yet at theorem strength (the P3→Markov step remains a sketch). The status is therefore: DAG structure is \textit{strongly motivated} by operational postulates, not \textit{forced} in the mathematical sense. ACT uses DAG + AND/OR because (a) AND/OR is the most parsimonious complete basis for binary combination, (b) the DAG naturally supports causal reasoning, and (c) the representation converged across three independent formalism attempts. Alternative parameterizations within the graphical structure are legitimate research directions.

\textbf{Combination assignment is principled but fallible.} The question "if I remove one parent, can $v$ still be achieved?" is derivable from $M_t$'s causal model — it's a principled assignment, not arbitrary. But the assignment can be wrong (false AND = pessimistic over-investment; false OR = optimistic under-investment), and should be updateable when evidence reveals a different structural relationship.

\textbf{Connection to Pearl's framework.} The strategy DAG shares structure with a causal Bayesian network: directed edges, propositional nodes, and a probabilistic factorization. In Regime A domains, where the agent performs genuine interventions and observes isolated outcomes, the DAG's edge credences approximate interventional probabilities, and Pearl's do-calculus applies to status propagation and plan evaluation. In Regimes B and C, the edge credences are the agent's working causal beliefs, updated from data of weaker identification strength. The DAG remains useful for planning --- the agent must act on \textit{some} causal model --- but the plan-confidence score $\hat P_\Sigma$ inherits the identification weaknesses of its constituent edges. An agent operating primarily in Regime C should treat $\hat P_\Sigma$ as a rough heuristic, not a calibrated probability. The regime-indexed interpretation (above) makes this gradient explicit rather than leaving it as a caveat.

\textbf{Depth penalties on calibration.} Beyond the confidence decay that deeper DAGs suffer ( \S\,\ref{chain-confidence-decay}), the two-edge strategic dynamics analysis (\S\,\ref{strategic-dynamics-derivation} (Props B.2-B.3)) shows that deeper edges are also harder to calibrate. Edge $k$ in a chain is tested only when all upstream edges succeed, so its effective correction rate is attenuated by $\prod_{j<k}\theta_j$ (the evidence-starvation effect). Deeper DAGs therefore face a double penalty: lower aggregate confidence AND slower convergence of edge credences toward truth. This reinforces the structural pressure toward shallow, observable strategies — deep plans require both high per-edge reliability and sustained observability at every intermediate level to remain calibratable.

\textbf{Edge-independence scope.} The AND/OR status propagation treats edge failures as independent. This is a \textbf{load-bearing assumption}, not a minor caveat — it enters every result that depends on $\hat P_\Sigma$.

\textit{What it buys:} Tractable $O(\lvert V\rvert + \lvert E\rvert)$ status propagation; single-parameter edges; clean persistence proofs (the sector condition transfers from per-edge credence to plan value via the Jacobian — Prop B.5 in \S\,\ref{strategic-dynamics-derivation}).

\textit{What it costs:} $\hat P_\Sigma$ systematically overestimates success under positively correlated failure (the dominant real-world case). The plan-confidence error $\delta_s = \hat P_\Sigma - \Phi$ proved persistent by B.5 tracks calibration \textit{within the independence model} — $\Phi$ is the AND/OR formula at true edge rates, not actual plan success ( \S\,\ref{strategic-calibration}). Gradient-based credit assignment ( \S\,\ref{credit-assignment-boundary}) inherits the same bias: the residual $(y_G - \hat P_\Sigma)$ conflates per-edge miscalibration with omitted correlation structure. Chain confidence decay ( \S\,\ref{chain-confidence-decay}) notes that positive correlation makes the independent formula conservative, understating actual fragility.

\textit{Mitigation:} Agents can model correlation structure explicitly by adding common-cause nodes to $\Sigma_t$ — a shared infrastructure dependency becomes a node whose children are the edges it correlates. This shifts the burden from the propagation formula (which remains independent-edge) to the DAG construction (which now represents the correlation). The independence assumption then holds \textit{within the augmented DAG}, at the cost of a larger graph.

\end{definition}

\begin{proposition}[Directed Separation]
\label{directed-separation}
The epistemic update function $f_M$ is goal-blind: it processes incoming events without reference to the agent's objectives or strategy. The purposeful update $f_G$ depends on the updated epistemic state. Action couples all substates. This directed asymmetry — epistemic update is independent of purpose; purposeful update depends on epistemic state — is the structural backbone of the theory.



\medskip\noindent\textbf{Formal Expression.}


% [Derived (directed-separation, from complete-agent-state + scope condition)]

\textbf{The update functions:}

$$M_{\tau^+} = f_M(M_{\tau^-}, e_\tau) \qquad \text{(no } G_t \text{ argument)}$$

$$G_{\tau^+} = f_G(G_{\tau^-}, M_{\tau^+}, e_\tau) \qquad \text{(depends on updated } M_t \text{)}$$

\textbf{The policy:}

$$a_t = \pi(M_t, G_t) \qquad \text{(couples all substates)}$$

The three lines encode the full coupling structure:
\begin{itemize}
  \item $f_M$ determines how the agent updates beliefs — independently of what it wants
  \item $f_G$ determines how the agent revises purpose — in light of what it now believes
  \item $\pi$ determines what the agent does — based on both what it knows and what it wants
\end{itemize}


% [Scope Condition (directed-separation-scope)]

The claim "$f_M$ has no $G_t$ argument" requires that the epistemic update is \textbf{goal-blind conditional on the realized event}. This holds when:

\begin{enumerate}
  \item The observation mechanism $h$ may be action-dependent ( \S\,\ref{scope-condition} allows this), but $f_M$ processes whatever event arrives without reference to why the agent sought that event
  \item The agent does not use its goals to filter, weight, or interpret observations differently — no goal-dependent attention thresholds or confirmation bias baked into $f_M$
\end{enumerate}


If the agent's goals influence the \textit{observation mechanism} (goal-directed sensing, attention allocation, query selection), the \textbf{event that arrives} depends on $G_t$ through $\pi \to a_t \to e_\tau$. But $f_M$ still processes the event goal-blindly. The directed separation is about the \textbf{processing} of events, not the \textbf{selection} of events.


\medskip\noindent\textbf{Architectural classification.}


% [Scope Condition (directed-separation-architecture)]

Whether directed separation holds is determined by the agent's \textbf{processing topology} — specifically, whether $G_t$ is causally upstream of $f_M$ in the agent's internal processing graph. This is a structural property of the architecture, not a tunable parameter.

\begin{adjustbox}{max width=\textwidth}
\small
\begin{tabular}{llll}
\toprule
Class & Topology & Directed separation & Examples \\
\midrule
\textbf{1. Modular} & Separate estimator and planner, connected through state-estimate interface & Holds by construction — estimator has no causal path from $G_t$ & Kalman filter + LQR; modular RL with separate world model; military intelligence separated from operations \\
\textbf{2. Fully merged} & Single mechanism handles both epistemic and strategic processing & Fails by construction — $G_t$ is causally upstream of every computation & Transformer LLM (attention processes goals and observations together); potentially human cognition (motivated reasoning) \\
\textbf{3. Partially modular} & Some shared infrastructure, some separate pathways & Holds for modular stages, fails for merged stages & Biological cortex (shared sensory areas, separate prefrontal); hybrid AI with separate preprocessing \\
\bottomrule
\end{tabular}
\end{adjustbox}


\textbf{Operationalization.} The degree of coupling in partially modular architectures (Class 3) can be quantified as:


$$\kappa_{\text{processing}} = \frac{I(G_t \,;\, M_{\tau^+} \mid e_\tau,\, M_{\tau^-})}{H(G_t \mid e_\tau,\, M_{\tau^-})}$$

where $I(\cdot;\cdot\mid\cdot)$ is conditional mutual information and $H(\cdot\mid\cdot)$ is conditional entropy. The conditioning on $M_{\tau^-}$ is essential: without it, prior correlation between goals and model state (which exists even in modular agents) inflates the measure. The quantity captures \textit{extra} goal information entering the epistemic update beyond what was already in the prior model — information that flows through shared causal paths in the processing infrastructure (paths that bypass the event $e_\tau$).

\begin{itemize}
  \item $\kappa_{\text{processing}} = 0$: Class 1 (modular). No information about $G_t$ reaches $M_{\tau^+}$ except through $e_\tau$.
  \item $\kappa_{\text{processing}} \approx 1$: Class 2 (fully merged). Nearly all goal information is available to the epistemic update.
  \item $0 < \kappa_{\text{processing}} < 1$: Class 3 (partially modular). The value depends on the architecture's interface design.
\end{itemize}


\textbf{Distribution dependence.} $\kappa_{\text{processing}}$ is a distribution-dependent measure: it quantifies how much goal-information actually flows through the shared pathways under a given distribution of tasks, goals, and events. It does not directly measure whether pathways \textit{exist} — that is the architectural classification (Class 1/2/3), which is structural and distribution-independent. A Class 1 agent has $\kappa = 0$ under ALL distributions (no pathway exists). A Class 2 agent has high $\kappa$ under most distributions (pathways exist and are used). A Class 3 agent's $\kappa$ varies with the task distribution — the same hybrid architecture may exhibit low coupling on familiar tasks (where the modular stages handle most processing) and high coupling on novel tasks (where goal-conditioned downstream reasoning dominates). The classification is the primary tool; the operationalization is a diagnostic for Class 3 agents where the degree of coupling is architecturally ambiguous.

\textbf{Why the classification is not a smooth parameter.} The architectural boundary between "has a separable perception module" and "processes everything through goal-conditioned attention" is discrete. Within the modular class, $\kappa \approx 0$ regardless of task. Within the fully merged class, $\kappa$ is high regardless of prompt design. Only in the partially modular class is $\kappa$ genuinely variable and worth parameterizing. This replaces the earlier $\kappa$-as-scalar framing (see \texttt{msc/spike-kappa-topology-insight.md}) which treated coupling as a smoothly tunable quantity.

\textbf{Implications for theory scope:}
\begin{itemize}
  \item \textbf{Class 1}: Section II's results apply exactly. The sequential orient cascade is the correct analysis.
  \item \textbf{Class 2}: Requires coupled formulation from the start — $X_{\tau^+} = f_X(X_{\tau^-}, e_\tau)$ without decomposition. This is the scope of \texttt{03-logogenic-agents/}.
  \item \textbf{Class 3}: The sequential cascade is an approximation. Approximation quality depends on $\kappa_{\text{processing}}$ and requires per-architecture error analysis.
\end{itemize}



\medskip\noindent\textbf{Epistemic Status.}


\textit{Conditional} on the scope condition above. The conditional claim (IF epistemic update is goal-blind, THEN the separation holds) is exact. Whether a particular agent satisfies the condition is determined by its processing architecture (Class 1/2/3).

The architectural classification: \textbf{robust qualitative}. The three classes are structurally distinct (modular vs merged vs partial), well-motivated by examples across domains, and supported by a formal operationalization ($\kappa_{\text{processing}}$). The classification replaces the earlier $\kappa$-as-scalar framing, which is documented in \texttt{msc/spike-kappa-topology-insight.md}. The operationalization of $\kappa_{\text{processing}}$ as conditional mutual information is well-defined but typically not computable in closed form for real architectures — it serves as a conceptual anchor rather than a practical measurement tool.


\medskip\noindent\textbf{Discussion.}


\textbf{This is a genuine scope restriction, not a footnote.} An LLM agent's prompt includes the task objective, which shapes how it interprets code, documentation, and error messages. Its $f_M$ is goal-conditioned in practice: the agent reading code with the goal "fix the auth bug" processes the same code differently than one with "add logging." The epistemic update and purposeful evaluation are entangled in the attention mechanism.

\textbf{When the approximation is good:}
\begin{itemize}
  \item Goal-conditioning affects \textit{attention} (which events to seek) more than \textit{interpretation} (how to process events that arrive)
  \item The agent has strong epistemic discipline (updates beliefs based on evidence quality, not goal alignment)
  \item The epistemic update is architecturally separated from goal evaluation (e.g., separate model-update and planning modules)
\end{itemize}


\textbf{When the approximation is poor:}
\begin{itemize}
  \item The agent exhibits confirmation bias (interpreting ambiguous evidence in goal-consistent ways)
  \item Goal-conditioning is deeply embedded in the processing architecture (attention-based models where the query includes intent)
  \item The agent's observation channel is strategically controlled by an adversary who knows the agent's goals
\end{itemize}


\textbf{What directed separation buys the theory.} Section I's $M_t$-side quantities — $\delta$, $\eta^*$, $\mathcal{T}$, the persistence condition — remain well-defined on $M_t$ regardless of whether directed separation holds. What directed separation provides is the \textit{clean factorized update}: $M_t$ updates independently, then $G_t$ updates in light of the new $M_t$, and the orient cascade resolves sequentially. Without directed separation, the $M_t$ dynamics depend on $G_t$, the update becomes a coupled system, and the sequential orient cascade becomes an approximation of a simultaneous fixed-point problem. The theory still applies — the quantities are well-defined — but the modular Section I → Section II lift becomes a coupled analysis.

\textbf{The deeper question.} Goal-conditioned epistemic dynamics — where $f_M$ depends on $G_t$ — is the formal territory of motivated reasoning, confirmation bias, and wishful thinking. A future extension would model these as departures from directed separation: coupling terms in $f_M$ that create richer (and more fragile) dynamics. The current theory treats this as out of scope, which is honest but leaves the most human-like and LLM-like agents as approximate fits.

\end{proposition}

\begin{definition}[Satisfaction Gap]
\label{satisfaction-gap}
The satisfaction gap measures whether the agent's objective is achievable: the distance between what the objective requires and what the best available policy can deliver. It answers "is this goal feasible given what I know and what I can do?"



\medskip\noindent\textbf{Formal Expression.}



$$A_O(M_t;\, \Pi, N_h) = \sup_{\pi \in \Pi} V_O(M_t, \pi;\, N_h)$$

The \textbf{objective attainability} — the best achievable value given current beliefs $M_t$, available policy class $\Pi$, and horizon $N_h$.


$$\delta_{\text{sat}} = V_{O_t}^{\min} - A_O(M_t;\, \Pi, N_h)$$

where $V_{O_t}^{\min}$ is the minimum trajectory value that counts as "objective met" — a threshold set by the objective itself (for constraints: all satisfied; for utility: a minimum acceptable level).

\begin{itemize}
  \item $\delta_{\text{sat}} > 0$: The objective is \textbf{unmet} under the best available policy, current model, and horizon.
  \item $\delta_{\text{sat}} \leq 0$: The objective is \textbf{attainable} in principle.
\end{itemize}


\textbf{Disambiguation.} $\delta_{\text{sat}} > 0$ does NOT automatically mean the goal is wrong. It means the goal is unmet given $(M_t, \Pi, N_h)$. The positive signal has multiple possible causes:

\begin{adjustbox}{max width=\textwidth}
\small
\begin{tabular}{lll}
\toprule
Cause & Fix & How to distinguish \\
\midrule
Goal is genuinely infeasible & Revise $O_t$ & Persists across $M_t$, $\Pi$, $N_h$ improvements \\
Policy class too narrow & Expand $\Pi$ (structural adaptation of $\Sigma_t$) & $\delta_{\text{sat}}$ decreases when richer policies are tried \\
Horizon too short & Extend $N_h$ & $\delta_{\text{sat}}$ decreases with longer planning horizon \\
Model is wrong about feasibility & Improve $M_t$ (reduce $\delta_{\text{epistemic}}$) & $\delta_{\text{sat}}$ changes when $M_t$ is corrected \\
Objectives jointly infeasible & Revise $O_t$ or relax constraints & Individual terminal satisfaction gaps are zero but AND-node fails; cross-terminal tradeoff is required \\
\bottomrule
\end{tabular}
\end{adjustbox}


Objective revision is the \textbf{last resort}, not the first response to unmet goals. The orient cascade ( \S\,\ref{orient-cascade}) formalizes this ordering.


\medskip\noindent\textbf{Epistemic Status.}


\textit{Exact as a definition — convention-relative as a diagnostic.} The satisfaction gap is a mathematical definition — the difference between a threshold and a supremum over a function class. The definition is precise; the \textit{computation} of $A_O$ is generally intractable (it requires optimization over a policy class), but the quantity is well-defined. However, the \textit{value} of $\delta_{\text{sat}}$ depends on three external parameters: the continuation convention ($\pi_{\text{cont}}$), the horizon ($N_h$), and the scalarization of the objective ($V_{O_t}$). These are not derived by ACT — they are choices the analyst makes. The satisfaction gap is therefore an intrinsic architectural diagnostic \textit{given} a measurement convention, not an absolute property of the agent.

\textbf{Convention dependence.} $A_O$ inherits the continuation convention from the value object ( \S\,\ref{value-object}). Under the canonical default ($\pi_{\text{cont}} = \pi_{\text{current}}$), $A_O = \sup_{\pi \in \Pi} V_O(M_t, \pi; N_h)$ is the \textbf{best first-action attainability} — the supremum is over the first action only, with continuation fixed at $\pi_{\text{current}}$. This is a \textit{one-step-improvement} quantity, not a full-policy optimum. Under the Bellman convention ($\pi_{\text{cont}} = \pi^*$), $A_O$ becomes the full-policy optimal value — a stronger but harder-to-compute quantity. The distinction matters: $\delta_{\text{sat}}$ and $\delta_{\text{regret}}$ mean different things under different conventions, and analyses using different conventions should not be compared directly. ACT adopts the one-step improvement convention as the canonical default; all diagnostics are comparable across analyses of the same agent only when computed under the same convention.


\medskip\noindent\textbf{Discussion.}


\textbf{Why two gap measures, not one.} An earlier formulation used a single $\delta_{\text{objective}}$ for goal-related mismatch. This conflates two distinct situations: "the goal is too hard" and "the strategy is too weak." When the agent is optimally pursuing an infeasible goal, $\delta_{\text{objective}}$ is large but there's no strategy to improve — the problem is the goal, not the plan. The satisfaction gap ($\delta_{\text{sat}}$) and control regret ( \S\,\ref{control-regret}) separate these cases, enabling the right corrective action.

\textbf{The disambiguation table is load-bearing.} Without it, an agent facing $\delta_{\text{sat}} > 0$ might immediately revise its objective when the real problem is an inadequate model or a too-narrow policy class. The table encodes the diagnostic procedure: check $M_t$ adequacy first (maybe the goal IS feasible but the model doesn't know it), then check $\Pi$ and $N_h$, and only then consider revising $O_t$.

**Dependence on $M_t$.** $A_O$ is computed from $M_t$, not from the true environment state $\Omega_t$. The agent's assessment of attainability could be wrong — an achievable goal might look unachievable with a bad model, or vice versa. Improving $M_t$ (reducing $\delta_{\text{epistemic}}$) brings the agent's attainability assessment closer to reality. This is why the orient cascade puts epistemic update before attainability evaluation.

\end{definition}

\begin{definition}[Control Regret]
\label{control-regret}
Control regret measures the gap between the best achievable performance and the agent's current performance — how much the agent is leaving on the table by following its current policy rather than the best available one. This is the signal for strategy revision.



\medskip\noindent\textbf{Formal Expression.}



$$\delta_{\text{regret}} = A_O(M_t;\, \Pi, N_h) - V_O(M_t, \pi_{\text{current}};\, N_h) \geq 0$$

Always non-negative: the current policy cannot outperform the best in its class.

\begin{itemize}
  \item $\delta_{\text{regret}} \approx 0$: The agent is doing the best it can within current $(\Pi, N_h, M_t)$. If $\delta_{\text{sat}} > 0$ simultaneously, the problem is not the current strategy — it's either the goal, the capability ($\Pi$, $N_h$), or the model ($M_t$). See \S\,\ref{satisfaction-gap}'s disambiguation.
  \item $\delta_{\text{regret}} \gg 0$: There's room for improvement without changing $O_t$. → Revise $\Sigma_t$.
\end{itemize}



\medskip\noindent\textbf{Epistemic Status.}


\textit{Exact as a definition — convention-relative as a diagnostic.} Like the satisfaction gap, this is a mathematical definition — a difference between two values of the same functional. The quantity is well-defined; computing it requires evaluating $A_O$ (generally intractable) and $V_O$ under the current policy (tractable in simulation, approximate in practice). Like $\delta_{\text{sat}}$, this diagnostic is convention-relative: it changes with $\pi_{\text{cont}}$, $N_h$, $\Pi$, and the scalarization $V_{O_t}$. Results computed under different conventions are not comparable. See \S\,\ref{satisfaction-gap} Epistemic Status for the full convention-dependence discussion.


\medskip\noindent\textbf{Discussion.}


\textbf{The diagnostic power of the two-gap system.} The satisfaction gap and control regret together encode a 2×2 diagnostic:

\begin{adjustbox}{max width=\textwidth}
\small
\begin{tabular}{lll}
\toprule
 & $\delta_{\text{sat}} \leq 0$ (attainable) & $\delta_{\text{sat}} > 0$ (unmet) \\
\midrule
$\delta_{\text{regret}} \approx 0$ (near-optimal) & \textbf{Success}: goal achievable, policy good & \textbf{Capability limit}: optimally pursuing an unmet goal → check $M_t$, $\Pi$, $N_h$, then consider revising $O_t$ \\
$\delta_{\text{regret}} \gg 0$ (suboptimal) & \textbf{Strategy problem}: goal achievable, policy poor → revise $\Sigma_t$ & \textbf{Both}: goal hard AND strategy weak → revise $\Sigma_t$ first, then reassess $\delta_{\text{sat}}$ \\
\bottomrule
\end{tabular}
\end{adjustbox}


This diagnostic is what makes the orient cascade ( \S\,\ref{orient-cascade}) actionable: each cell prescribes a different corrective action.

**Control regret as the signal for $\Sigma_t$ revision.** When $\delta_{\text{regret}}$ is high, the agent knows it could do better with a different strategy. The \textit{specific} corrections — which edges to revise, which branches to prune, which alternatives to add — come from the strategic calibration residual ( \S\,\ref{strategic-calibration}), which localizes the regret to specific parts of $\Sigma_t$.

\textbf{Regret approaching zero when optimally failing.} This is the key insight motivating the two-gap split. A single $\delta_{\text{objective}}$ would show "large gap" for both "bad strategy, achievable goal" and "good strategy, impossible goal." The first warrants strategy revision; the second warrants goal revision (after ruling out $M_t$/$\Pi$/$N_h$ inadequacy). Without the split, the agent cannot distinguish these cases and may waste effort optimizing a strategy that's already near-optimal for an infeasible goal.

\end{definition}

\begin{definition}[Strategic Calibration]
\label{strategic-calibration}
The strategic calibration residual measures whether the strategy's causal model is correct: are the edges in $\Sigma_t$ accurate predictors of how much value each step actually produces? This is the fine-grained diagnostic that localizes control regret to specific parts of the strategy.



\medskip\noindent\textbf{Formal Expression.}



For each edge $(i, j)$ in $\Sigma_t$ with credence $p_{ij}$, the \textbf{edge residual}:

$$r_{ij} = \mathbb{E}[\Delta V_O \mid \text{edge } (i,j) \text{ traversed},\, M_t] - \Delta V_O^{\text{observed}}$$

where $\Delta V_O$ is the change in $V_O(M_t, \pi;\, N_h)$ attributable to completing step $j$ — as predicted by $\Sigma_t$ versus as observed.

The \textbf{strategic calibration residual} aggregates across active edges:

$$\delta_{\text{strategic}} = \left(\sum_{(i,j) \in \text{active}} w_{ij} \cdot r_{ij}^2 \right)^{1/2}$$

where $w_{ij}$ weights edges by importance (e.g., criticality to the current plan's critical path).

\textbf{Conditioning.} The edge residual $r_{ij}$ is meaningful only when:
\begin{itemize}
  \item The edge was actually traversed (the agent attempted the step)
  \item $M_t$ is adequate (so the observed $\Delta V_O$ is meaningful, not noise)
  \item The agent followed $\Sigma_t$'s prescription for step $j$ (execution fidelity — otherwise the residual conflates bad plan with bad execution)
\end{itemize}


Without the execution fidelity condition, a positive residual could mean "the plan is wrong" or "the agent didn't follow the plan." These require different corrections ($\Sigma_t$ revision vs. execution improvement).


\medskip\noindent\textbf{Epistemic Status.}


\textit{Discussion-grade.} The edge residual concept is well-motivated: each edge predicts a value increment, and comparing prediction to observation is standard calibration. But the specific aggregation ($L^2$ norm with importance weights) is a reasonable first pass, not a derived result. The weighting scheme ($w_{ij}$ by criticality) is sensible but ungrounded. The conditioning requirements (especially execution fidelity) make this quantity hard to estimate in practice — the agent must know whether it followed its own plan, which requires a level of self-monitoring that many agents lack.

This is a \textbf{second-order inference} — it requires accumulating evidence over multiple edge traversals. It is inherently slower to evaluate than $\delta_{\text{epistemic}}$ (which updates on every observation) or $\delta_{\text{sat}}$ (which can be evaluated from $M_t$ and $\Sigma_t$ alone).


\medskip\noindent\textbf{Discussion.}


\textbf{Connection to \S\,\ref{strategy-persistence-schema}.} There are two distinct strategic mismatch quantities:

\begin{enumerate}
  \item \textbf{Plan-confidence error} $\delta_s = \hat P_\Sigma - \Phi$ — the scalar gap between the agent's plan-confidence score and the independence-model plan value at true edge parameters. Computable from status propagation alone, without credit assignment. The sector condition transfers to $\delta_s$ (Prop B.5 in \S\,\ref{strategic-dynamics-derivation}). Note: $\Phi$ is the AND/OR propagation formula evaluated at true edge rates — it is NOT actual plan success probability under correlated failure ( \S\,\ref{strategy-dag}, edge-independence caveat). $\delta_s$ tracks calibration \textit{within the independence model}, not calibration to strategic reality.
\end{enumerate}


\begin{enumerate}
  \item \textbf{Strategic calibration residual} $\delta_{\text{strategic}}$ (this segment) — an $L^2$ aggregation of per-edge value-increment residuals requiring credit assignment to compute.
\end{enumerate}


These are related (both measure strategy-reality mismatch) but not interchangeable. $\delta_s$ is the proven persistence target; $\delta_{\text{strategic}}$ provides finer-grained diagnostics but its persistence properties remain open, pending the credit-assignment machinery in \S\,\ref{credit-assignment-boundary}. The correction function for both is edge-credence revision ( \S\,\ref{edge-update-via-gain}); the disturbance is environmental changes that alter edge-traversal outcomes.

\textbf{Typing as value-increment residuals.} Each edge predicts a scalar (value increment), not a full state transition. This is the most tractable typing because it connects directly to the value object $V_O$ and allows aggregation across heterogeneous step types (a military advance and a logistics delivery produce different state changes but both produce value increments measurable on the same scale).

\textbf{Credit-assignment problem.} The edge residual $r_{ij}$ subtracts "predicted value increment" from "observed value change." These are different types: the prediction comes from $\Sigma_t$'s internal causal model, while the observation comes from empirical measurement of $\Delta V_O$. The subtraction is meaningful only when the observed value change can be \textit{attributed to the specific edge} — a credit-assignment problem the current formulation does not address. For edges with a single parent, attribution is straightforward. For multi-parent AND/OR nodes, the observed $\Delta V_O$ at the child reflects the combined effect of all parent edges, and decomposing it into per-edge contributions requires additional structure (e.g., Shapley-value decomposition, or sequential observation of parent completions). In the absence of clean attribution, $r_{ij}$ measures "how well did the overall prediction work?" rather than "how well did this specific edge predict?" — still useful for aggregate calibration, but not sufficient for targeted edge revision.

\end{definition}

\begin{proposition}[Orient Cascade]
\label{orient-cascade}
For actuated agents, epistrophe (the corrective phase of the cycle) expands into a multi-step cascade. The resolution order is forced by information dependency: epistemic update first, then attainability assessment, then strategy evaluation, then (if needed) objective revision. Each step's input depends on the output of prior steps. The ordering is not a design choice — it's a consequence of which quantities require which others.



\medskip\noindent\textbf{Formal Expression.}


% [Derived (orient-cascade, from information dependency between mismatch types)]

\begin{enumerate}
  \item **Reduce $\delta_{\text{epistemic}}$** — understand reality.
\end{enumerate}

Update $M_t$ via \S\,\ref{mismatch-signal} and \S\,\ref{update-gain}. Prerequisite for all purposeful evaluation, because $M_t$ appears in every subsequent formula.

\begin{enumerate}
  \item **Evaluate $\delta_{\text{sat}}$** — is the goal achievable?
\end{enumerate}

Compute $A_O(M_t; \Pi, N_h)$ using the updated $M_t$. Requires adequate $M_t$ to assess attainability ( \S\,\ref{satisfaction-gap}).

\begin{enumerate}
  \item **Evaluate $\delta_{\text{regret}}$** — is the policy suboptimal?
\end{enumerate}

Compare $A_O$ to $V_O(M_t, \pi_{\text{current}}; N_h)$ ( \S\,\ref{control-regret}). This step applies regardless of $\delta_{\text{sat}}$'s sign — the 2×2 diagnostic ( \S\,\ref{control-regret}) requires both quantities to distinguish four cases:
\begin{itemize}
  \item $\delta_{\text{sat}} \leq 0$, $\delta_{\text{regret}} \approx 0$: \textbf{success} — goal attainable, policy near-optimal.
  \item $\delta_{\text{sat}} \leq 0$, $\delta_{\text{regret}} \gg 0$: \textbf{strategy problem} — goal attainable, policy poor → revise $\Sigma_t$.
  \item $\delta_{\text{sat}} > 0$, $\delta_{\text{regret}} \gg 0$: \textbf{both} — goal hard AND strategy weak → revise $\Sigma_t$ first (cheaper than revising $O_t$), then reassess $\delta_{\text{sat}}$.
  \item $\delta_{\text{sat}} > 0$, $\delta_{\text{regret}} \approx 0$: \textbf{capability limit} — already doing the best available; proceed to step 5.
\end{itemize}


\begin{enumerate}
  \item **If $\delta_{\text{regret}}$ high, evaluate $\delta_{\text{strategic}}$** — is the plan's causal model wrong?
\end{enumerate}

Examine edge residuals. Requires adequate $M_t$ and evidence of suboptimal execution ( \S\,\ref{strategic-calibration}).

\begin{enumerate}
  \item **If $\delta_{\text{sat}} > 0$ persists across $\Sigma_t$ revisions** — revise $O_t$.
\end{enumerate}

The cascade's ordering ensures objective revision is the last resort, not the first response to unmet goals. The agent reaches this step only in the capability-limit quadrant ($\delta_{\text{sat}} > 0$, $\delta_{\text{regret}} \approx 0$), or after strategy revision fails to close the gap.

\textbf{Derivation.} Each step's input depends on prior steps' outputs:
\begin{itemize}
  \item You cannot evaluate strategy quality with a broken reality model (step 3 requires step 1)
  \item You cannot distinguish "bad strategy" from "infeasible goal" without both $\delta_{\text{sat}}$ and $\delta_{\text{regret}}$ (step 3 requires step 2)
  \item You should not revise the objective until you've verified that improving $\Sigma_t$ cannot close the gap (step 5 requires steps 3-4)
\end{itemize}


The ordering is forced by information dependency.


\medskip\noindent\textbf{Epistemic Status.}


The cascade \textbf{ordering} is \textit{exact}: it is a logical consequence of which quantities appear in which formulas. Steps 1-2 (epistemic update, attainability assessment) rest on well-typed quantities ( \S\,\ref{mismatch-signal}, \S\,\ref{satisfaction-gap}) and exact derivation. Steps 3-5 (control regret, strategic calibration, objective revision) depend on \S\,\ref{strategic-calibration}, which is discussion-grade — the credit-assignment problem and execution-fidelity requirement are acknowledged but unresolved ( \S\,\ref{strategic-calibration}, Epistemic Status). The ordering of all five steps is forced by information dependency (each step's input depends on prior steps' output). The \textit{content} of steps 3-5 — what exactly the agent computes and whether the quantities are estimable in practice — inherits strategic-calibration's discussion-grade status. What is NOT derived is the \textit{timing} — how long the agent should spend on each step before proceeding.


\medskip\noindent\textbf{Discussion.}


**$G_t$ complexity bounded by $M_t$ capacity.** $\Sigma_t$'s evaluable complexity is bounded by $M_t$'s ability to observe which strategy edges are intact. An agent with poor model sufficiency ($S(M_t) \ll 1$) cannot meaningfully evaluate a complex $\Sigma_t$ — the strategic calibration residual requires adequate $M_t$ to distinguish "edge prediction wrong" from "observation too noisy to tell."

This creates a \textbf{virtuous cycle}: better $M_t$ → richer evaluable $\Sigma_t$ → better-directed action → faster $M_t$ improvement. And a \textbf{vicious one}: degraded $M_t$ → forced strategy simplification → cruder action → further $M_t$ degradation. The vicious cycle is the strategic analog of the death spiral described in the persistence condition ( \S\,\ref{persistence-condition}) — the agent loses the capacity to maintain complex plans, which reduces the quality of its actions, which further degrades its model.

\textbf{Connection to Boyd's OODA.} The orient cascade is the formal analog of Boyd's "Orient" — not just model updating, but the structured interaction between reality-understanding and strategy. Boyd's insight was that Orient is the critical step, not Decide. The cascade provides a mathematical mechanism consistent with this insight: Orient resolves the information dependencies that make Decide meaningful. An agent that skips to Decide (strategy revision) without adequate Orient (model update + attainability assessment) will revise its strategy based on stale or incorrect beliefs. Whether the dependency structure in the cascade captures the actual cognitive process Boyd described is an empirical question.

\textbf{Timescale structure.} The cascade implies a natural timescale ordering for the different update types:

$$\nu_{\text{epistemic}} \gg \nu_{\text{edge-update}} \gg \nu_{\gamma\text{-reclassify}} \gg \nu_{\text{prune/graft}} \gg \nu_{O\text{-revision}}$$

Weight updates are frequent (every observation). Combination-type reclassification is rare (needs strong structural evidence). Pruning/grafting is rarer still (abandon or create causal hypotheses). Objective revision is rarest (change what you want, not how you get it). This ordering is an empirical observation for many agent populations, consistent with the cascade but not derived from it.

\end{proposition}

\begin{proposition}[Observability Dominance]
\label{observability-dominance}
Unobservable strategy edges cannot be updated — the gain principle drives their update rate to zero. This means the agent's effective strategy is limited to the parts it can observe, regardless of the nominal confidence in unobservable paths. Observability dominates nominal confidence in determining which strategies are epistemically alive.



\medskip\noindent\textbf{Formal Expression.}


% [Derived (observability-dominance, from update-gain + strategy-dag)]

For a path $P$ through $\Sigma_t$, the \textbf{path observability}:

$$\text{obs}(P) = \min_{v \in P} \sigma_v$$

where $\sigma_v$ is the observability of node $v$ — how well the agent can determine whether $v$ has been achieved. The weakest link determines the path's observability.

\textbf{Observability-adjusted confidence:}

$$\text{conf}_{\text{obs}}(P) = \text{conf}(P) \cdot \text{obs}(P)$$

When $\sigma_v \approx 0$ for any node $v$ on the path: by \S\,\ref{update-gain}, $\eta_{\text{edge}} = U_{\text{edge}} / (U_{\text{edge}} + U_{\text{obs}}) \to 0$ as $U_{\text{obs}} \to \infty$. The edges connecting to $v$ are \textbf{frozen at their prior} — the agent cannot update them regardless of what happens. The path is epistemically dead.


\medskip\noindent\textbf{Epistemic Status.}


\textit{Robust qualitative.} The mechanism (high observation noise → low gain → frozen edges) is a direct consequence of \S\,\ref{update-gain}'s uncertainty ratio. The specific functional form ($\text{conf}_{\text{obs}} = \text{conf} \cdot \text{obs}$) is a first-order approximation — the actual relationship between observability and effective confidence is more complex (it depends on how many observations are accumulated, the prior strength, and the noise structure). The qualitative prediction (low observability → frozen beliefs → ineffective strategy) is robust.


\medskip\noindent\textbf{Discussion.}


\textbf{Observability as the gateway to learning.} An agent choosing between a strong-but-blind path (high $\text{conf}(P)$, low $\text{obs}(P)$) and a weak-but-visible path (lower $\text{conf}(P)$, high $\text{obs}(P)$) should prefer the visible one. After one attempt: the visible path yields large $\eta_{\text{edge}}$ — the agent quickly learns whether it works and can redirect. The blind path yields tiny $\eta_{\text{edge}}$ — the agent is still guessing after $n$ attempts. Observability enables learning; opacity prevents it.

\textbf{Unobservable regions are absorbing.} Once significant strategy investment operates through unobservable nodes: frozen beliefs → no mismatch signal → no reason to revise → the agent cannot learn and cannot recognize that it cannot learn. Escape requires external shock, proactive observability investment (instrumenting previously unmonitored nodes), or another agent whose observations cover the blind spot ( \S\,\ref{communication-gain}).

**Connection to \#code-quality-as-observation-infrastructure (cross-component reference — see \texttt{02-tst-core/}).** In the software domain, code quality directly determines $\sigma_v$ — well-structured code with good tests makes strategy steps (features, refactors, deployments) observable. Poor code quality reduces observability, freezing the developer's causal beliefs about what changes will accomplish. This makes code quality a strategic concern, not just an aesthetic one.

\textbf{Optimal decomposition depth.} Finer decomposition (more intermediate nodes) provides earlier failure detection — detect a problem at step $k$ rather than discovering it at step $n$. But finer decomposition also increases the number of uncertain edges ( \S\,\ref{chain-confidence-decay}). The optimal decomposition depth balances incremental confirmation against compound decay: decompose as finely as observation channels allow, but no finer.

\textbf{Quantitative content from the two-edge case.} The analysis in \S\,\ref{strategic-dynamics-derivation} (Props B.2-B.3) instantiates this result for the minimal multi-edge chain $A \to B \to G$, giving the qualitative claims above precise mathematical form.

*Observable intermediate $B$.* When the agent can observe whether $B$ was achieved, each edge gets independent Bayesian updates. The per-edge sector parameters are $\alpha_1 = 1/(n_1+1)$ and $\alpha_2 = \theta_1/(n_2+1)$, where $\theta_1$ is the true success probability of the first edge. The overall sector parameter is $\alpha_\Sigma = \min(\alpha_1, \alpha_2)$ — a weakest-link result. The $\theta_1$ factor in $\alpha_2$ is the \textbf{evidence-starvation effect}: downstream edge 2 can only be tested when upstream edge 1 succeeds (with probability $\theta_1$), so its effective correction rate is attenuated by $\theta_1$. For a depth-$d$ chain, this generalizes to $\alpha_k = \prod_{j<k}\theta_j / (n_k+1)$ — the deepest edge faces exponential attenuation. This is derived for the two-edge case and conjectured (with clear inductive mechanism) for depth-$d$.

*Unobservable intermediate $B$.* When the agent observes only the terminal outcome $y_G$, per-edge identification fails entirely. The marginal Bayesian update has a systematic bias: the zero-correction-at-truth property (A1) is violated with bias $O(1/n)$. The agent's point-estimate updates for each edge are downward-biased because success always credits both edges fully ($\alpha_k \to \alpha_k + 1$), but failure distributes blame fractionally via proportional attribution. The proportional-blame update turns out to be exactly the marginal Bayesian point estimate — not a heuristic — but it discards the posterior correlation that failure introduces between the edge beliefs.

The consequence is stronger than "frozen edges": unobservable intermediates don't just prevent updating — they make per-edge identification impossible, forcing the agent to plan-level aggregation. If the agent tracks $\hat{\Phi} = p_1 p_2$ as a single Beta on the plan's overall success probability, the single-edge sector condition applies with $\alpha_{\Sigma,\text{plan}} = 1/(n_\Phi + 1)$. But diagnostic resolution is lost: the agent knows the plan is failing but cannot localize which step needs revision.

\textit{The observability investment tradeoff.} Making $B$ observable yields a quantifiable improvement in $\alpha_\Sigma$: from the plan-level rate $1/(n_\Phi+1)$ to the per-edge weakest-link rate $\min(1/(n_1+1), \theta_1/(n_2+1))$. The improvement is positive whenever $\theta_1 > 1/2$ and experience is distributed similarly across edges. This gives the observability-dominance principle concrete economic content: the value of instrumenting an intermediate node is the difference in sector parameters, which translates directly to persistence margin.

\end{proposition}

\begin{hypothesis}[Edge Update via Gain]
\label{edge-update-via-gain}
The uncertainty-ratio gain principle ( \S\,\ref{update-gain}) extends from epistemic updates to strategy-edge updates: edge credences revise in proportion to the ratio of edge uncertainty to observation noise, modulated by the identifiability of the causal link ( \S\,\ref{edge-update-causal-validity}). This gives a principled, conservative update rule that avoids overreacting to single observations and degrades gracefully when causal identification is weak.



\medskip\noindent\textbf{Formal Expression.}


% [Hypothesis (edge-update-via-gain)]

Edge credences update via:

$$p_{ij}^{\text{new}} = p_{ij}^{\text{old}} + \eta_{\text{edge}} \cdot \left(\text{signal}(o_t, i, j) - p_{ij}^{\text{old}}\right)$$

where:
\begin{itemize}
  \item $\text{signal}(o_t, i, j) \in [0, 1]$: evidential content of observation $o_t$ about the causal link $i \to j$
  \item $\eta_{\text{edge}} = U_{\text{edge}} / (U_{\text{edge}} + U_{\text{obs}})$: update gain
\end{itemize}


with:
\begin{itemize}
  \item $U_{\text{edge}}$: uncertainty about this specific causal link. If $p_{ij} \sim \text{Beta}(\alpha_{ij}, \beta_{ij})$: $U_{\text{edge}} = \text{Var}[\text{Beta}] = \alpha\beta / ((\alpha + \beta)^2(\alpha + \beta + 1))$
  \item $U_{\text{obs}}$: observation noise on the channel confirming this link. $U_{\text{obs}} \propto 1/\sigma_j$ (inverse of observability of node $j$)
\end{itemize}


\textbf{Beta-Bernoulli instantiation.} For binary observations (success/failure of step $j$):
\begin{itemize}
  \item Observe success: $\alpha_{ij} \to \alpha_{ij} + 1$
  \item Observe failure: $\beta_{ij} \to \beta_{ij} + 1$
  \item Point estimate: $p_{ij} = \alpha / (\alpha + \beta)$
  \item Effective gain: $\eta_{\text{edge}} = 1/(n+1)$ where $n = \alpha + \beta$
\end{itemize}


This is the standard Bayesian conjugate update, yielding $\Delta\hat p = (y - \hat p)/(n+1)$. The gain $1/(n+1)$ is the exact conjugate update rate — not a literal substitution into $\eta = U_{\text{edge}}/(U_{\text{edge}} + U_{\text{obs}})$, which is a structural principle (conservative updating proportional to relative uncertainty), not a universal algebraic formula. The Beta-Bernoulli gain satisfies the same principle: it decreases as posterior certainty increases, trusting accumulated evidence over individual observations. But the algebraic derivation is conjugate analysis, not Kalman-style variance ratios, because the Bernoulli observation model has different noise structure than the Gaussian case where the variance-ratio formula is exact. The sector-condition analysis in \S\,\ref{strategic-dynamics-derivation} (Props B.1-B.4) uses the exact Beta-Bernoulli gain $1/(n+1)$ directly.


\medskip\noindent\textbf{Epistemic Status.}


\textit{Hypothesis.} The extension of the gain principle from $M_t$ to $\Sigma_t$ edges is structurally motivated: the same uncertainty-ratio logic applies wherever an agent updates a belief in response to noisy evidence. The update \textit{form} is validated for Beta-Bernoulli edges across four topologies (\S\,\ref{strategic-dynamics-derivation}, Props B.1-B.4). The signal function — how to compute $\text{signal}(o_t, i, j)$ — is now characterized at the theory level:

\begin{itemize}
  \item \textbf{Theory-level resolution.} \S\,\ref{credit-assignment-boundary} shows that persistence does not require per-edge credit assignment (Prop B.5), that any signal function with directional fidelity suffices for persistence (Level 1), and that exact attribution is \#P-hard in general (Level 3). The gradient-based candidate $\text{signal}_k = p_k + J_k \cdot (y_G - \hat P_\Sigma) / \lVert\mathbf{J}\rVert^2$ satisfies directional fidelity for monotone DAGs and is computable in $O(\lvert V\rvert + \lvert E\rvert)$.
  \item \textbf{What remains domain-specific.} The choice among signal function implementations (gradient attribution, proportional blame, belief propagation, exact Bayesian) depends on the domain's observability structure and computational budget. This is engineering, analogous to how the gain \textit{principle} ($\eta^* = U_M/(U_M + U_o)$) is theory while the gain \textit{estimator} (Kalman, TD-learning, intuition) is engineering.
  \item \textbf{Remaining open questions.} (1) Continuous outcomes with multi-parent nodes — signal function candidates exist but are not verified. (2) Interaction with $M_t$ updates �� edge updates use observations already processed for $M_t$; whether this creates statistical dependencies that bias edge updates is not analyzed. (3) The gradient candidate inherits $\hat P_\Sigma$'s overestimation bias under correlated failures ( \S\,\ref{strategy-dag}) — see \S\,\ref{credit-assignment-boundary} Working Notes.
\end{itemize}



\medskip\noindent\textbf{Discussion.}


\textbf{Connection to \S\,\ref{observability-dominance}.} When $\sigma_j \approx 0$: $U_{\text{obs}} \to \infty$, $\eta_{\text{edge}} \to 0$. The edge is frozen at its prior. Unobservable links cannot be updated — this is the mechanism that makes unobservable paths epistemically dead.

\textbf{Connection to \S\,\ref{strategic-calibration}.} The edge update rule is the correction function that the strategic calibration residual calls for: persistently positive $r_{ij}$ (predicted value increment exceeds observed) should reduce $p_{ij}$; persistently negative $r_{ij}$ should increase it. The gain principle provides the update magnitude — how much to adjust given the evidence strength.

\textbf{Conservative by design.} The uncertainty ratio ensures the agent doesn't overreact to single observations: a well-established edge ($\alpha + \beta$ large, $U_{\text{edge}}$ small) requires strong evidence to revise; a newly hypothesized edge ($\alpha + \beta$ small, $U_{\text{edge}}$ large) is easily moved by evidence. This mirrors the epistemic gain's behavior for $M_t$ — the agent trusts accumulated experience over individual observations.

\end{hypothesis}

\begin{scope}[Causal Validity of Edge Updates]
\label{edge-update-causal-validity}
The gain-based edge update ( \S\,\ref{edge-update-via-gain}) revises edge credences $p_{ij}$ --- causal efficacy estimates whose identification strength varies with the data regime ( \S\,\ref{strategy-dag}). This segment scopes where the update yields credences that approximate the interventional quantity $P(j \mid do(i), M_t)$, where it yields partially identified estimates, and where it yields associational proxies.



\medskip\noindent\textbf{Formal Expression.}


% [Scope Condition (edge-update-causal-validity)]


\medskip\noindent\textbf{Where the agent has direct intervention.}


By \S\,\ref{strategy-dag}, leaf action nodes are propositions about the agent's own actions: "action $a$ succeeds at $\tau_v$." When the agent executes an action leaf, it performs a genuine $do(\cdot)$ operation. The edge from that leaf to its child carries credence $p_{ij} = \text{Cr}(j \text{ advances} \mid do(i), M_t)$, and the execution-observation pair $(do(i), o_j)$ is interventional data for that edge.

However, interventional data does not automatically yield clean causal identification. By \S\,\ref{loop-interventional-access}, the loop provides data \textit{generated under intervention} — but between the intervention and a usable causal estimate stand coverage, within-step confounding, delay, and partial observability. The following conditions determine when the interventional data is strong enough for valid edge revision:

\textbf{(C1) The parent is an action leaf under the agent's control.} The agent directly executed the action. This makes the data interventional in character — the agent chose the action, the environment responded. For condition leaves (observable states the agent doesn't control) and internal nodes (propositional combinations achieved indirectly), $do(i)$ is not directly available. See "Indirect edges" below for the weaker identification available at those positions.

\textbf{(C2) The outcome is attributable.} The agent can distinguish whether $j$ advanced specifically because of $do(i)$, or for other concurrent reasons. This is the credit-assignment problem identified in \S\,\ref{strategic-calibration}. It is trivially satisfied for single-parent nodes (one possible cause) and for well-isolated interventions. It is violated when multiple parent edges of $j$ fire concurrently.

\textbf{(C3) Execution conditions vary.} The agent does not systematically execute $i$ only under conditions that independently favor $j$'s success. If it does, the observed success rate carries selection bias: $P(j \mid \text{chose to execute } i) \neq P(j \mid do(i))$ because the decision to execute correlates with favorable conditions. This is mitigated when the agent varies execution contexts across episodes, or when external factors (CI pipelines, scheduled operations) force execution regardless of conditions.

C1 establishes that the data is interventional. C2 and C3 determine whether the interventional signal can be cleanly extracted. All three are satisfied simultaneously in Regime A domains; they degrade together in Regime B and C domains.


\medskip\noindent\textbf{Three regimes.}


These conditions partition domains into admissibility regimes, paralleling \S\,\ref{ciy-observational-proxy}:

\begin{adjustbox}{max width=\textwidth}
\small
\begin{tabular}{lllll}
\toprule
Regime & C1 & C2 & C3 & Causal validity of leaf-edge updates \\
\midrule
\textbf{A: Intervention-rich} & Agent controls leaf actions & Good isolation (one action at a time) & Conditions vary (CI, diverse contexts) & \textbf{Strong.} Updates approximate interventional. \\
\textbf{B: Partial intervention} & Agent acts but with coordination constraints & Concurrent actions blur attribution & Self-selection likely & \textbf{Moderate.} Updates carry optimistic bias. \\
\textbf{C: Observation-only} & Agent did not act (condition leaves, passive monitoring) & Attribution impossible & Confounding dominant & \textbf{Weak.} Updates reflect association, not causation. \\
\bottomrule
\end{tabular}
\end{adjustbox}



\medskip\noindent\textbf{Indirect edges.}


For edges between non-leaf nodes, or edges whose parent is a condition node, the agent does not directly intervene on the parent. Instead, the agent's interventions at the leaves propagate upward through the DAG. The edge $(i, j)$ where $i$ is an internal node receives \textit{indirect} interventional evidence: the agent intervened on leaves below $i$, observed that $i$ was (or wasn't) achieved, and then observed whether $j$ advanced.

This indirect evidence is weaker for two reasons:
\begin{enumerate}
  \item \textbf{Compounding attribution}: the agent must attribute $j$'s outcome to edge $(i, j)$ after already attributing $i$'s achievement to the leaf-level interventions. Each attribution step introduces uncertainty.
  \item \textbf{Confounding from below}: $i$'s achievement depends on multiple leaf actions and condition states. Even if each leaf intervention is clean, their combined effect on $i$ may be confounded.
\end{enumerate}


The identification strength for indirect edges decreases with depth in the DAG — deeper edges are farther from the agent's direct interventions and have more confounding pathways.


\medskip\noindent\textbf{Identifiability-adjusted gain.}


% [Hypothesis (identifiability-coefficient)]

The update gain should be adjusted by the agent's confidence in causal attribution:

$$\eta_{\text{edge}}^{\text{adj}} = \eta_{\text{edge}} \cdot \iota_{ij}$$

where $\iota_{ij} \in [0, 1]$ is the \textbf{identifiability coefficient} — the agent's estimate of how cleanly the observed outcome can be attributed to edge $(i, j)$ specifically.

\begin{itemize}
  \item $\iota_{ij} = 1$: clean attribution (leaf-originating edge, single parent, isolated execution in Regime A).
  \item $\iota_{ij} \approx 0$: no attribution possible (deep internal edge, many concurrent causes, Regime C).
\end{itemize}


For leaf-originating edges in Regime A: $\iota_{ij} \approx 1$. For internal edges at depth $d$: $\iota_{ij}$ decreases with $d$ (each level of indirect inference degrades attribution). The precise functional form is domain-dependent.

This unifies two sources of frozen edges:
\begin{enumerate}
  \item Low \textbf{observability} ($\sigma_v \approx 0$ from \S\,\ref{observability-dominance}): the node's outcome is hard to \textit{measure}.
  \item Low \textbf{identifiability} ($\iota_{ij} \approx 0$): the outcome is measurable but can't be \textit{attributed} to this edge.
\end{enumerate}


Both drive $\eta_{\text{edge}} \to 0$ and produce the same effect: the edge is frozen at its prior.


\medskip\noindent\textbf{Epistemic Status.}


\textit{Conditional} on the three conditions C1–C3 and the DAG position of the edge. Max attainable: conditional (the conditions are genuine restrictions, not removable assumptions).

The restriction to leaf-originating edges for strong identification: \textbf{derived} from \S\,\ref{strategy-dag}'s node definitions. Only action-leaf nodes correspond to operations the agent directly performs. This is a structural property of the DAG, not an empirical claim.

The three-regime classification: \textbf{robust qualitative}. The regimes parallel the CIY admissibility regimes in \S\,\ref{ciy-observational-proxy}. The underlying logic is Pearl's interventional/observational distinction applied to strategy-edge updates.

The identifiability coefficient $\iota_{ij}$: \textbf{hypothesis}. The concept is sound (discounting by attribution confidence is standard in causal inference), but the specific form — a scalar multiplier on the gain — is a first-order approximation.

The depth-dependent degradation for indirect edges: \textbf{robust qualitative}. Each level of indirect inference adds an attribution step, and uncertainty compounds. The specific degradation rate is domain-dependent.


\medskip\noindent\textbf{Discussion.}


\textbf{Why this gap matters.} Without causal validity conditions, \S\,\ref{edge-update-via-gain}'s status is ambiguous: it might be updating credences that claim to be interventional using evidence that is merely associational. The conditions make explicit what's needed for the update to preserve the interventional semantics of \S\,\ref{strategy-dag}'s edge credences.

\textbf{Connection to \S\,\ref{observability-dominance}.} Observability gates whether the agent can \textit{detect} an outcome. Identifiability gates whether the agent can \textit{attribute} an outcome to a cause. Both are prerequisites for learning. The combined effective gain:

$$\eta_{\text{eff}} = \frac{U_{\text{edge}}}{U_{\text{edge}} + U_{\text{obs}}} \cdot \iota_{ij}$$

captures both gates in a single quantity. When either gate closes ($U_{\text{obs}} \to \infty$ or $\iota_{ij} \to 0$), effective gain goes to zero.

\textbf{Connection to the signal function.} The identifiability coefficient is one component of the unspecified signal function flagged in \S\,\ref{edge-update-via-gain}'s working notes. The full signal function $\text{signal}(o_t, i, j)$ decomposes into: (a) what outcome was observed ($o_t$), (b) how attributable is it to edge $(i, j)$ ($\iota_{ij}$), and (c) what does the attributed outcome imply about the edge's causal strength. This segment addresses (b); (a) and (c) remain open.

\textbf{Software as Regime A.} In software development, the agent runs a specific test (leaf action) and observes the result. C1 is satisfied (the agent ran the test), C2 is satisfied (the test targets one behavior), C3 is satisfied (CI runs regardless of conditions). Leaf-originating edges in software strategies have $\iota \approx 1$. This is the maximally favorable domain for causal edge updates.

\textbf{Optimal decomposition depth revisited.} \S\,\ref{observability-dominance} notes that finer decomposition (more intermediate nodes) provides earlier failure detection but adds uncertain edges via \S\,\ref{chain-confidence-decay}. This segment adds another cost of depth: deeper edges have lower $\iota_{ij}$, making them harder to learn causally. The optimal decomposition depth balances three factors: (a) observability of intermediate nodes, (b) confidence decay through chains, and (c) identifiability degradation with depth.

\end{scope}

\begin{discussion}[Credit Assignment Boundary]
\label{credit-assignment-boundary}
The strategy-revision loop requires assigning credit for observed outcomes to specific edges in the strategy DAG — decomposing "the plan partially worked" into "step 3 failed, step 5 was irrelevant, step 1 succeeded." This is ACT's version of RL's temporal credit assignment problem. This segment characterizes the boundary between tractable and intractable cases, states what the theory requires of any credit-assignment scheme, and identifies what the theory can guarantee without solving the problem at all.



\medskip\noindent\textbf{Formal Expression.}



\medskip\noindent\textbf{The Credit Assignment Problem for Strategy DAGs.}


% [Discussion (credit-assignment-problem)]

Given strategy DAG $\Sigma_t = (V, E, p, \gamma)$, an observed outcome at the root (and possibly at some intermediate nodes), produce per-edge signals $\text{signal}(o_t, i, j)$ for each edge $(i,j) \in E$ such that the edge update

$$p_{ij}^{\text{new}} = p_{ij} + \eta_{\text{edge}} \cdot (\text{signal}(o_t, i, j) - p_{ij})$$

drives credences toward truth. The problem is trivial when all intermediates are observable (each edge updates from its own observation) and genuinely hard when intermediates are unobservable and the DAG has shared structure.


\medskip\noindent\textbf{Default Signal Function (Gradient-Based Attribution).}



ACT's default signal function for edge updates, analogous to $\eta^* = U_M/(U_M + U_o)$ being the default gain:

$$\text{signal}_k(o_t) = p_k + \frac{J_k \cdot (y_G - \hat P_\Sigma)}{\lVert\mathbf{J}\rVert^2}$$

where $\mathbf{J} = \nabla_\mathbf{p} P_\Sigma$ is the plan-value gradient (computable from status propagation in $O(\lvert V\rvert + \lvert E\rvert)$), $y_G$ is the observed root outcome, and $\hat P_\Sigma$ is the current plan-confidence score. The signal attributes the plan-level residual $(y_G - \hat P_\Sigma)$ to each edge in proportion to its sensitivity $J_k = \partial P_\Sigma / \partial p_k$.

\textbf{Properties:}
\begin{itemize}
  \item \textbf{Directional fidelity (B1):} Satisfies $\mathbb{E}[\text{signal}_k - p_k] \propto J_k(\Phi - \hat P_\Sigma) \propto J_k \cdot \delta_s$. Since $J_k \geq 0$ for monotone AND/OR DAGs and $\delta_s = \Phi - \hat P_\Sigma$ has the correct sign, the expected signal pushes each edge toward truth.
  \item \textbf{Sector parameter:} $\alpha_s = \eta_{\text{edge}}$ for componentwise corrections (Prop B.5b); $\alpha_s = \eta_{\text{edge}} / \kappa(\mathbf{J})^2$ for coupled corrections.
  \item \textbf{Computational cost:} $O(\lvert V\rvert + \lvert E\rvert)$ — the same forward pass that computes $\hat P_\Sigma$ also yields $\mathbf{J}$.
  \item \textbf{Relationship to RL:} This is the ACT analog of REINFORCE — the Jacobian $\mathbf{J}$ is the score function, and the signal attributes the plan-level return to each edge via the score-weighted residual.
\end{itemize}


\textbf{Correlated-failure caveat.} The residual $(y_G - \hat P_\Sigma)$ decomposes into per-edge miscalibration \textit{plus} omitted correlation structure. When edge failures are correlated (shared infrastructure, common-mode risks — the dominant real-world case per \S\,\ref{strategy-dag}), $\hat P_\Sigma$ systematically overestimates success, making the residual systematically negative on failure. The gradient signal then blames individual edges for what is actually model-class misspecification (missing correlation nodes in the DAG). In these settings, the signal still has directional fidelity \textit{on average} (it pushes edges downward when the plan is overconfident, which is the correct direction), but the per-edge attribution is contaminated by the correlation structure the DAG doesn't represent. Agents in adversarial or common-mode-risk environments should treat gradient-based attribution as a coarse diagnostic and invest in making correlation structure explicit (adding common-cause nodes to $\Sigma_t$) rather than relying on per-edge signals to identify the problem.

Domains with richer observation structure can do better (Thompson sampling, full belief propagation, domain-specific attribution). The gradient-based signal is the \textit{concrete Level 1 default} — the minimum viable credit-assignment scheme that satisfies the theory's requirements.


\medskip\noindent\textbf{What the Theory Can Guarantee Without Solving Credit Assignment.}


Three results hold independently of any specific credit-assignment scheme:

\textbf{1. Persistence is credit-assignment-free.} Proposition B.5 in \S\,\ref{strategic-dynamics-derivation} shows that the sector condition transfers from per-edge credence space to the value-residual diagnostic $\delta_{\text{strategic}}$ via the Jacobian $\mathbf{J} = \nabla_\mathbf{p} P_\Sigma$. The Jacobian is computable from status propagation in $O(\lvert V\rvert + \lvert E\rvert)$ — no outcome decomposition required. The persistence guarantee (whether the strategy can be maintained) does not depend on the agent's ability to attribute outcomes to edges.

\textbf{2. The diagnostic framework is plan-level.} The satisfaction gap ( \S\,\ref{satisfaction-gap}), control regret ( \S\,\ref{control-regret}), and the orient cascade ordering ( \S\,\ref{orient-cascade}) operate on aggregate value, not per-edge quantities. They tell the agent \textit{whether} the strategy is failing (and whether the failure is feasibility vs. optimality vs. calibration), without requiring per-edge attribution.

\textbf{3. Observability-dominance identifies the tractable edges.} \S\,\ref{observability-dominance} determines which edges have nonzero observability — only these can receive informative signals. Edges with zero observability are frozen regardless of the credit-assignment scheme. The tractable boundary is the observable subgraph of $\Sigma_t$.


\medskip\noindent\textbf{The Tractable Cases.}


% [Discussion (tractable-credit-assignment)]

Credit assignment is solved (exact, polynomial-time) when:

\begin{adjustbox}{max width=\textwidth}
\small
\begin{tabular}{lll}
\toprule
Condition & Why tractable & Update rule \\
\midrule
\textbf{All intermediates observable} & Each edge has its own observation; updates decouple & Beta-Bernoulli per edge (Prop B.2) \\
\textbf{Binary outcomes, independent edges, linear chain} & Marginal Bayesian update = proportional blame & Prop B.3 (with plan-level fallback for unobservable) \\
\textbf{Tree DAG, observable leaves} & No shared descendants; message passing is exact & Belief propagation (standard) \\
\bottomrule
\end{tabular}
\end{adjustbox}



\medskip\noindent\textbf{The Intractable Cases: Three Independent Barriers.}


% [Discussion (intractable-credit-assignment)]

Exact per-edge attribution in general AND/OR DAGs with partial observability faces three independent barriers (full analysis in \texttt{msc/spike-credit-assignment-boundaries.md}):

\textbf{1. Computational intractability (\#P-hardness).} The "contribution of edge $k$ to the observed outcome" has the form of a Shapley value over a cooperative game defined by the AND/OR propagation. Since AND/OR DAGs can represent any monotone Boolean function (including weighted threshold functions), and Shapley value computation for weighted voting games is \#P-complete (Deng and Papadimitriou, 1994), exact attribution is \#P-hard. \textit{Caveat:} the reduction is to \textit{exact} computation; approximate Shapley values are computable in polynomial time with sampling.

\textbf{2. Information-theoretic underdetermination.} When intermediates are unobservable, per-edge attribution is \textit{underdetermined}, not just hard. The identifiable subspace has dimension bounded by the number of observable nodes:

$$\dim(\mathcal{I}(\mathcal{V}_{\text{obs}})) \leq \lvert\mathcal{V}_{\text{obs}}\rvert$$

When $\lvert\mathcal{V}_{\text{obs}}\rvert < \lvert E\rvert$ (fewer observable nodes than edges), some directions in $\boldsymbol\theta$-space are fundamentally unresolvable from the available data. Any attribution in the unidentifiable directions relies on prior beliefs, not evidence.

\textbf{3. The posterior correlation barrier.} Even for approximately identifiable cases, any factored representation (independent Beta posteriors per edge) necessarily discards the correlation introduced by failure at multi-parent nodes. The exact posterior complexity grows exponentially with the number of observed failures. The factored representation is an approximation by construction — coupled corrections are inherent to the problem, not an artifact of a bad algorithm.


\medskip\noindent\textbf{The Design Requirement.}


% [Discussion (credit-assignment-design-requirement)]

The theory does not prescribe a specific credit-assignment scheme. It states what any scheme must satisfy for the persistence guarantees to hold:

\textbf{Minimal requirement (from \S\,\ref{gain-sector-bridge}):} The per-edge signal function must have \textbf{directional fidelity} — the expected update for each edge must point toward the true credence:

$$\mathbb{E}[(\text{signal}(o_t, i, j) - p_{ij}) \cdot (p_{ij} - \theta_{ij})] \leq 0$$

(the expected correction is non-positively correlated with the current error). This is the per-component version of condition B1 from the bridge theorem. Any signal function satisfying this produces sector-satisfying corrections that transfer losslessly to value space (Prop B.5b, componentwise case).

\textbf{Sufficient condition for persistence:} Per-component directional fidelity + bounded gain ($\eta_{\text{edge}} > 0$). The theory guarantees persistence when these hold, regardless of how the signals are computed.

\textbf{What's NOT required:} Exact attribution, unbiased estimation, minimum-variance estimation, or optimality of any kind. The persistence guarantee is robust to approximation — a sloppy but directionally correct signal function still produces bounded strategic mismatch. The \textit{quality} of the approximation affects the \textit{tightness} of the persistence bound (how close $R^*_\Sigma$ is to zero), not whether persistence holds at all.


\medskip\noindent\textbf{The Hierarchy of Credit Assignment Quality.}


\begin{adjustbox}{max width=\textwidth}
\small
\begin{tabular}{llll}
\toprule
Level & Requirement & What it buys & Cost \\
\midrule
\textbf{0} (none) & Plan-level tracking only & Persistence guarantee (Prop B.5) & No per-edge diagnostics \\
\textbf{1} (directional) & Directional fidelity per edge & Persistence + rough per-edge diagnostics & Gradient computation $O(\lvert V\rvert + \lvert E\rvert)$ \\
\textbf{2} (approximate) & Proportional blame / expectation propagation & Persistence + per-edge diagnostics (with bias) & Factor-graph inference \\
\textbf{3} (exact) & Full Bayesian posterior & Persistence + optimal per-edge calibration & \#P-hard (general case) \\
\bottomrule
\end{tabular}
\end{adjustbox}


ACT's formal guarantees require only Level 0. Practical agents need at least Level 1 for adaptive behavior — and the default signal function (above) provides a concrete Level 1 scheme. Level 2 is the sweet spot for most applications. Level 3 is a mathematical ideal that is computationally unattainable in the general case.


\medskip\noindent\textbf{Epistemic Status.}


\textit{Mixed.} The default signal function (gradient-based attribution) is a \textit{formulation} — a concrete, well-motivated representational choice analogous to the gain principle's $\eta^*$ formula. It satisfies directional fidelity for monotone AND/OR DAGs (derivable from Jacobian non-negativity). The boundary characterization (tractable cases, intractability barriers, design requirement) is \textit{discussion-grade}, with the intractability argument at sketch level and the design requirement derived from the bridge theorem ( \S\,\ref{gain-sector-bridge}, \S\,\ref{strategic-dynamics-derivation} Prop B.5).

Max attainable: \textit{conditional} — with a formal intractability reduction, the boundary characterization could be promoted. The design requirement is already exact (it follows from the bridge theorem).


\medskip\noindent\textbf{Discussion.}


\textbf{ACT characterizes the structure, not the algorithm.} The theory's contribution to credit assignment is not a solution but a characterization: when it's trivial (observable intermediates), when it's hard (general partial observability), what any solution must satisfy (directional fidelity), and what guarantees hold without it (persistence, plan-level diagnostics). This is the right level of ambition for a theory of adaptive systems — the specific algorithm is domain-specific engineering; the structural characterization is universal.

\textbf{The analogy to the gain principle.} \S\,\ref{update-gain} characterizes the optimal gain structure ($\eta^* = U_M/(U_M + U_o)$) without prescribing how $U_M$ and $U_o$ are estimated. A Kalman filter computes them exactly; an RL agent approximates them; a human intuits them. The gain \textit{principle} is theory; the gain \textit{estimator} is engineering. Credit assignment has the same structure: the \textit{requirement} (directional fidelity) is theory; the \textit{implementation} (proportional blame, gradient attribution, belief propagation) is engineering.

\textbf{The observability lever.} The most powerful insight from this analysis is that credit assignment is primarily an \textit{observability design problem}, not an algorithm design problem. An agent that designs its strategy with observable intermediates (instrumented plans, staged rollouts, checkpoints) sidesteps the intractability entirely. This connects to \S\,\ref{observability-dominance}'s practical guidance: invest in making intermediate states observable, because unobservable regions are both epistemically dead (frozen credences) and computationally intractable (no efficient attribution).

\textbf{OKRs as observability-by-design.} The Objectives and Key Results framework is a direct organizational instantiation of this principle. The OKR discipline converts a deep, partially-unobservable strategy DAG (objective → vague initiatives → daily actions) into one with explicitly observable intermediate nodes (objective → measurable Key Results → tracked initiatives). In ACT terms: Key Results are intermediate nodes with $\sigma_v \approx 1$ by construction, making credit assignment between actions and objectives componentwise (Prop B.2) rather than intractable.

OKR failure modes map to ACT predictions:

\begin{adjustbox}{max width=\textwidth}
\small
\begin{tabular}{llll}
\toprule
OKR Failure & ACT Analog & Formal Quantity & Consequence \\
\midrule
\textbf{Vanity metrics} (measurable but irrelevant) & Observable node not causally connected to objective & High $\sigma_v$, low $p_{ij}$ & Edge updates, but the edge doesn't lead where the agent thinks — correction effort is wasted \\
\textbf{Too many Key Results} & Wide OR-node exploration-gating & $\alpha_\Sigma \propto 1/k$ (Prop B.4) & Correction capacity diluted across alternatives; persistence threshold harder to meet \\
\textbf{Lagging indicators} & Evidence starvation by delay & $\nu_{\text{obs}} \ll \rho$ & By the time the KR is measured, the correction window has passed; mismatch accumulated beyond $R_\Sigma$ \\
\textbf{Goodhart's Law} (metric becomes the goal) & Terminal-condition misalignment with $O_t$ & $V_{O_t}(\tau) < V_{O_t}^{\min}$ despite terminals achieved & Well-formedness constraint ( \S\,\ref{strategy-dag}) violated; agent optimizes the intermediate rather than the objective \\
\bottomrule
\end{tabular}
\end{adjustbox}


This is not an analogy — it is a domain instantiation. The same formal machinery that predicts when strategic persistence holds also predicts when OKRs work: when Key Results are genuinely causally connected to the Objective ($p_{ij}$ is high and calibrated), few enough to monitor effectively ($k$ is small), and measured on a timescale that permits correction (observation rate exceeds environment drift rate). The TST bridge to software team dynamics runs through exactly this connection.

\end{discussion}

\begin{formulation}[Structural Change as Parametric Limit]
\label{structural-change-as-parametric-limit}
In the probabilistic DAG, "structural" changes to $\Sigma_t$ are continuous operations on edge weights and node sets — not a separate mechanism. Pruning is a credence dropping below threshold; grafting is a new causal hypothesis initialized at a prior. This dissolves the sharp line between parametric update (adjusting weights) and structural change (adding/removing edges).



\medskip\noindent\textbf{Formal Expression.}



The six operations on $\Sigma_t$, ordered from most to least frequent:

\begin{adjustbox}{max width=\textwidth}
\small
\begin{tabular}{lll}
\toprule
Operation & What changes & Trigger \\
\midrule
Reweighting & Edge credence $p_{ij}$ & New observation about the link ( \S\,\ref{edge-update-via-gain}) \\
$\gamma$ reclassification & Node combination type AND↔OR & Strong structural evidence that combination semantics changed \\
Pruning & Remove failed branch ($p_{ij} \to \approx 0$) & Credence drops below viability threshold \\
Grafting & Add new branch ($0 \to p_{ij}$) & Discovery of a new possible path (initialized at prior) \\
Objective revision & Change terminal nodes & Feasibility failure or opportunity ( \S\,\ref{satisfaction-gap}) \\
Full restructure & Replace entire $\Sigma_t$ & Catastrophic failure ( \S\,\ref{structural-adaptation-necessity}) \\
\bottomrule
\end{tabular}
\end{adjustbox}


A healthy agent does continuous strategic maintenance (reweight, occasionally prune and graft) and rarely reaches catastrophic restructuring. Full restructure is the strategic analog of \S\,\ref{structural-adaptation-necessity}'s model-class change — the rare, expensive event when the entire representational structure must be replaced.


\medskip\noindent\textbf{Epistemic Status.}


\textit{Robust qualitative.} The continuity claim (structural change as parametric limit) is a property of the probabilistic representation: in a space where edges carry real-valued credences, adding or removing an edge is a boundary event, not a discontinuity. The frequency ordering of operations is an empirical pattern, not a derived result — it's consistent with the observation that deeper changes (restructuring > pruning > reweighting) require more evidence and are more costly.


\medskip\noindent\textbf{Discussion.}


\textbf{Connection to TF-10's destruction-creation.} \S\,\ref{structural-adaptation-necessity} describes the rare case when the entire model class must be replaced — the agent's representational framework is fundamentally inadequate. In the strategy DAG, this corresponds to full restructure: the causal theory encoded in $\Sigma_t$ is so wrong that incremental revision (reweighting, pruning) cannot fix it. The DAG must be rebuilt from a different starting point.

The key insight is that this extreme case is the \textit{limit} of the continuous process, not a separate mechanism. An agent that maintains healthy strategic hygiene (regular reweighting, timely pruning of failing branches, proactive grafting of alternatives) will rarely need full restructure. An agent that neglects maintenance — letting failing branches persist, ignoring negative evidence — accumulates strategic debt until catastrophic restructuring becomes unavoidable.

\textbf{Pruning and grafting thresholds.} When should the agent prune (remove an edge with very low credence) rather than just keeping it at low $p_{ij}$? The answer involves the cognitive cost of maintaining edges in $\Sigma_t$ — each edge consumes representational capacity and evaluation time. For agents with bounded representational capacity (LLMs with finite context windows), pruning is necessary to keep $\Sigma_t$ within capacity. The threshold is domain-dependent and connects to the open question of cognitive cost of $\Sigma_t$ ( \S\,\ref{strategy-dimension} Working Notes).

\end{formulation}

\begin{definition}[Strategic Tempo]
\label{strategic-tempo}
The effective rate at which an agent acquires useful revisions to its strategy $\Sigma_t$ --- the sum of per-edge correction capacities across the strategy DAG.



\medskip\noindent\textbf{Formal Expression.}



$$\mathcal T_\Sigma = \sum_{(i,j) \in E} \nu_{ij} \cdot \eta_{\text{edge},ij}$$

where:
\begin{itemize}
  \item $(i,j)$ indexes the edges of the strategy DAG ( \S\,\ref{strategy-dag})
  \item $\nu_{ij}$ is the effective observation rate for edge $(i,j)$ --- how often the agent obtains evidence about the causal link $i \to j$
  \item $\eta_{\text{edge},ij}$ is the per-edge update gain ( \S\,\ref{edge-update-via-gain})
\end{itemize}


\textbf{Parallel with epistemic tempo.} The definition mirrors \S\,\ref{adaptive-tempo}'s $\mathcal{T} = \sum_k \nu^{(k)} \cdot \eta^{(k)*}$, replacing observation channels with strategy edges and epistemic gain with edge gain. The structural parallel is exact: both are sums of (rate $\times$ quality) products over correction channels.

\textbf{Key difference: endogenous edge rates.} Epistemic tempo's channel rates $\nu^{(k)}$ are largely exogenous --- the environment generates observations at its own pace. Strategic tempo's edge rates $\nu_{ij}$ are \textit{endogenous}: they depend on the agent's action policy (which edges get tested) and on upstream success (downstream edges are tested only when upstream edges fire). This endogeneity is the source of the structural differences between epistemic and strategic persistence.


\medskip\noindent\textbf{Consistency verification.}


The definition is consistent with the four verified topologies from \S\,\ref{strategic-dynamics-derivation}:

**Case B.1 (single edge $A \to G$).** One edge, $\nu = \nu_{AG}$, $\eta_{\text{edge}} = 1/(n+1)$. $\mathcal T_\Sigma = \nu_{AG}/(n+1)$. The sector parameter $\alpha_\Sigma = 1/(n+1)$ is the per-observation correction quality; $\mathcal T_\Sigma = \nu \cdot \alpha_\Sigma$, matching the epistemic tempo pattern exactly.

**Case B.2 (two-edge AND chain, $A \to B \to G$, $B$ observable).** Two edges. Edge 1 is tested at rate $\nu_1 = \nu$ (every execution). Edge 2 is tested only when edge 1 succeeds: $\nu_2 = \nu \cdot \theta_1$. The bottleneck edge has $\alpha_\Sigma = \min(1/(n_1+1),\; \theta_1/(n_2+1))$. $\mathcal T_\Sigma = \nu/(n_1+1) + \nu\theta_1/(n_2+1)$, consistent with depth-gated attenuation.

**Case B.3 (two-edge AND chain, $B$ unobservable).** Per-edge tempo is ill-defined (the marginal point estimate is biased). Plan-level tempo is well-defined: $\mathcal T_{\Sigma,\text{plan}} = \nu/(n_\Phi + 1)$, treating $\hat\Phi = p_1 p_2$ as a single tracked quantity.

**Case B.4 (two-arm OR node, $\varepsilon$-greedy).** Edge 1 tested at rate $\nu_1 = \nu(1-\varepsilon)$, edge 2 at rate $\nu_2 = \nu\varepsilon$. $\mathcal T_\Sigma = \nu(1-\varepsilon)/(n_1+1) + \nu\varepsilon/(n_2+1)$. Action selection directly controls the rate allocation --- exploration-gated, not depth-gated.


\medskip\noindent\textbf{Structural decomposition.}


\textbf{AND-chains: depth-gated (geometric attenuation).} In a chain of depth $d$ with edge success probabilities $\theta_k$, the effective observation rate for edge $k$ is:

% [Derived (Conditional on independent edges)]

$$\nu_k = \nu \cdot \prod_{j < k} \theta_j$$

Each additional depth level attenuates by a factor $\theta_k < 1$. For a uniform chain ($\theta_k = \theta$, $n_k = n$ for all $k$):

$$\mathcal{T}_\Sigma = \frac{\nu}{n+1} \sum_{k=1}^{d} \theta^{k-1} = \frac{\nu}{n+1} \cdot \frac{1 - \theta^d}{1 - \theta}$$

This converges to $\nu / ((n+1)(1-\theta))$ as $d \to \infty$ --- total strategic tempo is bounded even for arbitrarily deep chains. The marginal tempo contribution of edge $k$ decays as $\theta^{k-1}$, falling below any fixed threshold at depth $d^*$ ( \S\,\ref{strategy-complexity-cost}). Deep AND-chains have low $\mathcal T_\Sigma$ at their leaves regardless of how fast the agent acts --- the evidence-starvation effect identified in \S\,\ref{strategic-dynamics-derivation}.

\textbf{OR-nodes: exploration-gated.} At an OR-node with $m$ alternatives under $\varepsilon$-exploration, the rate allocated to alternative $l$ is:


$$\nu_l = \begin{cases} \nu(1 - \varepsilon + \varepsilon/m) & l = l^* \text{ (greedy arm)} \\ \nu \cdot \varepsilon/m & l \neq l^* \text{ (exploratory arms)} \end{cases}$$

The bottleneck is the least-explored alternative. Pure greedy ($\varepsilon = 0$) gives $\nu_l = 0$ for non-greedy arms, making those edges permanently uncorrectable.


\medskip\noindent\textbf{Regime adjustment via identifiability.}


By \S\,\ref{edge-update-causal-validity}, the effective gain on each edge is modulated by the identifiability coefficient $\iota_{ij} \in [0, 1]$:


$$\mathcal T_\Sigma = \sum_{(i,j) \in E} \nu_{ij} \cdot \iota_{ij} \cdot \eta_{\text{edge},ij}$$

where $\iota_{ij} = 1$ for directly intervened edges with attributable outcomes (Regime A in \S\,\ref{edge-update-causal-validity}), $\iota_{ij} < 1$ for edges requiring observational proxy or partial identification, and $\iota_{ij} = 0$ for edges that are causally unidentifiable. This ensures the tempo accounts for the quality of causal evidence, not just its quantity.


\medskip\noindent\textbf{Per-edge persistence.}


% [Derived (from persistence-condition applied per edge)]

For $\Sigma_t$ to persist, every edge must maintain bounded mismatch. The bottleneck condition is:

$$\forall (i,j) \in E: \quad \nu_{ij} \cdot \iota_{ij} \cdot \eta_{\text{edge},ij} > \frac{\rho_{\Sigma,ij}}{R_{\Sigma,ij}}$$

This is the per-edge analog of \S\,\ref{per-dimension-persistence}'s per-dimension condition for $M_t$. The aggregate relationship between $\mathcal T_\Sigma$ and the average correction rate $\alpha_\Sigma$ is:

$$\alpha_\Sigma \leq \frac{\mathcal T_\Sigma}{\lvert E\rvert} \leq \mathcal T_\Sigma$$

(minimum $\leq$ average $\leq$ sum). Consequently, $\mathcal T_\Sigma > \lvert E\rvert \cdot \rho_\Sigma / R_\Sigma$ is \textit{necessary} for persistence but not sufficient --- the persistence condition is bottleneck-limited by the weakest edge, not governed by the aggregate.


\medskip\noindent\textbf{Epistemic Status.}


The definition itself is axiomatic --- it names a quantity by analogy with \S\,\ref{adaptive-tempo}. The consistency verification with the four cases from \S\,\ref{strategic-dynamics-derivation} is \textit{derived} (conditional on Beta-Bernoulli dynamics). The AND-chain geometric attenuation is \textit{derived} (conditional on independent edge outcomes). The OR-node exploration gating and identifiability adjustment are \textit{hypotheses} in the general DAG case, though verified for the specific topologies above. The bottleneck-limited persistence observation is \textit{derived} from \S\,\ref{per-dimension-persistence}'s result applied to the strategy substate.

Max attainable: conditional. Currently conditional because the general DAG case (mixed AND/OR topologies, correlated edges) has not been verified.


\medskip\noindent\textbf{Discussion.}


\textbf{Connection to \S\,\ref{per-dimension-persistence}.} The per-edge persistence condition inherits the same structure as the per-dimension epistemic result: the weak edge is the bottleneck. Scalar $\mathcal T_\Sigma$ overestimates effective strategic adaptation for the same reason scalar $\mathcal T$ overestimates epistemic adaptation --- it averages over heterogeneous correction capacities.

\textbf{Three-way tradeoff.} Strategic tempo competes with both epistemic tempo and exploitation for the agent's finite action capacity. Each action that tests a strategy edge (improving $\mathcal T_\Sigma$) is an action not spent gathering epistemic information (improving $\mathcal T$) or pursuing the current best action (exploitation). The allocation is addressed in \S\,\ref{exploit-explore-deliberate} — the extended deliberation threshold is derived, but the broader three-way framing is discussion-grade. Deliberation (internal computation) supplements $\mathcal T_\Sigma$ via an internal channel distinct from action-generated evidence.

\textbf{Software as Regime A.} In software development, tests are genuine interventions with attributable outcomes ($\iota_{ij} \approx 1$). This makes software agents naturally high-$\mathcal T_\Sigma$ --- a structural advantage identified but not yet formalized in the TST domain ( \texttt{02-tst-core/}).

\end{definition}

\begin{formulation}[Cognitive Cost of Strategy]
\label{strategy-complexity-cost}
The complexity cost of maintaining an explicit strategy $\Sigma_t$, formulated via minimum description length and the information bottleneck principle --- connecting DAG structure to the maintenance term $C_{\text{maintain}}$ in the explicit strategy condition ( \S\,\ref{explicit-strategy-condition}).



\medskip\noindent\textbf{Formal Expression.}



\medskip\noindent\textbf{Strategy description length.}



The minimum description length of a strategy DAG $\Sigma_t = (V, E, p, \gamma)$ ( \S\,\ref{strategy-dag}) decomposes as:

$$\operatorname{DL}(\Sigma_t) = \operatorname{DL}_{\text{struct}}(G) + \operatorname{DL}_{\text{param}}(p \mid G)$$

where:
\begin{itemize}
  \item $\operatorname{DL}_{\text{struct}}(G)$: bits to encode the DAG topology --- node identities, edge connectivity, AND/OR labels $\gamma$. Scales as $O(\lvert E\rvert \log \lvert V\rvert)$ for sparse DAGs.
  \item $\operatorname{DL}_{\text{param}}(p \mid G)$: bits to encode the edge credences given the topology. For Beta-distributed credences, each edge requires $O(\log n_{ij})$ bits where $n_{ij} = \alpha_{ij} + \beta_{ij}$ is the effective sample size.
\end{itemize}


The total scales as $O(\lvert E\rvert \log \lvert V\rvert)$ for moderate-precision credences, growing linearly in the number of edges and logarithmically in the number of nodes.


\medskip\noindent\textbf{Strategy IB objective.}



The optimal strategy complexity balances parsimony against decision-relevance, by analogy with \S\,\ref{information-bottleneck}'s compression-prediction trade-off for $M_t$:

$$\Sigma_t^* = \arg\min_{\Sigma_t} \left[\operatorname{DL}(\Sigma_t) - \beta_\Sigma \cdot I(\Sigma_t;\, \pi^* \mid M_t)\right]$$

where:
\begin{itemize}
  \item $\operatorname{DL}(\Sigma_t)$: description length (complexity cost)
  \item $I(\Sigma_t;\, \pi^* \mid M_t)$: mutual information between the strategy and the optimal policy, conditioned on the current model --- how much the strategy helps the agent choose good actions beyond what the model already provides
  \item $\beta_\Sigma > 0$: trade-off parameter analogous to \S\,\ref{information-bottleneck}'s $\beta$
\end{itemize}


When $\beta_\Sigma$ is low (high maintenance cost relative to decision value), the agent prefers simple strategies. When $\beta_\Sigma$ is high (strategy is cheap to maintain relative to its decision value), the agent can afford complex plans. The explicit strategy condition ( \S\,\ref{explicit-strategy-condition}) is the binary threshold: $\beta_\Sigma$ large enough that \textit{any} $\Sigma_t$ is worth maintaining.


\medskip\noindent\textbf{Maximum useful chain depth.}


% [Derived (Conditional on Beta-Bernoulli, per-edge persistence)]

From \S\,\ref{strategic-tempo}'s per-edge persistence condition, an AND-chain of depth $d$ with per-edge observation rate $\nu$, true success probability $\theta$ per edge, and effective sample size $n$ per edge persists only if the deepest edge satisfies:

$$\nu \cdot \theta^{d-1} \cdot \frac{1}{n+1} > \frac{\rho_\Sigma}{R_\Sigma}$$

Solving for the maximum depth at which persistence is achievable:

$$d^* = 1 + \left\lfloor \frac{\log\bigl(\frac{\nu}{(n+1)\rho_\Sigma / R_\Sigma}\bigr)}{\log(1/\theta)} \right\rfloor$$

When $\nu / ((n+1)\rho_\Sigma / R_\Sigma) \leq 1$, even $d = 1$ fails --- the agent cannot maintain a single edge under these conditions.

\textbf{Interpretation.} Beyond depth $d^*$, evidence starvation makes edges uncorrectable faster than the environment invalidates them. The agent accumulates strategic mismatch on deep edges regardless of how fast it acts at the top of the chain.

\textbf{Quantitative illustration} ($\theta = 0.8$, $\nu = 1$):

\begin{adjustbox}{max width=\textwidth}
\small
\begin{tabular}{lll}
\toprule
$n$ & $\rho_\Sigma / R_\Sigma$ & $d^*$ \\
\midrule
10 & 0.01 & 10 \\
10 & 0.1 & 0 \\
100 & 0.01 & 5 \\
100 & 0.1 & 0 \\
\bottomrule
\end{tabular}
\end{adjustbox}


High evidence requirements ($n$ large) and volatile environments ($\rho_\Sigma / R_\Sigma$ large) severely limit useful chain depth.


\medskip\noindent\textbf{Triple depth penalty.}


Deep AND-chains suffer three independent penalties that compound:

\begin{enumerate}
  \item \textbf{Confidence decay} ( \S\,\ref{chain-confidence-decay}): aggregate confidence $\prod p_k$ decays geometrically with depth. The plan is \textit{less likely to succeed}.
  \item \textbf{Evidence starvation} ( \S\,\ref{strategic-dynamics-derivation}): effective observation rate $\nu_k = \nu \cdot \prod_{j < k}\theta_j$ decays geometrically. The plan is \textit{harder to calibrate}.
  \item \textbf{Cognitive cost} (this segment): each additional depth level adds $O(\log \lvert V\rvert)$ bits to description length. The plan is \textit{more expensive to maintain}.
\end{enumerate}


All three are multiplicative in depth, making deep sequential strategies exponentially costly along three independent dimensions.


\medskip\noindent\textbf{Enriched explicit strategy condition.}



The maintenance cost $C_{\text{maintain}}$ from \S\,\ref{explicit-strategy-condition} decomposes as:

$$C_{\text{maintain}} = C_{\text{represent}} + C_{\text{revise}} + C_{\text{monitor}}$$

where:
\begin{itemize}
  \item $C_{\text{represent}} \propto \operatorname{DL}(\Sigma_t)$: cognitive cost of holding the strategy in working memory (proportional to description length)
  \item $C_{\text{revise}} \propto \sum_{(i,j)} \nu_{ij} \cdot c_{\text{update}}$: cost of processing edge updates (proportional to strategic tempo $\mathcal T_\Sigma$ times per-update cost)
  \item $C_{\text{monitor}} \propto \lvert\{(i,j) : \iota_{ij} < 1\}\rvert$: cost of monitoring edges with partial identifiability (the agent must do extra causal reasoning for non-trivial edges)
\end{itemize}


This decomposition makes the \S\,\ref{explicit-strategy-condition}'s maintenance term concrete: each component maps to a quantity defined elsewhere in the theory.


\medskip\noindent\textbf{Complexity compression operations.}


% [Discussion (complexity-compression)]

The IB objective suggests three compression operations, corresponding to structural changes from \S\,\ref{structural-change-as-parametric-limit}:

\begin{enumerate}
  \item \textbf{Edge pruning} (operation 3 in \S\,\ref{structural-change-as-parametric-limit}): remove edges with $\eta_{\text{edge},ij} \cdot I_{\text{edge},ij} < c_{\text{bit}}$, where $I_{\text{edge},ij}$ is the decision-relevance of that edge and $c_{\text{bit}}$ is the per-bit maintenance cost. Edges that contribute less decision value than their representational cost are candidates for removal.
  \item \textbf{Node merging} (reducing $\lvert V\rvert$): collapse intermediate nodes that serve no decision-distinguishing function. This reduces $\operatorname{DL}_{\text{struct}}$ by a factor proportional to the reduction in $\lvert V\rvert$.
  \item \textbf{Depth truncation} at $d^*$: prune all edges beyond the maximum useful depth. This is not optimization but necessity --- edges beyond $d^*$ cannot maintain bounded mismatch.
\end{enumerate}



\medskip\noindent\textbf{Epistemic Status.}


The description length formulation is a \textit{formulation} --- it applies standard MDL to the strategy DAG, which is a representational choice not a derived necessity. The IB objective is \textit{formulation/discussion-grade} --- it names the right trade-off by analogy with \S\,\ref{information-bottleneck} but the mutual information term $I(\Sigma_t;\, \pi^* \mid M_t)$ is not operationalized. The maximum useful depth $d^*$ is \textit{derived} conditional on Beta-Bernoulli dynamics and the per-edge persistence condition from \S\,\ref{strategic-tempo}. The triple depth penalty is an \textit{observation} combining results from three independent segments. The enriched maintenance decomposition is \textit{formulation}. The compression operations are \textit{discussion-grade}.

Max attainable: conditional. The DL formulation is standard; the IB objective and compression operations would require empirical validation or formal optimality proofs to advance beyond discussion-grade. The depth bound could reach exact status for specific edge models.


\medskip\noindent\textbf{Discussion.}


\textbf{Connection to \S\,\ref{explicit-strategy-condition}.} The enriched maintenance decomposition gives the cost inequality quantitative content: an agent can now \textit{compute} whether a proposed strategy is worth maintaining by estimating its description length, revision cost, and monitoring burden, and comparing against the exploration/repair cost of operating without it.

\textbf{LLM context windows as DL constraint.} For language-constituted agents ( \texttt{03-logogenic-agents/}), the strategy must fit in the context window. A context window of $W$ tokens imposes $\operatorname{DL}(\Sigma_t) \leq W \cdot \log_2(\lvert\text{vocab}\rvert)$ as a hard constraint. This makes the IB trade-off non-optional: the agent \textit{must} compress its strategy, and the depth bound $d^*$ becomes a context-window-limited quantity. A 128K-token context window may support a 500-edge DAG encoded in natural language; a 4K-token window may support only a 15-edge sketch.

\textbf{Strategy simplification pressure.} The triple depth penalty creates systematic pressure toward shallow, wide (OR-heavy) strategies over deep, sequential (AND-heavy) ones. This aligns with the structural pressure identified in \S\,\ref{chain-confidence-decay}'s Discussion, now grounded in three independent mechanisms rather than one.

\end{formulation}

\begin{remark}[Proposed-schema: Strategy Persistence Schema]
\label{strategy-persistence-schema}
The sector-condition mathematics ( \S\,\ref{sector-condition-stability}, \S\,\ref{sector-condition-derivation}) proves bounded mismatch for any system with: a mismatch state, a correction function satisfying sector bounds, and bounded disturbance. This mathematics is domain-agnostic — it doesn't care whether the mismatch is epistemic or strategic. If the strategic update dynamics can be shown to satisfy the same structural conditions, strategy persistence follows by the same result.



\medskip\noindent\textbf{Formal Expression.}


% [Proposed Schema (strategy-persistence-schema)]

\textbf{If} strategic update dynamics satisfy:
\begin{itemize}
  \item \textbf{(SA1)} Zero correction at zero strategic mismatch: when the mismatch state is zero, no revision occurs
  \item \textbf{(SA2')} Local sector condition on strategic correction: the correction function is appropriately bounded
  \item \textbf{(SA3)} Sufficient exploration (OR-nodes only): the action selection policy allocates correction capacity to all OR alternatives at a rate exceeding the strategic disturbance-to-reserve ratio
  \item Bounded strategic disturbance: the rate at which the environment invalidates causal links is bounded
\end{itemize}


\textbf{Then} $\Sigma_t$ persists iff:

$$\alpha_\Sigma > \frac{\rho_\Sigma}{R_\Sigma}$$

where $\alpha_\Sigma$ is the strategic correction rate, $\rho_\Sigma$ is the strategic disturbance rate, and $R_\Sigma$ is the strategic reserve (tolerance for strategic mismatch before performance degrades catastrophically).

\textbf{Which mismatch state?} The schema applies to any mismatch state for which conditions (SA1)-(SA3) can be verified. Two candidates exist:

\begin{itemize}
  \item \textbf{Plan-confidence error} $\delta_s = \hat P_\Sigma - \Phi$: the scalar difference between the agent's plan-confidence score and the independence-model plan value at true edge parameters. This is the mismatch for which persistence IS proved (Props B.1-B.5 in \S\,\ref{strategic-dynamics-derivation}). It is computable from status propagation without credit assignment.
  \item \textbf{Strategic-calibration residual} $\delta_{\text{strategic}}$: the per-edge value-increment residual aggregation defined in \S\,\ref{strategic-calibration}. This is the mismatch the orient cascade ( \S\,\ref{orient-cascade}) uses for edge-level revision. Persistence of $\delta_{\text{strategic}}$ remains \textbf{open} and requires the credit-assignment machinery in \S\,\ref{credit-assignment-boundary}.
\end{itemize}


The verified instances below all use per-edge credence error $\boldsymbol\delta_c = (\hat p_k - \theta_k)$ or the plan-level surrogate $\delta_s$. They do not verify the schema for $\delta_{\text{strategic}}$ directly.


\medskip\noindent\textbf{Epistemic Status.}


\textit{Sketch, with verified instances.} This is a \textbf{result schema}, not a proven result in the general case. The mathematical template (sector conditions → bounded mismatch) is derived ( \S\,\ref{sector-condition-derivation}). What was missing was instantiation — showing that specific strategic update dynamics satisfy the template's preconditions. Four cases have now been verified (full derivations in \S\,\ref{strategic-dynamics-derivation}):

\begin{enumerate}
  \item \textbf{Single edge, Beta-Bernoulli} ($A \to G$): Sector condition satisfied globally with $\alpha_\Sigma = 1/(n+1)$. The bound is tight (expected correction is exactly linear). (A1) satisfied. Persistence condition: $1/(n+1) > \rho_\Sigma / R_\Sigma$.
\end{enumerate}


\begin{enumerate}
  \item \textbf{Two-edge chain, observable intermediate} ($A \to B \to G$, $B$ observable): Sector condition satisfied globally with $\alpha_\Sigma = \min(1/(n_1+1), \theta_1/(n_2+1))$ — a weakest-link result. Correction function is diagonal (no cross-edge coupling). (A1) satisfied. The $\theta_1$ factor in edge 2's rate is the evidence-starvation effect.
\end{enumerate}


\begin{enumerate}
  \item \textbf{Two-edge chain, unobservable intermediate} ($A \to B \to G$, $B$ not observable): Per-edge sector condition \textbf{fails} — the marginal Bayesian update violates (A1) with bias $O(1/n)$. But plan-level tracking (treating $\hat{\Phi} = p_1 p_2$ as a single Beta) recovers the sector condition with $\alpha_{\Sigma,\text{plan}} = 1/(n_\Phi + 1)$, at the cost of per-edge diagnostic resolution.
\end{enumerate}


\begin{enumerate}
  \item **Two-arm OR-node, $\varepsilon$-greedy** ($A_1, A_2 \to G$, $G$ is OR): Sector condition satisfied with $\alpha_\Sigma = \min((1-\varepsilon)/(n_1+1),\; \varepsilon/(n_2+1))$ — an \textbf{exploration-gated} weakest-link, not depth-gated as in AND chains. (SA3) required: minimum exploration rate $\varepsilon > \rho_\Sigma(n_{\max}+1)/R_\Sigma$. Pure greedy ($\varepsilon = 0$) \textbf{fails} the sector condition. With optimal equal-rate allocation, $\alpha_\Sigma = 1/(n_1 + n_2 + 2)$ — the correction capacity is split across alternatives.
\end{enumerate}


The schema is no longer purely hypothetical. The sector parameter for strategic dynamics is the edge update gain $\eta_{\text{edge}}$ — the same quantity that governs epistemic persistence. The structural parallel between epistemic and strategic persistence is not an analogy but a mathematical identity at the sector-framework level.


\medskip\noindent\textbf{Discussion.}


\textbf{What's needed to promote this from schema to result.}

\begin{enumerate}
  \item \textbf{Strategic mismatch state}: partially resolved. Prop B.5 in \S\,\ref{strategic-dynamics-derivation} shows the sector condition transfers from per-edge credence error to \textbf{plan-confidence error} $\delta_s = \hat P_\Sigma - \Phi$ — the scalar difference between the agent's plan-confidence score and the independence-model plan value at true edge parameters (note: $\Phi$ is NOT actual plan success probability under correlated failure — see \S\,\ref{strategy-dag} edge-independence caveat). For linear correction (Beta-Bernoulli), the transfer is exact ($\alpha_s = \alpha_c$); for nonlinear correction, $\alpha_s \geq \alpha_c / \kappa(\mathbf{J})^2$. \textbf{However}, $\delta_s$ is distinct from the \textbf{strategic-calibration residual} $\delta_{\text{strategic}}$ defined in \S\,\ref{strategic-calibration}, which is an $L^2$ aggregation of per-edge value-increment residuals requiring credit assignment to compute. Persistence of $\delta_s$ (plan-level tracking) is proved; persistence of $\delta_{\text{strategic}}$ (per-edge diagnostics) remains open and requires the credit-assignment machinery in \S\,\ref{credit-assignment-boundary}.
\end{enumerate}


\begin{enumerate}
  \item ~~\textbf{Strategic correction function}: needs to satisfy the sector condition.~~ \textbf{Resolved} for Beta-Bernoulli edges. Props B.1-B.4 in \S\,\ref{strategic-dynamics-derivation} verify the sector condition for four topologies (single edge, two-edge AND observable/unobservable, two-arm OR).
\end{enumerate}


\begin{enumerate}
  \item \textbf{Strategic disturbance}: The rate at which the environment invalidates causal links in $\Sigma_t$. \textbf{Still open} as a formalized quantity — currently a domain parameter ($\rho_\Sigma$), analogous to how $\rho$ for epistemic disturbance is a domain parameter in \S\,\ref{persistence-condition}.
\end{enumerate}


\begin{enumerate}
  \item ~~\textbf{Sector condition verification}: the critical mathematical step.~~ \textbf{Resolved} for four topologies. See \S\,\ref{strategic-dynamics-derivation}.
\end{enumerate}


\begin{enumerate}
  \item ~~\textbf{Credit assignment / signal function}: needed for edge updates.~~ \textbf{Characterized at the theory level.} \S\,\ref{credit-assignment-boundary} shows persistence does not require credit assignment (Prop B.5), establishes directional fidelity as the minimal requirement, and provides a gradient-based default signal function. The specific update algorithm is domain engineering, not theory — the same way the gain \textit{estimator} is domain engineering while the gain \textit{principle} ($\eta^* = U_M/(U_M + U_o)$) is theory. Caveat: the default gradient signal inherits $\hat P_\Sigma$'s overestimation bias under correlated failures ( \S\,\ref{strategy-dag}, \S\,\ref{credit-assignment-boundary}).
\end{enumerate}


\textbf{The structural parallel is genuine.} This schema extends \textit{structural persistence} (see Persistence in \texttt{LEXICON.md}) from the epistemic substate to the strategy substate — asking whether the strategy correction machinery can outpace the rate at which the environment invalidates the agent's causal theory. It inherits the same limitation: structural persistence of $\Sigma_t$ does not address operational persistence (how close $\lVert\delta_{\text{strategic}}\rVert$ is to $R_\Sigma$) or continuity persistence (whether the agent's strategic identity coheres through time). The persistence condition for $M_t$ ( \S\,\ref{persistence-condition}) says: adaptive tempo must exceed the ratio of disturbance to critical mismatch. If the same mathematics applies to $\Sigma_t$, then strategy persistence requires strategic tempo to exceed the ratio of strategic disturbance to critical strategic mismatch. The strategic analog of "the environment changes faster than the agent can learn" is "the world invalidates plans faster than the agent can revise them." Both lead to the same catastrophic outcome: the system cannot maintain bounded mismatch and begins to degrade.

\textbf{What this would buy the theory.} If promoted to a result, strategy persistence would:
\begin{itemize}
  \item Provide a formal criterion for "when does a strategy remain viable?"
  \item Connect strategic failure modes to the same mathematical framework as epistemic failure modes
  \item Enable quantitative comparison: is the bottleneck epistemic persistence (model can't keep up with reality changes) or strategic persistence (plans can't keep up with requirement changes)?
  \item Ground the organizational intuition that plans need to be revised faster than the situation changes
\end{itemize}


\end{remark}

\begin{discussion}[Three-Way Resource Allocation]
\label{exploit-explore-deliberate}
At each decision point, an actuated agent with explicit strategy $\Sigma_t$ faces a three-way allocation of its finite cycle budget across exploit (take the currently-best action), explore (take an information-gathering action), and deliberate (pause acting to revise the model or strategy internally).



\medskip\noindent\textbf{Formal Expression.}



\medskip\noindent\textbf{The three activities.}



An actuated agent's cycle budget decomposes into three activities:

\begin{enumerate}
  \item \textbf{Exploit}: select and execute the action maximizing expected value under the current model and strategy. Earns immediate value $Q_O(M_t, a; \pi_{\text{cont}}, N_h)$.
\end{enumerate}


\begin{enumerate}
  \item \textbf{Explore}: select and execute an action maximizing causal information yield. Earns future model improvement via $\text{CIY}(a; M_t)$.
\end{enumerate}


\begin{enumerate}
  \item \textbf{Deliberate}: pause external action to run the orient cascade ( \S\,\ref{orient-cascade}) --- internal model refinement ($\Delta\eta^*$), strategy revision ($\Delta\alpha_\Sigma$), or objective reassessment. Earns improved future action quality.
\end{enumerate}


\textbf{The key distinction:} Exploit and explore both acquire new data from the environment. Deliberation acquires no new data --- it applies additional computation to existing data and model state. Deliberation's information ceiling is the gap between the total information already in the agent's observation history and the amount incorporated into $M_t$. For an agent with perfect inference, this gap is zero. Deliberation's value is therefore a property of the agent's computational limitations, not of the environment.


\medskip\noindent\textbf{Extended deliberation threshold.}


% [Derived (Conditional on GA-4, deliberation-drift assumption from CREF:deliberation-cost:CREF)]

Extending \S\,\ref{deliberation-cost} to incorporate strategic deliberation, the deliberation threshold becomes:

$$\Delta\tau^* = \arg\max_{\Delta\tau \geq 0} \left[\underbrace{\Delta\eta^*(\Delta\tau) \cdot \lVert\delta_{\text{post}}\rVert}_{\text{epistemic benefit}} + \underbrace{\Delta V_\Sigma(\Delta\tau)}_{\text{strategic benefit}} - \underbrace{\rho_{\text{delib}} \cdot \Delta\tau}_{\text{drift cost}}\right]$$

where:
\begin{itemize}
  \item $\Delta\eta^*(\Delta\tau)$ is the improvement in model update gain from $\Delta\tau$ of internal simulation ( \S\,\ref{deliberation-cost})
  \item $\lVert\delta_{\text{post}}\rVert$ is the post-deliberation mismatch magnitude
  \item $\Delta V_\Sigma(\Delta\tau)$ is the improvement in strategy value from deliberation: the expected increase in $V_O$ from revising $\Sigma_t$ during the pause
  \item $\rho_{\text{delib}}$ is the local mismatch drift rate during inaction (GA-4)
\end{itemize}


\textbf{First-order condition:}

$$\frac{\partial \Delta\eta^*}{\partial \Delta\tau} \cdot \lVert\delta_{\text{post}}\rVert + \frac{\partial \Delta V_\Sigma}{\partial \Delta\tau} = \rho_{\text{delib}}$$

Stop deliberating when the marginal joint improvement rate (epistemic plus strategic) drops below the mismatch drift rate.

\textbf{Reduction to existing results.} When $\Delta V_\Sigma = 0$ (no strategy revision benefit), this reduces to the deliberation threshold in \S\,\ref{deliberation-cost}. When $\Delta\tau^* = 0$ (deliberation never beneficial), the agent acts immediately.


\medskip\noindent\textbf{On the structure of the allocation.}


% [Discussion (allocation-structure)]

\textbf{Two-stage decomposition.} The three-way allocation can be decomposed into two nested decisions: (1) how long to deliberate before acting, (2) conditional on acting, which action to take (solved by \S\,\ref{ciy-unified-objective}). This decomposition is natural because exploit and explore are both external actions ($\in \mathcal{A}$), while deliberation is internal processing. However, this decomposition is a \textit{modeling convenience}, not a structural necessity. A unified objective over $\mathcal{A} \cup \{\text{deliberate}\}$ with temporal discount $\gamma$ is equally valid:

$$\bar{\pi}^* = \arg\max_{\bar{a} \in \mathcal{A} \cup \{\text{deliberate}\}} \bar{Q}(\bar{a}; M_t, \gamma)$$

where $\bar{Q}(\text{deliberate}) = \gamma \cdot \mathbb{E}[V_{\text{act}}(f_{\text{delib}}(M_t))]$ and $\bar{Q}(a) = r(a) + \gamma \cdot \mathbb{E}[V_{\text{act}}(M_{t+1})]$.

The two-stage decomposition is useful because it allows reusing the CIY-unified objective for the action-selection step, but it is not the unique or forced structure.

\textbf{Additive benefit decomposition.} The objective function above decomposes deliberation benefit into additive epistemic and strategic terms. This assumes the two benefits are approximately independent. Under directed separation ( \S\,\ref{directed-separation}), $M_t$ updates independently of $G_t$, which makes the decomposition structurally motivated for Class 1 agents. However, strategic benefit depends on the improved model ($\Delta V_\Sigma(\Delta\tau, M_t + \Delta M_t(\Delta\tau))$), creating an interaction that the additive form ignores. The additive form is a linearization that holds when $\Delta M_t$ is small relative to $M_t$. For Class 2 agents, epistemic and strategic deliberation are coupled, and the additive decomposition is a convenience.


\medskip\noindent\textbf{Control regret as deliberation ceiling.}


% [Discussion (deliberation-ceiling)]

When $\delta_{\text{regret}} \approx 0$ ( \S\,\ref{control-regret}), strategic deliberation has no value: the agent is already executing the best available strategy. This is the formal reason why low control regret suppresses deliberation. The ceiling on strategic deliberation value is $\delta_{\text{regret}}$ --- the most the agent could gain by finding a better strategy.

However, this ceiling applies only to the \textit{strategy-revision} channel of deliberation. Deliberation may also:
\begin{itemize}
  \item Reveal that the \textit{objective} should change (step 5 of the orient cascade), producing value not bounded by $\delta_{\text{regret}}$
  \item Reduce \textit{uncertainty about} strategy quality, enabling more confident exploitation even when the strategy doesn't change
\end{itemize}



\medskip\noindent\textbf{Qualitative regime descriptions.}


% [Discussion (qualitative-regimes)]

The three activities have natural dominance conditions that follow from the structure of the problem:

\textbf{Exploit dominates} when the strategy is near-optimal ($\delta_{\text{regret}} \approx 0$), model uncertainty is low ($\lambda(M_t) \approx 0$), and the environment penalizes pauses ($\rho_{\text{delib}}$ high). The agent knows what to do and must act now. This is Boyd's implicit guidance and control.

\textbf{Explore dominates} when model uncertainty is high ($U_M$ large), and the model needs correction via external evidence that internal simulation cannot provide. Exploration acquires new data; deliberation only processes existing data.

\textbf{Deliberate dominates} when the strategy needs revision ($\delta_{\text{regret}} \gg 0$ or $\delta_{\text{strategic}} \gg 0$), the environment is stable during pauses ($\rho_{\text{delib}}$ low), and the agent has unprocessed information or structural hypotheses it can evaluate internally. This regime requires that internal computation is cheaper or more efficient than external probing --- otherwise exploration dominates deliberation.

\textbf{Boundary conditions.} As $\rho_{\text{delib}} \to \infty$: $\Delta\tau^* \to 0$, and the allocation collapses to the binary exploit/explore tradeoff. As $\lambda \to 0$ and $\Delta V_\Sigma \to 0$: the allocation collapses to pure exploitation. As $\delta_{\text{regret}} \to 0$: strategic deliberation value vanishes, and the allocation reduces to \S\,\ref{deliberation-cost}'s epistemic think/act threshold combined with \S\,\ref{ciy-unified-objective}'s exploit/explore balance.

These regime descriptions are qualitative observations about the problem structure, not derived predictions. They restate what the constituent objectives already imply (high uncertainty favors information acquisition, low cost favors deliberation) without adding quantitative content. The theory does not provide the thresholds at which one regime transitions to another --- these are domain-specific.


\medskip\noindent\textbf{Epistemic Status.}


\textit{Discussion-grade.} The extended deliberation threshold is \textit{derived} (conditional on GA-4 and the deliberation-drift assumption from \S\,\ref{deliberation-cost}) and is the one genuinely formal contribution of this segment. The regime descriptions are \textit{discussion-grade} --- qualitatively correct observations that restate the problem structure. The two-stage decomposition is a \textit{formulation choice}, not a derived result. The additive benefit decomposition is a \textit{linearization} that is structurally motivated under directed separation but not derived.

**Simulation check (drifting multi-armed bandit, \texttt{msc/sim-three-way-tradeoff.py}):** The three-way allocation with perfect foresight outperformed binary exploit/explore by 2--6\% across five configurations. The oracle almost never chose deliberation (0--8 out of 200 steps). Exploration earned 3.5x more information per step than deliberation. A unified objective (single argmax with temporal discount) outperformed two-stage UCB by 8--13\%. The simulation is UNFAVORABLE to deliberation (no strategy structure, no irreversible actions, small action space), so these results are a lower bound on deliberation's value in complex domains. The simulation confirms that deliberation adds little in simple settings and that the two-stage decomposition is not necessary.

Max attainable: \textit{discussion-grade} for the overall framing. The extended deliberation threshold could reach \textit{conditional} (it inherits its status from \S\,\ref{deliberation-cost}). The regime descriptions are unlikely to advance beyond discussion-grade without domain-specific instantiation that provides quantitative thresholds.


\medskip\noindent\textbf{Discussion.}


\textbf{Why this is Section II content, not Section I.} Section I's deliberation-cost handles the binary think/act threshold using only $M_t$ and $\eta^*$. The three-way allocation is inherently a Section II result because it requires:
\begin{itemize}
  \item $\Sigma_t$ to exist (the strategy that deliberation can improve)
  \item $\delta_{\text{regret}}$ to be defined (the signal that motivates strategic deliberation)
  \item The orient cascade (the internal structure of what "deliberation" does for an actuated agent)
\end{itemize}


Without $G_t$, there is no strategic deliberation mode --- the agent can only improve $M_t$ or act, which is the Section I binary.

\textbf{Deliberation as computation, not a third activity.} The deepest distinction is not between three "activities" but between two resource types: \textit{data acquisition} (external action, covering both exploit and explore) and \textit{computation on existing data} (deliberation). Exploration and exploitation both involve acting in the environment and receiving observations; they differ in the objective governing action selection. Deliberation involves no new observations --- only additional processing of what the agent already has.

This reframing explains why deliberation has diminishing returns (the information in existing data is finite), why it depends on agent architecture (a perfect Bayesian has zero deliberation value), and why exploration generally dominates deliberation for information acquisition (new data beats reprocessing old data). Deliberation's comparative advantage is in domains where:
\begin{itemize}
  \item Internal computation is cheaper than external probing (bits per unit cost, not bits per timestep)
  \item The action space is too large for systematic exploration (planning/search)
  \item Actions are irreversible and error costs are high (deliberation as risk management)
  \item The model has latent relational structure exploitable by additional computation
\end{itemize}


\textbf{Deliberation as internal exploration.} \S\,\ref{deliberation-cost}'s Open Question 3 asks whether internal simulation can generate genuine CIY. If so, the boundary between "explore" and "deliberate" is permeable: an agent running mental simulations that surface model inconsistencies is performing a kind of exploration without external action. This makes deliberation relatively more attractive for agents with high-fidelity internal models and less attractive for agents whose internal models are too crude to generate meaningful surprises.

\textbf{Connection to \S\,\ref{strategy-complexity-cost}.} Deliberation duration required to revise $\Sigma_t$ grows with strategy complexity. The description length $\operatorname{DL}(\Sigma_t)$ from \S\,\ref{strategy-complexity-cost} provides a lower bound on the deliberation time needed for meaningful revision.

\textbf{Connection to \S\,\ref{strategic-tempo}.} Strategic tempo $\mathcal{T}_\Sigma$ is the rate of useful strategy revision from \textit{external} evidence. Deliberation provides an \textit{internal} channel for strategy revision that supplements $\mathcal{T}_\Sigma$.

\textbf{Domain instantiations:}

\begin{adjustbox}{max width=\textwidth}
\small
\begin{tabular}{llll}
\toprule
Domain & Exploit & Explore & Deliberate \\
\midrule
Military (Boyd) & Execute the mission & Reconnaissance, probing attacks & Staff planning, war-gaming \\
RL agent & Greedy action & $\varepsilon$-exploration, UCB & MCTS planning, model-based rollouts \\
Software developer & Write code along current plan & Try unfamiliar approach, spike & Architecture review, read docs \\
Organization & Execute strategy & A/B test, pilot program & Strategic planning offsite \\
AI coding agent & Execute planned edit & Read unfamiliar code, run tests & Plan approach, analyze requirements \\
\bottomrule
\end{tabular}
\end{adjustbox}


\end{discussion}

\section{III. Agentic Composites}

\begin{quote}\itshape
Scope: multiple agents interacting through a shared environment, or equivalently, the internal structure of composite agents. The composition postulate ( \#composition-consistency) requires that the theory apply at every level of description where the scope condition is met. This section develops what "applies at every level" means formally (the composition closure criterion, which requires an additional contraction assumption beyond Section I's sector condition), what happens when composition is imperfect, and what the dynamics of inter-agent interaction look like.
\end{quote}

\begin{quote}\itshape
Correlated observations as default; independence as the special case requiring justification. Adversarial dynamics are one end of a teleological unity spectrum, not a separate theory.
\end{quote}

\begin{quote}\itshape
Two sources: the simulation-validated adversarial dynamics from TFT (TF-11/Appendix F, track-b simulations), and the composition spike (\texttt{msc/spike-agent-composition.md}) which derives the requirement for composition consistency from the scope condition's level-independence.
\end{quote}

\begin{scope}[Multi-Agent Scope]
\label{multi-agent-scope}
Section III applies wherever multiple agents satisfying the scope condition interact through a shared environment. Each agent observes, acts, and faces uncertainty; their actions affect each other's environments. This is the general case for organizations, teams, ecosystems, and adversarial encounters. Independence (agents whose actions don't affect each other) is the special case requiring justification.



\medskip\noindent\textbf{Formal Expression.}


% [Scope (multi-agent-scope)]

A multi-agent system consists of $N$ agents $\{A_1, \ldots, A_n\}$, each satisfying the scope condition ( \S\,\ref{scope-condition}), interacting through a shared environment with state $\Omega_t \in \mathcal{S}_{env}$:

\begin{itemize}
  \item Each agent $A_i$ has state $X_t^{(i)} = (M_t^{(i)}, G_t^{(i)})$
  \item Each agent observes: $o_t^{(i)} = h^{(i)}(\Omega_t, a_t^{(\neg i)}, \xi_t^{(i)})$ — observations may depend on other agents' actions
  \item Each agent acts: $a_t^{(i)} = \pi^{(i)}(X_t^{(i)})$
  \item The environment evolves: $\Omega_{t+1} = T(\Omega_t, a_t^{(1)}, \ldots, a_t^{(n)}, \omega_t)$
\end{itemize}


The coupling is through the environment: agent $i$'s actions enter agent $j$'s observation function and the shared environment transition. Agents may also communicate directly (a special case of action-observation coupling with a dedicated channel).


\medskip\noindent\textbf{Observation decomposition and routing.}



Each agent's observation decomposes into environmental and inter-agent components:

$$o_t^{(i)} = \left(o_{\text{env},t}^{(i)},\; \{m_{ji,t}\}_{j \in \mathcal N_t(i)}\right)$$

where:
\begin{itemize}
  \item $o_{\text{env},t}^{(i)} = h_\text{env}^{(i)}(\Omega_t, \xi_t^{(i)})$: direct environmental observation (no inter-agent content)
  \item $\mathcal N_t(i) \subseteq \{1, \ldots, N\} \setminus \{i\}$: the \textbf{communication neighborhood} — which agents send messages to $i$ at time $t$
  \item $m_{ji,t} = c_t^{(j \to i)}(X_t^{(j)})$: message from $j$ to $i$, determined by the sender's full state and the communication protocol
\end{itemize}



The \textbf{routing structure} $R_t = (\mathcal N_t, \{c_t^{(j \to i)}\})$ specifies:
\begin{itemize}
  \item The \textbf{topology} $\mathcal N_t$: who communicates with whom
  \item The \textbf{protocol} $c_t^{(j \to i)}$: the rule governing what class of information flows from $j$ to $i$
\end{itemize}


Note: the protocol $c_t^{(j \to i)}$ is a \textit{rule} specifying the channel, not the specific content of any message. Individual messages reflect the sender's state $X_t^{(j)}$ — including their individual goals — through the sender's policy. What the routing structure governs is the \textit{infrastructure}: which channels exist and what kind of information they carry.


Routing is \textbf{goal-blind} when neither the topology nor the protocol depends on the composite's goal state:

$$\mathcal N_t \perp G_t^c \qquad \text{and} \qquad c_t^{(j \to i)} \perp G_t^c \quad \forall\, j, i$$

This means the communication infrastructure does not change based on what the composite is trying to achieve. Individual messages naturally reflect individual agents' goals through their policies — this is action, not routing. The routing condition is about the \textit{structure} of information flow, not the \textit{content} of individual messages.

\textbf{Goal-dependent routing} occurs when either the topology or the protocol varies with $G_t^c$. Examples: activating crisis-specific communication channels, changing intelligence-sharing protocols based on the current mission, reassigning reporting chains based on the operational objective.


\medskip\noindent\textbf{Epistemic Status.}


\textit{Axiomatic.} This is a scope definition — it describes the class of systems Section III addresses. The only substantive choice is that coupling goes through the shared environment rather than through direct state modification. This follows from the agent boundary assumption ( \S\,\ref{agent-environment}): agents affect each other by affecting the environment, not by directly altering each other's internal states.


\medskip\noindent\textbf{Discussion.}


\textbf{Correlated observations as default.} When agents share an environment, their observations are generically correlated — they see aspects of the same reality. Independence (uncorrelated observations) requires the agents to observe non-overlapping aspects of the environment, which is the special case. Most multi-agent settings of interest involve substantial observation correlation.

\textbf{The adversarial case is one end of a spectrum.} Agents whose objectives conflict are not a separate theory — they are multi-agent systems with negative teleological unity ( \S\,\ref{unity-dimensions}). The same composition and coordination machinery applies; the dynamics are just different (destabilizing rather than stabilizing).

\textbf{Inter-agent communication as a special observation channel.} Messages from $j$ to $i$ are actions by $j$ and observations by $i$. The routing structure formalizes the infrastructure: who talks to whom ($\mathcal N_t$) and under what protocol ($c_t^{(j \to i)}$). The sender controls the content (unlike passive environmental observation), which introduces strategic manipulation ( \S\,\ref{communication-gain}). The gain from inter-agent communication enters the distributed tempo ( \S\,\ref{communication-gain}, Working Notes).

\textbf{The routing/content distinction matters for directed separation.} Individual messages reflect senders' goals — that's just action through policy. The directed-separation question at the composite level ( \S\,\ref{directed-separation-under-composition}) is about the \textit{routing structure}: does the infrastructure change based on the composite's goals? Goal-blind routing preserves directed separation; goal-dependent routing breaks it.

\end{scope}

\begin{formulation}[Composition Closure Criterion]
\label{composition-closure}
We define a group of interacting agents as a valid composite macro-agent when its closed-loop dynamics approximately commute with coarse-graining — that is, when projecting micro-states to macro-states and then running macro-dynamics yields approximately the same result as running micro-dynamics and then projecting.



\medskip\noindent\textbf{Formal Expression.}


Let a system consist of $N$ sub-agents interacting in a shared environment with state space $\mathcal S_{\text{env}}$. The micro-state, micro-observations, and micro-actions are:

$$X_{\text{micro}, t} = \{ (M_{i,t}, G_{i,t}) \}_{i=1}^N \in \mathcal X_{\text{micro}}$$

$$o_{\text{micro}, t} = \{ o_{i,t} \}_{i=1}^N \in \mathcal O_{\text{micro}}$$

$$a_{\text{micro}, t} = \{ a_{i,t} \}_{i=1}^N \in \mathcal A_{\text{micro}}$$

The coupled micro-dynamics form an action-observation loop:

$$X_{\text{micro}, t} \xrightarrow{\pi_{\text{micro}}} a_{\text{micro}, t} \xrightarrow{E} (\Omega_{t+1}, o_{\text{micro}, t+1}) \xrightarrow{f_{\text{micro}}} X_{\text{micro}, t+1}$$

We constrain our search to an admissible class of projections $\Lambda \in \mathcal P_{\text{adm}}$ mapping micro to macro, and an admissible class of macro-dynamics $(\pi_c, E_c, f_c) \in \mathcal M_{\text{adm}}$:
\begin{itemize}
  \item $\Lambda_x : \mathcal X_{\text{micro}} \to \mathcal X_c = (M_c, G_c)$
  \item $\Lambda_o : \mathcal O_{\text{micro}} \to \mathcal O_c$
  \item $\Lambda_a : \mathcal A_{\text{micro}} \to \mathcal A_c$
  \item $\Lambda_\Omega : \mathcal S_{\text{env}} \to \mathcal S_{\text{env}, c}$
\end{itemize}


Let $\mathcal D_{\text{micro}}$ be the distribution of reachable trajectories generated entirely by the true micro-system over horizon $H$.

*[Definition (Composition Closure)]* We define the minimal achievable closure defect $\varepsilon^*$ over the admissible classes as:

$$ \varepsilon^* = \inf_{\Lambda \in \mathcal P_{\text{adm}},\, (\pi_c, E_c, f_c) \in \mathcal M_{\text{adm}}} \big\lVert (\varepsilon_x, \varepsilon_a, \varepsilon_o) \big\rVert $$

Where the expected component errors evaluated over true micro-trajectories $\tau \sim \mathcal D_{\text{micro}}$ are:
\begin{itemize}
  \item $\varepsilon_x = \mathbb E_\tau \Big[ \frac{1}{H} \sum_{t=1}^H \big\lVert \Lambda_x\big(f_{\text{micro}}(X_{\text{micro}, t}, o_{\text{micro}, t+1})\big) - f_c\big(\Lambda_x(X_{\text{micro}, t}), \Lambda_o(o_{\text{micro}, t+1})\big) \big\rVert_\mathcal{X} \Big]$
  \item $\varepsilon_a = \mathbb E_\tau \Big[ \frac{1}{H} \sum_{t=1}^H \big\lVert \Lambda_a\big(\pi_{\text{micro}}(X_{\text{micro}, t})\big) - \pi_c\big(\Lambda_x(X_{\text{micro}, t})\big) \big\rVert_\mathcal{A} \Big]$
  \item $\varepsilon_o = \mathbb E_\tau \Big[ \frac{1}{H} \sum_{t=1}^H \big\lVert \Lambda_o\big(E_{\text{obs}}(\Omega_t, a_{\text{micro}, t})\big) - E_{c, \text{obs}}\big(\Lambda_\Omega(\Omega_t), \Lambda_a(a_{\text{micro}, t})\big) \big\rVert_\mathcal{O} \Big]$
\end{itemize}


A set of agents forms a meaningful composite agent when $\varepsilon^* \leq \varepsilon_{\text{max}}$.


\medskip\noindent\textbf{Admissibility constraints on macro-dynamics.}



$\mathcal M_{\text{adm}}$ is the class of macro-dynamics $(\pi_c, E_c, f_c)$ satisfying:

\textbf{(A1) ACT agent structure.} The macro-state decomposes as $X_c = (M_c, G_c)$. The macro-update is recursive:

$$X_{c,t+1} = f_c(X_{c,t}, o_{c,t+1})$$

The macro-policy is state-dependent: $a_{c,t} = \pi_c(X_{c,t})$. This ensures the macro-agent has the same structural components as a single ACT agent ( \S\,\ref{agent-environment}, \S\,\ref{recursive-update}, \S\,\ref{action-selection}).

\textbf{(A2) Macro-mismatch.} A mismatch signal is well-defined:

$$\delta_{c,t} = o_{c,t} - \hat{o}_{c,t}(M_c, a_{c,t-1})$$

where $\hat{o}_{c,t}$ is the macro-model's prediction of the next macro-observation. This ensures the macro-agent can detect prediction errors — the foundation of adaptation ( \S\,\ref{mismatch-signal}).

\textbf{(A3) Macro-tempo.} The macro-update has well-defined adaptive tempo:

$$\mathcal T_c = \sum_k \nu_c^{(k)} \cdot \eta_c^{(k)*}$$

where $k$ indexes the macro-agent's observation channels, $\nu_c^{(k)}$ is the event rate, and $\eta_c^{(k)*}$ is the optimal gain ( \S\,\ref{adaptive-tempo}, \S\,\ref{update-gain}). This ensures the persistence condition is formulable at the macro level.

\textbf{(A4) Bounded macro-correction.} The macro-correction function satisfies the sector condition ( \S\,\ref{sector-condition-stability}):

$$\delta_c^T F_c(\mathcal T_c, \delta_c) \geq \alpha_c \lVert \delta_c \rVert^2 \quad \text{for } \lVert \delta_c \rVert \leq R_c$$

with $\alpha_c > 0$ (positive correction rate) and $R_c > 0$ (finite reserve). This ensures the macro-agent's corrections work — they reduce mismatch rather than amplifying it, within the macro-model's capacity.

\textbf{What (A1)-(A4) prevent.} Without these constraints, the infimum over $\mathcal M_{\text{adm}}$ is trivially zero (any dynamical system can curve-fit micro-trajectories over a finite horizon). The constraints force the macro-dynamics to be a genuine ACT agent — with decomposed state, prediction errors, adaptive tempo, and stable correction — not an arbitrary function approximator. $\varepsilon^*$ then measures the irreducible cost of representing multiple agents as one ACT agent.

\textbf{What (A1)-(A4) do NOT require.} Directed separation ( \S\,\ref{directed-separation}) is not part of the admissibility constraints. It is an additional structural property that some composites have and others don't ( \S\,\ref{directed-separation-under-composition}). Results that depend on directed separation declare it as an additional assumption. Similarly, strategy structure ($G_c = (O_c, \Sigma_c)$ with a DAG) is not required — simpler goal representations are admissible.


\medskip\noindent\textbf{Admissibility constraints on projections.}



The admissible class of projections $\mathcal P_{\text{adm}}(\epsilon_I, L)$ consists of projections $\Lambda = (\Lambda_x, \Lambda_o, \Lambda_a, \Lambda_\Omega)$ satisfying three conditions:

\textbf{(P1) Information preservation.} The projection retains at least a fraction $(1 - \epsilon_I)$ of the predictive mutual information:

$$I\big(\Lambda_x(X_{\text{micro},t});\; \Lambda_o(o_{\text{micro},t+1}) \mid \Lambda_a(a_{\text{micro},t})\big) \geq (1 - \epsilon_I) \cdot I\big(X_{\text{micro},t};\; o_{\text{micro},t+1} \mid a_{\text{micro},t}\big)$$

The left side is the predictive information the macro-state has about next macro-observations, conditioned on the macro-action. The right side is the same quantity at the micro level. The parameter $\epsilon_I \in [0, 1)$ controls how much predictive power may be sacrificed by projection. This is the information bottleneck (Tishby et al. 1999) applied to the projection map.

\textbf{(P2) Lipschitz continuity.} Each component of $\Lambda$ is Lipschitz-continuous with constant $L$:

$$\lVert \Lambda_x(X) - \Lambda_x(X') \rVert_{\mathcal X_c} \leq L \cdot \lVert X - X' \rVert_{\mathcal X_{\text{micro}}} \quad \forall\, X, X' \in \mathcal X_{\text{micro}}$$

(and analogously for $\Lambda_o, \Lambda_a, \Lambda_\Omega$, each with its own Lipschitz constant). Without Lipschitz regularity, bounded closure defect does not imply bounded trajectory error — the bridge lemma requires this or something equivalent. When the projection is $L$-Lipschitz, the trajectory error bound becomes $L \cdot \varepsilon^* / \alpha_c$.

\textbf{(P3) Dimensionality reduction.} The macro-state space has strictly lower dimension than the micro-state space:

$$\dim(\mathcal X_c) < \dim(\mathcal X_{\text{micro}})$$

This prevents the trivial identity projection, which achieves $\varepsilon^* = 0$ but is not a genuine abstraction.

\textbf{Why three conditions.} No single condition suffices alone. (P1) prevents degenerate projections (a constant map $\Lambda_x = c$ has zero mutual information). (P2) prevents pathological projections (discontinuous maps that amplify micro-errors into unbounded macro-errors). (P3) prevents trivial projections (the identity map satisfies (P1) and (P2) but achieves nothing). Together they constrain $\mathcal P_{\text{adm}}$ to projections that are informative, regular, and genuinely reductive.

**The $(\epsilon_I, L)$ parameters are part of the problem specification, not derived quantities.** The natural choice for many applications is $L = 1$ (non-expansive projection) and $\epsilon_I$ chosen as the minimum information loss compatible with genuine dimensionality reduction.

\textbf{Independence from (A1)-(A4).} The macro-dynamics admissibility (A1)-(A4) partially constrains the projection — (A1) requires $\Lambda_x$ to preserve the $M_c / G_c$ decomposition, and (A4) implicitly requires regularity through the sector condition. But (A1)-(A4) do not specify how much predictive information the projection must retain. A macro-agent with a very coarse projection can satisfy (A1)-(A4) while being a poor representation of the micro-system. The information-preservation condition (P1) fills this gap. See \texttt{msc/spike-projection-admissibility.md} §5 for the full independence analysis.


\medskip\noindent\textbf{Bridge lemma: closure defect to trajectory error.}


% [Derived (bridge-lemma, from sector-condition-derivation + A4 + contraction assumption)]

If the macro-dynamics satisfy (A4) \textbf{and} the sector condition on the correction function implies contraction of the full update map $f_c(\cdot, o)$ in its state argument, then bounded closure defect implies bounded trajectory error.

\textbf{Setup.} Let $\tilde X_{c,t} = \Lambda_x(X_{\text{micro},t})$ be the projected micro-state ("true" macro-state) and $X_{c,t}$ be the macro-state evolved by macro-dynamics. Define trajectory error $e_t = X_{c,t} - \tilde X_{c,t}$.

\textbf{Per-step error evolution.} At step $t$, decompose the next-step error:

$$e_{t+1} = \underbrace{f_c(\tilde X_{c,t} + e_t, o_{c,t+1}) - f_c(\tilde X_{c,t}, o_{c,t+1})}_{\text{propagation of accumulated error}} + \underbrace{f_c(\tilde X_{c,t}, o_{c,t+1}) - \tilde X_{c,t+1}}_{\text{new closure error } \leq \varepsilon_x}$$

The first term measures how the macro-update propagates existing error. \textbf{Contraction assumption:} The sector condition on the correction function (A4) implies that the full update map $f_c(\cdot, o)$ is contracting in its state argument — that is, the correction dominates the state update, so differences in state input shrink through the update. This holds when the correction is the primary mechanism by which the state responds to itself (as opposed to new information from $o$), which is the normal regime for adaptive agents. It can fail during structural adaptation when the state changes discontinuously, or when the observation-driven component of the update amplifies state differences. In discrete time, the sector condition with correction rate $\alpha_c$ and event rate $\nu_c$ gives a per-step contraction factor:

$$\lambda = 1 - \alpha_c / \nu_c \quad (< 1 \text{ when } \alpha_c < \nu_c)$$

The condition $\alpha_c < \nu_c$ means the agent doesn't fully correct in a single step — always true for realistic systems. Under this contraction:

$$\lVert e_{t+1} \rVert \leq \lambda \lVert e_t \rVert + \varepsilon_x$$

\textbf{Bound.} This is a standard linear recurrence. Starting from $e_0 = 0$:

$$\lVert e_t \rVert \leq \varepsilon_x \sum_{k=0}^{t-1} \lambda^k = \varepsilon_x \frac{1 - \lambda^t}{1 - \lambda}$$

As $t \to \infty$:

$$\limsup_{t \to \infty} \lVert e_t \rVert \leq \frac{\varepsilon_x}{1 - \lambda} = \frac{\varepsilon_x \nu_c}{\alpha_c}$$

Since the closure defect $\varepsilon^*$ is the per-step error and $\varepsilon^* \nu_c$ is the closure error rate (per-step error × steps per unit time), this bound has the same structure as $\rho / \alpha$ from \S\,\ref{persistence-condition} — a ratio of disturbance rate to correction rate.

\textbf{Condition for meaningful composition.} The bound must fit within the sector-condition region:

$$\frac{\varepsilon^* \nu_c}{\alpha_c} < R_c \quad \Longleftrightarrow \quad \varepsilon^* < \frac{\alpha_c R_c}{\nu_c}$$

This parallels the persistence condition: the closure error rate must not exceed the macro-agent's adaptive reserve. If it does, the macro-description diverges from micro-reality.

\textbf{Why this is the persistence condition in disguise.} The trajectory error $e_t$ evolves under the same dynamics as mismatch $\delta_t$ — contraction from the correction function, perturbation from an external source. For mismatch, the source is environmental disturbance $\rho$. For trajectory error, the source is closure defect $\varepsilon^* \nu_c$. The Lyapunov argument is identical (cf. Prop A.1 in \S\,\ref{sector-condition-derivation}); only the labels change. The sector condition (A4) does double duty.


\medskip\noindent\textbf{Epistemic Status.}


\textit{Conditional.} Max attainable: conditional (formulation choice).

The closure defect $\varepsilon^*$ is well-defined given (A1)-(A4), (P1)-(P3), and specified norms. The macro-dynamics admissibility (A1)-(A4) and the projection admissibility (P1)-(P3) are independent formulation choices — (A1)-(A4) specify what "ACT-shaped macro-dynamics" means; (P1)-(P3) specify what "informative, regular, genuinely reductive projections" means. Both are needed for $\varepsilon^*$ to be an infimum over a non-trivial, non-degenerate class. The projection admissibility conditions have a proposed definition with one exact instantiation (two-Kalman case, \texttt{msc/spike-projection-admissibility.md}); general-case computability of (P1) remains open.

The bridge lemma is \textit{conditional} — not fully derived from (A4) alone. The argument requires that the sector condition on the correction function implies contraction of the full update map $f_c(\cdot, o)$ in its state argument. This is an additional assumption beyond (A4), not a consequence of it. The gap: (A4) constrains only the correction component of $f_c$, but the full update also includes observation-driven terms that may amplify state differences. For estimation-type agents (Kalman), contraction is verified exactly. For general agents, contraction remains an independent assumption that must be checked per-domain. The discrete-time argument itself (standard linear recurrence, no continuous-time gap) is exact once contraction is granted.

The norm choices ($\lVert\cdot\rVert_\mathcal{X}$, $\lVert\cdot\rVert_\mathcal{A}$, $\lVert\cdot\rVert_\mathcal{O}$, and the combination norm) are load-bearing. For estimation-type agents, the Mahalanobis norm (weighted by inverse prediction-error covariance) is the natural choice — verified exactly in the two-Kalman case. For general domains, norm specification remains open.


\medskip\noindent\textbf{Discussion.}


This criterion replaces intuitive questions about "where the boundary of an agent is" with a functional test: does a macroscopic ACT description preserve the underlying micro-dynamics well enough to remain predictive and capable? The core requirement is an \textbf{approximate dynamical homomorphism} — the macro-dynamics approximately commute with the projection.

\textbf{Relationship to \S\,\ref{composition-consistency}.} The Section I postulate requires that ACT's machinery be scale-invariant — predictions at different levels of description must be compatible. This segment operationalizes "compatible" as "bounded closure defect under admissible coarse-graining." The admissibility constraints ensure the macro-description is genuinely ACT-shaped, so the same persistence condition, the same tempo framework, and the same mismatch dynamics apply at the macro level with macro-level parameters.

\textbf{The sector condition does double duty.} (A4) ensures the macro-agent's corrections work (structural persistence), AND it ensures the macro-description tracks micro-reality (bridge lemma). Both uses address \textit{structural persistence} in the sense of \texttt{LEXICON.md} — the capacity of the correction machinery, not the current operating point or the composite's identity through time. This is the central insight: an ACT agent that is structurally persistent in its own right also persists as a faithful representation of its constituents, because both require the same thing — contracting correction dynamics under bounded perturbation.

\textbf{Deriving composite (A4) from sub-agent properties.} If each sub-agent satisfies the sector condition with parameters $(\alpha_i, R_i)$, and coordination costs are bounded by $\Delta\mathcal T_i^{\text{cost}}$ per agent ( \S\,\ref{team-persistence}), then the composite satisfies (A4) with:

$$\alpha_c \geq \min_i (\alpha_i - \Delta\mathcal T_i^{\text{cost}})$$

$$R_c \leq \min_i R_i$$

This is a weakest-link bound (conservative but clean). Cooperative coupling ( \S\,\ref{team-persistence}) can improve $\alpha_c$ beyond this bound by reducing effective disturbance. The key implication: (A4) is \textit{verifiable} from micro-level properties — compute each sub-agent's sector-condition parameters, estimate coordination costs, and check whether the composite has positive correction rate. No need to compute $f_c$ directly. See \texttt{msc/working-composition-admissibility.md} §6.2 for the derivation.

\textbf{Connection to team-persistence.} \S\,\ref{team-persistence} derives persistence conditions for sub-agents in a cooperative-adversarial network. This segment provides the macro-level complement: the conditions under which the composite itself is a valid ACT agent. Together they close the loop: sub-agents persist individually (team-persistence) AND the composite is a meaningful macro-agent (composition closure with admissibility).

\textbf{Two-Kalman instantiation.} The simplest nontrivial worked case: two Kalman filters tracking correlated scalar random walks with correlation $\rho_{\text{corr}}$, no communication. The natural projection keeps the state estimates and discards the covariance states (means-only projection, dimension 2 from micro-dimension 4). At steady state, the closure defect is exactly $\varepsilon^* = 0$ for \textbf{all} values of $\rho_{\text{corr}}$ — the means-only projection perfectly represents the micro-dynamics because the Kalman gains converge to constants and the discarded covariance state carries no information. The "cost of independence" (the estimation performance lost by not exploiting cross-correlations) is a \textbf{performance gap} $\Delta_{\text{perf}} \approx 2\rho_{\text{corr}}^2 q^2 r / (S^*)^2$ (quadratic in $\rho_{\text{corr}}$ for small correlations), not a closure defect — it measures suboptimality relative to a joint filter, not failure to track the micro-dynamics. The composition-closure framework diagnoses representability ("is this group a coherent ACT agent?"), not optimality ("is this group as good as a centralized agent?"). See \texttt{msc/spike-composition-correlated-kalman.md} for the full derivation, (A1)-(A4) verification, and the first genuine $\varepsilon^* > 0$ case (purposeful agents with Beta-Bernoulli strategy updates). All three (P1)-(P3) conditions are verified exactly: $\epsilon_I = 0$ (no information loss at steady state, since the discarded covariance state is constant), $L = 1$ (the projection is a coordinate map), and $\dim(\mathcal X_c) = 2 < 4 = \dim(\mathcal X_{\text{micro}})$.

\textbf{Norm specification for estimation-type agents.} The two-Kalman case identifies the Mahalanobis norm as the natural choice for agents whose primary function is state estimation. The state norm weights by the inverse prediction-error covariance: $\lVert X_c - X_c' \rVert^2 = (\hat\omega_c - \hat\omega_c')^T (P_{\text{pred}}^*)^{-1} (\hat\omega_c - \hat\omega_c')$. The observation norm weights by the inverse innovation covariance $S^{-1} = (P_{\text{pred}}^* + R)^{-1}$. The general principle: norms should weight by inverse uncertainty, so that differences in well-estimated components count more than differences in poorly-estimated ones. This is the norm the Kalman filter implicitly uses — the Kalman gain minimizes the expected squared Mahalanobis distance from truth. For domains without a natural covariance structure (discrete states, non-Gaussian models), the Euclidean norm remains the default. See \texttt{msc/spike-projection-admissibility.md} §4 for derivation and §4.5 for the general-case pattern.

\end{formulation}

\begin{sketch}[Composite Tempo Inequality]
\label{tempo-composition}
The adaptive tempo of a composite agent is bounded from above by the sum of its sub-agents' tempos. The gap between aggregate potential and realized composite tempo is the coordination overhead — tempo consumed by internal reconciliation rather than external mismatch correction.



\medskip\noindent\textbf{Formal Expression.}


Let $\mathcal T_i$ be the adaptive tempo of sub-agent $i$ within a composite group of $N$ agents. Let $\mathcal T_c$ be the adaptive tempo of the composite macro-agent $A_c$ defined by an admissible coarse-graining $\Lambda$.

*[Derived (Sub-additive Tempo, sketch)]* For any composite agent $A_c$ with minimal closure defect $\varepsilon^* \geq 0$, the composite tempo is bounded by the sum of individual tempos:
$$ \mathcal{T}_c \leq \sum_{i=1}^N \mathcal{T}_i $$

*[Definition (Coordination Overhead)]* We define the \textbf{coordination overhead penalty} $C_{\text{coord}}$ as the difference between aggregate potential and realized macro-tempo:
$$ C_{\text{coord}} := \Big( \sum_{i=1}^N \mathcal{T}_i \Big) - \mathcal{T}_c $$


\medskip\noindent\textbf{Epistemic Status.}


\textit{Sketch.} Max attainable: exact (with a complete proof). The inequality is structurally motivated: composition cannot create corrective capacity out of nothing, and internal reconciliation consumes tempo that doesn't contribute to external mismatch correction. The $\varepsilon^* \to C_{\text{coord}}$ connection is sketched: closure defect as internal disturbance rate gives $C_{\text{coord}} \geq \varepsilon^* \nu_c$, and the inequality $\mathcal T_c \leq \sum \mathcal T_i$ follows. \textbf{Caveat:} this connection depends on the bridge lemma in \S\,\ref{composition-closure}, which requires a contraction assumption on the macro-update map — an assumption that is structurally motivated by (A4) but not formally derived from it (see \S\,\ref{composition-closure}, Epistemic Status). The "follows directly" language below should be read as "follows from the bridge lemma under its contraction assumption." The remaining gap: whether $\varepsilon^* \nu_c$ is a tight lower bound or whether additional coordination costs contribute beyond the closure-defect mechanism.


\medskip\noindent\textbf{Derivation sketch.}


\begin{enumerate}
  \item Tempo is a rate of mismatch correction ($\nu \cdot \eta^*_{\text{eff}}$).
  \item In the ideal case ($\varepsilon^* = 0$), every "tick" of an individual's ACT loop contributes directly to the macro ACT loop without friction. Thus $\mathcal T_c = \sum \mathcal T_i$ and $C_{\text{coord}} = 0$.
  \item When $\varepsilon^* > 0$, closure fails partially. Sub-agents spend a fraction of their individual tempo correcting \textit{internal} mismatches generated by other sub-agents (e.g., Alice changes an API, Bob must update his $M_t$ to understand it), or take actions that counteract each other.
  \item The macro-agent $A_c$ does not gain adaptive advantage against the \textit{external} environment from these internal reconciliation cycles. These cycles are consumed by $C_{\text{coord}}$.
\end{enumerate}



\medskip\noindent\textbf{The closure-defect to coordination-overhead connection.}


% [Derived (coordination-overhead-bound, from composition-closure bridge lemma)]

The closure defect $\varepsilon^*$ introduces internal mismatch at rate $\varepsilon^* \nu_c$ (per-step error × macro event rate). This internal mismatch must be corrected by the macro-agent's correction capacity — tempo spent on internal correction is tempo NOT spent on external mismatch correction.

The macro-agent faces total effective disturbance:

$$\rho_{\text{eff}} = \rho_{\text{ext}} + \varepsilon^* \nu_c$$

The fraction of macro-tempo consumed by closure correction:

$$C_{\text{coord}} \geq \varepsilon^* \nu_c$$

This is the coordination overhead: a lower bound equal to the closure error rate. When the macro-agent's correction of internal mismatch is less efficient than its correction of external mismatch (because internal mismatches may be harder to diagnose — they look like environmental change from the macro perspective), the actual overhead exceeds this bound.

\textbf{Consequences:}
\begin{itemize}
  \item The realized external tempo: $\mathcal T_c^{\text{ext}} = \mathcal T_c - C_{\text{coord}} \leq \sum \mathcal T_i - \varepsilon^* \nu_c$
  \item Adding an agent increases $\sum \mathcal T_i$ but may also increase $\varepsilon^*$ (larger composite → harder to approximate as a single macro-agent). Brooks's Law occurs when $\Delta\varepsilon^* \nu_c > \Delta\mathcal T_i$ — the closure defect increase from adding a member exceeds their tempo contribution.
  \item The persistence condition for the composite becomes: $\sum \mathcal T_i - \varepsilon^* \nu_c > \rho_{\text{ext}} / \lVert\delta_{\text{critical}}\rVert$ — explicitly separating internal overhead from external challenge.
\end{itemize}



\medskip\noindent\textbf{Discussion.}


\textbf{Equality conditions.} Strict equality ($\mathcal T_c = \sum \mathcal T_i$, i.e. $C_{\text{coord}} = 0$) requires:
\begin{enumerate}
  \item \textbf{Orthogonal Routing:} The information dependency DAG exactly matches organizational boundaries (no costly cross-talk required).
  \item \textbf{Perfect Shared Intent:} No tempo is lost to internal negotiation or conflicting pursuit of the macro-objective.
  \item \textbf{No Net Macro-Information Loss:} Observations may be redundant, but they do not result in wasted tempo after fusion, nor is critical macro-information lost during coordination.
\end{enumerate}


These are stated as sufficient conditions; whether they are also necessary is open.

\textbf{Brooks's Law.} Adding more agents (increasing $\sum \mathcal T_i$) only increases the composite tempo $\mathcal T_c$ if the corresponding increase in closure defect $\varepsilon^*$ doesn't let $C_{\text{coord}}$ dominate. The model provides a formal analog of Brooks's Law: if communication overhead ($C_{\text{coord}}$) grows faster than aggregate capability ($\sum \mathcal T_i$) when adding agents, then adding people to a late project makes it later. Whether this specific mechanism (coordination overhead consuming tempo) is the dominant cause in practice is an empirical question.

\textbf{Connection to \S\,\ref{team-persistence}.} The persistence condition for the composite agent requires $\mathcal T_c > \rho_{\text{ext}} / \Vert\delta_{\text{critical}}\Vert$. Since $\mathcal T_c = \sum \mathcal T_i - C_{\text{coord}}$, high coordination overhead can push the composite below the persistence threshold, causing it to disintegrate as a coherent entity — even though each sub-agent individually persists.

\end{sketch}

\begin{hypothesis}[Directed Separation Under Composition]
\label{directed-separation-under-composition}
When individual agents satisfy directed separation ( \S\,\ref{directed-separation}), does the composite macro-agent also satisfy it? The answer depends on whether the composite's internal information routing — which observations reach which sub-agents — is itself goal-dependent. Two cases arise, corresponding to the first two classes in \S\,\ref{directed-separation}'s architectural classification.



\medskip\noindent\textbf{Formal Expression.}


% [Hypothesis (directed-separation-under-composition, extending directed-separation to composites)]

Consider $N$ agents $A_1, \ldots, A_N$, each satisfying directed separation individually: $M_{i,\tau^+} = f_M^{(i)}(M_{i,\tau^-}, e_{i,\tau})$ with no $G_t^{(i)}$ argument.

\textbf{The question.} Directed separation is about \textbf{processing}, not \textbf{selection} ( \S\,\ref{directed-separation}, scope condition). A single agent's goals affect which events it seeks (through $\pi$), but $f_M$ processes whatever event arrives without reference to $G_t$. At the composite level, the analogous question is: does the composite's routing structure $R_t$ ( \S\,\ref{multi-agent-scope}) depend on the composite's goals $G_t^c$?

Note: sub-agents' goal-driven actions shape the environment, which other sub-agents observe. This is NOT a violation of directed separation — it is the same mechanism directed separation explicitly allows at the single-agent level: goals influence events through action, but processing of realized events is goal-blind. Agent $B$ observing environmental changes caused by $A$'s goal-driven behavior is $B$ processing a realized event goal-blindly. The observations carry information about $A$'s goals, but that is a property of the event's content, not a failure of goal-blind processing. Similarly, individual messages reflecting individual agents' goals is action through policy, not a routing-structure dependence on $G_t^c$.


\medskip\noindent\textbf{Case 1: Goal-blind routing.}


% [Hypothesis (Case 1)]

If the routing structure satisfies $R_t \perp G_t^c$ ( \S\,\ref{multi-agent-scope}, goal-blind routing) — neither the communication topology $\mathcal N_t$ nor the protocol $c_t^{(j \to i)}$ depends on the composite's goals — then:

\begin{itemize}
  \item Each sub-agent processes observations goal-blindly (individual directed separation)
  \item The routing is goal-blind (by construction)
  \item Therefore $f_M^c$ is $G_t^c$-independent
\end{itemize}


Directed separation \textbf{survives} at the composite level. Examples: military command structures with doctrinal communication protocols, software teams with defined code-review processes, multi-agent AI systems with protocol-specified message passing.


\medskip\noindent\textbf{Case 2: Goal-dependent routing.}


% [Hypothesis (Case 2)]

If $R_t$ depends on $G_t^c$ — either the topology $\mathcal N_t$ changes (different reporting chains activated depending on the mission) or the protocol $c_t^{(j \to i)}$ changes (different intelligence products shared depending on the objective) — then the composite's effective observation function has a goal argument:

$$o_c = h^c(\Omega, a_{\text{micro}}, G_t^c, \xi)$$

Even if each sub-agent's $f_M^{(i)}$ processes $o_i$ goal-blindly, the \textbf{set of observations reaching each sub-agent} depends on $G_t^c$ through the routing function. Directed separation \textbf{fails} at the composite level. The composite is analogous to a Class 2 (merged) architecture: goal content shapes the information pathway, not through individual interpretation but through collective routing.


\medskip\noindent\textbf{Epistemic Status.}


\textit{Conditional} on the routing structure of the composite. Max attainable: conditional (the two cases are genuine architectural alternatives). Previously discussion-grade; upgraded after routing formalization in \S\,\ref{multi-agent-scope} and architectural classification promotion in \S\,\ref{directed-separation} (2026-04-01).

\textbf{Foundation status.}

\begin{enumerate}
  \item The \textbf{architectural classification} (Class 1/2/3) is in \S\,\ref{directed-separation}'s Formal Expression with formal operationalization ($\kappa_{\text{processing}}$). Status: robust qualitative.
  \item The \textbf{routing structure} $R_t$ is defined in \S\,\ref{multi-agent-scope} with a formal goal-independence condition ($R_t \perp G_t^c$). The definition decomposes into topology independence and protocol independence. Whether this captures all relevant ways that composite information flow can depend on goals is an open question.
  \item The \textbf{admissible coarse-graining} $\Lambda$ from \S\,\ref{composition-closure} now has specified admissibility constraints: (A1)-(A4) for macro-dynamics and (P1)-(P3) for projections. The remaining gaps are validation (exercising the machinery on a purposeful multi-agent example) and computability (P1 requires conditional mutual information, tractable only for linear-Gaussian), not missing definitions.
\end{enumerate}


The logic of each case is sound given the routing definition. Case 1: goal-blind routing + goal-blind processing = goal-blind composite. Case 2: goal-dependent routing means the composite's observation function depends on $G_t^c$, regardless of individual processing. Both follow from directed separation's scope condition.

The claim that the architectural classification lifts to composition is \textbf{structurally motivated} by \S\,\ref{composition-consistency} (the theory must give consistent answers at every level of description). This is an argument from theoretical coherence, not a derivation.


\medskip\noindent\textbf{Discussion.}


\textbf{What this is and what it isn't.} This segment provides a two-case taxonomy for directed separation under composition. It is not yet a derivation — the remaining caveat (admissibility placeholders) prevents that. It IS a useful classification that identifies which composites fall within Section III's scope (Case 1) and which require a coupled formulation (Case 2).

\textbf{Most composites of interest are Case 1.} Fixed communication structures are the norm in designed multi-agent systems: military doctrine specifies information flow regardless of the current mission; software development processes define code review, CI, and deployment pipelines independent of the feature being built; multi-agent AI protocols specify message formats and channels. Goal-dependent routing (Case 2) is rarer: crisis management (where communication structure changes with threat type), ad hoc teams assembled for specific objectives, and attention-based multi-agent architectures where query routing depends on the shared goal.

\textbf{Connection to logogenic agents.} LLM-based agents are individually Class 2 (goal-conditioned attention). Whether composites of LLM agents are Case 1 or Case 2 depends on whether the inter-agent communication protocol is fixed or goal-dependent. A multi-agent LLM system with fixed API contracts between agents is Case 1 at the composite level even though each agent is individually Class 2. The internal architecture of each agent and the composition architecture are separate questions.

\textbf{Adversarial implications.} Goal-dependent routing (Case 2) makes the composite's objective partially observable through its communication patterns. An adversary who can observe which information is being routed where can infer the composite's current goal. This creates a meta-strategic incentive for fixed routing: it preserves not just epistemic hygiene but operational security.

\end{hypothesis}

\begin{discussion}[Unity Dimensions]
\label{unity-dimensions}
The quality of a composite agent's composition can be decomposed along four substantially independent dimensions: epistemic (shared model), teleological (shared objective), strategic (coordinated action), and perceptual (shared observations). These dimensions predict the closure defect $\varepsilon^*$ ( \S\,\ref{composition-closure}) — high unity along all four predicts low $\varepsilon^*$.



\medskip\noindent\textbf{Formal Expression.}



For a composite agent $A_c$ composed of sub-agents $\{A_1, \ldots, A_n\}$:

\textbf{Epistemic unity} $U_M$ — how much of the reality model is shared:

$$U_M = \frac{I(M_t^{(1)}; \ldots; M_t^{(n)})}{H(M_t^{(1)}, \ldots, M_t^{(n)})}$$

The fraction of total model information that is shared (multi-information / total-correlation ratio). $U_M = 1$ for identical models; $U_M = 0$ for independent models.

\textbf{Teleological unity} $U_O$ — how aligned are the objectives:

$$U_O^{(i,j)} = \text{corr}\!\left(V_{O_t^{(i)}}(\tau),\; V_{O_t^{(j)}}(\tau)\right)$$

over trajectories the composite encounters. $+1$ for identical objectives; $-1$ for perfectly opposed; $0$ for orthogonal. The composite teleological unity is an aggregation over all pairs. The scalar ranges from fully cooperative to fully adversarial per objective dimension.

\textbf{Strategic unity} $U_\Sigma$ — how coordinated is the joint policy:

% [Discussion]

$$U_\Sigma = 1 - \frac{D_{\text{KL}}(\pi^c_{\text{actual}} \Vert \pi^c_{\text{optimal}})}{D_{\text{KL}}(\pi^c_{\text{independent}} \Vert \pi^c_{\text{optimal}})}$$

where $\pi^c_{\text{optimal}}$ is the jointly optimal policy. $U_\Sigma = 1$ when actual matches optimal; $U_\Sigma = 0$ when actual matches independent (no coordination). Requires knowing the jointly optimal policy, which is itself a strong assumption.

\textbf{Perceptual unity} $U_{\text{obs}}$ — how much of the observation stream is shared:

The fraction of total observation information that reaches all sub-agents. Full perceptual unity means all agents observe the same signals; zero means private observations only. Enables epistemic convergence without explicit model-sharing.


\medskip\noindent\textbf{Epistemic Status.}


\textit{Discussion-grade.} Max attainable: empirical. The four dimensions are qualitatively motivated: they correspond to the four components of agent state ($M_t$, $O_t$, $\Sigma_t$, and the observation channel). The specific metrics proposed above are sketches — the information-theoretic formulations ($U_M$, $U_\Sigma$) are well-defined in principle but would require specifying distributions and distance measures for any practical computation. The claim that these dimensions are substantially independent is a hypothesis, not derived — epistemic unity may enable strategic unity (shared models allow coordination without explicit planning), so independence is approximate at best.


\medskip\noindent\textbf{Discussion.}


\textbf{Clausewitz's three gaps.} These dimensions map to the gaps identified by Clausewitz (systematized by Bungay in \textit{The Art of Action}):

\begin{adjustbox}{max width=\textwidth}
\small
\begin{tabular}{lll}
\toprule
Clausewitz Gap & Unity Dimension & Formal Quantity \\
\midrule
Knowledge gap & Epistemic unity ($U_M$) & $1 - U_M$: fraction of model not shared \\
Alignment gap & Teleological unity ($U_O$) & $1 - U_O$: objective misalignment \\
Effects gap & Strategic + Perceptual unity & $1 - U_\Sigma$ + observation routing costs \\
\bottomrule
\end{tabular}
\end{adjustbox}


The mapping is not perfect — Clausewitz's "effects gap" blends action coordination with observation feedback — but it provides 200+ years of organizational evidence for the qualitative decomposition.

\textbf{Connection to closure defect.} The unity dimensions serve as predictors of the component closure errors in \S\,\ref{composition-closure}: high $U_M$ predicts low $\varepsilon_x$ (state update error), high $U_O$ predicts low $\varepsilon_a$ (action error, since aligned goals produce compatible actions), high $U_{\text{obs}}$ predicts low $\varepsilon_o$ (observation error). The mapping from unity to closure error is not yet formalized.

\textbf{What each dimension's absence costs.}

\begin{itemize}
  \item Low $U_M$: prediction conflicts → uncoordinated actions based on contradictory beliefs. Internal mismatch component from model disagreement.
  \item Low $U_O$: strategic friction → sub-agents pursue conflicting sub-goals. Effort wasted or counterproductive.
  \item Low $U_\Sigma$: redundancy and gaps → two agents fix the same bug while a critical one goes unnoticed.
  \item Low $U_{\text{obs}}$: information silos → critical signals observed by one agent but not actionable by the composite.
\end{itemize}


\end{discussion}

\begin{definition}[Shared Intent]
\label{shared-intent}
When sub-agents within a composite must coordinate, they face a communication problem: transmitting the full objective $O_t$ and strategy $\Sigma_t$ is expensive (high bandwidth, high latency), but acting without any shared purpose wastes coordination potential. The Information Bottleneck ( \S\,\ref{information-bottleneck}) applied to inter-agent communication predicts an optimal compression: transmit enough of $G_t$ to align behavior, not more.



\medskip\noindent\textbf{Formal Expression.}



Let $G_t^{\text{full}}$ be the source agent's complete purposeful state $(O_t, \Sigma_t)$. Let $G_t^{\text{shared}}$ be the compressed representation communicated to partners. The shared intent is the IB-optimal compression:

$$G_t^{\text{shared}} = \arg\min_{G_s} \left[ I(G_t^{\text{full}}; G_s) - \beta \cdot I(G_s; a_t^{\text{coordinated}}) \right]$$

where $a_t^{\text{coordinated}}$ is the jointly optimal action and $\beta$ controls the complexity-relevance tradeoff. At high $\beta$, the agent communicates more detail (approaching full model sharing). At low $\beta$, communication is minimal (approaching independent action).

The shared intent is the \textit{minimal sufficient statistic} of the sender's purposeful state for predicting the jointly optimal coordination behavior.


\medskip\noindent\textbf{Epistemic Status.}


\textit{Discussion-grade.} Max attainable: conditional (conditional on the IB framework being appropriate for inter-agent communication). The application of IB to inter-agent communication is structurally motivated — IB compresses optimally given a relevance criterion, and coordination relevance is the natural criterion — but the specific formulation assumes: (1) the sender knows the jointly optimal action (which requires knowing other agents' states), (2) the compression is lossless in the IB sense (real communication introduces noise, delay, and misinterpretation), (3) the $\beta$ parameter is fixed rather than dynamically adjusted. These are strong assumptions. The qualitative prediction (communicate purpose before plans before models) is more robust than the specific IB formulation.


\medskip\noindent\textbf{Discussion.}


\textbf{What gets compressed out.} The IB compression preferentially preserves:
\begin{enumerate}
  \item Terminal objectives (what the agent is trying to achieve) — these are compact and change slowly
  \item High-level strategy (which approach, not which specific steps) — moderate size, moderate change rate
  \item Strategic details (specific edge credences in $\Sigma_t$) — large, change fast, low coordination value
\end{enumerate}


**Connection to cognitive cost of $\Sigma_t$.** For agents with bounded communication capacity (bandwidth-limited channels, finite context windows), the DAG must be summarized for transmission. A 500-node strategy DAG cannot be shared in full; the IB compression identifies which structural features of the DAG matter for coordination.

\textbf{Organizational communication patterns.} Commander's intent in military doctrine is an empirical instantiation: the commander communicates \textit{what} to achieve and \textit{why}, not \textit{how}. This is IB-optimal if objectives change slowly (low $\nu_O$) and strategies change fast (high $\nu_\Sigma$) — communicating objectives gives a long shelf life per bit transmitted.

\end{definition}

\begin{hypothesis}[Auftragstaktik Principle]
\label{auftragstaktik-principle}
For a composite agent with limited communication bandwidth, the optimal allocation prioritizes sharing objectives over strategies over models. This captures the structural insight of Auftragstaktik (mission-type tactics): investing communication bandwidth in shared purpose (teleological unity) while accepting lower epistemic and strategic unity, granting sub-agents autonomy to adapt locally. The model predicts the same priority ordering that military doctrine discovered empirically; whether the mechanism (IB-optimal bandwidth allocation) is the actual reason Auftragstaktik works is an open question.



\medskip\noindent\textbf{Formal Expression.}


% [Hypothesis (auftragstaktik-principle)]

Let a composite agent's total inter-agent communication bandwidth be $B = B_O + B_\Sigma + B_M$, allocated across objective sharing ($B_O$), strategy coordination ($B_\Sigma$), and model synchronization ($B_M$).

The hypothesis: the allocation that maximizes composite tempo $\mathcal{T}_c$ (or equivalently, minimizes coordination overhead $C_{\text{coord}}$) prioritizes:

$$B_O > B_\Sigma > B_M$$

when:
\begin{itemize}
  \item Objectives change slowly relative to strategies: $\nu_O \ll \nu_\Sigma$
  \item Strategies change slowly relative to models: $\nu_\Sigma \ll \nu_M$
  \item Sub-agents have sufficient local adaptive capacity: each $\mathcal{T}_i > \rho_i^{\text{local}} / \Vert\delta_{\text{critical}}^i\Vert$
\end{itemize}


The priority ordering follows from the IB framework ( \S\,\ref{shared-intent}): the bits with the longest shelf life and highest coordination value per bit should be transmitted first. Objectives change slowly and enable autonomous coordination (sub-agents who share objectives can independently choose compatible strategies). Models change fast and provide diminishing coordination value (two agents with the same model but different objectives still conflict).


\medskip\noindent\textbf{Epistemic Status.}


\textit{Discussion-grade.} Max attainable: empirical. The priority ordering is a qualitative prediction grounded in the IB framework and supported by extensive military-organizational evidence (Bungay's analysis of Clausewitz, Wehrmacht doctrine, modern mission command). But it is not derived — the IB optimization would need to be solved explicitly with realistic cost functions to confirm the ordering, and the conditions under which the ordering reverses (e.g., when model synchronization is critical because the situation is genuinely ambiguous) are not characterized. The empirical evidence is strong but comes primarily from one domain (military command); generalization to software teams, AI agent swarms, and other settings is plausible but unverified.


\medskip\noindent\textbf{Discussion.}


\textbf{When the ordering reverses.} The prioritization $B_O > B_\Sigma > B_M$ assumes sub-agents can independently construct adequate local models. When the environment is genuinely ambiguous and local observations are insufficient (fog of war, novel codebase, unprecedented market conditions), model synchronization may be worth more than objective sharing — sub-agents who share the same wrong model at least err consistently, which is sometimes better than each having a different wrong model.

\textbf{Bungay's evidence.} In \textit{The Art of Action}, Bungay documents that organizations consistently fail by inverting this priority: they over-invest in controlling \textit{how} subordinates act (strategy sharing, $B_\Sigma$) rather than ensuring subordinates understand \textit{why} (objective sharing, $B_O$). The result: subordinates who follow instructions precisely but cannot adapt when conditions change, because they lack the teleological context to improvise intelligently.

\textbf{The software team instantiation.} A well-functioning development team has:
\begin{itemize}
  \item High $B_O$: clear product goals, understood by all (sprint goals, feature objectives)
  \item Moderate $B_\Sigma$: architectural decisions shared, implementation details autonomous
  \item Low $B_M$: each developer builds their own mental model of the code they touch; full codebase understanding is neither expected nor efficient
\end{itemize}


When this inverts (micromanagement of implementation details, unclear product goals), team tempo drops — consistent with the Auftragstaktik prediction.

\textbf{Connection to Conway's Law.} Conway's Law (system structure mirrors communication structure) is a consequence: when $B_\Sigma$ is low and $B_O$ is high, sub-agents coordinate through shared objectives rather than explicit action coordination, producing systems whose boundaries reflect objective decomposition rather than communication channels.

\end{hypothesis}

\begin{proposition}[Team Persistence]
\label{team-persistence}
Teams persist where individuals cannot through two physically distinct cooperative mechanisms: communication (allies share observations that improve correction) and action (allies act in the shared environment to reduce disturbance at its source). These mechanisms enter the persistence condition at different points — tempo and disturbance respectively — and a given cooperative interaction contributes through one mechanism or the other, not both.



\medskip\noindent\textbf{Formal Expression.}



\medskip\noindent\textbf{Distributed Tempo.}



Agent $i$'s effective tempo includes contributions from both direct observation and communication from allies:

$$\mathcal{T}_i = \underbrace{\sum_k \nu_i^{(k)} \eta_i^{(k)*}}_{\text{direct observation tempo}} + \underbrace{\sum_{j \in \mathcal{N}(i)} \nu_{ji}^{\text{comm}} \, \eta_{ji}^*}_{\text{communication tempo}}$$

where $\nu_{ji}^{\text{comm}}$ is the rate of communication events from agent $j$ to agent $i$, and $\eta_{ji}^*$ is the communication gain ( \S\,\ref{communication-gain}). Faster team adaptation comes not only from faster individual sensing but from faster, more reliable knowledge transfer.

\textbf{Channel independence caveat.} Both sums are additive, inheriting the channel-independence assumption from \S\,\ref{adaptive-tempo}: each channel and each communication source contributes non-redundant correction capacity. When allies report correlated information (overlapping observations, shared intelligence sources, redundant status reports), the communication tempo overcounts. The additive formula is an upper bound; the redundancy penalty depends on the mutual information between communication sources. See \S\,\ref{adaptive-tempo} for the single-agent version of this caveat.


\medskip\noindent\textbf{Cooperative-Adversarial Disturbance Decomposition.}



The disturbance rate experienced by agent $i$ decomposes into environment, adversarial, and cooperative components:

$$\rho_i = \rho_{i,\text{env}} + \sum_{j \in \mathcal{A}_i} \gamma_{j \to i}^{\text{adv}} \, \mathcal{T}_j - \sum_{j \in \mathcal{C}_i} \gamma_{j \to i}^{\text{coop}} \, \mathcal{T}_j$$

where $\mathcal A_i$ is the set of agents adversarially coupled to $i$, $\mathcal C_i$ is the set cooperatively coupled, and the $\gamma$ coefficients capture coupling effectiveness (as in \S\,\ref{adversarial-destabilization}).

\textbf{The cooperative term is negative — but through action, not communication.} Allies reduce agent $i$'s effective disturbance by \textit{acting in the shared environment} to prevent or mitigate disturbance at its source. Examples: an ally stabilizes a shared resource, neutralizes a threat before it reaches $i$, or absorbs environmental variation through their own actions. The mechanism is causal coupling through the shared environment, not information transfer.

\textbf{Separation from communication tempo.} The communication tempo term (above) captures allies \textit{telling} agent $i$ things that improve its correction. The cooperative disturbance term captures allies \textit{doing} things that reduce the disturbance $i$ faces. These are physically distinct: communication improves $i$'s correction function (better $\alpha_i$, higher $\mathcal{T}_i$); cooperative action reduces the disturbance $\rho_i$ that $i$ must correct against. A single cooperative event contributes through one channel or the other. An ally's message about a threat enters through communication tempo; an ally's action that eliminates the threat enters through disturbance reduction. Counting a single event in both terms would double-count the benefit and make the persistence threshold systematically optimistic.

\textbf{Effective disturbance rate.} The decomposition can yield $\rho_i < 0$ when cooperative coupling dominates both environment disturbance and adversarial coupling. The sector-condition analysis ( \S\,\ref{sector-condition-stability}) assumes non-negative disturbance (GA-2). Define:


$$\rho_i^{\text{eff}} = \max(\rho_i, \, 0)$$

When $\rho_i^{\text{eff}} = 0$, the agent's cooperative network fully absorbs all disturbance — the persistence condition is trivially satisfied and mismatch decays to zero. This is an idealized limit; in practice, $\rho_i^{\text{eff}} > 0$ because cooperative coupling is imperfect and environment disturbance is never fully preempted. All downstream uses of $\rho_i$ in the persistence and reserve conditions should be read as $\rho_i^{\text{eff}}$.


\medskip\noindent\textbf{Team Persistence Condition.}


% [Derived (team-persistence, from sector-condition-stability, persistence-condition)]

Applying the sector-condition framework ( \S\,\ref{sector-condition-stability}) with $\rho_i^{\text{eff}}$, agent $i$ persists iff:

$$\frac{\rho_i^{\text{eff}}}{\alpha_i} < R_i$$

Substituting the decomposition (the $\max(\cdot, 0)$ in $\rho_i^{\text{eff}}$ is omitted to expose the three levers; the condition is trivially satisfied when the numerator is non-positive):

$$\frac{\rho_{i,\text{env}} + \sum_j \gamma_{j \to i}^{\text{adv}} \mathcal{T}_j - \sum_j \gamma_{j \to i}^{\text{coop}} \mathcal{T}_j}{\alpha_i} < R_i$$

This reveals three distinct levers for team persistence:

\begin{enumerate}
  \item **Increase $\alpha_i$** (individual correction efficiency) — better models, better gain calibration
  \item \textbf{Increase cooperative tempo} ($\gamma^{\text{coop}} \mathcal T_j$) — more reliable allies, faster communication channels
  \item \textbf{Reduce adversarial coupling} ($\gamma^{\text{adv}} \mathcal T_j$) — better deception detection, reduced exposure to adversarial actions
\end{enumerate}



\medskip\noindent\textbf{Coordination Overhead Threshold.}


% [Discussion — Coordination Threshold]

Communication channels have costs: time to compose and parse messages, bandwidth limitations, synchronization requirements. These costs reduce the agent's effective tempo by diverting capacity from direct adaptation. Let $\Delta \mathcal T_i^{\text{cost}}(j)$ represent the tempo-equivalent coordination cost of maintaining the channel with $j$ — the reduction in $i$'s direct observation tempo caused by the overhead, in units of $[t^{-1}]$.

The net benefit of adding agent $j$ to $i$'s communication network is positive only when:

$$\nu_{ji}^{\text{comm}} \, \eta_{ji}^* > \Delta \mathcal{T}_i^{\text{cost}}(j)$$

Both sides have units $[t^{-1}]$: the LHS is communication tempo gained, the RHS is direct-adaptation tempo lost to coordination overhead. This implies a natural team-size limit: adding members increases communication tempo with diminishing returns (as $U_{\text{src}}$ and $U_o$ accumulate across diverse sources) while coordination costs grow, potentially superlinearly. The optimal team size occurs where the marginal communication tempo equals the marginal coordination cost.


\medskip\noindent\textbf{Epistemic Status.}


Conditional on the communication-gain hypothesis ( \S\,\ref{communication-gain}). The distributed tempo definition is a \textit{formulation} — a representational choice extending \S\,\ref{adaptive-tempo} to the multi-agent case. The disturbance decomposition is a \textit{formulation} — the additive structure and the sign convention are modeling choices, not derivations. The persistence condition is \textit{derived} from the sector-condition framework ( \S\,\ref{sector-condition-stability}) given the decomposition: the derivation is exact under the same assumptions (GA-2, GA-3 applied to $\rho_i^{\text{eff}}$). The coordination overhead threshold is \textit{discussion-grade} — qualitatively clear but the claim about diminishing returns and superlinear costs is asserted, not derived.

Max attainable: \textit{robust-qualitative} for the persistence condition (it inherits the sector-condition's robustness but the decomposition is a modeling choice). The coordination threshold could reach \textit{conditional} with a concrete cost model.


\medskip\noindent\textbf{Discussion.}


\textbf{Compositional analog of \S\,\ref{persistence-condition}.} Like the single-agent persistence condition, this segment addresses \textit{structural persistence} (see Persistence in \texttt{LEXICON.md}) — whether the composite correction machinery can outpace the effective disturbance rate. It does not address operational persistence (whether any particular sub-agent is near its boundary) or continuity persistence (whether the team maintains coherent identity through personnel changes). A team can be structurally persistent as a composite while individual members are operationally fragile or while the team's continuity is interrupted by turnover. The single-agent persistence condition says an agent persists when $\mathcal{T} > \rho / \Vert\delta_{\text{critical}}\Vert$. This segment extends that condition to agents embedded in a cooperative-adversarial network. The formal structure is identical — the sector-condition machinery applies unchanged — but the \textit{inputs} ($\mathcal T_i$ and $\rho_i$) now include inter-agent terms. This is consistent with \S\,\ref{composition-consistency}: the same dynamical laws apply at every level of description; what changes between levels is which channels contribute to tempo and which sources contribute to disturbance.

\textbf{Why teams can persist where individuals cannot.} Two distinct mechanisms combine. First, communication tempo raises $\mathcal{T}_i$ — allies provide observations that improve correction. Second, cooperative action lowers $\rho_i$ — allies act in the environment to reduce disturbance at its source. An individual agent facing $\rho_{i,\text{env}} > \alpha_i R_i$ fails the persistence condition. Adding cooperative allies can either raise the numerator's denominator (tempo) or lower the numerator directly (disturbance), or both — through physically distinct mechanisms.

\textbf{Timescale separation and \S\,\ref{composition-consistency}.} The distributed tempo definition presumes that communication events and direct observation events are comparable — they enter additively into $\mathcal T_i$. This requires that the communication timescale is not so slow relative to the environment dynamics that communicated information is stale on arrival. When communication latency approaches $1/\rho_{i,\text{env}}$, the effective $\eta_{ji}^*$ degrades (the observation uncertainty $U_{o,ji}$ increases with staleness), naturally suppressing the communication tempo contribution.

\textbf{Complement to \S\,\ref{adversarial-destabilization}.} That segment characterizes when an adversary can push an agent past its stability boundary. This segment characterizes the cooperative counterpart: when allies can pull an agent back from instability. The $\gamma$ coefficients have the same structure — coupling effectiveness — but opposite sign in the disturbance decomposition.

\end{proposition}

\begin{theorem}[Adversarial Tempo Advantage]
\label{adversarial-tempo-advantage}
Under adversarial coupling where one agent's actions contribute to the other's disturbance rate, the steady-state mismatch ratio scales superlinearly with the tempo ratio.



\medskip\noindent\textbf{Formal Expression.}


% [Derived (adversarial-tempo-advantage, from mismatch-dynamics steady state and adversarial-destabilization coupling model)]

\textbf{Setup.} Two agents $A$ and $B$ with adaptive tempos $\mathcal T_A$ and $\mathcal T_B$, coupled via the same model as \S\,\ref{adversarial-destabilization}:

$$\rho_A = \rho_{\text{base}} + \gamma_B \cdot \mathcal{T}_B, \qquad \rho_B = \rho_{\text{base}} + \gamma_A \cdot \mathcal{T}_A$$

where $\gamma_A, \gamma_B > 0$ are the coupling effectivenesses and $\rho_{\text{base}}$ is the background disturbance rate (taken equal for both agents for clarity; the asymmetric case is a straightforward generalization).

\textbf{Steady-state mismatch ratio.} From the linear mismatch dynamics ( \S\,\ref{mismatch-dynamics}), $\Vert\delta\Vert_{ss} = \rho / \mathcal{T}$. Substituting the coupled disturbance rates:

$$\Vert\delta_A\Vert_{ss} = \frac{\rho_{\text{base}} + \gamma_B \cdot \mathcal{T}_B}{\mathcal{T}_A}, \qquad \Vert\delta_B\Vert_{ss} = \frac{\rho_{\text{base}} + \gamma_A \cdot \mathcal{T}_A}{\mathcal{T}_B}$$

Taking the ratio:

$$\frac{\Vert\delta_B\Vert_{ss}}{\Vert\delta_A\Vert_{ss}} = \frac{\rho_{\text{base}} + \gamma_A \cdot \mathcal{T}_A}{\mathcal{T}_B} \cdot \frac{\mathcal{T}_A}{\rho_{\text{base}} + \gamma_B \cdot \mathcal{T}_B}$$

% [Result (adversarial-tempo-advantage)]

$$\frac{\Vert\delta_B\Vert_{ss}}{\Vert\delta_A\Vert_{ss}} = \frac{(\rho_{\text{base}} + \gamma_A \cdot \mathcal{T}_A) \cdot \mathcal{T}_A}{(\rho_{\text{base}} + \gamma_B \cdot \mathcal{T}_B) \cdot \mathcal{T}_B}$$


\medskip\noindent\textbf{Coupling-Dominant Limit.}


In the coupling-dominant regime ($\gamma_A \cdot \mathcal T_A \gg \rho_{\text{base}}$ and $\gamma_B \cdot \mathcal T_B \gg \rho_{\text{base}}$), the base disturbance becomes negligible:

$$\frac{\Vert\delta_B\Vert_{ss}}{\Vert\delta_A\Vert_{ss}} \to \frac{\gamma_A \cdot \mathcal{T}_A^2}{\gamma_B \cdot \mathcal{T}_B^2} = \frac{\gamma_A}{\gamma_B} \cdot \left(\frac{\mathcal{T}_A}{\mathcal{T}_B}\right)^2$$

For symmetric coupling ($\gamma_A = \gamma_B$):

$$\frac{\Vert\delta_B\Vert_{ss}}{\Vert\delta_A\Vert_{ss}} \to \left(\frac{\mathcal{T}_A}{\mathcal{T}_B}\right)^2$$

The exponent is $b = 2$: a \textbf{squared} tempo advantage. A 2:1 tempo ratio yields a 4:1 mismatch ratio, not 2:1.


\medskip\noindent\textbf{Derivation.}


From \S\,\ref{mismatch-dynamics}, the steady-state condition $d\Vert\delta\Vert/dt = 0$ gives $\Vert\delta\Vert_{ss} = \rho / \mathcal{T}$.

\begin{enumerate}
  \item Substitute the coupling model into $B$'s steady state: $\Vert\delta_B\Vert_{ss} = (\rho_{\text{base}} + \gamma_A \cdot \mathcal T_A) / \mathcal T_B$.
\end{enumerate}


\begin{enumerate}
  \item Substitute the coupling model into $A$'s steady state: $\Vert\delta_A\Vert_{ss} = (\rho_{\text{base}} + \gamma_B \cdot \mathcal T_B) / \mathcal T_A$.
\end{enumerate}


\begin{enumerate}
  \item Form the ratio:
\end{enumerate}


$$\frac{\Vert\delta_B\Vert_{ss}}{\Vert\delta_A\Vert_{ss}} = \frac{(\rho_{\text{base}} + \gamma_A \cdot \mathcal{T}_A) / \mathcal{T}_B}{(\rho_{\text{base}} + \gamma_B \cdot \mathcal{T}_B) / \mathcal{T}_A} = \frac{(\rho_{\text{base}} + \gamma_A \cdot \mathcal{T}_A) \cdot \mathcal{T}_A}{(\rho_{\text{base}} + \gamma_B \cdot \mathcal{T}_B) \cdot \mathcal{T}_B}$$

\begin{enumerate}
  \item In the coupling-dominant limit, $\rho_{\text{base}} \to 0$ relative to $\gamma \cdot \mathcal{T}$:
\end{enumerate}


$$\frac{\Vert\delta_B\Vert_{ss}}{\Vert\delta_A\Vert_{ss}} \to \frac{\gamma_A \cdot \mathcal{T}_A \cdot \mathcal{T}_A}{\gamma_B \cdot \mathcal{T}_B \cdot \mathcal{T}_B} = \frac{\gamma_A}{\gamma_B} \cdot \left(\frac{\mathcal{T}_A}{\mathcal{T}_B}\right)^2$$

$\square$


\medskip\noindent\textbf{Regime 2: Stochastic Coupling ($b = 3/2$).}


% [Derived (stochastic-tempo-advantage, from Model S steady state and coupling model)]

When the adversarial coupling enters as zero-mean noise rather than persistent drift, the disturbance model is Model S (GA-2S). The coupling model: $\sigma_B = \sigma_{\text{base}} + \gamma_A \cdot \mathcal{T}_A$ (the adversary's tempo increases the unpredictability of disturbances, not their systematic direction).

From the Model S steady state (Prop A.1S, scalar $n = 1$, linear $\alpha = \mathcal{T}$):

$$\lVert\delta\rVert_{\text{rms}} = \frac{\sigma}{\sqrt{2\mathcal{T}}}$$

Substituting the coupled noise scales:

$$\lVert\delta_B\rVert_{\text{rms}} = \frac{\sigma_{\text{base}} + \gamma_A \mathcal{T}_A}{\sqrt{2\mathcal{T}_B}}, \qquad \lVert\delta_A\rVert_{\text{rms}} = \frac{\sigma_{\text{base}} + \gamma_B \mathcal{T}_B}{\sqrt{2\mathcal{T}_A}}$$

Taking the ratio:

$$\frac{\lVert\delta_B\rVert_{\text{rms}}}{\lVert\delta_A\rVert_{\text{rms}}} = \frac{(\sigma_{\text{base}} + \gamma_A \mathcal{T}_A)\sqrt{\mathcal{T}_A}}{(\sigma_{\text{base}} + \gamma_B \mathcal{T}_B)\sqrt{\mathcal{T}_B}}$$

In the coupling-dominant limit ($\gamma \mathcal{T} \gg \sigma_{\text{base}}$, symmetric coupling $\gamma_A = \gamma_B$):

$$\frac{\lVert\delta_B\rVert_{\text{rms}}}{\lVert\delta_A\rVert_{\text{rms}}} \to \frac{\gamma_A \mathcal{T}_A \sqrt{\mathcal{T}_A}}{\gamma_B \mathcal{T}_B \sqrt{\mathcal{T}_B}} = \frac{\gamma_A}{\gamma_B} \cdot \left(\frac{\mathcal{T}_A}{\mathcal{T}_B}\right)^{3/2}$$

The exponent is $b = 3/2$. $\square$

\textbf{Why 3/2, not 2.} The $1/\sqrt{\alpha}$ scaling of Model S (vs. $1/\alpha$ for Model D) removes one half-power from the denominator. The numerator contributes $\mathcal{T}_A^1$ from the coupling; the denominator contributes $\mathcal{T}_B^{1/2}$ from the noise averaging. Combined with the $A$-side denominator ($1/\mathcal{T}_A^{1/2}$): $\mathcal{T}_A^{3/2}/\mathcal{T}_B^{3/2}$.


\medskip\noindent\textbf{Summary of Regime-Dependent Exponents.}


\begin{adjustbox}{max width=\textwidth}
\small
\begin{tabular}{lllll}
\toprule
Regime & Coupling type & Dominance & Exponent $b$ & Source \\
\midrule
1 & Deterministic drift (Model D) & Coupling-dominant & $2$ & Derived above \\
2 & Stochastic noise (Model S) & Coupling-dominant & $3/2$ & Derived above \\
3 & Either & Non-coupling-dominant & $\to 1$ (det.) or $\to 1/2$ (stoch.) & Asymptotic limit \\
\bottomrule
\end{tabular}
\end{adjustbox}


\textbf{Regime 3 (non-coupling-dominant).} When $\rho_{\text{base}} >rsim \gamma \cdot \mathcal{T}$ (or $\sigma_{\text{base}} >rsim \gamma \cdot \mathcal{T}$), the base disturbance dominates and the coupling terms become a perturbation. The mismatch ratio degrades toward $\mathcal T_A / \mathcal T_B$ (linear, $b = 1$) for Model D, or toward $(\mathcal T_A / \mathcal T_B)^{1/2}$ for Model S.

The simulation validation across all three regimes is in \S\,\ref{adversarial-exponent-regimes}.


\medskip\noindent\textbf{Epistemic Status.}


Both coupling-dominant exponents are \textit{exact} conditional on their respective disturbance models. The squared law ($b = 2$) is \textit{exact} under Model D (deterministic bounded disturbance, GA-2) with coupling-dominant conditions. The $3/2$ law ($b = 3/2$) is \textit{exact} under Model S (stochastic disturbance, GA-2S) with coupling-dominant conditions, derived from the $1/\sqrt{\alpha}$ steady-state scaling (Prop A.1S in \S\,\ref{sector-condition-derivation}). Both derivations are straightforward algebra from the respective steady-state formulas and the coupling model. The coupling model itself is an \textit{assumption} — the same one used in \S\,\ref{adversarial-destabilization}.

The non-coupling-dominant limits ($b \to 1$ for Model D, $b \to 1/2$ for Model S) are derived asymptotically. The smooth transition between regimes is confirmed by simulation ( \S\,\ref{adversarial-exponent-regimes}) but the interpolation formula is empirical. The transition between regimes is smooth, not sharp.

Max attainable: exact conditional on the disturbance model and coupling model. The result is as strong as its assumptions; no additional work changes the epistemic status without changing the dynamical model.


\medskip\noindent\textbf{Discussion.}


\textbf{Superlinearity is the key result.} The naive expectation — twice as fast yields twice the advantage — is wrong under adversarial coupling. The mechanism is that the faster agent both (a) corrects its own mismatch faster and (b) generates disturbance for the opponent faster. These two effects multiply, producing the squared exponent. Speed advantage is not additive; it compounds.

\textbf{Relationship to \S\,\ref{adversarial-destabilization}.} The steady-state mismatch ratio quantifies how much worse the slower agent does \textit{while both agents persist}. The destabilization threshold ( \S\,\ref{adversarial-destabilization}) marks where the slower agent fails entirely — its correction mechanism breaks down. Below the threshold, this segment's mismatch ratio applies. Above it, \S\,\ref{adversarial-destabilization}'s Lyapunov divergence takes over. The two results are complementary: this one gives the score; that one gives the game-ending condition.

\textbf{Regime dependence is operationally significant.} Whether an adversary's tempo increase produces systematic drift (positional maneuvering, API changes, doctrinal initiative) or unpredictable noise (feints, randomized attacks, market volatility) determines the scaling law. The distinction is not academic — $b = 2$ vs. $b = 3/2$ means a 3:1 tempo ratio yields 9:1 vs. 5.2:1 mismatch ratio. The model predicts that consistent, directional pressure is more effective per unit of tempo than unpredictable disruption.

\textbf{Formal analog of OODA-loop observations.} The squared scaling is consistent with Boyd's observation that getting inside the opponent's decision cycle has disproportionate effects. The theory identifies a specific mechanism (multiplicative interaction of correction speed and disturbance generation) and a specific condition (coupling-dominant regime) under which this disproportionality holds. Whether this mechanism is the dominant one in actual adversarial interactions is an empirical question, not a mathematical one.

\end{theorem}

\begin{hypothesis}[Communication Gain]
\label{communication-gain}
When an agent incorporates information from another agent (rather than from direct observation), the optimal update gain extends the uncertainty ratio with additional terms for source quality and alignment uncertainty.



\medskip\noindent\textbf{Formal Expression.}


% [Hypothesis (communication-gain)]

$$\eta_{ji}^* = \frac{U_{M_i}}{U_{M_i} + U_{o,ji} + U_{\text{src},j} + U_{\text{align},ji}}$$

where:
\begin{itemize}
  \item $U_{M_i}$: agent $i$'s model uncertainty (same as \S\,\ref{update-gain})
  \item $U_{o,ji}$: \textbf{communication channel noise} — latency, ambiguity, compression loss, bandwidth limitations of the channel between $j$ and $i$
  \item $U_{\text{src},j}$: \textbf{source quality uncertainty} — $i$'s uncertainty about $j$'s model calibration and domain competence
  \item $U_{\text{align},ji}$: \textbf{alignment uncertainty} — $i$'s uncertainty about whether $j$'s communications serve $i$'s interests or $j$'s potentially conflicting objectives
\end{itemize}


When all additional terms are zero (perfect channel, calibrated and aligned source): $\eta_{ji}^* \to 1$ (full trust). When any term is large: $\eta_{ji}^* \to 0$ (ignore the signal).

\textbf{Connection to single-agent case.} When $j$ is the environment (direct observation): $U_{\text{src}} = U_{\text{align}} = 0$, recovering \S\,\ref{update-gain}'s standard form $\eta^* = U_M / (U_M + U_o)$.


\medskip\noindent\textbf{Epistemic Status.}


\textit{Hypothesis.} The additive denominator treats all uncertainty sources as independent, zero-mean noise — a structural heuristic, not a strict variance derivation. This is appropriate for $U_{o,ji}$ (channel noise) and $U_{\text{src},j}$ (miscalibration), which are typically unstructured. For $U_{\text{align},ji}$ (deception), additivity is conservative: it correctly drives $\eta_{ji}^*$ toward zero when alignment uncertainty is high, but misses the adversary's \textit{actual} strategy — presenting as trustworthy to exploit high $\eta_{ji}^*$. The additive model captures the \textit{defender's} response to detected misalignment; it does not model the \textit{attacker's} optimization over the defender's trust dynamics.

All four uncertainty terms must be expressed in a \textbf{common predictive-dispersion scale} before summation — the same units as $U_{M_i}$ (variance of the predictive distribution over the observed quantity). When hard to estimate directly, a conservative approximation: set $U_{\text{src}} + U_{\text{align}}$ to the empirical variance of $j$'s past prediction residuals as observed by $i$, minus the known channel noise $U_{o,ji}$.


\medskip\noindent\textbf{Discussion.}


\textbf{The denominator terms have different natures.} $U_{o,ji}$ is a property of the \textit{channel} — improvable by infrastructure. $U_{\text{src},j}$ is a property of the \textit{source} — improvable by $j$ improving its model, or estimable by $i$ through calibration tracking. $U_{\text{align},ji}$ is a property of the \textit{relationship} — the game-theoretic variable.

\textbf{Trust calibration as a meta-model.} Agent $i$'s estimates of $U_{\text{src},j}$ and $U_{\text{align},ji}$ constitute a \textbf{trust meta-model} — a model of models. This meta-model is itself subject to ACT's full apparatus: it has mismatch (trust prediction errors), should be updated with appropriate gain (not overreacting to single disagreements), and can be structurally inadequate ( \S\,\ref{structural-adaptation-necessity} — the agent's trust model class may not capture the actual reliability structure of its sources).

\textbf{Risk-asymmetric trust.} The Bayesian posterior on source reliability gives the best \textit{estimate}, but the \textit{decision} about how much to trust should be risk-weighted. Trusting a deceptive agent (HILP — high impact, low probability) can cause catastrophic model corruption ( \S\,\ref{adversarial-destabilization}, effects spiral). Mild miscalibration toward a reliable source (LIHP) causes small ongoing inefficiency. For high-stakes interactions, use a conservative quantile of the trust posterior rather than the mean — require more evidence before granting high trust.

\textbf{Trust transitivity.} When agent $i$ has no direct experience with agent $k$, but trusted intermediary $j$ provides an assessment, the transitive trust question arises. A Bayesian mixture model discounts the recommendation by the intermediary's own reliability:

$$P_i(\theta_k \mid s_j) \propto \left[r_{ji} \cdot P(s_j \mid \theta_k) + (1 - r_{ji}) \cdot P_0(s_j)\right] \cdot P_i(\theta_k)$$

where $r_{ji}$ is $i$'s reliability estimate of $j$ and $\theta_k$ is $k$'s true alignment. When $r_{ji} \to 0$, the posterior collapses to the prior (no update); when $r_{ji} \to 1$, the full informative likelihood applies. This model gives a principled three-phase trust formation: prior → transitive update → direct experience (which eventually swamps the prior).

\end{hypothesis}

\begin{proposition}[Adversarial Destabilization]
\label{adversarial-destabilization}
When two agents are coupled such that one's praxis contributes to the other's disturbance rate, the faster agent can generate aporia in the target faster than the target's epistrophe can resolve it — driving the target outside its invariant region and causing the correction mechanism to break down entirely.



\medskip\noindent\textbf{Formal Expression.}


% [Derived (adversarial-destabilization, from sector-condition-stability)]

\textbf{Setup.} Let both agents $A$ and $B$ satisfy the sector condition (A1, A2', A3 from \S\,\ref{sector-condition-derivation}) with respective parameters $(\alpha_A, R_A, \rho_{A,\text{base}})$ and $(\alpha_B, R_B, \rho_{B,\text{base}})$. The coupling:

% [Assumption (Coupling Model)]

$$\rho_B = \rho_{B,\text{base}} + \gamma_A \cdot \mathcal{T}_A$$

where $\gamma_A > 0$ represents the effectiveness of $A$'s actions at disrupting $B$'s environment.

\textbf{Result.} $B$'s ultimately bounded radius under coupled dynamics is:

$$R^*_B = \frac{\rho_{B,\text{base}} + \gamma_A \cdot \mathcal{T}_A}{\alpha_B}$$

$B$ diverges (exits its sector-condition region) when $R^*_B > R_B$, i.e., when:

$$\gamma_A \cdot \mathcal{T}_A > \alpha_B R_B - \rho_{B,\text{base}} = \Delta\rho^*_B$$

That is:

$$\mathcal{T}_A > \frac{\Delta\rho^*_B}{\gamma_A}$$

Symmetrically, $B$ destabilizes $A$ when $\mathcal{T}_B > \Delta\rho^*_A / \gamma_B$.

The adversarial outcome depends on whether either agent can push the other past its stability limit. $\square$

\textbf{Interpretation.} "Getting inside the opponent's OODA loop" has a precise Lyapunov characterization: Agent $A$ destabilizes Agent $B$ when $A$'s praxis, multiplied by coupling effectiveness, generates aporia in $B$ faster than $B$'s epistrophe can resolve it — specifically, when $A$'s tempo times coupling exceeds $B$'s adaptive reserve $\Delta\rho^*_B$. This captures:

\begin{itemize}
  \item \textbf{Asymmetric coupling} ($\gamma_A \neq \gamma_B$): an agent with lower tempo but higher coupling effectiveness can still win.
  \item \textbf{Finite reserves}: an agent with very high $\mathcal{T}$ but operating near its model-class limit ($\Delta\rho^*$ small) is vulnerable despite high tempo.
  \item \textbf{Structural collapse}: when $R^*_B > R_B$, the failure mode is not merely "large mismatch" but "correction mechanism breakdown" — connecting to \S\,\ref{structural-adaptation-necessity}.
\end{itemize}



\medskip\noindent\textbf{Corollary: The Effects Spiral.}


When Agent $B$ is driven past its stability boundary ($R^*_B > R_B$), and $B$'s degrading model causes $B$'s actions to become erratic in a way that increases $A$'s coupling effectiveness ($\gamma_A$ increases with $\Vert\delta_B\Vert$), the result is a positive-feedback Lyapunov instability:

% [Discussion — Mechanism Schematic]

$$\Vert\delta_B\Vert \uparrow \;\Rightarrow\; B\text{'s actions become erratic} \;\Rightarrow\; \gamma_A \uparrow \;\Rightarrow\; \rho_B \uparrow \;\Rightarrow\; \Vert\delta_B\Vert \uparrow$$

With $\gamma_A$ now an increasing function of $\Vert\delta_B\Vert$, the disturbance term in $B$'s dynamics grows superlinearly. $\dot{V}_B > 0$ and increasing — mismatch accelerates away from the stability region. The spiral terminates only when $B$ undergoes structural adaptation ( \S\,\ref{structural-adaptation-necessity} — changing the model class) or ceases to function as an adaptive agent entirely.


\medskip\noindent\textbf{Epistemic Status.}


The destabilization threshold is \textit{exact} under the coupling model (which treats $\mathcal{T}_A$ as exogenous). The coupling model itself is an \textit{assumption} — it decouples the agents rather than modeling the fully coupled dynamical system where both agents' mismatch states co-evolve. The analysis therefore characterizes the \textit{destabilization threshold} (the conditions under which $A$ \textit{can} push $B$ past its stability boundary) rather than the full transient dynamics. This is a worst-case bound, treating $A$ as operating at its steady-state tempo.

The effects spiral (corollary) is \textit{discussion-grade} — the positive-feedback mechanism is qualitatively clear, but formalizing the $\gamma_A(\Vert\delta_B\Vert)$ functional form and proving instability under it requires specifying how an agent's degrading model affects its action quality, which the theory does not yet formalize.

A full coupled Lyapunov analysis with a joint function $V(\delta_A, \delta_B)$ would capture mutual feedback effects but requires specifying how each agent's mismatch state affects the other's disturbance in real time — an open extension.


\medskip\noindent\textbf{Discussion.}


\textbf{Destabilization vs. steady-state ratio.} The destabilization threshold is a failure of \textit{structural persistence} (see Persistence in \texttt{LEXICON.md}) — the point where the correction machinery can no longer outpace the adversarially amplified disturbance. An adversary does not need to attack operational persistence (pushing the target near its boundary) or continuity persistence (disrupting identity) directly; destroying structural persistence is sufficient, because without it, operational persistence degrades to zero and the agent ceases to function regardless of its continuity stance. The linear analysis in \S\,\ref{mismatch-dynamics} gives the steady-state mismatch ratio under coupling: a quantitative result about how much worse $B$ does. This segment gives the qualitative result: under what conditions does $B$ \textit{fail entirely}, not merely fall behind. The linear analysis tells you the score; the Lyapunov analysis tells you when the game is over.

\textbf{Connection to \S\,\ref{adversarial-tempo-advantage}.} The simulation results show the tempo advantage is superlinear (exponent $\approx 2$ in pure adversarial regimes). This Lyapunov result explains WHY: the destabilization threshold creates a phase transition — below it, $B$ persists (possibly with degraded performance); above it, $B$'s correction mechanism collapses entirely, and the effects spiral accelerates the collapse.

\end{proposition}

\begin{observation}[Adversarial Exponent Regimes]
\label{adversarial-exponent-regimes}
The adversarial tempo advantage exponent — the power $b$ in $\lVert\delta_B\rVert / \lVert\delta_A\rVert \sim (\mathcal T_A / \mathcal T_B)^b$ — is not a single number. It depends on two structural features of the disturbance: whether the adversarial coupling enters as deterministic drift (Model D) or stochastic noise (Model S), and whether the coupling dominates the base disturbance rate. Three regimes, with the coupling-dominant exponents now derived analytically from the respective disturbance models.



\medskip\noindent\textbf{Formal Expression.}


% [Derived (adversarial-exponent-regimes, from Model D/S steady states + coupling model; validated by simulation)]

\textbf{Regime 1: Model D (deterministic drift), coupling-dominant.} When adversarial coupling enters as a persistent directional disturbance ($\rho_B = \rho_{\text{base}} + \gamma \cdot \mathcal T_A$, GA-2) and coupling dominates ($\gamma \cdot \mathcal T_B \gg \rho_{\text{base}}$):

$$b = 2 \qquad \text{(simulation: 1.999)}$$

Derived from the Model D steady state $\lVert\delta\rVert_{ss} = \rho/\mathcal{T}$ (Prop A.1). See \S\,\ref{adversarial-tempo-advantage}.

\textbf{Regime 2: Model S (stochastic noise), coupling-dominant.} When adversarial coupling enters through the noise scale of zero-mean perturbations ($\sigma_B = \sigma_{\text{base}} + \gamma \cdot \mathcal T_A$, GA-2S) and coupling dominates:

$$b = \frac{3}{2} \qquad \text{(simulation: 1.481)}$$

Derived from the Model S steady state $\lVert\delta\rVert_{\text{rms}} = \sigma_w/\sqrt{2\mathcal{T}}$ (Prop A.1S). The $1/\sqrt{\mathcal{T}}$ scaling (vs. $1/\mathcal{T}$ for Model D) removes one half-power from the denominator, reducing the exponent from 2 to 3/2. See \S\,\ref{adversarial-tempo-advantage}.

\textbf{Regime 3: Non-coupling-dominant.} When base disturbance is comparable to or exceeds the adversarial coupling ($\rho_{\text{base}} >rsim \gamma \cdot \mathcal T_B$):

$$b \to 1.0 \text{ (Model D)} \quad \text{or} \quad b \to 0.5 \text{ (Model S)}$$

The exponent degrades smoothly as the base-to-coupling ratio increases. The asymptotic limits are derived (they reflect the $1/\mathcal{T}$ or $1/\sqrt{\mathcal{T}}$ scaling without the coupling numerator); the smooth interpolation is empirical.

\begin{adjustbox}{max width=\textwidth}
\small
\begin{tabular}{lll}
\toprule
$\rho_{\text{base}} / (\gamma \cdot \mathcal T_B)$ & Exponent (deterministic) & Exponent (stochastic) \\
\midrule
0.002 & 1.999 & 1.481 \\
0.20 & 1.877 & 1.101 \\
2.0 & 1.445 & 0.791 \\
6.3 & 1.213 & 0.577 \\
\bottomrule
\end{tabular}
\end{adjustbox}



\medskip\noindent\textbf{Epistemic Status.}


\textit{Exact conditional on disturbance model.} The coupling-dominant exponents are derived, not empirical: $b = 2$ follows from the Model D steady state (Prop A.1) and the coupling model; $b = 3/2$ follows from the Model S steady state (Prop A.1S) and the coupling model. The simulation results (6 variants, multiple parameter sweeps) now serve as validation of the derived exponents, not as their epistemic foundation. The non-coupling-dominant limits ($b \to 1$, $b \to 1/2$) are derived asymptotically; the smooth interpolation between coupling-dominant and non-coupling-dominant is empirical. What remains empirical is whether a given real adversarial interaction is better modeled as Model D or Model S — that is a domain question, not a theory question.


\medskip\noindent\textbf{Discussion.}


\textbf{The disturbance model determines the exponent.} The mismatch dynamics ( \S\,\ref{mismatch-dynamics}) now distinguish two disturbance models: Model D (bounded deterministic, GA-2) with steady-state $\rho/\mathcal{T}$, and Model S (stochastic zero-mean, GA-2S) with steady-state $\sigma_w/\sqrt{2\mathcal{T}}$. The different steady-state scaling is the root cause of the different exponents. This resolves the ambiguity that previously existed in the single-$\rho$ formulation.

\textbf{Why the squared law held for the coupling-dominance sweep.} In Variant A, the coupling enters as deterministic drift: $\rho_B = \rho_{\text{base}} + \gamma \cdot \mathcal T_A$, and the steady state is $\Vert\delta_B\Vert = \rho_B / \mathcal T_B$. The ratio $\Vert\delta_B\Vert / \Vert\delta_A\Vert$ in the coupling-dominant limit gives $(\mathcal T_A / \mathcal T_B)^2$ directly.

\textbf{Nonlinear correction creates thresholds, not lower exponents.} For saturating, sigmoid, and breakdown correction functions under deterministic drift, the issue is not a reduced exponent but a catastrophic divergence when $\rho$ exceeds the correction capacity ($\rho > \mathcal{T} \cdot R$). This is exactly the persistence threshold failure ( \S\,\ref{persistence-condition}), observed directly in simulation.

\textbf{Domain interpretation.} Whether a given opponent's tempo increase causes deterministic drift or stochastic noise depends on the domain:
\begin{itemize}
  \item Military: an opponent who maneuvers faster creates systematic positional change (drift, $b \approx 2$)
  \item Market: a competitor who acts unpredictably creates noise in signals ($b \approx 1.5$)
  \item Software: a fast-changing API creates systematic drift in the codebase state (drift)
  \item Adversarial ML: an opponent who varies attack vectors increases observation noise ($b \approx 1.5$)
\end{itemize}


\end{observation}

\begin{observation}[Observation Noise Gates Adversarial Advantage]
\label{observation-gates-advantage}
Observation noise collapses the adversarial tempo advantage. When agents observe their mismatch through a noisy channel, the faster agent's additional corrections become noisy, partially offsetting its tempo advantage. The optimal gain ( \S\,\ref{update-gain}) partially restores the advantage but cannot fully recover it.



\medskip\noindent\textbf{Formal Expression.}


% [Observation (observation-gates-advantage, from track-b Variant E)]

In a two-agent adversarial system with observation noise $\sigma_{\text{obs}}$ added to each agent's mismatch signal:

\begin{adjustbox}{max width=\textwidth}
\small
\begin{tabular}{lll}
\toprule
$\sigma_{\text{obs}}$ & Exponent (fixed $\eta$) & Exponent (optimal $\eta^*$) \\
\midrule
0.00 & 1.04 & 1.04 \\
0.10 & 1.00 & 0.97 \\
0.20 & 0.92 & 0.94 \\
0.50 & 0.60 & 0.63 \\
1.00 & 0.18 & 0.40 \\
\bottomrule
\end{tabular}
\end{adjustbox}


At $\sigma_{\text{obs}} = 1.0$ (10x the process noise), the fixed-gain adversarial exponent drops from $\sim 1.0$ to $\sim 0.2$ — tempo advantage nearly vanishes. The Riccati-optimal gain restores it to $\sim 0.4$, more than doubling the advantage but not recovering the noise-free level.

\textbf{The mechanism.} When observation noise is high, each correction step adds noise to the mismatch estimate. The faster agent makes more corrections per unit time, each noisy, partially offsetting the benefit of higher tempo. The optimal gain mitigates this by reducing $\eta$ to match the noise level — correcting less aggressively but more accurately.


\medskip\noindent\textbf{Epistemic Status.}


\textit{Empirical.} Max attainable: derived (the mechanism is analytically tractable via Riccati analysis of noisy AR(1) processes). The observation that noise degrades advantage is confirmed by simulation. The optimal gain's partial restoration is consistent with the uncertainty ratio principle ( \S\,\ref{update-gain}: $\eta^* = U_M / (U_M + U_o)$). The quantitative degradation curve ($b$ vs. $\sigma_{\text{obs}}$) is empirical at these parameters; a general analytical expression would require solving the coupled noisy-AR(1) system.


\medskip\noindent\textbf{Discussion.}


\textbf{Observation quality gates tempo advantage.} Boyd insisted that the quality of Orient (observation processing) matters more than raw OODA speed. The simulation results show a formal analog of this pattern: faster tempo with noisy observations ($\sigma_{\text{obs}}$ high) gives nearly zero advantage over a slower agent with equally noisy observations. The tempo advantage is gated by observation quality — consistent with Boyd's emphasis, though the model captures a specific mechanism (noisy correction steps) rather than the full richness of Orient processing.

\textbf{The optimal gain helps most in the moderate-noise regime.} At $\sigma_{\text{obs}} = 0.05$ (observation noise half of process noise), the optimal gain cuts steady-state mismatch by 52\% compared to fixed gain. At very high noise, the improvement is less dramatic in absolute terms but more important relatively (0.40 vs. 0.18 exponent).

\textbf{Practical implication.} An agent facing an adversary with superior tempo should invest in degrading the adversary's observation quality rather than trying to match their speed. Conversely, an agent with superior tempo should protect its observation channels — the tempo advantage is only as good as the observation quality that supports it.

\textbf{Connection to code quality.} In the software domain ( \#code-quality-as-observation-infrastructure — cross-component reference, see \texttt{02-tst-core/}), code quality IS observation infrastructure. A well-structured codebase provides low-noise observations (clear tests, readable code, explicit interfaces). A poorly structured codebase adds observation noise to every development cycle, degrading the developer's effective tempo regardless of how fast they work.

\end{observation}

\begin{theorem}[Per-Dimension Persistence]
\label{per-dimension-persistence}
The scalar persistence condition overestimates adaptive capacity when the agent's correction gain varies across dimensions. The weak dimension is the bottleneck — it dominates the aggregate mismatch regardless of performance on strong dimensions. The correct condition is per-dimension, with the form depending on whether the disturbance is deterministic (Model D) or stochastic (Model S).



\medskip\noindent\textbf{Formal Expression.}


% [Result (per-dimension-persistence)]

For an agent with $d$-dimensional mismatch $\delta_t \in \mathbb{R}^d$, diagonal correction gain $\eta = \text{diag}(\eta_1, \ldots, \eta_d)$, and per-dimension disturbance:


\medskip\noindent\textbf{Model D: Deterministic Per-Dimension Threshold.}


Under bounded disturbance $|w_k(t)| \leq \rho_k$ (GA-2, per dimension), the per-dimension steady-state mismatch is:

$$|\delta_k|_{ss} = \frac{\rho_k}{\alpha_k}$$

\textbf{Persistence requires} $\alpha_k > \rho_k / R_k$ \textbf{for each dimension}, or in linear operational form:

$$\mathcal{T}_k > \frac{\rho_k}{\delta_{\text{critical},k}} \quad \text{for each dimension } k$$

This is the deterministic worst-case bound — exact under bounded disturbance by the same Lyapunov argument as Prop A.1, applied per dimension.


\medskip\noindent\textbf{Model S: Stochastic Per-Dimension Steady State.}


Under stochastic disturbance $w_{k,t} \sim N(0, \rho_k^2)$ (GA-2S, per dimension), the discrete AR(1) process $\delta_{k,t+1} = (1 - \eta_k)\delta_{k,t} + w_{k,t}$ has stationary distribution $\delta_k \sim N(0, \rho_k^2/(2\eta_k - \eta_k^2))$, giving:

$$E[|\delta_k|] = \frac{\rho_k}{\sqrt{2\eta_k - \eta_k^2}} \cdot \sqrt{\frac{2}{\pi}}$$

which for small $\eta_k$ approximates:

$$E[|\delta_k|] \approx \frac{\rho_k}{\sqrt{2\eta_k}} \cdot \sqrt{\frac{2}{\pi}}$$

\textbf{Persistence requires} (from $E[|\delta_k|] < \delta_{\text{critical},k}$, small-$\eta_k$ approximation):

$$\eta_k > \frac{\rho_k^2}{\pi \cdot \delta_{\text{critical},k}^2}$$

Or as a probability bound ($P(|\delta_k| > \delta_{\text{critical},k}) < \epsilon$, using the Gaussian tail):

$$\eta_k > \frac{\rho_k^2 \cdot z_{1-\epsilon}^2}{2 \cdot \delta_{\text{critical},k}^2}$$

where $z_{1-\epsilon}$ is the Gaussian quantile. Note: the stochastic threshold is quadratic in $\rho_k/\delta_{\text{critical},k}$ (not linear as in Model D), reflecting the $1/\sqrt{\alpha}$ scaling.


\medskip\noindent\textbf{Common Structure.}


The aggregate $L_2$ mismatch $\lVert\delta\rVert = \sqrt{\sum_k \delta_k^2}$ is dominated by the dimension with the largest $\rho_k / \eta_k$ ratio (Model S) or $\rho_k / \alpha_k$ ratio (Model D). The qualitative conclusion — the weak dimension is the bottleneck — holds for both models.


\medskip\noindent\textbf{Epistemic Status.}


\textit{Exact conditional on disturbance model.} Both per-dimension forms are now derived from their respective disturbance models:

\begin{enumerate}
  \item \textbf{Model D threshold} ($\mathcal{T}_k > \rho_k/\delta_{\text{critical},k}$) follows from Prop A.1 applied per dimension — the same Lyapunov argument with bounded disturbance, restricted to each coordinate. This is exact under GA-2 + GA-3.
\end{enumerate}


\begin{enumerate}
  \item \textbf{Model S steady state} ($E[|\delta_k|] = \rho_k/\sqrt{2\eta_k - \eta_k^2} \cdot \sqrt{2/\pi}$) follows from the AR(1) stationary distribution under Gaussian disturbance (GA-2S). The stochastic persistence threshold ($\eta_k > \rho_k^2/(\pi \cdot \delta_{\text{critical},k}^2)$) is derived from this formula. The 4-significant-figure simulation match validates Model S, not Model D.
\end{enumerate}


The previously noted "regime mixing" is resolved: the two threshold forms belong to different disturbance models. The Model D threshold is linear in $\rho_k$; the Model S threshold is quadratic. The 72\% scalar overestimate and weak-dimension bottleneck are structural properties that hold under both models.


\medskip\noindent\textbf{Discussion.}


\textbf{Scalar tempo overestimates.} In a 3D system with gains $\eta = (0.15, 0.03, 0.03)$ and disturbances $\rho = (0.20, 0.20, 0.02)$:

\begin{adjustbox}{max width=\textwidth}
\small
\begin{tabular}{lllll}
\toprule
Dimension & $\eta_k$ & $\rho_k$ & $\rho_k / \eta_k$ & $E[\Vert\delta_k\Vert]$ \\
\midrule
1 (well-tracked) & 0.15 & 0.20 & 1.33 & 0.303 \\
2 (weak) & 0.03 & 0.20 & 6.67 & 0.656 \\
3 (unimportant) & 0.03 & 0.02 & 0.67 & 0.066 \\
\bottomrule
\end{tabular}
\end{adjustbox}


Scalar prediction: $\rho / \mathcal{T} = 0.284 / 0.21 = 1.35$. Actual $\Vert\delta\Vert_{L_2} = 0.785$. Overestimate: 72\%. Dimension 2 alone accounts for 84\% of the $L_2$ mismatch.

\textbf{Isotropic allocation dominates.} Equalizing the same total gain budget ($\eta = 0.07$ per dimension) reduces $\Vert\delta\Vert_{L_2}$ from 0.785 to 0.685 — a 13\% improvement — because it reduces the bottleneck effect on the weak dimension.

\textbf{Adversarial exploitation.} An adversary who identifies the target's weak dimension can concentrate disturbance there. Targeted attack (80\% on the weak dimension) amplifies the mismatch ratio by 17\% (from 2.70 to 3.15). The real danger is structural: if the weak dimension's mismatch exceeds its critical threshold ($R_{\text{max}}$), correction fails on that dimension while the aggregate $\Vert\delta\Vert_{L_2}$ may still look manageable. Per-dimension monitoring is essential.

\textbf{Implications for the persistence condition.} Like the scalar result, per-dimension persistence addresses \textit{structural persistence} (see Persistence in \texttt{LEXICON.md}) — whether the correction machinery on each dimension can outpace that dimension's disturbance rate. An agent can be structurally persistent on every dimension while still being operationally fragile on one (near its per-dimension $R_k$ boundary). The scalar persistence condition ( \S\,\ref{persistence-condition}) remains correct as a \textit{necessary} condition: if the aggregate tempo is insufficient, the agent fails. But it is not \textit{sufficient} — an agent can satisfy the scalar condition while failing on a single dimension. The per-dimension condition has two forms: Model D ($\mathcal{T}_k > \rho_k/\delta_{\text{critical},k}$, exact under bounded disturbance) and Model S ($\eta_k > \rho_k^2/(\pi \cdot \delta_{\text{critical},k}^2)$, exact under Gaussian disturbance). Both predict per-dimension failure correctly; the choice depends on the disturbance character in the domain.

\textbf{Connection to multi-agent systems.} The per-dimension result has a direct multi-agent analog: in a composite agent, each sub-agent's contribution to composite tempo may be strong in some dimensions and weak in others. The composite's persistence requires coverage across all relevant dimensions — a team of specialists who each handle one dimension well composes better than a team of generalists who are mediocre at everything, provided the dimension assignment matches.

\end{theorem}

\section{Appendices: Details}

\begin{quote}\itshape
Supporting material: derivations, sketches, simulation results, and operationalization procedures backing the main theory claims.
\end{quote}

\begin{derivation}[Sector Condition Stability — Lyapunov Derivation]
\label{sector-condition-derivation}
Complete Lyapunov derivations of bounded mismatch and adaptive reserve for the sector-condition results stated in \S\,\ref{sector-condition-stability}.



\medskip\noindent\textbf{Motivation.}


\S\,\ref{mismatch-dynamics} hypothesizes the linear ODE $d\Vert\delta\Vert/dt = -\mathcal{T}\Vert\delta\Vert + \rho$ as a first-order approximation. The linear form yields clean closed-form results but commits to a specific functional relationship between mismatch magnitude and correction rate. True correction dynamics are almost certainly nonlinear — exhibiting saturation at large mismatch, threshold effects near zero, and structural breakdown when the model class is exhausted.

A Lyapunov approach proves persistence and stability under much weaker assumptions: any correction dynamics satisfying qualitative monotonicity properties (the sector condition). The results below are strictly more general — the linear case is recovered where the sector bounds coincide.


\medskip\noindent\textbf{Setup.}


Let $\delta(t) \in \mathbb{R}^n$ be the mismatch vector — the difference between the model's predictions and reality across $n$ observable dimensions. The vector treatment connects to per-dimension tempo analysis ( \S\,\ref{per-dimension-persistence}).

The mismatch dynamics are:


$$\frac{d\delta}{dt} = -F(\mathcal{T}, \delta) + w(t)$$

where:

\begin{itemize}
  \item $F(\mathcal{T}, \delta): \mathbb R_+ \times \mathbb{R}^n \to \mathbb{R}^n$ is the \textbf{correction function} — how the agent's adaptive process reduces mismatch. It maps to the same space as $\delta$ (so that the inner product $\delta^T F$ in the sector condition is well-defined). This subsumes the update gain $\eta^*$ ( \S\,\ref{update-gain}), event rate $\nu$, and the structure of the update rule.
  \item $w(t)$ is the \textbf{disturbance} — new mismatch introduced by environmental change, with $\Vert w(t)\Vert \leq \rho$ (bounded disturbance rate).
\end{itemize}


The linear case from \S\,\ref{mismatch-dynamics} has $F(\mathcal{T}, \delta) = \mathcal{T} \cdot \delta$.


\medskip\noindent\textbf{Assumptions on $F$.}


We require only qualitative properties, not a specific functional form:


\medskip\noindent\textbf{(A1) Zero Correction at Zero Mismatch.}


% [Assumption A1]

$$F(\mathcal{T}, 0) = 0$$

No correction is applied when the model perfectly matches reality.


\medskip\noindent\textbf{(A2') Local Sector Condition.}


There exists a region $\mathcal B_R = \{\delta : \Vert\delta\Vert \leq R\}$ and $\alpha > 0$ such that (following the sector-condition framework of Lur'e):

% [Assumption A2' (sector-condition)]

$$\delta^T F(\mathcal{T}, \delta) \geq \alpha \Vert\delta\Vert^2 \quad \forall \delta \in \mathcal{B}_R$$

The correction function always points "inward" (reducing mismatch), and its magnitude is bounded below relative to $\Vert\delta\Vert^2$. The linear case has $\alpha = \mathcal{T}$. A saturating correction has $\alpha$ decreasing for large $\Vert\delta\Vert$. A threshold correction has $\alpha = 0$ for small $\Vert\delta\Vert$.

The local form allows the correction to break down outside $\mathcal B_R$ — the structural adaptation regime of \S\,\ref{structural-adaptation-necessity}.

\textbf{Grounding of GA-3.} The sector condition (A2') is not an irreducible assumption for well-designed agents. \S\,\ref{gain-sector-bridge} shows that gain-based updating ( \S\,\ref{update-gain}) with directional fidelity produces correction functions satisfying A2', with $\alpha = \eta^* \cdot c_{\min}$. For gradient-based agents, A2' is equivalent to local strong convexity of the loss function, with $\alpha = \eta \cdot \mu$ where $\mu$ is the strong convexity modulus. The Lyapunov proofs below are unchanged — they operate downstream of A2' regardless of whether it is assumed or derived.


\medskip\noindent\textbf{(A3) Tempo Monotonicity.}


For fixed $\delta$, $\delta^T F(\mathcal{T}, \delta)$ is monotone increasing in $\mathcal{T}$. Higher tempo means faster correction.

\textbf{Connection to ACT parameters.} The sector parameter $\alpha$ is determined by the adaptive tempo $\mathcal{T}$ ( \S\,\ref{adaptive-tempo}) and the structure of the correction function. In the linear case, $\alpha = \mathcal{T} = \sum_k \nu^{(k)} \cdot \eta^{(k)*}$. In nonlinear cases, $\alpha$ represents the \textit{worst-case} correction efficiency within the valid region — the minimum ratio of correction power to mismatch magnitude. The radius $R$ represents the model class capacity: how large a mismatch can grow before the correction mechanism fails (i.e., before the sector condition ceases to hold), at which point structural adaptation ( \S\,\ref{structural-adaptation-necessity}) becomes necessary.


\medskip\noindent\textbf{Candidate Lyapunov Function.}



$$V(\delta) = \frac{1}{2}\Vert\delta\Vert^2$$

Positive definite, radially unbounded, continuously differentiable. Its level sets $V = c$ are spheres of radius $\sqrt{2c}$ in mismatch space.


\medskip\noindent\textbf{Proposition A.1: Bounded Mismatch (Single Agent).}


\textbf{Statement.} Under (A1), (A2'), (A3), with bounded disturbance $\Vert w(t)\Vert \leq \rho$, the mismatch $\delta(t)$ is \textbf{ultimately bounded}: there exists $R^* > 0$ such that $\Vert\delta(t)\Vert \leq R^*$ for all sufficiently large $t$, provided $R^* < R$ (the ultimately bounded region fits within the sector-condition region).

\textbf{Proof.}

Compute $\dot{V}$ along trajectories:

% [Derived (Proof Step)]

$$\dot{V} = \delta^T \dot{\delta} = \delta^T[-F(\mathcal{T}, \delta) + w(t)]$$

% [Derived (Proof Step)]

$$= -\delta^T F(\mathcal{T}, \delta) + \delta^T w(t)$$

By (A2'): $\delta^T F(\mathcal{T}, \delta) \geq \alpha\Vert\delta\Vert^2$

By Cauchy-Schwarz: $\delta^T w(t) \leq \Vert\delta\Vert \cdot \Vert w(t)\Vert \leq \rho\Vert\delta\Vert$

Therefore:

% [Derived (Proof Step)]

$$\dot{V} \leq -\alpha\Vert\delta\Vert^2 + \rho\Vert\delta\Vert = -\Vert\delta\Vert(\alpha\Vert\delta\Vert - \rho)$$

$\dot{V} < 0$ whenever $\Vert\delta\Vert > \rho/\alpha$ \textbf{and} $\Vert\delta\Vert \leq R$ (where A2' holds).

Define $R^* = \rho/\alpha$.

**Invariance of $\mathcal B_R$.** When $R^* < R$ (the persistence condition), the ball $\mathcal B_R$ is \textit{positively invariant}: any trajectory starting in $\mathcal B_R$ remains in $\mathcal B_R$ for all future time. At the boundary $\Vert\delta\Vert = R$, the sector condition holds and $\dot{V} = -\Vert\delta\Vert(\alpha\Vert\delta\Vert - \rho) = -R(\alpha R - \rho) < 0$ (since $\alpha R > \rho$ by the persistence condition). Trajectories at the boundary point inward. Therefore $\mathcal B_R$ is invariant, and the sector condition applies to the entire future trajectory of any trajectory starting inside $\mathcal B_R$.

\textbf{Ultimate boundedness.} Within $\mathcal B_R$, the Lyapunov function is strictly decreasing for $\Vert\delta\Vert > R^*$ and may increase for $\Vert\delta\Vert < R^*$. All trajectories starting in $\mathcal B_R$ are ultimately bounded by $R^*$ — they are driven into $\mathcal B_{R^*}$ and remain in a neighborhood of it (with possible oscillation at the boundary due to the disturbance).

\textbf{Initial condition requirement.} The result requires the trajectory to start inside $\mathcal B_R$ (or to be brought inside by some external mechanism). A trajectory starting outside $\mathcal B_R$ is not covered — the sector condition does not hold there, and the correction function may fail to reduce mismatch. This is precisely the structural adaptation regime of \S\,\ref{structural-adaptation-necessity}: an agent whose mismatch exceeds $R$ has exhausted its model class capacity and needs structural change to re-enter the region where parametric correction works.

The agent persists (from within $\mathcal B_R$) iff $R^* < R$, i.e., iff $\alpha > \rho/R$. $\square$

\textbf{Interpretation.} The ultimately bounded region has radius $R^* = \rho/\alpha$. In the linear case, $\alpha = \mathcal{T}$, recovering \S\,\ref{persistence-condition}'s steady-state result $R^* = \rho/\mathcal{T}$ exactly. But Proposition A.1 holds for \textit{any} correction function satisfying the sector condition, not just the linear one.

\textbf{The persistence threshold, generalized.} The agent persists (mismatch remains bounded within the model class capacity) iff $\rho/\alpha < R$. If the correction function breaks down (A2' fails) before $R^*$ is reached, the agent may diverge. This IS \S\,\ref{structural-adaptation-necessity}'s trigger: when $\rho/\alpha > R$ (the environment demands more correction than the model class can provide), parametric adaptation fails and structural change is required.


\medskip\noindent\textbf{Proposition A.2: Stability Margin (Adaptive Reserve).}


\textbf{Statement.} Under the conditions of A.1, the agent can tolerate a sudden increase in disturbance rate of:

% [Derived (adaptive-reserve)]

$$\Delta\rho^* = \alpha R - \rho$$

without mismatch diverging (where $R$ is the radius of the sector-condition region from A2'). Beyond this, $R^*$ exceeds $R$ and the correction function may fail.

\textbf{Proof.} After a shock, the new disturbance rate is $\rho + \Delta\rho$. The new ultimately bounded radius is $(\rho + \Delta\rho)/\alpha$. This remains within the valid region iff $(\rho + \Delta\rho)/\alpha \leq R$, i.e., $\Delta\rho \leq \alpha R - \rho$. $\square$

\textbf{Interpretation.} $\Delta\rho^*$ is the agent's \textbf{adaptive reserve} — how much additional environmental volatility it can absorb before its model breaks down. This is a single number characterizing an agent's robustness to shock:

\begin{itemize}
  \item An agent operating well below capacity ($\rho \ll \alpha R$) has a large reserve — it is \textbf{robust}.
  \item An agent near its limit ($\rho \approx \alpha R$) has a small reserve — it is \textbf{fragile}.
\end{itemize}


\begin{adjustbox}{max width=\textwidth}
\small
\begin{tabular}{lll}
\toprule
Domain & Large $\Delta\rho^*$ (robust) & Small $\Delta\rho^*$ (fragile) \\
\midrule
Control & Kalman filter on slow target & Same filter on erratic target \\
Biology & Organism in stable niche & Same organism under climate change \\
Organization & Well-capitalized firm, stable market & Startup in volatile market \\
Military & Force with operational depth & Force at culmination point \\
\bottomrule
\end{tabular}
\end{adjustbox}



\medskip\noindent\textbf{Proposition A.1S: Bounded Mismatch Under Stochastic Disturbance.}


\textbf{Statement.} Under (A1), (A2'), (A3), with stochastic disturbance (GA-2S: $w(t)$ is zero-mean with $\mathbb{E}[\lVert w(t)\rVert^2] = \sigma_w^2$), the mismatch $\delta(t)$ satisfies:

% [Derived (stochastic-bounded-mismatch, from Itô-Lyapunov analysis)]

$$\mathbb{E}[\lVert\delta(t)\rVert^2] \leq \lVert\delta(0)\rVert^2 e^{-2\alpha t} + \frac{n\sigma_w^2}{2\alpha}$$

for all $t$, where $n = \dim(\delta)$. The agent persists in the mean-square sense iff:

$$\alpha > \frac{n\sigma_w^2}{2R^2}$$

\textbf{Proof.}

The SDE form of the mismatch dynamics under stochastic disturbance is:

$$d\delta = -F(\mathcal{T}, \delta)\,dt + \sigma_w\,dW_t$$

where $W_t$ is a standard $n$-dimensional Wiener process.

Apply Itô's formula to $V(\delta) = \frac{1}{2}\lVert\delta\rVert^2$:

% [Derived (Proof Step)]

$$dV = \delta^T(-F)\,dt + \delta^T \sigma_w\,dW_t + \frac{1}{2}\sigma_w^2 n\,dt$$

The last term is the Itô correction: $\frac{1}{2}\text{tr}(\sigma_w^2 I_n) = \frac{n}{2}\sigma_w^2$.

Taking expectations (the Itô integral $\delta^T \sigma_w\,dW_t$ has zero expectation):

% [Derived (Proof Step)]

$$\frac{d}{dt}\mathbb{E}[V] = \mathbb{E}[\delta^T(-F)] + \frac{n}{2}\sigma_w^2$$

By (A2'): $\delta^T F \geq \alpha\lVert\delta\rVert^2 = 2\alpha V$. Therefore:

% [Derived (Proof Step)]

$$\frac{d}{dt}\mathbb{E}[V] \leq -2\alpha\,\mathbb{E}[V] + \frac{n}{2}\sigma_w^2$$

This is a linear ODE in $\mathbb{E}[V]$ with solution:

$$\mathbb{E}[V(t)] \leq V(0)\,e^{-2\alpha t} + \frac{n\sigma_w^2}{4\alpha}(1 - e^{-2\alpha t})$$

Since $V = \frac{1}{2}\lVert\delta\rVert^2$, the steady-state mean-square mismatch is:

$$\mathbb{E}[\lVert\delta\rVert^2]_{ss} = \frac{n\sigma_w^2}{2\alpha}$$

The RMS steady-state mismatch is:

% [Derived (stochastic-steady-state)]

$$R^*_S \;=\; \lVert\delta\rVert_{\text{rms}} = \sigma_w\sqrt{\frac{n}{2\alpha}}$$

Persistence requires $R^*_S < R$, giving $\alpha > n\sigma_w^2 / (2R^2)$. $\square$

\textbf{Interpretation.} Model D (Prop A.1) gives $R^* = \rho/\alpha$, scaling as $1/\alpha$. Model S gives $R^*_S = \sigma_w\sqrt{n/(2\alpha)}$, scaling as $1/\sqrt{\alpha}$. Doubling the correction efficiency halves the deterministic steady-state mismatch but only reduces the stochastic steady-state by a factor of $\sqrt{2} \approx 1.41$. Correction is less effective against noise than against drift. This difference in scaling propagates into the adversarial exponent regimes ( \S\,\ref{adversarial-exponent-regimes}): $b = 2$ under Model D, $b = 3/2$ under Model S.

\textbf{Tail bound.} At steady state, Markov's inequality gives $P(\lVert\delta\rVert > R) \leq n\sigma_w^2/(2\alpha R^2)$. The agent stays within $R$ with probability $\geq 1 - \epsilon$ provided $\alpha \geq n\sigma_w^2/(2\epsilon R^2)$. For the linear case (Ornstein-Uhlenbeck), the stationary distribution is Gaussian and exact tail probabilities are available. The Markov bound is the general result holding under (A2') alone.


\medskip\noindent\textbf{Summary of Results.}


\begin{adjustbox}{max width=\textwidth}
\small
\begin{tabular}{llll}
\toprule
Result & What it proves & Assumptions & Linear case recovery \\
\midrule
\textbf{A.1} (Bounded Mismatch) & $R^* = \rho/\alpha$ & (A1), (A2'), bounded $\rho$ (GA-2) & $\alpha = \mathcal{T}$ gives $R^* = \rho/\mathcal{T}$ \\
\textbf{A.1S} (Stochastic Bounded Mismatch) & $R^*_S = \sigma_w\sqrt{n/(2\alpha)}$ & (A1), (A2'), stochastic $w$ (GA-2S) & $\alpha = \mathcal{T}$ gives $R^*_S = \sigma_w\sqrt{n/(2\mathcal{T})}$ \\
\textbf{A.2} (Stability Margin) & $\Delta\rho^* = \alpha R - \rho$ & Same as A.1 & $R \to \infty$ for linear (always stable if $\mathcal{T} > 0$) \\
\bottomrule
\end{tabular}
\end{adjustbox}



\medskip\noindent\textbf{Epistemic Status.}


The setup and assumptions are \textit{definitions} — they specify what we mean by "correction function" and "disturbance." Propositions A.1 and A.2 are \textit{exact} — they follow from the assumptions via standard Lyapunov theory (Khalil 2002, Chapters 4 and 9). Proposition A.1S is \textit{exact} under an implicit strengthening of A2': the Ito-Lyapunov proof uses the sector bound on the entire trajectory, but the Wiener process can push $\delta$ outside $\mathcal B_R$ where the local sector condition does not hold. The proof is valid when A2' holds in a sufficiently large region (or globally), or when a stopping-time argument bounds the contribution of excursions. For practical purposes the tail bound (Markov inequality) quantifies the excursion probability and the gap is higher-order when $R^*_S \ll R$. This is a standard subtlety in applied stochastic Lyapunov theory, not specific to ACT. The assumptions themselves (sector condition, bounded disturbance) are \textit{empirical claims} about the qualitative behavior of real correction dynamics. The sector-condition framework originates with Lur'e (1957); the Lyapunov stability results are standard. The application to adaptive agents under the ACT framework is new but the mathematics is not.


\medskip\noindent\textbf{Discussion.}


\textbf{Key value.} The persistence threshold and adaptive reserve are no longer contingent on the linear hypothesis in \S\,\ref{mismatch-dynamics}. They hold for any correction dynamics satisfying the sector condition — a mild qualitative assumption that says "correction points inward with at least baseline efficiency $\alpha$." This is a significant epistemic upgrade: from \textit{hypothesis-dependent} to \textit{robust under qualitative assumptions}.

\textbf{What the proofs do NOT illuminate.} (1) Quantitative steady-state values — Lyapunov gives \textit{bounds}, not exact values; the linear analysis remains necessary for quantitative predictions. (2) Convergence rates — standard Lyapunov tells you stable/unstable, not how fast. (3) Optimal gain structure — \S\,\ref{update-gain} comes from estimation theory, not stability theory. (4) Model sufficiency — the \S\,\ref{information-bottleneck} framework is information-theoretic, complementary to but independent of stability analysis.

\textbf{Two disturbance models.} Proposition A.1 assumes bounded disturbance (GA-2: $\lVert w(t)\rVert \leq \rho$); Proposition A.1S assumes stochastic disturbance (GA-2S: $\mathbb{E}[\lVert w(t)\rVert^2] = \sigma_w^2$). These are not approximations to each other — they capture structurally different environments. Model D (bounded) covers persistent directional change: an adversary who maneuvers, an API that drifts, a climate that shifts. Model S (stochastic) covers unpredictable fluctuations around a stable mean: market noise, sensor noise, random perturbations. The choice of model is an empirical question for each domain; the theory provides both tools. Neither model handles heavy-tailed environmental shocks (financial crises, ecological catastrophes, strategic surprise) where $\lVert w(t)\rVert$ can exceed any finite bound with non-negligible probability. Extreme tail events are better understood as triggers for structural adaptation ( \S\,\ref{structural-adaptation-necessity}) rather than disturbances to be absorbed parametrically.

\end{derivation}

\begin{derivation}[Gain-Sector Bridge — Proofs and Verification]
\label{gain-sector-derivation}
Complete derivations for the results stated in \S\,\ref{gain-sector-bridge}: gain-based updating produces correction functions satisfying the sector condition (GA-3) whenever the update rule has directional fidelity. Covers the Kalman case (scalar and matrix), the general bridge theorem, the gradient-descent equivalence, and numerical verification.



\medskip\noindent\textbf{Motivation.}


\S\,\ref{sector-condition-derivation} proves persistence and stability downstream of the sector condition (GA-3): $\delta^T F(\mathcal{T}, \delta) \geq \alpha \lVert\delta\rVert^2$. The sector parameter $\alpha$ governs the persistence bound ($R^* = \rho/\alpha$), the adaptive reserve ($\Delta\rho^* = \alpha R - \rho$), and the adversarial scaling exponent. But whether the correction function $F$ induced by the gain principle ( \S\,\ref{update-gain}) actually satisfies GA-3 was not derived — GA-3 was the theory's softest structural joint.

This appendix closes the gap. For optimal Bayesian updates (Kalman, conjugate families), the bridge is exact: the sector parameter equals the gain. For gradient-based agents, GA-3 is equivalent to local strong convexity of the loss function. The load-bearing assumption shifts from "the correction function satisfies the sector condition" (opaque) to "the update rule has directional fidelity" (transparent, checkable).



\medskip\noindent\textbf{Proposition B.1: Scalar Kalman Sector Condition.}


\textbf{Statement.} For a scalar linear-Gaussian system ($n = m = 1$, $H = 1$) with Kalman filter gain $K = P^- / (P^- + R_{\text{obs}})$, the correction function $F(e) = K \cdot e$ satisfies the sector condition with:

% [Derived (scalar-Kalman-sector, from Kalman update)]

$$\alpha = K = \frac{P^-}{P^- + R_{\text{obs}}} = \eta^*$$

The sector parameter equals the Kalman gain, which equals the uncertainty-ratio gain from \S\,\ref{update-gain}. The bound is tight (the correction is linear).

\textbf{Proof.} The scalar Kalman update for state estimate $\hat x$ given innovation $\delta = o - \hat x$ is $\hat x_\text{post} = \hat x_\text{prior} + K \delta$, where the gain is $K = P^- / (P^- + R_\text{obs})$ with $P^-$ the prior variance and $R_\text{obs}$ the observation noise variance.

The state-space estimation error $e = x - \hat{x}$ evolves under the correction step as:

$$e_{\text{post}} = e_{\text{prior}} - K(He_{\text{prior}} + \varepsilon) = (1 - K) e_{\text{prior}} - K\varepsilon$$

With $H = 1$, the expected correction function is $F(e) = K \cdot e$. The sector product is:

$$e \cdot F(e) = K \cdot e^2$$

Since $P^- > 0$ and $R_\text{obs} > 0$, we have $K \in (0, 1)$. The sector condition $e \cdot F(e) \geq \alpha \cdot e^2$ holds with $\alpha = K$. The bound is tight because the correction is exactly linear. $\square$

\textbf{Steady-state sector parameter.} At steady state (random walk dynamics, $A = 1$), the prior variance satisfies the algebraic Riccati equation. Solving $K^2 R_{\text{obs}} + QK - Q = 0$ (where $Q$ is the process noise variance) gives:

% [Derived (steady-state scalar sector parameter)]

$$\alpha_{ss} = K_{ss} = \frac{-Q + \sqrt{Q^2 + 4 Q R_{\text{obs}}}}{2 R_{\text{obs}}}$$

Limiting behavior: when $Q \gg R_{\text{obs}}$ (fast dynamics, clean observations), $K_{ss} \to 1$; when $Q \ll R_{\text{obs}}$ (slow dynamics, noisy observations), $K_{ss} \approx \sqrt{Q / R_{\text{obs}}}$.

\textbf{Connection to ACT quantities.} The adaptive tempo for a single observation channel at rate $\nu$ is $\mathcal{T} = \nu \cdot K_{ss}$, and the sector parameter in the continuous-time framework is $\alpha = \mathcal{T}$. The bridge is trivial in the scalar case: the gain IS the sector parameter.



\medskip\noindent\textbf{Proposition B.2: Matrix Kalman Sector Condition.}


\textbf{Statement.} For a linear-Gaussian system with state $x \in \mathbb{R}^n$, observations $o \in \mathbb{R}^m$, observation matrix $H$, and Kalman gain $K = P^- H^T (H P^- H^T + R_{\text{obs}})^{-1}$, the correction function $F(e) = KH \cdot e$ satisfies the sector condition in the $(P^-)^{-1}$-weighted inner product, restricted to the observable subspace $\mathcal{O} = \operatorname{range}(H^T)$:

% [Derived (matrix-Kalman-sector, from Kalman covariance reduction)]

$$e^T (P^-)^{-1} KH \, e = e^T H^T S^{-1} H \, e \geq \alpha \cdot e^T (P^-)^{-1} e$$

where $S = H P^- H^T + R_{\text{obs}}$ is the innovation covariance, and:

$$\alpha = 1 - \lambda_{\max}(P_{t\vert t} \, P_{t\vert t-1}^{-1})$$

restricted to observable directions. In unobservable directions ($e \in \ker(H)$), $\alpha = 0$.

\textbf{Proof.} Define the $(P^-)^{-1}$-weighted inner product: $\langle u, v \rangle_{P^-} = u^T (P^-)^{-1} v$. In this inner product the sector product becomes:

$$\langle e, KH \, e \rangle_{P^-} = e^T (P^-)^{-1} KH \, e$$

Substituting $K = P^- H^T S^{-1}$:

$$e^T (P^-)^{-1} P^- H^T S^{-1} H \, e = e^T H^T S^{-1} H \, e$$

Since $S \succ 0$, the matrix $H^T S^{-1} H$ is symmetric positive semidefinite, with null space $\ker(H)$. So $\langle e, KH \, e \rangle_{P^-} \geq 0$ for all $e$, with equality iff $He = 0$.

For the sector bound, we need $e^T H^T S^{-1} H \, e \geq \alpha \cdot e^T (P^-)^{-1} e$. From the Kalman covariance identity $KH = I - P (P^-)^{-1}$ (where $P = P_{t\vert t}$ is the posterior covariance), the eigenvalues of $KH$ in the $(P^-)^{-1}$-weighted sense are $1 - \lambda_i(P(P^-)^{-1})$, where $\lambda_i$ denotes generalized eigenvalues.

Since $P \preceq P^-$ (the posterior is no less certain than the prior), $\lambda_i(P(P^-)^{-1}) \leq 1$, so all eigenvalues of $KH$ are non-negative. In observable directions, $P \prec P^-$ strictly, giving $\alpha > 0$. In unobservable directions, $P = P^-$ (no information gained), so $\alpha = 0$. $\square$

\textbf{Fully observable diagonal case.} When $H = I$, $A = I$, $Q = \operatorname{diag}(q_1, \ldots, q_n)$, $R_{\text{obs}} = \operatorname{diag}(r_1, \ldots, r_n)$, the problem decouples into $n$ scalar problems (each governed by Prop B.1), and:

% [Derived (diagonal-Kalman-sector)]

$$\alpha = \min_{i=1}^n K_i = \min_{i=1}^n \frac{-q_i + \sqrt{q_i^2 + 4 q_i r_i}}{2 r_i}$$

The bottleneck is the dimension with the worst signal-to-noise ratio, consistent with \S\,\ref{per-dimension-persistence}.

\textbf{The weighted-norm subtlety.} The Lyapunov proofs in \S\,\ref{sector-condition-derivation} use the Euclidean norm ($V = \frac{1}{2}\lVert e\rVert^2$), while the matrix Kalman sector condition holds in the $(P^-)^{-1}$-weighted norm. For fully observable systems with bounded condition number $\kappa(P^-)$, the norms are equivalent:

$$\alpha_{\text{Euclidean}} \geq \alpha_{\text{weighted}} / \kappa(P^-)$$

The Lyapunov results remain valid with this quantitative adjustment. The weighted Lyapunov function $V(e) = \frac{1}{2} e^T (P^-)^{-1} e$ yields the tighter bound directly.



\medskip\noindent\textbf{Proposition B.3: Bridge Theorem.}


\textbf{Statement.} Given the gain-based update $M_t = M_{t-1} + \eta^* \cdot g(\delta_t)$ ( \S\,\ref{update-gain}), with $H$ mapping state-space corrections to observation-space mismatch reduction, the induced correction function $F(\delta) = \eta^* \cdot H \, g(\delta)$ satisfies the sector condition (GA-3) with parameter $\alpha > 0$ whenever:

\textbf{(B1) Directional fidelity.} The composite mismatch transform preserves the mismatch-reducing direction:

% [Derived (bridge-theorem, from update-gain + B1)]

$$\delta^T H \, g(\delta) \geq c \lVert\delta\rVert^2 \quad \text{for } \lVert\delta\rVert \leq R$$

for some $c > 0$. The sector parameter is:

$$\alpha = \eta^* \cdot c_{\min}, \qquad c_{\min} = \inf_{\lVert\delta\rVert \leq R} \frac{\delta^T H \, g(\delta)}{\lVert\delta\rVert^2}$$

\textbf{Proof.} The gain-based update changes the model state by $\Delta M = \eta^* \cdot g(\delta)$. The induced change in predicted observation is $H \cdot \Delta M = \eta^* \cdot H \, g(\delta)$, so the expected correction function (the mismatch reduction per step) is:

$$F(\delta) = \eta^* \cdot H \, g(\delta)$$

The sector product is:

$$\delta^T F(\delta) = \eta^* \cdot \delta^T H \, g(\delta)$$

By B1: $\delta^T H \, g(\delta) \geq c_{\min} \lVert\delta\rVert^2$. Since $\eta^* > 0$ (the agent updates), we have:

$$\delta^T F(\delta) \geq \eta^* \cdot c_{\min} \lVert\delta\rVert^2 = \alpha \lVert\delta\rVert^2 \qquad \square$$

\textbf{When B1 holds by construction.} For optimal Bayesian updates (Kalman, conjugate families, exponential families), the posterior update minimizes expected loss, ensuring the correction aligns with the mismatch. B1 holds by optimality, not by assumption. The gain principle $\eta^* = U_M / (U_M + U_o)$ determines only the magnitude; the direction is guaranteed by the update rule's optimality. For gradient descent on a convex loss, B1 follows from gradient monotonicity (the defining property of convexity).



\medskip\noindent\textbf{Proposition B.4: Gradient Equivalence.}


\textbf{Statement.} For any agent updating via gradient descent on a differentiable loss $L$ with learning rate $\eta$ and minimizer $M^*$ (where $\nabla L(M^*) = 0$):

% [Derived (sector-convexity-equivalence)]

$$\text{GA-3 holds with } (\alpha, R) \iff L \text{ is locally } (\alpha/\eta)\text{-strongly convex on } \mathcal{B}_R(M^*)$$

The sector parameter factors as $\alpha = \eta \cdot \mu$, where:

$$\mu = \inf_{\lVert\delta\rVert \leq R} \lambda_{\min}(\nabla^2 L(M^* + \delta))$$

is the strong convexity modulus. The basin radius $R$ is the largest ball around $M^*$ where $\nabla^2 L$ remains positive definite.

\textbf{Proof.} The correction function for gradient descent is $F(\delta) = \eta \cdot \nabla L(M^* + \delta)$ (continuous-time form absorbing event rate $\nu$ into $\mathcal{T} = \nu \cdot \eta$). The sector condition requires:

$$\delta^T F(\delta) = \eta \cdot \delta^T \nabla L(M^* + \delta) \geq \alpha \lVert\delta\rVert^2$$

Dividing by $\eta > 0$:

$$\delta^T \nabla L(M^* + \delta) \geq \frac{\alpha}{\eta} \lVert\delta\rVert^2$$

A differentiable function $L$ is $\mu$-strongly convex on $\mathcal{S}$ iff the gradient monotonicity condition holds (Nesterov 2004, Theorem 2.1.10): for all $x, y \in \mathcal{S}$,

$$(\nabla L(x) - \nabla L(y))^T (x - y) \geq \mu \lVert x - y\rVert^2$$

Setting $x = M^* + \delta$, $y = M^*$, and using $\nabla L(M^*) = 0$:

$$\nabla L(M^* + \delta)^T \delta \geq \mu \lVert\delta\rVert^2$$

This is identical to the sector condition divided by $\eta$. The equivalence follows:

($\Rightarrow$) GA-3 gives $\delta^T \nabla L(M^* + \delta) \geq (\alpha/\eta) \lVert\delta\rVert^2$. Since $\nabla L(M^*) = 0$, gradient monotonicity holds with $\mu = \alpha/\eta$.

($\Leftarrow$) $\mu$-strong convexity gives $\delta^T \nabla L(M^* + \delta) \geq \mu \lVert\delta\rVert^2$. Multiplying by $\eta$: $\delta^T F(\delta) \geq \eta\mu \lVert\delta\rVert^2$. So GA-3 holds with $\alpha = \eta\mu$. $\square$

**The $\alpha$ factorization.** The decomposition $\alpha = \eta \cdot \mu$ separates agent design ($\eta$, responsiveness) from environment structure ($\mu$, loss curvature). The adaptive tempo maps as $\alpha \propto \mathcal{T}$: higher gain or steeper curvature means faster mismatch correction.



\medskip\noindent\textbf{Loss Function Classification.}


The gradient equivalence (Prop B.4) classifies which loss functions satisfy GA-3, and with what parameters:

\begin{adjustbox}{max width=\textwidth}
\small
\begin{tabular}{lllll}
\toprule
Loss class & Strong convexity & $\alpha$ & Basin $R$ & Status \\
\midrule
Quadratic (Kalman, linear regression) & Global: $\mu = \lambda_{\min}(H)$ & $\eta \cdot \lambda_{\min}(H)$ & $\infty$ & Exact \\
Exponential family (natural params) & Global: $\mu = \lambda_{\min}(\text{Fisher})$ & $\eta \cdot \lambda_{\min}(\text{Fisher})$ & $\infty$ & Exact \\
L2-regularized convex loss & Global: $\mu \geq \lambda$ & $\eta \cdot \lambda$ (floor) & $\infty$ & Exact \\
Unregularized logistic & Global convex, not strongly & $\eta \cdot \mu(R)$, decaying & $\infty$ (but $\alpha \to 0$) & Local \\
Quasi-convex & $\delta^T \nabla L \geq 0$ but ratio $\to 0$ & $\eta \cdot \mu(R) \to 0$ & Finite effective & Local \\
Non-convex within basin & Local: $\mu > 0$ in basin & $\eta \cdot \mu_\text{local}$ & Distance to nearest saddle & Local \\
Non-convex beyond basin & Fails: $\lambda_{\min}(\nabla^2 L) < 0$ & N/A & N/A & Fails \\
\bottomrule
\end{tabular}
\end{adjustbox}


\textbf{Key observations:}

\begin{enumerate}
  \item \textit{All exponential family models in natural parameter form satisfy GA-3 globally.} This subsumes Kalman, Beta-Bernoulli, Gaussian, and Poisson posteriors. The Hessian is the Fisher information matrix, which is positive definite in the interior of the natural parameter space.
\end{enumerate}


\begin{enumerate}
  \item \textit{L2 regularization transforms locally strongly convex losses into globally strongly convex ones.} The regularization parameter $\lambda$ provides a floor: $\alpha \geq \eta \lambda$ everywhere.
\end{enumerate}


\begin{enumerate}
  \item \textit{Non-convex losses satisfy GA-3 within each basin of attraction.} The basin radius $R$ is the distance from the local minimum to the nearest inflection surface (where $\lambda_{\min}(\nabla^2 L) = 0$). Beyond the basin, the sector condition fails — this IS the structural-adaptation trigger from \S\,\ref{structural-adaptation-necessity}.
\end{enumerate}




\medskip\noindent\textbf{Failure Modes.}


The bridge fails in precisely five cases. These characterize the exact conditions under which the gain-based update does NOT produce a correction function satisfying the sector condition.


\medskip\noindent\textbf{FM-1: Directional Infidelity.}


The mismatch transform $g$ rotates the correction away from the mismatch direction, so $\delta^T H g(\delta) \leq 0$.

\textit{Example.} If $\delta \in \mathbb{R}^2$ and $g(\delta) = R_{90}\delta$ (a 90-degree rotation), then $\delta^T g(\delta) = 0$ for all $\delta$. The correction is perpendicular to the mismatch and never reduces it.

\textit{When it occurs.} Pathological parameterizations where the model-space update direction is misaligned with observation-space mismatch. For optimal Bayesian updates, B1 holds by construction — the posterior minimizes expected loss.


\medskip\noindent\textbf{FM-2: Gain Collapse.}


$\eta^* \to 0$ while $\rho > 0$, so $\alpha \to 0$ and the persistence condition $\alpha > \rho/R$ eventually fails. An agent accumulating experience in a non-stationary environment will eventually have gain too small to track changes. This is not a bridge failure but a persistence-condition failure — the bridge holds ($\alpha = \eta^* \cdot c > 0$ for finite experience), but $\alpha$ falls below $\rho/R$. See \S\,\ref{update-gain}.


\medskip\noindent\textbf{FM-3: Nonlinear Saturation.}


The correction function $g$ saturates at large $\lVert\delta\rVert$, so the sector ratio $\delta^T g(\delta) / \lVert\delta\rVert^2$ decays. For a saturating function such as $g(\delta) = \tanh(\delta)$, the local sector parameter is:

$$\alpha(R) = \eta^* \cdot \frac{\tanh(R)}{R}$$

For small $R$, $\alpha \approx \eta^*$ (linear regime). For large $R$, $\alpha \approx \eta^* / R$ (saturated regime). The sector condition holds locally — the valid region shrinks as the nonlinearity strengthens. This is exactly what the local form A2' in \S\,\ref{sector-condition-derivation} anticipates.


\medskip\noindent\textbf{FM-4: Unobservable Directions.}


When $\ker(H) \neq \{0\}$ (fewer observations than state dimensions), the correction has no effect on mismatch components in $\ker(H)$. The sector condition holds only in the observable subspace $\mathcal{O} = \operatorname{range}(H^T)$. This is a structural limitation: the agent cannot correct errors it cannot observe. See \S\,\ref{observability-dominance}.

For persistence in the full state space, the system must be \textit{detectable}: all unstable modes must be observable. Stable unobservable modes decay naturally without correction.


\medskip\noindent\textbf{FM-5: Model Misspecification.}


The model class does not contain the truth, so the gradient direction is systematically wrong. B1 fails because the correction aims at the wrong target. The sector parameter degrades proportionally to the misspecification. This is the \S\,\ref{structural-adaptation-necessity} trigger: the correction function fails not because the gain is wrong but because the correction \textit{direction} is wrong.



\medskip\noindent\textbf{Simulation Results Summary.}


Six numerical experiments validate the gradient equivalence (Prop B.4) and the loss-function classification. Full code and traces are in \texttt{msc/spike-gain-sector-bridge-nonlinear.md}.


\medskip\noindent\textbf{Experiment 1: Quadratic Loss (Global Strong Convexity).}


$L(w) = \frac{1}{2}(w - w^*)^T H(w - w^*)$ with $w \in \mathbb{R}^5$, Hessian eigenvalues $\{0.5, 1.0, 2.0, 3.0, 5.0\}$, $\eta = 0.1$. The sector ratio $\delta^T F / \lVert\delta\rVert^2$ at every gradient-descent step fell within $[\eta\mu, \eta L] = [0.05, 0.50]$, confirming Prop B.4 exactly. No step violated the bound.


\medskip\noindent\textbf{Experiment 2: Logistic Regression (Position-Dependent Convexity).}


Logistic regression with $d = 2$, $N = 5000$ samples. Unregularized: sector ratio positive along the entire trajectory, with local $\mu$ ranging from 0.076 (at MLE) to higher values far from it. With L2 regularization ($\lambda = 0.1$): sector ratio bounded below by $\eta\lambda = 0.03$ at every point, confirming the global floor.


\medskip\noindent\textbf{Experiment 3: Gaussian Mixture (Non-Convex, Basin Structure).}


Two-component Gaussian mixture. Within basin ($r < 3.1$): Hessian positive definite, sector ratio positive, gradient descent converges to correct optimum. At basin boundary ($r \approx 3.1$): Hessian eigenvalue crosses zero. Beyond basin ($r > 3.5$): sector ratio negative, gradient descent converges to wrong optimum. Basin boundary confirmed as sharp transition, matching the \S\,\ref{structural-adaptation-necessity} trigger.


\medskip\noindent\textbf{Experiment 4: Local Sector Parameter Measurement.}


Same Gaussian mixture. Sector parameter $\alpha(r)$ decreases monotonically from 181.4 (at $r = 0.5$) to 125.8 (at $r = 3.0$), then drops to zero at $r \approx 3.1$. The decay is moderate within the basin (69\% retention at the boundary); the transition at the boundary is sharp.


\medskip\noindent\textbf{Experiment 5: Exponential Family (Poisson, Natural Parameters).}


Poisson in natural parameter space. Sector ratio positive at every step, ranging from 0.00177 to 0.00240 (with $\eta = 0.001$). Confirms global strong convexity of exponential family losses in natural parameterization.


\medskip\noindent\textbf{Experiment 6: Quasi-Convex and Multi-Modal Losses.}


\textit{Quasi-convex} ($L(w) = -1/(1+w^2)$): sector ratio always positive but decays as $O(1/w^4)$. No uniform lower bound exists — $\alpha(R) \to 0$, making persistence bounds vacuous for large $R$.

\textit{Multi-modal} ($L(w) = \sin(w) + 0.01w^2$): sector condition holds within each basin ($\lvert\delta\rvert < 3.0$), fails at basin boundary (gradient reversal at local maximum). Basin radius $\approx$ distance from local minimum to nearest local maximum.



\medskip\noindent\textbf{Epistemic Status.}


\textbf{Propositions B.1 and B.2} are \textit{exact} — they follow from standard Kalman filter algebra. The only subtlety is the weighted norm in the matrix case (B.2), which is a quantitative refinement, not a qualitative gap. Both are standard results in estimation theory reframed in sector-condition language.

\textbf{Proposition B.3} (Bridge Theorem) is a \textit{conditional derivation}: exact under B1 (directional fidelity). B1 holds by construction for optimal Bayesian updates (the posterior minimizes expected loss, ensuring the correction aligns with the mismatch). For approximate update rules, B1 is a design condition, not a global assumption. The condition is transparent and checkable for specific systems.

\textbf{Proposition B.4} (Gradient Equivalence) is an \textit{exact mathematical equivalence} — the bidirectional implication is proved, not assumed. The loss-function classification follows directly. The simulation experiments confirm the theoretical predictions across all six loss classes tested; no violations were observed where the theory predicts the sector condition holds.

\textbf{Max attainable:} \textit{conditional} for the bridge (B1 is inherent — pathological update rules exist), \textit{exact} for the gradient equivalence. The condition cannot be removed: there exist correction functions that violate the sector condition (FM-1 through FM-5).

\textbf{What remains open.} (1) Non-gradient agents (PID controllers, rule-based systems, human judgment): GA-3 remains an empirical claim for these agent classes. (2) Time-varying $\alpha$: when the gain adapts (Adam, RMSprop, Kalman transient), $\alpha(t)$ varies; the Lyapunov analysis extends to time-varying $\alpha(t) \geq \underline{\alpha} > 0$ via standard results (Khalil 2002, Chapter 8). (3) Stochastic gradients: SGD satisfies the sector condition \textit{in expectation}; the per-step noise enters as effective disturbance in the Prop A.1S framework.



\medskip\noindent\textbf{Discussion.}


\textbf{The formal chain is closed for well-designed agents.} The prediction chain becomes:

$$\text{gain principle} + \text{B1 (directional fidelity)} \;\xrightarrow{\text{Prop B.3 (exact)}}\; \text{sector condition (GA-3)} \;\xrightarrow{\text{Props A.1, A.1S, A.2 (exact)}}\; \text{persistence, reserve, scaling}$$

For optimal Bayesian updates, B1 holds by construction, making the full chain a derivation. For gradient agents, B1 is equivalent to local strong convexity — a well-characterized property.

\textbf{GA-3 is grounded, not floating.} The assumption load shifts from "the correction function satisfies the sector condition" (what correction function?) to "the update rule has directional fidelity" (the correction at least points toward reality). For the important agent classes — Kalman filters, conjugate posteriors, exponential families, gradient descent on convex losses — this is automatic.

\textbf{Basin boundary = structural adaptation trigger.} For gradient agents with non-convex losses, the basin radius $R$ is the convexity radius of the loss landscape. When mismatch exceeds $R$, the correction function reverses direction. This is the \S\,\ref{structural-adaptation-necessity} trigger with a precise geometric characterization: structural adaptation is needed when the agent crosses an inflection surface of its loss landscape.

\end{derivation}

\begin{derivation}[Recursive Update — Uniqueness Derivation]
\label{recursive-update-derivation}
Derivation showing that $M_{\tau^+} = f(M_{\tau^-}, e_\tau)$ is the \textit{unique} update form consistent with directed time, partial observability, and state completeness. Not merely one option, but the only one.



\medskip\noindent\textbf{Setup.}


We work within ACT's scope ( \S\,\ref{scope-condition}): an agent coupled to an environment $\Omega$ through observation and action channels, with residual uncertainty.

**Universe of information at event time $\tau$.** The following information exists (in the broadest ontological sense) at the moment event $e_\tau$ occurs:

\begin{adjustbox}{max width=\textwidth}
\small
\begin{tabular}{ll}
\toprule
Information & Description \\
\midrule
$\Omega_\tau$ & The environment state \\
$\mathcal{C}_{\tau^-}$ & The complete interaction history ( \S\,\ref{chronica}) up to (but not including) $e_\tau$ \\
$\{M_{\tau'}\}_{\tau' \leq \tau^-}$ & The agent's prior internal states, culminating in $M_{\tau^-}$ \\
$e_\tau$ & The current event (observation arriving or action completing) \\
$\{e_{\tau'}\}_{\tau' > \tau}$ & Future events (not yet occurred) \\
\bottomrule
\end{tabular}
\end{adjustbox}


The question: of these, which can the update $M_{\tau^+}$ depend on?


\medskip\noindent\textbf{The Three Constraints.}


\textbf{Constraint 1 — Arrow of time ( \S\,\ref{causal-structure} postulate).} Events are temporally ordered and this ordering is irreversible. An update occurring at time $\tau$ cannot depend on events that have not yet occurred:

$$M_{\tau^+} \text{ cannot depend on } \{e_{\tau'}\}_{\tau' > \tau}$$

This is a physical constraint — the most primitive one. In a classical universe, information from the future is simply not available. Even if the agent can \textit{predict} future events, those predictions are part of $M_{\tau^-}$ (they are internal computations, not future information).

\textbf{Constraint 2 — Partial observability ( \S\,\ref{scope-condition}).} The agent cannot access $\Omega_\tau$ directly. Its only interface with the environment is through the event $e_\tau$, which is a lossy function of $\Omega_\tau$ (via \S\,\ref{observation-function}):

$$M_{\tau^+} \text{ cannot depend on } \Omega_\tau \text{ except through } e_\tau$$

This is a scope constraint. If the agent could access $\Omega$ directly, the residual uncertainty condition in \S\,\ref{scope-condition} would be trivially violable.

\textbf{Constraint 3 — State completeness ( \S\,\ref{agent-model}).} $M_{\tau^-}$ is the agent's \textit{complete} internal state just before event $e_\tau$. There is no information about the agent's past that is available to the update mechanism but not encoded in $M_{\tau^-}$:

$$M_{\tau^+} \text{ cannot depend on } \mathcal{C}_{\tau^-} \text{ or } \{M_{\tau'}\}_{\tau' < \tau^-} \text{ except through } M_{\tau^-}$$

This constraint does the most interesting work and deserves careful examination (see Discussion below).


\medskip\noindent\textbf{The Derivation.}


\textbf{Result (Recursive Update Uniqueness).} Under Constraints 1–3, the model update at event time $\tau$ must have the form

$$M_{\tau^+} = f(M_{\tau^-}, e_\tau)$$

for some function $f: \mathcal{M} \times \mathcal{E} \to \mathcal{M}$. No other update form is consistent with the three constraints.

\textbf{Derivation.} Consider the most general possible update. The updated state $M_{\tau^+}$ is a function of \textit{all accessible information}:

$$M_{\tau^+} = F(\text{accessible information at } \tau)$$

We characterize the accessible information by eliminating what is not accessible.

\textbf{(i) Eliminate future events.} By C1 (arrow of time), $\{e_{\tau'}\}_{\tau' > \tau}$ is not accessible.

After this elimination, the candidate dependency set is:
$$\{\Omega_\tau,\; \mathcal{C}_{\tau^-},\; \{M_{\tau'}\}_{\tau' \leq \tau^-},\; e_\tau\}$$

\textbf{(ii) Eliminate direct environment access.} By C2 (partial observability), the agent cannot access $\Omega_\tau$ except through the event $e_\tau$. Any information from $\Omega_\tau$ that reaches the agent does so through $e_\tau$ — already in the dependency set.

After this elimination:
$$\{\mathcal{C}_{\tau^-},\; \{M_{\tau'}\}_{\tau' \leq \tau^-},\; e_\tau\}$$

**(iii) Reduce past information to $M_{\tau^-}$.** By C3 (state completeness), $M_{\tau^-}$ is the agent's complete internal state. Every element of $\mathcal{C}_{\tau^-}$ and every prior model state $M_{\tau'}$ ($\tau' < \tau^-$) that could influence the update can do so \textit{only through} its effect on $M_{\tau^-}$. The agent's internal state evolves through a sequence of updates; the cumulative effect of all prior events is exactly $M_{\tau^-}$. The raw events that produced this state are no longer separately available — they were "consumed" by the update mechanism and their information (to the extent it was retained) is now encoded in $M_{\tau^-}$.

Could the agent maintain a separate log of raw events outside of $M$? It could — but that log *is part of $M$*. Whatever information the agent retains in any form — model parameters, cached data, raw event buffers, metadata — is by definition part of its complete internal state $M_{\tau^-}$. If something is available to the update mechanism and not in $M_{\tau^-}$, then $M_{\tau^-}$ was not the complete state — contradicting C3.

After this elimination:
$$\{M_{\tau^-},\; e_\tau\}$$

Therefore:
$$M_{\tau^+} = F(M_{\tau^-}, e_\tau) \equiv f(M_{\tau^-}, e_\tau)$$

This is the unique form: no information beyond $(M_{\tau^-}, e_\tau)$ is accessible under the three constraints, so no update form depending on anything else is realizable. $\square$

\textbf{Corollary (Between-events dynamics).} Between events, no new event $e$ arrives. The same argument applies with $e_\tau$ removed from the accessible set:

$$\frac{dM}{d\tau} = g(M_\tau)$$

The agent's internal evolution between events (prediction, decay, internal simulation) depends only on the current state. $\square$

\textbf{Corollary (Serial special case).} When observations and actions alternate at a uniform rate on a single channel, each event $e_t$ is the pair $(o_t, a_{t-1})$. The update becomes:

$$M_t = f(M_{t-1}, o_t, a_{t-1})$$

This is the familiar discrete-time form. $\square$


\medskip\noindent\textbf{Information-Set Formalization.}


For readers who prefer a measure-theoretic framing:

The agent's \textbf{information set} at time $\tau$ is the sigma-algebra $\mathcal{I}_\tau^{agent}$ — the collection of events (in the probability-theoretic sense) about which the agent can condition its update.

\begin{itemize}
  \item \textbf{C1} restricts $\mathcal{I}_\tau^{agent} \subseteq \sigma(\{e_{\tau'} : \tau' \leq \tau\} \cup \{\Omega_\tau\} \cup \{M_{\tau'} : \tau' \leq \tau^-\})$ — no future information.
  \item \textbf{C2} further restricts: $\sigma(\Omega_\tau) \setminus \sigma(e_\tau)$ is not in $\mathcal{I}_\tau^{agent}$ — the agent cannot condition on aspects of $\Omega_\tau$ not captured by $e_\tau$.
  \item \textbf{C3} further restricts: $\sigma(\{e_{\tau'} : \tau' < \tau\} \cup \{M_{\tau'} : \tau' < \tau^-\}) \subseteq \sigma(M_{\tau^-})$ from the agent's perspective.
\end{itemize}


After all three restrictions: $\mathcal{I}_\tau^{agent} = \sigma(M_{\tau^-}, e_\tau)$.

By the Doob–Dynkin lemma, any $\sigma(M_{\tau^-}, e_\tau)$-measurable random variable is a (Borel) function of $(M_{\tau^-}, e_\tau)$. Therefore $M_{\tau^+} = f(M_{\tau^-}, e_\tau)$ for some measurable $f$. $\square$


\medskip\noindent\textbf{Attempts to Break the Result.}


Before trusting the proof, seven counterexample attacks:


\medskip\noindent\textbf{Attack 1: Simultaneous events.}


Two events arrive at exactly the same time: $e_\tau^{(1)}$ and $e_\tau^{(2)}$. The update has three arguments: $f(M_{\tau^-}, e_\tau^{(1)}, e_\tau^{(2)})$.

\textbf{Verdict:} Not deep — \S\,\ref{event-driven-dynamics} defines events as atomic. If we allow bundled events, the form holds with $e_\tau$ as a set. Reveals that "event" needs careful definition, but the form is preserved.


\medskip\noindent\textbf{Attack 2: Continuous environmental influence.}


An agent embedded in a physical system experiences continuous forces (gravity, temperature, electromagnetic fields). These aren't "events" in \S\,\ref{event-driven-dynamics}'s sense; they're continuous signals. The true dynamics would be $dM/d\tau = g(M_\tau, o(\tau))$ where $o(\tau)$ is a continuous observation stream.

\textbf{Verdict:} Genuine limitation of the event-driven formulation. The between-events corollary $dM/d\tau = g(M_\tau)$ holds only when the agent is truly isolated between events. For continuous coupling, the analogous result is the general state-space representation $\dot{M} = g(M, u)$ from control theory — arrived at by the same three constraints. The event-driven version is a special case for digital/sampled systems.


\medskip\noindent\textbf{Attack 3: The C3 circularity.}


C3 defines $M$ as the agent's complete internal state. Any apparent counterexample is dissolved by expanding $M$. Consider: an agent has a "model" (neural net weights) and a "replay buffer" (stored raw events). C3 says $M = (\text{weights}, \text{buffer})$. The model space is just larger than you thought.

\textbf{Verdict:} The deepest objection. The proof essentially: (1) Define $M$ to be everything the agent has. (2) Observe the update can only use what the agent has. (3) Therefore $f(M_{\tau^-}, e_\tau)$. The real content is the \textit{analytical commitment}: by defining $M$ as complete, we commit to Markovian analysis, which then makes \S\,\ref{model-sufficiency} the right quality metric. See Epistemic Status below.


\medskip\noindent\textbf{Attack 4: Shared state between agents.}


Agents A and B share a common memory bank (shared database). The clean resolution is the multi-agent framework: the shared memory is part of the \textit{composite} system's state, and each agent's interaction with it is mediated by events (reads and writes). Not a true counterexample but highlights that C3 requires careful delineation of agent boundaries.


\medskip\noindent\textbf{Attack 5: External randomness not in $e_\tau$.}


Hardware thermal noise used in the update. The stochastic case $M_{\tau^+} \sim P(\cdot \mid M_{\tau^-}, e_\tau)$ is a special case of $f$ where $f$ is a randomized function. The \textit{form} — dependence on exactly $(M_{\tau^-}, e_\tau)$ — is preserved. The result statement should explicitly allow stochastic $f$.


\medskip\noindent\textbf{Attack 6: Time-dependent updates.}


Could $f$ depend on the timestamp $\tau$ itself? Yes — consistently. The event $e_\tau$ in \S\,\ref{event-driven-dynamics} carries a timestamp: $e_\tau = (\text{type}, \text{channel}, \text{payload}, \tau)$. So time-dependence enters through $e_\tau$. Alternatively, the agent may maintain an internal clock as part of $M_{\tau^-}$. Either way, $f(M_{\tau^-}, e_\tau)$ accommodates time-dependence.


\medskip\noindent\textbf{Attack 7: Agents that store full history.}


An agent with $M_{\tau^-} \supseteq \mathcal{C}_{\tau^-}$ is entirely consistent. The model space $\mathcal{M}$ is simply large enough to include the raw history. The \S\,\ref{information-bottleneck} argues compression is \textit{wise} — but the recursive update form holds regardless of compression level.


\medskip\noindent\textbf{Epistemic Status.}


The result is correct but partly definitional. The three constraints have different epistemic characters:

\begin{adjustbox}{max width=\textwidth}
\small
\begin{tabular}{lll}
\toprule
Constraint & Character & Can it be violated? \\
\midrule
C1 (arrow of time) & Physical law & Not in a classical universe \\
C2 (partial observability) & Scope definition & Only by leaving ACT's scope \\
C3 (state completeness) & Analytical commitment & Not without redefining $M$ \\
\bottomrule
\end{tabular}
\end{adjustbox}


C1 and C2 do genuine eliminative work — they rule out update forms that depend on future events or on raw $\Omega$. These are non-trivial constraints.

C3 is a definitional commitment that produces the Markov structure. It cannot be "violated" because any violation is absorbed by expanding $M$. This is not a weakness — it's the nature of the claim. The result says: *the Markovian analysis is the only one consistent with C1 + C2 + the definition of $M$ as complete*. The alternative — an update that depends on something outside $M$ — is not "wrong" but rather means $M$ was misspecified.

\textbf{What the result says:} C1 eliminates a physically impossible class of updates (future-dependent). C2 eliminates a scope-excluded class ($\Omega$-dependent). After (1) and (2), the \textit{only remaining question} is how the past enters: through the full history $\mathcal{C}_{\tau^-}$ or through a compressed state $M_{\tau^-}$. C3 says the agent \textit{has} a complete state, and whatever that state is, it's all the agent has. The Markov form follows.

\textbf{What the result does NOT say:} That $M$ must be a lossy compression (the agent could store full history). That the Markov property is "natural" or "optimal" (it's a consequence of how $M$ is defined). That continuous-coupling systems are event-driven (the event framework is one abstraction; $\dot{M} = g(M, u)$ is the more general one, arrived at by the same three constraints).


\medskip\noindent\textbf{Discussion.}


\textbf{Recursion as a consequence of completeness.} The recursive form is not an assumption bolted on — it follows from the definition of $M_t$ as complete. The sufficiency of the recursive form is precisely what \S\,\ref{model-sufficiency} measures: when $S(M_t) = 1$, the recursive update loses nothing.

\textbf{What this opens.} The proof yields the \textit{form}. It immediately invites the follow-up questions that the rest of the theory addresses: What should $f$ preserve? → \S\,\ref{information-bottleneck} and \S\,\ref{model-sufficiency}. How should $f$ weight new information? → \S\,\ref{update-gain}. When is $\mathcal{M}$ itself inadequate? → \S\,\ref{structural-adaptation-necessity}.

\end{derivation}

\begin{sketch}[Multi-Timescale Stability]
\label{multi-timescale-stability}
When adaptive processes operate at $N$ nested timescales, composite stability requires each level to be stable given its slower levels, with sufficient timescale separation between adjacent pairs.



\medskip\noindent\textbf{Formal Expression.}




\medskip\noindent\textbf{The General $N$-Timescale System.}


The temporal nesting in \S\,\ref{temporal-nesting} creates a coupled multi-timescale system with $N$ levels. Singular perturbation theory provides tools to analyze such systems. Define a hierarchy of state variables:


$$x^{(1)}, \; x^{(2)}, \; \ldots, \; x^{(N)}$$

where $x^{(1)}$ is the fastest (e.g., mismatch at the reactive/parametric level) and $x^{(N)}$ is the slowest (e.g., architectural or meta-structural state). The coupled dynamics:


$$\dot{x}^{(k)} = \frac{1}{\epsilon_k} \, G^{(k)}\!\left(x^{(1)}, \ldots, x^{(N)}\right) + w^{(k)}(t)$$

where $\epsilon_1 \ll \epsilon_2 \ll \cdots \ll \epsilon_N$ encode the timescale separation and each $G^{(k)}$ may depend on the states at all levels.


\medskip\noindent\textbf{The Two-Timescale Special Case.}


The simplest nontrivial instance has $N = 2$:

\begin{itemize}
  \item Fast state $x^{(1)} = \delta$ (mismatch under parametric adaptation)
  \item Slow state $x^{(2)} = \mathcal{M}$ (model class, changing on a structural timescale)
\end{itemize}


$$\dot{x}^{(1)} = -F(\mathcal{T}, x^{(1)}; x^{(2)}) + w(t) \quad \text{(fast: parametric adaptation)}$$

$$\dot{x}^{(2)} = \epsilon \, G(x^{(1)}, x^{(2)}) \quad \text{(slow: structural adaptation)}$$

where $\epsilon \ll 1$ reflects the timescale separation and $F$ depends on $x^{(2)}$ (the correction function is determined by the current model class).


\medskip\noindent\textbf{Sketch of Approach (General Case).}


The standard singular perturbation result (Tikhonov's theorem, generalized) applies layer by layer: if level $k$ is stable for each fixed configuration of the slower levels $k+1, \ldots, N$ (each level has a stable attractor given the levels above it), and each successive slow manifold is itself stable, then the composite $N$-level system is stable.

\S\,\ref{temporal-nesting}'s convergence constraint $\nu_{n+1} \ll \nu_n$ is the condition ensuring sufficient timescale separation at each boundary — i.e., $\epsilon_k / \epsilon_{k+1} \ll 1$ for each $k$. When this separation is violated between any adjacent pair, the faster level's transients contaminate the slower level's dynamics, potentially destabilizing the composite system.


\medskip\noindent\textbf{Epistemic Status.}


This is a \textit{sketch}, not a complete result. The framework and approach are presented as a guide for future development. The claim that timescale separation ensures composite stability is a standard result in singular perturbation theory; the application to ACT's nested adaptive levels is new but follows the standard pattern.

Making it rigorous requires specifying the dynamics $G^{(k)}$ for levels deeper than parametric adaptation. \S\,\ref{structural-adaptation-necessity} gives the \textit{trigger condition} for structural change but not the \textit{dynamics} of how change at deeper levels proceeds. Specifying these would require theories of how agents search over model classes, modify their own architecture, or restructure their adaptive mechanisms — open problems in RL (architecture search, meta-learning), biology (evolutionary dynamics), and organizational theory (institutional change).


\medskip\noindent\textbf{Discussion.}


\textbf{The convergence constraint as stability condition.} The sketch suggests that \S\,\ref{temporal-nesting}'s convergence constraint is not merely a heuristic but a formal condition for composite stability across arbitrarily many timescales. This connects the empirical observation (don't let deeper-level changes happen too fast) to a stability-theoretic foundation.

\textbf{Applicability to LLM systems.} LLMs involve many parallel adaptive processes — pretraining (slowest), fine-tuning, LoRA-style adaptation, in-context learning, retrieval/RAG updates, tool-use feedback, and within-generation attention dynamics — without clean boundaries between "parametric" and "structural." The $N$-timescale framework accommodates this naturally: each mechanism operates at its characteristic rate, and the stability analysis requires only that adjacent timescales be sufficiently separated, regardless of how many levels exist or how they are labeled.

\end{sketch}

\begin{detail}[Linear ODE Approximation]
\label{linear-ode-approximation}
The full linear treatment of mismatch dynamics: scalar and vector forms, steady-state solutions under both disturbance models, transient behavior, convergence rate, validity conditions, breakdown modes, discrete-time connection, and adversarial coupling. This appendix collects the linear-case results that \S\,\ref{mismatch-dynamics} states as a hypothesis, derives them explicitly, and delineates exactly where the linear approximation is valid and where it fails.



\medskip\noindent\textbf{Formal Expression.}



\medskip\noindent\textbf{1. The scalar and vector forms.}


The mismatch vector $\delta(t) \in \mathbb{R}^n$ evolves under the general dynamics from \S\,\ref{sector-condition-derivation}:


$$\frac{d\delta}{dt} = -F(\mathcal{T}, \delta) + w(t)$$

The \textbf{linear approximation} sets $F(\mathcal{T}, \delta) = \mathcal{T} \cdot \delta$, giving:

% [Hypothesis (linear-vector-ode)]

$$\frac{d\delta}{dt} = -\mathcal{T} \cdot \delta + w(t)$$

This is exact when the correction function is linear in $\delta$ (Kalman filter, Beta-Bernoulli conjugate update). It is an approximation otherwise.

\textbf{Scalar form (equality).} For a scalar mismatch $\delta(t) \in \mathbb{R}$:

% [Derived (scalar-ode, from linear-vector-ode with $n = 1$)]

$$\frac{d\delta}{dt} = -\mathcal{T} \cdot \delta + w(t)$$

This is an equality: the scalar ODE is exactly the $n = 1$ case of the vector ODE.

\textbf{Norm form (inequality).} For vector $\delta \in \mathbb{R}^n$, take the time derivative of $\lVert\delta\rVert = (\delta^T \delta)^{1/2}$:

$$\frac{d\lVert\delta\rVert}{dt} = \frac{\delta^T \dot\delta}{\lVert\delta\rVert} = \frac{\delta^T(-\mathcal{T}\delta + w)}{\lVert\delta\rVert} = -\mathcal{T}\lVert\delta\rVert + \frac{\delta^T w}{\lVert\delta\rVert}$$

By Cauchy-Schwarz, $\delta^T w / \lVert\delta\rVert \leq \lVert w\rVert$. Under Model D (GA-2: $\lVert w(t)\rVert \leq \rho$):

% [Derived (norm-inequality, from Cauchy-Schwarz)]

$$\frac{d\lVert\delta\rVert}{dt} \leq -\mathcal{T} \cdot \lVert\delta\rVert + \rho$$

\textbf{This is an inequality, not an equality.} The Cauchy-Schwarz step is tight only when $w(t)$ is parallel to $\delta(t)$ (worst-case disturbance). The norm form stated in \S\,\ref{mismatch-dynamics} is this upper bound. For the scalar case ($n = 1$), Cauchy-Schwarz is automatically tight and the inequality becomes an equality.


\medskip\noindent\textbf{2. Steady-state solutions.}



\medskip\noindent\textbf{Model D (deterministic bounded disturbance, GA-2: $\lVert w(t)\rVert \leq \rho$).}


Setting $d\lVert\delta\rVert/dt = 0$ in the norm inequality gives the worst-case steady state:

% [Derived (model-d-steady-state, from norm inequality)]

$$\lVert\delta\rVert_{ss} = \frac{\rho}{\mathcal{T}}$$

This is the tight upper bound: steady-state mismatch equals the ratio of disturbance to correction. In the scalar case ($n = 1$), this is exact (the Ornstein-Uhlenbeck process with deterministic forcing converges to $\rho/\mathcal{T}$). In the vector case, it is the worst case over all disturbance directions.


\medskip\noindent\textbf{Model S (stochastic zero-mean disturbance, GA-2S: $d\delta = -\mathcal{T}\delta\,dt + \sigma_w\,dW_t$).}


The Ornstein-Uhlenbeck process has stationary variance derived via Ito-Lyapunov analysis (Prop A.1S in \S\,\ref{sector-condition-derivation}, specialized to $\alpha = \mathcal{T}$):

% [Derived (model-s-steady-state, from Ito-Lyapunov with linear correction)]

$$\mathbb{E}[\lVert\delta\rVert^2]_{ss} = \frac{n\sigma_w^2}{2\mathcal{T}}$$

The RMS steady-state mismatch is:

$$\lVert\delta\rVert_{\text{rms}} = \sigma_w\sqrt{\frac{n}{2\mathcal{T}}}$$

For the scalar case ($n = 1$): $\lVert\delta\rVert_{\text{rms}} = \sigma_w / \sqrt{2\mathcal{T}}$.

\textbf{Scaling difference.} Model D gives $\lVert\delta\rVert_{ss} \propto 1/\mathcal{T}$; Model S gives $\lVert\delta\rVert_{\text{rms}} \propto 1/\sqrt{\mathcal{T}}$. Doubling the correction rate halves the deterministic steady-state mismatch but only reduces the stochastic steady-state by a factor of $\sqrt{2} \approx 1.41$. Correction is less effective against noise than against drift.


\medskip\noindent\textbf{3. Transient solution.}


**Model D, constant $\rho$.** The linear ODE $d\lVert\delta\rVert/dt = -\mathcal{T}\lVert\delta\rVert + \rho$ with initial condition $\lVert\delta(0)\rVert = \lVert\delta_0\rVert$ has the standard first-order linear solution:

% [Derived (transient-model-d)]

$$\lVert\delta(t)\rVert = \lVert\delta_0\rVert\,e^{-\mathcal{T} t} + \frac{\rho}{\mathcal{T}}(1 - e^{-\mathcal{T} t})$$

Mismatch decays exponentially from initial conditions toward the steady state $\rho/\mathcal{T}$ with time constant $\tau = 1/\mathcal{T}$. After $k$ time constants, the transient has decayed by a factor of $e^{-k}$: 63\% after one time constant, 95\% after three, 99.3\% after five.

\textbf{Model S, Ornstein-Uhlenbeck.} The mean-square mismatch evolves as (from Prop A.1S with $\alpha = \mathcal{T}$):

% [Derived (transient-model-s)]

$$\mathbb{E}[\lVert\delta(t)\rVert^2] = \lVert\delta_0\rVert^2\,e^{-2\mathcal{T} t} + \frac{n\sigma_w^2}{2\mathcal{T}}(1 - e^{-2\mathcal{T} t})$$

The variance converges twice as fast as the mean (rate $2\mathcal{T}$ vs. $\mathcal{T}$) because the Lyapunov function $V = \frac{1}{2}\lVert\delta\rVert^2$ is quadratic — its dynamics have doubled exponential rate.


\medskip\noindent\textbf{4. When the linear approximation is valid.}


The linear case $F(\mathcal{T}, \delta) = \mathcal{T} \cdot \delta$ is a special case of the sector condition (GA-3) from \S\,\ref{sector-condition-stability} with:

% [Derived (linear-sector-parameters)]

$$\alpha = \mathcal{T}, \qquad R \to \infty$$

That is, the sector condition holds globally with lower bound $\alpha$ equal to the tempo, and the sector-condition region has infinite radius. In the notation of \S\,\ref{gain-sector-bridge}:

\begin{itemize}
  \item \textbf{Directional fidelity parameter:} $c_{\min} = 1$ (the correction points exactly at the mismatch).
  \item \textbf{Sector parameter:} $\alpha = \eta^* \cdot c_{\min} = \eta^* = \mathcal{T}/\nu$.
  \item \textbf{No upper saturation:} $c_{\max} = c_{\min} = 1$ (the sector bounds coincide).
\end{itemize}


\textbf{Consequence for persistence.} With $R \to \infty$, structural persistence ($\alpha > \rho/R$) is trivially satisfied for any $\mathcal{T} > 0$. The binding constraint is task adequacy alone: $\mathcal{T} > \rho/\lVert\delta_{\text{critical}}\rVert$ (Model D) or $\mathcal{T} > n\sigma_w^2 / (2\lVert\delta_{\text{critical}}\rVert^2)$ (Model S). This is why the linear operational forms in \S\,\ref{persistence-condition} exist: in the linear world, structural persistence is free.

\textbf{The linear approximation is exact when:}

\begin{enumerate}
  \item \textbf{Kalman filter} (scalar or matrix): the gain $K_t$ produces $F(\delta) = K_t H \delta$, which is linear. The sector parameter equals the Kalman gain: $\alpha = K$ (scalar) or $\alpha = \lambda_{\min}^+(KH)$ (matrix). See \S\,\ref{gain-sector-bridge}.
\end{enumerate}


\begin{enumerate}
  \item \textbf{Beta-Bernoulli conjugate update:} the correction $F(\delta) = \delta/(n+1)$ is linear with $\alpha = 1/(n+1) = \eta_{\text{edge}}$.
\end{enumerate}


\begin{enumerate}
  \item \textbf{Exponential family with natural parameters:} the natural-parameter update is linear in the sufficient statistic, giving $\alpha = \eta \cdot \lambda_{\min}(\text{Fisher})$.
\end{enumerate}


\begin{enumerate}
  \item \textbf{Gradient descent on quadratic loss:} $F(\delta) = \eta \nabla^2 L \cdot \delta$ is linear with $\alpha = \eta \cdot \lambda_{\min}(\nabla^2 L)$.
\end{enumerate}



\medskip\noindent\textbf{5. When the linear approximation breaks down.}


The linear form $F = \mathcal{T} \cdot \delta$ fails when the true correction function deviates from linearity. Three failure modes, each corresponding to a specific nonlinearity:

**Saturation at large $\lVert\delta\rVert$.** The correction mechanism is overwhelmed by large mismatch. The true correction satisfies $F(\delta) < \mathcal{T}\lVert\delta\rVert$ for large $\lVert\delta\rVert$ — correction is slower than the linear prediction. Simulation confirms: for a saturating function $g(\delta) = \delta / (1 + \lvert\delta\rvert/R)$, the sector parameter at the capacity boundary is $\alpha \approx \mathcal{T}/2$ (half the linear value). The linear approximation overstates persistence margins. The sector-condition framework ( \S\,\ref{sector-condition-stability}) handles this via finite $R$ and reduced $\alpha$.

**Threshold effects at small $\lVert\delta\rVert$.** Below a detection threshold $\varepsilon$, small mismatches go uncorrected: $F(\delta) \approx 0$ for $\lVert\delta\rVert < \varepsilon$. This creates a dead zone where the model drifts. The linear approximation (which predicts correction at all scales) misses this. The sector condition fails locally at small $\lVert\delta\rVert$ (the ratio $\delta^T F / \lVert\delta\rVert^2 \to 0$), but the dead zone is bounded and does not affect large-scale stability.

\textbf{Structural breakdown.} Beyond some critical mismatch, the correction rate drops to zero because the model class is no longer appropriate: the mismatch exceeds the model's representational capacity. $F(\delta) \approx 0$ for $\lVert\delta\rVert > R$. This is the structural adaptation trigger ( \S\,\ref{structural-adaptation-necessity}). The linear approximation, which predicts correction growing without bound, misses this entirely. The sector-condition framework captures it via the finite radius $R$ of the sector-condition region.


\medskip\noindent\textbf{6. Discrete-time connection.}


The continuous ODE is a fluid-limit approximation of the discrete event-driven dynamics ( \S\,\ref{event-driven-dynamics}). The discrete mismatch dynamics are:

$$\delta_{k+1} = \delta_k - \eta^* F_d(\delta_k) + w_k$$

In the linear case ($F_d(\delta) = \delta$), this becomes:

$$\delta_{k+1} = (1 - \eta^*)\delta_k + w_k$$

which is an AR(1) process with coefficient $\lambda = 1 - \eta^*$.

\S\,\ref{discrete-sector-condition} proves the formal connection:

\textbf{Model D.} The discrete and continuous steady states are \textit{identical} — the fluid-limit gap is exactly zero. Both give $R^* = \rho/\alpha$. The discrete form has $\alpha = (1 - \lvert\lambda\rvert)/\eta^* = (1 - \lvert 1 - \eta^*\rvert)/\eta^*$, which equals $1$ (hence $\alpha_{\text{continuous}} = \nu \cdot 1 = \mathcal{T}$) when $0 < \eta^* < 1$.

\textbf{Model S.} The discrete steady-state variance is $\sigma^2_{\text{step}} / (1 - \lambda^2_{\text{eff}})$, which recovers the continuous result $n\sigma_w^2/(2\mathcal{T})$ in the fluid limit ($\eta^* \to 0$, $\nu \to \infty$, $\nu\eta^* = \mathcal{T}$ fixed). The variance gap is $O((\eta^*)^2 c^2_{\max})$ — quantitatively small whenever $\eta^* c_{\max} \ll 1$.

\textbf{Fluid-limit validity.} The ODE approximation is formally justified by the fluid limit theorem in \S\,\ref{discrete-sector-condition} (conditional on Lipschitz regularity of $F_d$). The approximation is most accurate when $\eta^* \ll 1$ (the small-gain regime). It is least accurate during initial transients when $\eta^*$ is large, but this phase is short-lived and the transient error is bounded by $O(\eta^* c_{\max} / \nu^{1/2})$.


\medskip\noindent\textbf{7. Adversarial coupling in the linear case.}


When two agents are coupled, $A$'s praxis contributes to $B$'s disturbance:

% [Hypothesis (Linear Coupling Model)]

$$\rho_B = \rho_{B,\text{base}} + \gamma_A \cdot \mathcal{T}_A$$

Under the linear approximation, the steady-state mismatches (Model D, coupling-dominant: $\gamma_A \mathcal{T}_A \gg \rho_{B,\text{base}}$) are:

% [Derived (linear-adversarial-steady-state, from Model D steady state + coupling model)]

$$\lVert\delta_B\rVert_{ss} \approx \frac{\gamma_A \mathcal{T}_A}{\mathcal{T}_B}, \qquad \lVert\delta_A\rVert_{ss} \approx \frac{\gamma_B \mathcal{T}_B}{\mathcal{T}_A}$$

The ratio:

% [Derived (squared-tempo-law, from linear steady states)]

$$\frac{\lVert\delta_B\rVert_{ss}}{\lVert\delta_A\rVert_{ss}} = \frac{\gamma_A}{\gamma_B}\left(\frac{\mathcal{T}_A}{\mathcal{T}_B}\right)^2$$

Under symmetric coupling ($\gamma_A \approx \gamma_B$), tempo advantage squares: a 2:1 tempo ratio yields a 4:1 mismatch ratio; a 3:1 ratio yields 9:1.

\textbf{Why the squared law.} The exponent 2 arises because the $1/\mathcal{T}$ steady-state scaling appears in both numerator (faster agent generates more disturbance) and denominator (faster agent corrects better): one power from the coupling, one from the steady-state. Under Model S (stochastic coupling, $1/\sqrt{\mathcal{T}}$ scaling), the exponent reduces to $3/2$. See \S\,\ref{adversarial-exponent-regimes} for the full regime analysis.

\textbf{Simulation validation.} Variant A confirmed $b = 1.999$ (Model D, coupling-dominant). Variants C-D confirmed $b = 1.481$ (Model S, coupling-dominant). Both within numerical precision of the derived values. See \S\,\ref{simulation-results}.


\medskip\noindent\textbf{Epistemic Status.}


\textit{Exact for the linear case.} The linear ODE is a standard first-order linear system; the solutions (transient and steady-state, both Model D and Model S) are textbook results for the Ornstein-Uhlenbeck process and its deterministic analog. The sector-condition parameters ($\alpha = \mathcal{T}$, $R \to \infty$) are exact when the correction function is linear. The discrete-time connection is exact (zero gap for Model D, $O((\eta^*)^2)$ gap for Model S). The adversarial scaling law ($b = 2$ for Model D, $b = 3/2$ for Model S) is derived and simulation-validated.

\textbf{What is exact vs. what is an approximation.} The linear ODE is exact for systems with linear correction dynamics (Kalman, Beta-Bernoulli, quadratic loss). For nonlinear systems, the linear ODE is a local Taylor approximation valid near the steady state. The sector-condition framework ( \S\,\ref{sector-condition-stability}, \S\,\ref{sector-condition-derivation}) is the general treatment; this appendix derives what falls out when the sector bounds coincide.

**The norm form is an upper bound for $n > 1$.** The inequality introduced by Cauchy-Schwarz at the norm step is tight only in the worst case. For typical disturbances that are not persistently aligned with the mismatch direction, the actual norm trajectory lies below the upper bound. The scalar case ($n = 1$) is an equality. This distinction is load-bearing: downstream results that use $d\lVert\delta\rVert/dt = -\mathcal{T}\lVert\delta\rVert + \rho$ as an equality (e.g., the exact transient solution, the exact steady-state ratio) are strictly correct only for $n = 1$ or for worst-case disturbance. The Lyapunov analysis in \S\,\ref{sector-condition-derivation} avoids this issue by working with $V = \frac{1}{2}\lVert\delta\rVert^2$ directly, without passing through the norm derivative.

\textbf{Max attainable:} exact (the linear case is its own exact analysis; there is nothing to tighten).


\medskip\noindent\textbf{Discussion.}


\textbf{Why this appendix exists.} The linear ODE in \S\,\ref{mismatch-dynamics} is labeled as a hypothesis and described as a "first-order approximation." That is correct for the general case. But for agents with linear correction dynamics, the ODE is not an approximation at all -- it is exact. This appendix serves two roles: (1) providing the complete derivation that \S\,\ref{mismatch-dynamics} points to rather than carrying inline, and (2) making precise when the "approximation" is exact and when it genuinely introduces error.

\textbf{Pedagogical value.} The linear case is the entry point for understanding ACT's mismatch dynamics. The closed-form solutions (exponential convergence, ratio steady state, squared scaling law) build intuition that the general sector-condition results ( \S\,\ref{sector-condition-stability}) cannot provide -- Lyapunov gives bounds and existence, not quantitative trajectories. The linear treatment also connects directly to the AR(1) simulation framework ( \S\,\ref{simulation-results}) and classical control theory (PI controllers, Kalman filters).

\textbf{Speed-quality substitutability.} From $\mathcal{T} = \nu \cdot \eta^*$ (single channel): doubling event rate $\nu$ has the same effect on $\lVert\delta\rVert_{ss}$ as doubling update quality $\eta^*$. They are multiplicative when both improve: 50\% improvement in each yields $1.5 \times 1.5 = 2.25\times$, not $3\times$. This is the formal analog of Boyd's insight that Orient quality often matters more than raw OODA speed.

\textbf{The persistence threshold in the linear world.} From the steady state: $\lVert\delta\rVert_{ss} < \lVert\delta_{\text{critical}}\rVert$ iff $\mathcal{T} > \rho/\lVert\delta_{\text{critical}}\rVert$ (Model D). This is task adequacy only -- structural persistence is trivially satisfied because $R \to \infty$. The linear world makes persistence look easier than it is. For nonlinear agents, both structural and task-adequacy conditions must hold ( \S\,\ref{persistence-condition}).

\end{detail}

\begin{derivation}[Graph Structure Uniqueness]
\label{graph-structure-uniqueness}
Four operational postulates — directed temporal ordering, probabilistic uncertainty, state-local revisability, and observable intermediates — derive that the strategy representation is a directed acyclic graph (proved), with the Markov property conditional on causal sufficiency of the strategy.



\medskip\noindent\textbf{Formal Expression.}



\medskip\noindent\textbf{The Postulates.}


Four properties that a strategy representation must satisfy. Each is independently motivated from the adaptive-systems foundation.


\medskip\noindent\textbf{P1: Directed Temporal Ordering.}


% [Derived (from CREF:causal-structure:CREF)]

If component $A$ of the strategy causally produces component $B$, then $A$ temporally precedes $B$. The strategy representation must respect this directionality — edges point from causes to effects, from actions to outcomes, from prerequisites to goals.

This is a consequence of the temporal postulate ( \S\,\ref{causal-structure}): the arrow of time is constitutive, not incidental. Reversing a causal edge would mean effects precede causes, which is physically impossible.


\medskip\noindent\textbf{P2: Probabilistic Uncertainty.}


% [Derived (from Cox's theorem)]

The agent's uncertainty about whether each step of the strategy will succeed must be quantified by a measure satisfying Cox's axioms (consistency, universality, non-negativity). The unique such measure is probability (Cox 1946). The agent may use other representations internally (confidence scores, fuzzy logic), but these must be mappable to probability to be consistent.


\medskip\noindent\textbf{P3: State-Local Revisability.}


% [Derived (from CREF:chain-confidence-decay:CREF + bounded computation)]

When the agent observes evidence about one component of its strategy (e.g., "step 3 succeeded" or "prerequisite 2 is blocked"), it must be able to update its beliefs about that component and its consequences without recomputing the entire strategy from scratch.

\textbf{Why this is forced, not chosen:}

\textit{From fragility.} Additive log-confidence ( \S\,\ref{chain-confidence-decay}) means longer chains are exponentially less reliable. The agent will frequently encounter partial failures. Each partial failure requires re-evaluation of the affected portion of the strategy. If each re-evaluation requires full recomputation, the agent's planning tempo $T_\Sigma$ is catastrophically slow — potentially violating the strategy persistence condition.

\textit{From bounded computation.} The agent has finite computational resources (the IB constraint applies to planning as well as model maintenance). Full recomputation of a strategy with $N$ components costs $O(N)$ or worse. Local revision costs $O(\lvert\text{affected}\rvert)$, which can be much smaller.

\textit{From the persistence condition.} Strategy must be revised faster than the environment invalidates it. Local revision directly increases $T_\Sigma$ by reducing the per-update cost. An agent that must recompute everything on each update has lower $T_\Sigma$ and is more likely to fall below the persistence threshold.


\medskip\noindent\textbf{P4: Observable Intermediates.}


% [Derived (from CREF:chain-confidence-decay:CREF + monitoring requirement)]

The strategy representation must have internal checkpoints — observable states between the initial action and the final goal — that the agent can monitor to detect partial failure.

Without intermediates, the agent cannot detect chain failure until the final outcome. By the time the final outcome reveals failure, all intermediate actions are wasted. With intermediates, the agent can detect failure at step $k$ and revise, saving the cost of steps $k+1$ through $n$. The value of early detection grows with chain length, because longer chains fail more often (P2 + \S\,\ref{chain-confidence-decay}).


\medskip\noindent\textbf{The Derivation.}



\medskip\noindent\textbf{Step 1: P1 implies directed edges.}


Each component $X_i$ of the strategy has a set of direct causes $\text{Pa}(X_i)$ — the components whose outcomes directly influence $X_i$'s outcome. P1 requires that these causal relationships are directed: $\text{Pa}(X_i)$ temporally precedes $X_i$. This gives directed edges $\text{Pa}(X_i) \to X_i$.


\medskip\noindent\textbf{Step 2: P2 implies probability distributions on edges.}


Each edge carries uncertainty: $P(X_i \mid \text{Pa}(X_i))$. By P2 (Cox), this is a probability distribution. The joint distribution over all strategy components is some $P(X_1, \ldots, X_n)$.


\medskip\noindent\textbf{Step 3: P3 + P1 + causal sufficiency imply the Markov condition.}


% [Derived (Conditional on causal sufficiency of $\Sigma_t$)]

\textbf{Claim.} P3 (state-local revisability) combined with P1 (temporal ordering) and causal sufficiency of the strategy representation implies the local Markov property.

\textbf{The local Markov property.} Each variable $X_i$ is conditionally independent of its non-descendants given its parents:

$$X_i \perp \text{NonDesc}(X_i) \mid \text{Pa}(X_i)$$

\textbf{The argument has three parts:}

\textbf{(a) P3 requires some local neighborhood to suffice.} When the agent updates its belief about $X_i$, it must do so using a local neighborhood $\text{Ne}(X_i)$ rather than the entire graph. For the local update to be correct:

$$P(X_i \mid \text{Ne}(X_i), \text{rest}) = P(X_i \mid \text{Ne}(X_i))$$

P3 establishes that SOME such neighborhood exists. It does not specify which.

\textbf{(b) P1 identifies parents as the natural neighborhood.} In a directed graph respecting temporal ordering, causal influence on $X_i$ flows from temporally preceding nodes through directed edges. The nodes with direct causal influence on $X_i$ are precisely $\text{Pa}(X_i)$ — the parents. Non-parent predecessors can only affect $X_i$ through the parents (their influence is mediated). Non-descendants are either predecessors (influence mediated through parents) or temporally unrelated (no causal path to $X_i$). Therefore: in a temporally ordered directed graph, $\text{Pa}(X_i)$ is the minimal local neighborhood that captures all direct causal influence on $X_i$.

\textbf{(c) Causal sufficiency makes parents sufficient.} The above requires that all direct causes of $X_i$ ARE represented as parents in the graph — that there are no latent common causes between $X_i$ and its non-descendants. This is the causal sufficiency assumption.

For $\Sigma_t$ specifically, causal sufficiency is a reasonable assumption: the agent \textit{constructed} the strategy graph, so all nodes are explicit. There are no "hidden" strategy components — the agent chose every node. If environmental factors affect multiple strategy steps (shared infrastructure, weather, market conditions), they should appear as condition nodes in $\Sigma_t$ ( \S\,\ref{strategy-dag}, source constraint). If they are omitted, the strategy is model-inadequate — exactly the situation \S\,\ref{structural-adaptation-necessity} addresses — and the Markov condition fails for the explicit nodes. This is a model quality issue, not a structural limitation of the argument.

\textbf{Assembling (a)-(c).} P3 requires a sufficient local neighborhood. P1 identifies parents as the causal neighborhood. Causal sufficiency ensures parents capture all direct influences. Therefore:

$$P(X_i \mid \text{Pa}(X_i), \text{NonDesc}(X_i)) = P(X_i \mid \text{Pa}(X_i))$$

This IS the local Markov property.

\textbf{Equivalence to factorization.} For positive distributions (all configurations have nonzero probability — a mild regularity condition), the local Markov property is equivalent to the global Markov property and to the factorization property (Lauritzen 1996, Theorem 3.27; Verma 1993):

$$P(X_1, \ldots, X_n) = \prod_{i=1}^{n} P(X_i \mid \text{Pa}(X_i))$$

This is the defining property of a Bayesian network.

\textbf{What about non-parent local neighborhoods?} The earlier sketch noted that other sparse factorizations (e.g., factor graphs with message-passing) might also support local revision. This is true — factor-graph message passing achieves exact inference for trees and approximate inference for loopy graphs, using a different local neighborhood (the factor's scope, not the parents). But P1 (directed temporal ordering) selects the parent-conditional factorization specifically, because the causal direction matters: the agent needs to reason about consequences of interventions ($do(X_i)$), and Pearl's do-calculus is defined on the parent-conditional DAG factorization, not on undirected factor graphs. Directed local neighborhoods (parents) are strongly motivated by the combination of locality (P3) and causal directionality (P1); undirected alternatives sacrifice the causal semantics that \S\,\ref{strategy-dag} requires.

\textbf{Why non-descendants specifically.} Descendants of $X_i$ depend on $X_i$ by construction — they are downstream. Updating $X_i$ based on $\text{Pa}(X_i)$ does not require knowing descendants, because causal influence flows FROM $X_i$ TO descendants, not the reverse. Learning about descendants can provide evidence about $X_i$ (diagnostic reasoning / explaining away), but that evidence flows through the graph structure itself via belief propagation — it does not violate P3 because the update still propagates locally.


\medskip\noindent\textbf{Step 4: P1 + finite horizon implies acyclicity (proved).}


This is the strongest piece of the argument. See the dedicated section below.


\medskip\noindent\textbf{Step 5: Assembly.}


P1 (directed edges) + P2 (probabilistic) + P3 (Markov condition) + P4 (internal nodes) + finite horizon (acyclicity):

\textbf{The strategy representation must be representable as a directed acyclic graph with probability distributions at each node conditioned on its parents — a Bayesian network.}


\medskip\noindent\textbf{Acyclicity Derivation.}


% [Derived (from CREF:causal-structure:CREF + finite planning horizon)]

This resolves a former known fragility in the theory. Acyclicity of $\Sigma_t$ is derived, not assumed.

\textbf{Result.} For a strategy representation over a finite future horizon, temporal ordering forces acyclicity.

\textbf{Derivation.}

\begin{enumerate}
  \item Each node $X_i$ in $\Sigma_t$ represents a future event or state with temporal position $\tau_i > t$ (the future time at which the step occurs or the state is evaluated).
  \item Each edge $X_i \to X_j$ requires $\tau_i < \tau_j$ (P1: causes temporally precede effects).
  \item A cycle $X_i \to X_j \to \cdots \to X_i$ would require $\tau_i < \tau_j < \cdots < \tau_i$, which is impossible for a real-valued time index.
  \item Therefore the graph is acyclic. $\square$
\end{enumerate}


Formally: a finite set with a strict partial order (future events ordered by time) is representable as a DAG. This is a standard result in order theory — every finite partial order has a Hasse diagram, which is a DAG.

\textbf{The iteration objection resolved.} A strategy that says "try $A$, if fail try $B$, if fail try $A$ again" appears cyclic.

In the time-indexed representation:

$$A_1 \to \text{check}_1 \to B_1 \to \text{check}_2 \to A_2 \to \ldots$$

Each attempt is a distinct node at a distinct time. The apparent cycle is a linear chain in the unrolled view. Iteration "terminates" when either a node succeeds (remaining retry nodes become probability-zero), the agent exhausts its resource budget (a constraint truncating the chain), or the horizon ends. Any finite-horizon strategy, including those with "loops" in the informal sense, is acyclic when time-indexed.

\textbf{Scope.} This applies to $\Sigma_t$ (the agent's strategy over the future), not to $M_t$'s model of the environment. $M_t$ may include cyclic causal processes — feedback loops in the physical world, market dynamics, ecosystem interactions. The acyclicity is specific to the purposeful substate because $\Sigma_t$ represents planned future actions and the future is partially ordered by time. $M_t$'s model of environmental dynamics may need to represent cycles (via time-unrolled DBNs or other cyclic structures).

\textbf{Connection to Pearl.} Pearl's do-calculus is defined on DAGs. Extensions to cyclic structures exist (cyclic SCMs, equilibrium models) but are substantially more complex and lose some of do-calculus's clean properties. The temporal argument here shows that for strategy representations (future-looking plans), acyclicity is not a convenience restriction on Pearl's framework — it is a consequence of the temporal structure of planning.


\medskip\noindent\textbf{What Is Derived vs. What Is Chosen.}


\begin{adjustbox}{max width=\textwidth}
\small
\begin{tabular}{lll}
\toprule
Property & Motivating postulate & Strength \\
\midrule
Directed edges & Temporal ordering (P1, \S\,\ref{causal-structure}) & Proved \\
Probabilistic uncertainty & Cox's theorem (P2) & Proved \\
Acyclicity & Temporal ordering + finite horizon (P1) & Proved \\
Internal structure & Fragility + monitoring (P4, \S\,\ref{chain-confidence-decay}) & Derived \\
Markov factorization & Local revisability (P3) + temporal ordering (P1) + causal sufficiency & Conditional — holds for causally sufficient strategies \\
\textbf{DAG with Markov property} & \textbf{P1 + P2 + P3 + P4 + causal sufficiency} & \textbf{Conditional — causal sufficiency is the remaining assumption} \\
AND/OR parameterization & Boolean completeness + parsimony & Hypothesis (binary outcomes only) \\
Single-parameter edges & Parsimony / IB & Formulation choice \\
Specific node ontology & — & Formulation choice \\
\bottomrule
\end{tabular}
\end{adjustbox}


The dividing line: acyclicity and directed edges are proved; the full DAG-with-Markov-property is conditional on causal sufficiency. The parameterization (AND/OR, CPT form, edge semantics) is a formulation choice within the strongly motivated structure, motivated by parsimony and domain fit but not by mathematical necessity.


\medskip\noindent\textbf{Equivalence Class.}


\textbf{Within the DAG class.} Multiple DAGs can encode the same conditional independence relations, forming a Markov equivalence class identified by a CPDAG (completed partially directed acyclic graph). Two DAGs in the same equivalence class make identical probabilistic predictions but may differ in causal interpretation.

\textbf{Across representation types.} Factor graphs, junction trees, influence diagrams, and chain graphs are NOT simple presentational variants of DAGs:

\begin{itemize}
  \item \textbf{Factor graphs} and \textbf{junction trees} preserve factorization and inference structure without necessarily preserving directed causal semantics.
  \item \textbf{Influence diagrams} add decision and utility nodes — a richer object, not an equivalent one.
  \item \textbf{Chain graphs} can express independence models that are not representable as DAGs at all.
  \item \textbf{Markov equivalence} is a statement within DAG classes, not across all graphical model types.
\end{itemize}


The correct claim is narrow: for a given factorized distribution, DAG and factor-graph representations can compute the same marginals. But causal semantics (do-calculus) are DAG-specific and do not transfer to undirected or mixed representations without additional structure.

\textbf{ACT's choice.} ACT uses DAG + AND/OR because: (a) AND/OR is the most parsimonious complete basis for binary combination ( \S\,\ref{and-or-scope}), (b) the DAG naturally supports causal/interventional reasoning (Pearl's do-calculus), and (c) the representation converged across three independent formalism attempts.


\medskip\noindent\textbf{Epistemic Status.}


The acyclicity derivation is \textit{exact} — it follows from temporal ordering over a finite horizon via standard order theory. The individual postulates P1, P2, and P4 are each well-grounded (temporal structure, Cox's theorem, and chain fragility respectively).

The P3→Markov step is now \textit{conditional on causal sufficiency} of $\Sigma_t$. The argument has three parts: P3 requires some local neighborhood (locality), P1 identifies parents as the causal neighborhood (directionality), and causal sufficiency ensures parents capture all direct influences (no hidden common causes). Assembling these gives the local Markov property, which Lauritzen (1996, Theorem 3.27) proves equivalent to the factorization property for positive distributions. The causal sufficiency assumption is reasonable for agent-constructed strategies (all nodes are explicit) and fails precisely when the strategy is model-inadequate — the situation \S\,\ref{structural-adaptation-necessity} addresses.

The previous concern about alternative sparse factorizations (factor graphs, message-passing structures) is resolved: P1's directed temporal ordering selects the parent-conditional factorization specifically, because undirected alternatives sacrifice the causal semantics that \S\,\ref{strategy-dag} requires for interventional reasoning.

Max attainable: \textit{exact} for acyclicity (already there). \textit{Conditional} for the full DAG-with-Markov-property claim — conditional on causal sufficiency, which is a property of strategy quality rather than agent architecture. This is an honest upgrade from the previous "sketch" status.

The AND/OR restriction is a \textit{hypothesis} for binary outcomes, grounded in Boolean completeness and parsimony. For non-binary outcomes, it does not apply and richer parameterizations within the strongly motivated graphical structure are needed.

The analogy to Cox's theorem — consistency axioms force probability; operational postulates strongly motivate graphical structure — is suggestive but not established. Cox's theorem has a formal proof; this argument has a derivation sketch with one step that needs tightening. The analogy is worth stating as motivation but should not be relied on as evidence.


\medskip\noindent\textbf{Discussion.}


\textbf{The contribution of this analysis.} The dividing line between derived/conditional structure and chosen parameterization is the key result. The agent community (and future researchers) know exactly what is proved (acyclicity, directed edges), what is conditional (Markov property, contingent on causal sufficiency), and what is a formulation choice (the parameterization within the graph). Alternative parameterizations can be explored without abandoning the theoretical foundation.

\textbf{Acyclicity deserves emphasis.} The existing theory previously flagged "DAG acyclicity is an assumption, not forced" as a known fragility. The derivation above resolves this: acyclicity of $\Sigma_t$ follows from temporal ordering over a finite planning horizon. This is specific to the strategy — $M_t$'s environment model is not restricted to acyclic structures.

\textbf{The causal sufficiency assumption.} The Markov condition assumes causal sufficiency — no hidden common causes within the strategy. Since the agent constructs $\Sigma_t$, there are no "hidden" variables internal to the strategy (the agent chose all the nodes). But environmental common causes that affect multiple strategy steps may violate the Markov condition within $\Sigma_t$. If so, latent variable models or causal graphs with hidden nodes would be needed. This connects to the edge identifiability problem noted in \S\,\ref{strategy-dag}.

\end{derivation}

\begin{derivation}[Sector Condition Verification for Strategic Dynamics]
\label{strategic-dynamics-derivation}
Complete verification that specific strategy update dynamics satisfy the sector condition from \S\,\ref{sector-condition-derivation}, backing the schema proposed in \S\,\ref{strategy-persistence-schema}. Four cases are treated: a single-edge strategy, a two-edge AND-chain with observable intermediate, the same chain with unobservable intermediate, and a two-arm OR-node under exploratory action selection.



\medskip\noindent\textbf{Motivation.}


\S\,\ref{strategy-persistence-schema} proposes that the sector-condition mathematics (Props A.1, A.1S, A.2 in \S\,\ref{sector-condition-derivation}) extends from epistemic to strategic mismatch. The schema identifies three structural conditions (SA1, SA2', SA3) and conjectures the persistence form $\alpha_\Sigma > \rho_\Sigma / R_\Sigma$. This appendix instantiates the schema for concrete Beta-Bernoulli edge dynamics, verifying the conditions case by case. Each proposition establishes the sector parameter $\alpha_\Sigma$ for a specific DAG topology; the persistence condition then follows by direct substitution into Proposition A.1 of \S\,\ref{sector-condition-derivation}.



\medskip\noindent\textbf{Common Setup.}


Throughout, edge credences follow Beta-Bernoulli dynamics: edge $k$ has true success probability $\theta_k \in (0,1)$, agent belief $p_k \sim \text{Beta}(\alpha_k, \beta_k)$ with point estimate $\hat p_k = \alpha_k / n_k$ where $n_k = \alpha_k + \beta_k$, and conjugate update on observing $y_k \in \{0,1\}$:


$$\Delta \hat p_k = \frac{y_k - \hat p_k}{n_k + 1}, \qquad \eta_k = \frac{1}{n_k + 1}$$

The single-step update matches the gain-based form from \S\,\ref{edge-update-via-gain} with $\eta_{\text{edge}} = 1/(n_k + 1)$.

Strategic mismatch per edge is $\delta_k = \hat p_k - \theta_k$. The sector-condition framework ( \S\,\ref{sector-condition-derivation}) requires: (SA1) $F(\mathbf{0}) = \mathbf{0}$, (SA2') $\boldsymbol{\delta}^T \mathbf{F}(\boldsymbol{\delta}) \geq \alpha_\Sigma \lVert\boldsymbol{\delta}\rVert^2$ within a region $\mathcal B_R$.



\medskip\noindent\textbf{Proposition B.1: Single Edge ($A \to G$).}


\textbf{Setup.} One action $A$, one goal $G$, one edge with credence $\hat p$ and true probability $\theta$. Mismatch: $\delta_\Sigma = \hat p - \theta$ (scalar).

\textbf{Statement.} Under Beta-Bernoulli updating, the expected correction function satisfies the sector condition globally (all $\lvert\delta_\Sigma\rvert \leq 1$) with:

% [Derived (Conditional on Beta-Bernoulli model)]

$$\alpha_\Sigma = \frac{1}{n+1}$$

The bound is tight.

\textbf{Proof.}

The expected correction. Since $y \sim \text{Bernoulli}(\theta)$:

% [Derived (Proof Step)]

$$\mathbb{E}[\Delta \hat p] = \frac{\theta - \hat p}{n+1} = -\frac{\delta_\Sigma}{n+1}$$

Identifying $F_\Sigma(\delta_\Sigma) = -\mathbb{E}[\Delta\delta_\Sigma] = \delta_\Sigma/(n+1)$:

% [Derived (Proof Step)]

$$\delta_\Sigma \cdot F_\Sigma(\delta_\Sigma) = \frac{\delta_\Sigma^2}{n+1} = \frac{1}{n+1}\lVert\delta_\Sigma\rVert^2$$

So $\alpha_\Sigma = 1/(n+1)$. The correction is exactly linear, so the sector bound is tight (not merely a lower bound).

\textbf{Verification of (SA1).} At $\delta_\Sigma = 0$: $F_\Sigma(0) = 0$. Satisfied.

\textbf{Stochastic ultimate bound.} Decomposing $\delta_{\Sigma,t+1} = \delta_\Sigma \cdot n/(n+1) + (y - \theta)/(n+1)$ and computing $\mathbb{E}[\delta_{\Sigma,t+1}^2]$:

% [Derived (stochastic ultimate bound, single edge)]

$$\mathbb{E}[\delta_{\Sigma,t+1}^2] = \delta_\Sigma^2 \cdot \frac{n^2}{(n+1)^2} + \frac{\theta(1-\theta)}{(n+1)^2}$$

The expected Lyapunov derivative $\mathbb{E}[\Delta V] < 0$ whenever:

$$\lvert\delta_\Sigma\rvert > \sqrt{\frac{\theta(1-\theta)}{2n+1}}$$

The stochastic noise floor decreases as $O(1/\sqrt{n})$ --- the standard Bayesian posterior concentration rate. $\square$



\medskip\noindent\textbf{Proposition B.2: Two-Edge AND-Chain, Observable Intermediate ($A \to B \to G$, $B$ observed).}


\textbf{Setup.} Action $A$ leads to intermediate $B$, then to goal $G$. Both $y_B$ and $y_G$ are observed. Edge credences $\hat p_1, \hat p_2$ with true probabilities $\theta_1, \theta_2$. Plan confidence: $\hat P_\Sigma = \hat p_1 \hat p_2$ (AND propagation from \S\,\ref{and-or-scope}). Mismatch vector: $\boldsymbol{\delta}_\Sigma = (\delta_1, \delta_2)^T$.

The generative model:

$$y_B \sim \text{Bernoulli}(\theta_1), \qquad y_G \mid y_B = 1 \sim \text{Bernoulli}(\theta_2), \qquad y_G \mid y_B = 0 = 0$$

\textbf{Statement.} Under Beta-Bernoulli updating with observable intermediate, the expected correction function is diagonal and satisfies the sector condition globally with:

% [Derived (Conditional on Beta-Bernoulli model, observable intermediate)]

$$\alpha_\Sigma = \min\!\left(\frac{1}{n_1+1},\; \frac{\theta_1}{n_2+1}\right)$$

\textbf{Proof.}

Edge 1 is tested every trial: $\Delta \hat p_1 = (y_B - \hat p_1)/(n_1+1)$, giving $\mathbb{E}[\Delta \hat p_1] = -\delta_1/(n_1+1)$.

Edge 2 is tested only when $y_B = 1$ (probability $\theta_1$). When $y_B = 0$, edge 2 receives no information:

% [Derived (Proof Step)]

$$\mathbb{E}[\Delta \hat p_2] = \theta_1 \cdot \frac{\theta_2 - \hat p_2}{n_2+1} + (1 - \theta_1) \cdot 0 = -\frac{\theta_1 \delta_2}{n_2+1}$$

The expected correction function is diagonal:

$$\mathbf{F}(\boldsymbol{\delta}_\Sigma) = \begin{pmatrix} \delta_1/(n_1+1) \\ \theta_1 \delta_2/(n_2+1) \end{pmatrix}$$

\textbf{Verification of (SA1).} At $\boldsymbol{\delta}_\Sigma = \mathbf{0}$: $\mathbf{F}(\mathbf{0}) = \mathbf{0}$. Satisfied.

The sector product:

% [Derived (Proof Step)]

$$\boldsymbol{\delta}_\Sigma^T \mathbf{F}(\boldsymbol{\delta}_\Sigma) = \frac{\delta_1^2}{n_1+1} + \frac{\theta_1 \delta_2^2}{n_2+1} \geq \min\!\left(\frac{1}{n_1+1}, \frac{\theta_1}{n_2+1}\right)(\delta_1^2 + \delta_2^2)$$

The minimum is achieved by whichever edge has the lower effective correction rate. $\square$

\textbf{The evidence-starvation effect.} Edge 2's effective correction rate $\theta_1/(n_2+1)$ is attenuated by the upstream success probability $\theta_1$. When $\theta_1$ is small, edge 2 is starved of evidence --- it is tested rarely, learns slowly, and becomes the weakest link. This gives a precise mechanism for \S\,\ref{chain-confidence-decay}: downstream edges in AND-chains are harder to calibrate, with correction rate proportional to the product of all upstream success probabilities.

**Generalization to depth $d$.** For a chain of depth $d$ with all intermediates observable, edge $k$'s effective correction rate is $\alpha_k = \prod_{j=1}^{k-1}\theta_j / (n_k + 1)$, giving:

% [Derived (depth-$d$ chain sector parameter)]

$$\alpha_\Sigma = \min_{k=1}^{d}\; \frac{\prod_{j=1}^{k-1}\theta_j}{n_k+1}$$



\medskip\noindent\textbf{Proposition B.3: Two-Edge AND-Chain, Unobservable Intermediate ($A \to B \to G$, $B$ not observed).}


\textbf{Setup.} Same DAG as Proposition B.2, but only $y_G \in \{0,1\}$ is observed. The agent cannot distinguish "$A \to B$ failed" from "$B \to G$ failed" when $y_G = 0$.

\textbf{Statement.} (a) The per-edge sector condition fails: the marginal Bayesian update violates (SA1) with systematic bias $O(1/n)$. (b) Plan-level tracking (treating $\hat\Phi = \hat p_1 \hat p_2$ as a single Beta over the composite probability $\Phi = \theta_1\theta_2$) recovers the sector condition with:

% [Derived (Conditional on Beta-Bernoulli model, unobservable intermediate)]

$$\alpha_{\Sigma,\text{plan}} = \frac{1}{n_\Phi + 1}$$

\textbf{Proof of (a): Per-edge (SA1) violation.}

The joint posterior after success ($y_G = 1$) factors cleanly: $\alpha_k \to \alpha_k + 1$ for both edges. The joint posterior after failure ($y_G = 0$) includes the factor $(1 - \theta_1\theta_2)$, which does not decompose into separate functions of $\theta_1$ and $\theta_2$. Failure introduces posterior correlation.

The exact marginal Bayesian update for edge 1 on failure is:

% [Derived (Proof Step)]

$$\Delta \hat p_1^{(\text{fail})} = -\frac{\hat p_1 \, q_1}{n_1 + 1}$$

where $q_1 = (1 - \hat p_1)/(1 - \hat p_1 \hat p_2)$ is the posterior probability that edge 1 caused the failure (the proportional-blame weight). The expected update, combining success (probability $\theta_1\theta_2$) and failure (probability $1 - \theta_1\theta_2$):

% [Derived (Proof Step)]

$$\mathbb{E}[\Delta \hat p_1] = \theta_1\theta_2 \cdot \frac{1 - \hat p_1}{n_1+1} - (1-\theta_1\theta_2)\cdot\frac{\hat p_1\,q_1}{n_1+1}$$

Evaluating at truth ($\hat p_k = \theta_k$):

% [Derived (SA1 violation magnitude)]

$$\mathbb{E}[\Delta \hat p_1]\big\rvert_{\hat p = \theta} = -\frac{\theta_1(1-\theta_1)(1-\theta_2)}{n_1+1} \neq 0$$

The bias is negative (pushes credence below truth), of magnitude $O(1/n)$, and arises because success always increments $\alpha_1$ by a full unit while failure increments $\beta_1$ by the fractional weight $q_1 < 1$. The success case over-credits edge 1 because the agent cannot separate edge 1's contribution from edge 2's. By symmetry, $\hat p_2$ is also biased downward. $\square$

\textbf{Proof of (b): Plan-level recovery.}

If the agent tracks $\hat\Phi$ directly as a Beta posterior over $\Phi = \theta_1\theta_2$ using $y_G \sim \text{Bernoulli}(\Phi)$, the setup reduces to a single-edge problem (Proposition B.1) with mismatch $\delta_\Phi = \hat\Phi - \Phi$ and sector parameter $\alpha_{\Sigma,\text{plan}} = 1/(n_\Phi + 1)$. All single-edge results apply. $\square$

\textbf{Tradeoff.} Per-edge tracking provides diagnostic resolution (which edge is failing?) but yields biased, non-sector-conforming updates when the intermediate is unobservable. Plan-level tracking sacrifices diagnostic resolution but provides clean sector-condition satisfaction. This is a concrete instance of \S\,\ref{observability-dominance}: unobservable nodes freeze per-edge learning, making the aggregate the correct unit of analysis.

\textbf{Non-identifiability.} From $y_G$ observations alone, only $\Phi = \theta_1\theta_2$ is identifiable. The per-edge decomposition $(\theta_1, \theta_2)$ is underdetermined --- any pair with the same product produces identical observations. Per-edge identification requires per-edge observability (or structural constraints that break the symmetry).



\medskip\noindent\textbf{Proposition B.4: Two-Arm OR-Node ($A_1, A_2 \to G$, $\varepsilon$-greedy).}


\textbf{Setup.} Goal $G$ reachable via two alternative actions $A_1, A_2$ (an OR-node from \S\,\ref{and-or-scope}). True probabilities $\theta_1, \theta_2$. Agent selects one action per trial: with probability $1 - \varepsilon$ the greedy choice (higher $\hat p_k$), with probability $\varepsilon$ the other. Mismatch vector: $\boldsymbol{\delta}_\Sigma = (\delta_1, \delta_2)^T$.

\textbf{Statement.} (a) Under a greedy policy ($\varepsilon = 0$), the sector condition fails for the full mismatch vector. (b) Under $\varepsilon$-greedy action selection with $\varepsilon > 0$, the sector condition is satisfied with:

% [Derived (Conditional on Beta-Bernoulli model, $\varepsilon$-greedy OR-node)]

$$\alpha_\Sigma = \min\!\left(\frac{1-\varepsilon}{n_1+1},\; \frac{\varepsilon}{n_2+1}\right)$$

where arm 1 is the greedy choice.

\textbf{Proof of (a): Greedy failure.}

Under greedy selection, only the chosen arm updates. Assuming $\hat p_1 > \hat p_2$ (arm 1 is greedy):

$$\mathbf{F}(\boldsymbol{\delta}_\Sigma) = \begin{pmatrix}\delta_1/(n_1+1) \\ 0\end{pmatrix}$$

The sector product $\boldsymbol{\delta}_\Sigma^T \mathbf{F} = \delta_1^2/(n_1+1)$. For the sector condition to hold, we need $\delta_1^2/(n_1+1) \geq \alpha_\Sigma(\delta_1^2 + \delta_2^2)$. Setting $\delta_1 = 0$: the LHS is zero but $\alpha_\Sigma \delta_2^2 > 0$ for $\delta_2 \neq 0$. No $\alpha_\Sigma > 0$ works. $\square$

**Proof of (b): $\varepsilon$-greedy recovery.**

The expected correction (arm 1 greedy):

% [Derived (Proof Step)]

$$\mathbb{E}[\mathbf{F}(\boldsymbol{\delta}_\Sigma)] = \begin{pmatrix}(1-\varepsilon)\,\delta_1/(n_1+1) \\ \varepsilon\,\delta_2/(n_2+1)\end{pmatrix}$$

The sector product:

% [Derived (Proof Step)]

$$\boldsymbol{\delta}_\Sigma^T \mathbf{F} = \frac{(1-\varepsilon)\,\delta_1^2}{n_1+1} + \frac{\varepsilon\,\delta_2^2}{n_2+1} \geq \min\!\left(\frac{1-\varepsilon}{n_1+1},\;\frac{\varepsilon}{n_2+1}\right)\!(\delta_1^2 + \delta_2^2)$$

\textbf{Verification of (SA1).} At $\boldsymbol{\delta}_\Sigma = \mathbf{0}$: $\mathbf{F}(\mathbf{0}) = \mathbf{0}$. Satisfied. $\square$

\textbf{Optimal exploration rate.} Maximizing $\alpha_\Sigma$ over $\varepsilon$ by equalizing the two terms:

% [Derived (optimal exploration rate)]

$$\varepsilon^* = \frac{n_1+1}{n_1+n_2+2}, \qquad \alpha_\Sigma^* = \frac{1}{n_1+n_2+2}$$

For equal experience ($n_1 = n_2$): $\varepsilon^* = 1/2$ (equal allocation). For unequal experience, $\varepsilon^*$ allocates more trials to the arm with higher $n$ (lower gain) to equalize correction rates.

**The $k$-arm generalization.** For $k$ OR-alternatives with $\varepsilon$-uniform exploration (probability $\varepsilon/(k-1)$ per non-greedy arm):

% [Derived ($k$-arm OR-node sector parameter)]

$$\alpha_\Sigma = \min\!\left(\frac{1-\varepsilon}{n_{\text{greedy}}+1},\;\min_{j \neq \text{greedy}}\frac{\varepsilon/(k-1)}{n_j+1}\right)$$

With optimal equal-rate exploration and equal experience: $\alpha_\Sigma = 1/(k(n+1))$. The sector parameter scales as $1/k$ --- each additional alternative dilutes correction capacity.

\textbf{Condition (SA3).} For the sector condition to be satisfiable at all, the exploration rate must satisfy:

% [Derived (minimum exploration rate for persistence)]

$$\varepsilon > \frac{\rho_\Sigma(n_{\max}+1)}{R_\Sigma}$$

A purely greedy agent ($\varepsilon = 0$) violates the sector condition because the unchosen alternative's mismatch grows without correction. This is the structural justification for (SA3) in \S\,\ref{strategy-persistence-schema}: OR-nodes require deliberate exploration to maintain strategic persistence.



\medskip\noindent\textbf{Persistence Conditions.}


Substituting each proposition's $\alpha_\Sigma$ into the persistence threshold $\alpha_\Sigma > \rho_\Sigma / R_\Sigma$ from Proposition A.1 of \S\,\ref{sector-condition-derivation}:

\textbf{Single edge (B.1):}

% [Derived (single-edge persistence threshold)]

$$\frac{1}{n+1} > \frac{\rho_\Sigma}{R_\Sigma} \quad\iff\quad n < \frac{R_\Sigma}{\rho_\Sigma} - 1$$

The critical experience level $n^* = R_\Sigma/\rho_\Sigma - 1$ is the point at which gain collapse occurs: the posterior has concentrated so much that the agent can no longer track environmental drift.

\textbf{Two-edge observable (B.2):}

% [Derived (two-edge persistence threshold, observable)]

$$\min\!\left(\frac{1}{n_1+1},\;\frac{\theta_1}{n_2+1}\right) > \frac{\rho_\Sigma}{R_\Sigma}$$

This decomposes into two independent conditions: $n_1 < R_\Sigma/\rho_\Sigma - 1$ and $n_2 < \theta_1 R_\Sigma/\rho_\Sigma - 1$. The downstream edge has a stricter bound, tightened by $\theta_1$. Gain collapse hits downstream edges first.

\textbf{Two-edge unobservable (B.3):}

$$\frac{1}{n_\Phi+1} > \frac{\rho_\Sigma}{R_\Sigma} \quad\iff\quad n_\Phi < \frac{R_\Sigma}{\rho_\Sigma} - 1$$

Same form as single-edge, applied to the plan-level aggregate.

**Two-arm OR, $\varepsilon$-greedy (B.4):**

% [Derived (OR-node persistence threshold)]

$$\min\!\left(\frac{1-\varepsilon}{n_1+1},\;\frac{\varepsilon}{n_2+1}\right) > \frac{\rho_\Sigma}{R_\Sigma}$$

With optimal exploration and equal experience ($n_1 = n_2 = n$):

$$n < \frac{R_\Sigma}{2\rho_\Sigma} - 1$$

The critical experience is halved relative to the single-edge case. Each additional OR-alternative further reduces the critical level by $1/k$.



\medskip\noindent\textbf{Summary of Results.}


\begin{adjustbox}{max width=\textwidth}
\small
\begin{tabular}{llllll}
\toprule
Case & Topology & $\alpha_\Sigma$ & SA1 & SA3 & Weakest link \\
\midrule
\textbf{B.1} Single edge & $A \to G$ & $1/(n+1)$ & Yes & N/A & --- \\
\textbf{B.2} Two-edge, $B$ observable & $A \to B \to G$ (AND) & $\min(1/(n_1\!+\!1),\;\theta_1/(n_2\!+\!1))$ & Yes & N/A & Depth-gated (evidence starvation) \\
\textbf{B.3} Two-edge, $B$ unobservable & $A \to B \to G$ (AND) & $1/(n_\Phi\!+\!1)$ (plan-level) & Per-edge: No ($O(1/n)$ bias); Plan: Yes & N/A & Credit-assignment collapse \\
\textbf{B.4} Two-arm OR, $\varepsilon$-greedy & $A_1, A_2 \to G$ (OR) & $\min((1\!-\!\varepsilon)/(n_1\!+\!1),\;\varepsilon/(n_2\!+\!1))$ & Yes & Required & Exploration-gated \\
\textbf{B.5a} Credence→value (linear) & Any & $\alpha_s = \alpha_c$ (Jacobian cancels) & — & — & None (exact transfer) \\
\textbf{B.5b} Credence→value (nonlinear, componentwise) & Any & $\alpha_s = \alpha_c$ ($J_k \geq 0$ preserves bound) & — & — & None (exact transfer) \\
\textbf{B.5b} Credence→value (coupled) & Any & $\alpha_s \geq \alpha_c / \kappa(\mathbf{J})^2$ & — & — & Inter-edge coupling \\
\bottomrule
\end{tabular}
\end{adjustbox}


\textbf{Structural results across cases:}

\begin{itemize}
  \item \textbf{Gain as sector parameter.} In every case, $\alpha_\Sigma$ is determined by the edge update gain $\eta_k = 1/(n_k+1)$, modulated by observability ($\theta_1$ attenuation in B.2), credit-assignment (plan-level collapse in B.3), or action-selection policy ($\varepsilon$ in B.4). The structural parallel between epistemic and strategic persistence is not an analogy --- it is a mathematical identity at the sector-framework level.
  \item \textbf{Evidence starvation (AND-chains).} Downstream edges in AND-chains have correction rates attenuated by the product of upstream success probabilities. Depth increases fragility: $\alpha_\Sigma$ decays exponentially with chain depth.
  \item \textbf{Exploration gating (OR-nodes).} Unchosen OR-alternatives receive zero correction. The sector condition requires deliberate exploration (SA3) at a rate exceeding $\rho_\Sigma(n_{\max}+1)/R_\Sigma$. Pure greedy policies fail.
  \item \textbf{Gain-collapse threshold.} In all cases, increasing experience eventually violates persistence when the environment drifts. The critical experience level $n^*$ is inversely proportional to the drift rate $\rho_\Sigma$.
\end{itemize}




\medskip\noindent\textbf{Proposition B.5: Bridge from Credence Error to Value Residuals.}


The propositions above verify the sector condition for the credence-error mismatch state $\boldsymbol\delta_c = (\hat p_k - \theta_k)$. This proposition extends the result to \textbf{plan-confidence error} — the scalar difference between the agent's plan-confidence score and the independence-model plan value evaluated at true edge parameters. This is a natural plan-level mismatch that is distinct from (but related to) the per-edge value-increment residuals $\delta_{\text{strategic}}$ defined in \S\,\ref{strategic-calibration}.

**Relationship to $\delta_{\text{strategic}}$.** \S\,\ref{strategic-calibration} defines $\delta_{\text{strategic}}$ as an $L^2$ aggregation of per-edge value-increment residuals — a quantity that requires credit assignment to compute. Plan-confidence error $\delta_s$ defined here is a \textit{different} quantity: it is the error in the DAG's aggregate self-assessment, computable from status propagation without credit assignment. The two are related (both measure strategy-reality mismatch) but not identical. This proposition shows the sector condition holds for plan-confidence error; extending it to $\delta_{\text{strategic}}$ directly would require the credit-assignment machinery discussed in \S\,\ref{credit-assignment-boundary}.


\medskip\noindent\textbf{Setup.}


The plan value $\hat P_\Sigma = P_\Sigma(\mathbf{p})$ is a function of the edge credence vector $\mathbf{p} = (p_1, \ldots, p_m)$. The \textbf{independence-model reference value} is $\Phi = P_\Sigma(\boldsymbol\theta)$ — the AND/OR propagation formula evaluated at the true edge parameters $\boldsymbol\theta$. Note: $\Phi$ is NOT the actual probability of plan success when edge outcomes are correlated ( \S\,\ref{strategy-dag}, edge-independence caveat). It is the best the independence model can do with perfect edge knowledge. Under correlated failure, the actual plan success probability may be lower than $\Phi$; the gap is a model-class limitation of the AND/OR DAG representation, not an estimation error. What $\delta_s$ tracks is calibration \textit{within the independence model}, not calibration to reality. Define:

\begin{itemize}
  \item \textbf{Credence-error mismatch:} $\boldsymbol\delta_c = \mathbf{p} - \boldsymbol\theta \in \mathbb{R}^m$
  \item \textbf{Plan-confidence error:} $\delta_s = \hat P_\Sigma - \Phi = P_\Sigma(\mathbf{p}) - P_\Sigma(\boldsymbol\theta)$ (scalar)
\end{itemize}


The gradient of plan value with respect to credences is $\mathbf{J} = \nabla_{\mathbf{p}} P_\Sigma \in \mathbb{R}^m$ (a vector, since $P_\Sigma$ is scalar). By first-order Taylor expansion:

$$\delta_s \approx \mathbf{J}^T \boldsymbol\delta_c$$


\medskip\noindent\textbf{B.5a: Linear Correction (Beta-Bernoulli).}


% [Derived (value-residual sector condition, linear case)]

When the correction in credence space is linear — $\mathbf{F}_c(\boldsymbol\delta_c) = \boldsymbol\eta \odot \boldsymbol\delta_c$ where $\eta_k = 1/(n_k+1)$ and $\odot$ is elementwise product — consider the sector ratio directly. Since $\delta_s = \mathbf{J}^T \boldsymbol\delta_c$ and $F_s = \mathbf{J}^T \mathbf{F}_c$:

$$\delta_s \cdot F_s = (\mathbf{J}^T \boldsymbol\delta_c)^T (\mathbf{J}^T \mathbf{F}_c) = \boldsymbol\delta_c^T \mathbf{J} \mathbf{J}^T \mathbf{F}_c$$

For the linear case $\mathbf{F}_c = \eta_{\min} \boldsymbol\delta_c$ (using the worst-case uniform gain for the bound):

$$\delta_s \cdot F_s = \eta_{\min} \cdot \boldsymbol\delta_c^T \mathbf{J} \mathbf{J}^T \boldsymbol\delta_c = \eta_{\min} \cdot \lVert\mathbf{J}^T \boldsymbol\delta_c\rVert^2 = \eta_{\min} \cdot \delta_s^2$$

Therefore:

$$\frac{\delta_s \cdot F_s}{\delta_s^2} = \eta_{\min}$$

\textbf{The Jacobian cancels.} The sector parameter in value-residual space equals the sector parameter in credence space, regardless of DAG structure:

$$\alpha_s = \alpha_c = \eta_{\min} = \min_k \frac{1}{n_k + 1}$$

This holds because $\mathbf{J}\mathbf{J}^T$ appears in both numerator and denominator of the sector ratio, and the linear correction commutes with the Jacobian mapping.

\textbf{Interpretation.} The Jacobian $\mathbf{J}$ determines \textit{which} credence errors matter most for plan value (high-sensitivity edges contribute more to value residuals). But it does not change the \textit{rate} at which those errors are corrected, because the correction is proportional to the error regardless of the error's plan-value impact.


\medskip\noindent\textbf{B.5b: Nonlinear Componentwise Correction.}


% [Derived (from per-component sector condition + Jacobian non-negativity)]

When the correction is nonlinear but \textbf{componentwise} — each edge corrects independently, so $(\mathbf{F}_c)_k$ depends only on $(\boldsymbol\delta_c)_k$ — and the plan-value Jacobian is \textbf{non-negative} ($J_k \geq 0$ for all $k$), the transfer is lossless regardless of nonlinearity.

\textbf{The per-component sector condition} (from Props B.1-B.4): for each edge $k$,

$$(\boldsymbol\delta_c)_k \cdot (\mathbf{F}_c)_k \geq \alpha_c \cdot (\boldsymbol\delta_c)_k^2$$

which gives $(\mathbf{F}_c)_k / (\boldsymbol\delta_c)_k \geq \alpha_c$ (the correction-to-mismatch ratio exceeds $\alpha_c$ for each edge independently).

\textbf{Jacobian non-negativity} holds for all well-formed AND/OR DAGs: increasing any edge credence $p_k$ never decreases plan value $P_\Sigma$ (monotonicity of AND/OR propagation). Therefore $J_k = \partial P_\Sigma / \partial p_k \geq 0$ for all $k$.

\textbf{The transfer.} Since $J_k \geq 0$ preserves the inequality direction:

$$J_k \cdot (\mathbf{F}_c)_k \geq J_k \cdot \alpha_c \cdot (\boldsymbol\delta_c)_k \quad \text{for each } k$$

Summing:

$$\mathbf{J}^T \mathbf{F}_c = \sum_k J_k (\mathbf{F}_c)_k \geq \alpha_c \sum_k J_k (\boldsymbol\delta_c)_k = \alpha_c \cdot \mathbf{J}^T \boldsymbol\delta_c$$

Therefore $F_s \geq \alpha_c \cdot \delta_s$, giving:

$$\alpha_s = \alpha_c$$

\textbf{No condition-number penalty.} The Jacobian's anisotropy is irrelevant when the correction is componentwise and the Jacobian is non-negative. The non-negativity ensures that every edge's correction contributes in the right direction to the value correction, regardless of how unevenly the Jacobian weights different edges.

\textbf{When does componentwise correction hold?} For all cases in Props B.1, B.2, and B.4 — the standard Beta-Bernoulli update on each edge, including the stochastic (nonlinear) single-step realization. The correction IS nonlinear (the actual per-step update $(y - \hat p)/(n+1)$ is random, not proportional to $\delta_k$), but it is componentwise (each edge updates from its own observation independently).

\textbf{When does it fail?} When edge corrections are coupled — correcting one edge changes another's correction direction. This occurs for:
\begin{itemize}
  \item \textbf{Unobservable intermediates (B.3):} The marginal Bayesian update couples edge posteriors. The correction for edge 2 depends on the joint posterior over both edges, not just on $(\boldsymbol\delta_c)_2$ alone.
  \item \textbf{Shared observations:} If testing one arm gives partial information about another (correlated arms), corrections couple across edges.
\end{itemize}


For coupled corrections, the general Cauchy-Schwarz bound gives:

$$\alpha_s \geq \frac{\alpha_c}{\kappa(\mathbf{J})^2}$$

where $\kappa(\mathbf{J})$ is the condition number of the Jacobian (ratio of maximum to minimum singular values). This is the bound from the earlier analysis but it now applies only to the coupled case, not to all nonlinear corrections.

\textbf{The refined picture: nonlinearity is not the problem; inter-edge coupling is.} A nonlinear but componentwise correction transfers losslessly through any non-negative Jacobian. A coupled correction — even a linear one — incurs the condition-number penalty because the coupling can redirect correction away from the value-relevant direction.


\medskip\noindent\textbf{B.5c: Implications.}


Three regimes for the credence-to-value bridge:

\begin{adjustbox}{max width=\textwidth}
\small
\begin{tabular}{llll}
\toprule
Correction type & Edge coupling & Transfer & $\alpha_s$ \\
\midrule
Linear & Any & Exact (Jacobian cancels) & $\alpha_c$ \\
Nonlinear, componentwise & Independent edges & Exact ($J_k \geq 0$ preserves bound) & $\alpha_c$ \\
Nonlinear, coupled & Coupled edges & $\kappa(\mathbf{J})^2$ penalty & $\alpha_c / \kappa^2$ \\
\bottomrule
\end{tabular}
\end{adjustbox}


\textbf{Componentwise correction closes the gap for all verified cases.} Props B.1, B.2, and B.4 all use componentwise edge updates with non-negative Jacobian. The operational diagnostic $\delta_{\text{strategic}}$ inherits the sector condition from credence error with no penalty — even for the nonlinear (stochastic) single-step realization. The derivation validates the operational mismatch, not merely a surrogate.

\textbf{Inter-edge coupling — not nonlinearity — is the fragility.} The $\kappa^2$ penalty arises only when correcting one edge changes another's correction direction (B.3's unobservable case, correlated observations). Deep or unbalanced DAGs have high $\kappa(\mathbf{J})$ when corrections couple, but the coupling, not the depth or imbalance per se, is what degrades the persistence guarantee.

\textbf{Connection to credit assignment.} The Jacobian $\mathbf{J} = \nabla_\mathbf{p} P_\Sigma$ is computable from the DAG's status propagation formulas — it does not require solving the credit-assignment problem. The credit-assignment problem is about \textit{decomposing observed value changes into per-edge contributions} (needed for the update rule). The Jacobian bridge is about \textit{transferring a per-edge sector condition to the value-residual space} (needed for the persistence guarantee). These are different problems: the bridge works even when credit assignment is unsolved, because it only requires the sensitivity structure, not the causal attribution.



\medskip\noindent\textbf{Epistemic Status.}


\textit{Conditional on the Beta-Bernoulli model.} Propositions B.1, B.2, and B.4 are \textit{derived}: the sector-condition verification is exact algebra under the stated generative model. The persistence conditions follow by direct application of Proposition A.1 ( \S\,\ref{sector-condition-derivation}), which is independently established. Proposition B.3(a) (the SA1 violation) is also derived; B.3(b) reduces to B.1 by construction. Proposition B.5a (the credence-to-value bridge for linear correction) is \textit{exact} — the Jacobian cancellation is algebraic, independent of DAG structure. Proposition B.5b (the nonlinear transfer) is \textit{conditional on Jacobian regularity} — the condition-number bound requires $\sigma_{\min}(\mathbf{J}) > 0$, which holds for non-degenerate DAGs.

All results use the \textit{expected-value} sector condition. A full stochastic treatment (Foster-Lyapunov or supermartingale convergence) would give probability bounds rather than expected-value bounds. The expected-value analysis gives the correct asymptotic behavior (posterior concentration at $O(1/\sqrt{n})$) but does not prove almost-sure convergence. For the stationary Beta-Bernoulli case, almost-sure convergence is guaranteed by the standard Bayesian consistency theorem independently of this derivation.

The time-varying $\alpha_\Sigma$ issue remains: since $n_k$ increases with each observation, $\alpha_\Sigma$ decreases over time. The sector-condition framework assumes constant $\alpha$. The results here give an \textit{instantaneous} persistence check at the current experience level, not a trajectory guarantee. Experience discounting (exponential forgetting with factor $\lambda$) stabilizes $\alpha_\Sigma$ at approximately $1 - \lambda$, yielding the persistence requirement $\lambda < 1 - \rho_\Sigma/R_\Sigma$.

\textbf{What remains open:}

\begin{enumerate}
  \item \textit{General DAG topology.} Only linear AND-chains and isolated OR-nodes are verified. Mixed AND/OR DAGs may exhibit interaction effects between evidence starvation and exploration gating.
  \item \textit{Continuous outcomes.} The Beta-Bernoulli model gives conjugate, closed-form updates. Non-conjugate cases (continuous signals, partial observability) require approximate inference, and the sector condition must be verified for the approximation.
  \item \textit{Modified sector condition for biased correction.} The $O(1/n)$ bias in the unobservable case (B.3a) could be accommodated by a sector-condition variant tolerating asymptotically vanishing bias, potentially recovering per-edge results.
  \item \textit{Correlated edges.} All cases assume independent edges. Shared infrastructure or common-mode failures would introduce additional coupling.
  \item \textit{Adaptive exploration.} Proposition B.4 uses fixed $\varepsilon$. Adaptive strategies (UCB, Thompson sampling) allocate exploration based on current uncertainty and should yield tighter sector bounds.
\end{enumerate}




\medskip\noindent\textbf{Discussion.}


\textbf{The sector parameter is the edge update gain.} Across all four cases, $\alpha_\Sigma$ is a function of $\eta_k = 1/(n_k+1)$ --- the same quantity that governs epistemic persistence ( \S\,\ref{adaptive-tempo}, \S\,\ref{update-gain}). Strategic persistence and epistemic persistence are governed by the same mathematical machinery, not merely structurally parallel. This validates the schema proposed in \S\,\ref{strategy-persistence-schema} at the level of concrete dynamics.

\textbf{AND vs. OR: structural vs. behavioral fragility.} The evidence-starvation effect (B.2) is structural --- it depends on the DAG topology and upstream reliability, not on the agent's policy. The exploration-gating effect (B.4) is behavioral --- it depends on the action-selection policy. AND-node persistence is improved by increasing upstream reliability (making early steps more likely to succeed). OR-node persistence is improved by increasing exploration (testing alternatives more frequently). The distinction maps to the familiar depth-vs-breadth tradeoff in planning.

\textbf{Gain collapse as the dominant failure mode.} In every drifting-environment case, the persistence threshold takes the form $n < c/\rho_\Sigma$ for some constant $c$. Accumulated experience suppresses the update gain, eventually making the agent unable to track environmental change. This gives the gain-collapse phenomenon ( \S\,\ref{update-gain}) a precise quantitative threshold for each topology.

\textbf{The credence-to-value bridge decouples persistence from credit assignment.} Proposition B.5 shows that strategic persistence (in value-residual space) can be established from per-edge sector conditions (in credence space) without solving the credit-assignment problem. The bridge requires only the value-sensitivity Jacobian $\mathbf{J} = \nabla_\mathbf{p} P_\Sigma$, which is computable from the DAG's status propagation formulas in $O(\lvert V\rvert + \lvert E\rvert)$ time. Credit assignment — attributing observed value changes to specific edges — is needed for the \textit{update rule} (computing the signal function in \S\,\ref{edge-update-via-gain}), not for the \textit{persistence guarantee}. This means the persistence analysis can proceed even while the update rule remains incompletely specified.

\textbf{Inter-edge coupling is the true fragility, not nonlinearity or DAG depth.} Proposition B.5b shows that nonlinear componentwise corrections transfer losslessly through any non-negative Jacobian — the $\kappa(\mathbf{J})^2$ penalty arises only when corrections are \textit{coupled} across edges. This refines the earlier qualitative argument from \S\,\ref{chain-confidence-decay}: the structural pressure toward shallow, balanced strategies is not about depth per se but about the coupling that deep or complex structures tend to introduce. A deep AND-chain with fully observable intermediates (B.2) has componentwise corrections and transfers losslessly despite high depth. The same chain with unobservable intermediates (B.3) has coupled corrections and incurs the penalty. The distinction is observability-mediated coupling, not topological complexity.

\end{derivation}

\begin{derivation}[Discrete-Time Sector Condition]
\label{discrete-sector-condition}
Discrete-time analogs of Props A.1, A.1S, and A.2 via contraction mapping, closing the fluid-limit gap (GA-5) between the event-driven dynamics ( \S\,\ref{event-driven-dynamics}) and the continuous-time Lyapunov results in \S\,\ref{sector-condition-derivation}.



\medskip\noindent\textbf{Formal Expression.}



\medskip\noindent\textbf{Setup.}


The discrete mismatch dynamics at event step $k$ are:


$$\delta_{k+1} = \delta_k - \eta^* F_d(\delta_k) + w_k$$

where $\eta^*$ is the update gain ( \S\,\ref{update-gain}), $F_d$ is the discrete correction direction, and $w_k$ is the per-step disturbance. The continuous correction function $F(\mathcal{T}, \delta)$ from \S\,\ref{sector-condition-derivation} decomposes as $F = \nu \cdot \eta^* \cdot F_d$ at event rate $\nu$.


\medskip\noindent\textbf{(DA2') Discrete Sector Condition.}


% [Assumption DA2' (discrete-sector-condition)]

There exist constants $c_{\min} > 0$ and $c_{\max} < 2/\eta^*$ such that for all $\lVert\delta\rVert \leq R$:

$$c_{\min} \lVert\delta\rVert^2 \leq \delta^T F_d(\delta) \leq c_{\max} \lVert\delta\rVert^2$$

The \textbf{lower bound} ($c_{\min}$) is directional fidelity — the correction points inward, identical to the continuous sector condition (A2') from \S\,\ref{sector-condition-derivation} via \S\,\ref{gain-sector-bridge}.

The \textbf{upper bound} ($c_{\max} < 2/\eta^*$) is the \textbf{no-overshoot condition}: each correction step must not reverse the mismatch. This is the classical step-size condition $\eta^* < 2/L$ for gradient descent (where $L$ is the Lipschitz constant of the gradient). For Bayesian updates, this is satisfied by construction — the posterior lies between prior and data.


\medskip\noindent\textbf{Contraction factor.}


Under DA2', the per-step Lyapunov function $V_k = \frac{1}{2}\lVert\delta_k\rVert^2$ satisfies:

% [Derived (contraction, from DA2')]

$$\lVert\delta_{k+1} - 0\rVert^2 \leq \lambda^2 \lVert\delta_k\rVert^2 + 2\lVert\delta_k\rVert \lVert w_k\rVert + \lVert w_k\rVert^2$$

where:

$$\lambda = \max(\lvert 1 - \eta^* c_{\min}\rvert,\; \lvert 1 - \eta^* c_{\max}\rvert)$$

DA2' ensures $\lvert\lambda\rvert < 1$: the lower bound gives $1 - \eta^* c_{\min} < 1$, and the upper bound gives $1 - \eta^* c_{\max} > -1$.


\medskip\noindent\textbf{Proposition DA.1: Bounded Mismatch (Deterministic).}


\textbf{Statement.} Under DA2' with bounded per-step disturbance $\lVert w_k\rVert \leq \rho_{\text{step}}$, the mismatch is ultimately bounded:

% [Derived (DA.1, discrete bounded mismatch)]

$$R^*_D = \frac{\rho_{\text{step}}}{1 - \lvert\lambda\rvert}$$

\textbf{Proof.} The contraction mapping yields:

$$\lVert\delta_{k+1}\rVert \leq \lvert\lambda\rvert \lVert\delta_k\rVert + \rho_{\text{step}}$$

This is an affine contraction with $\lvert\lambda\rvert < 1$. By the Banach fixed-point theorem, all trajectories starting in $\mathcal B_R$ converge to the ball of radius $R^*_D = \rho_{\text{step}}/(1 - \lvert\lambda\rvert)$, provided $R^*_D < R$. $\square$

\textbf{Recovery of continuous result.} In the fluid limit ($\eta^* \to 0$, $\nu \to \infty$, $\nu \eta^* = \mathcal{T}$ fixed), $\lambda \to 1 - \eta^* c_{\min}$ and $\rho_{\text{step}} \to \rho/\nu$. Then:

$$R^*_D = \frac{\rho/\nu}{\eta^* c_{\min}} = \frac{\rho}{\nu \eta^* c_{\min}} = \frac{\rho}{\alpha}$$

recovering Prop A.1 exactly. The discrete-to-continuous gap for Model D steady state is \textbf{zero}.


\medskip\noindent\textbf{Proposition DA.2: Adaptive Reserve (Discrete).}


\textbf{Statement.} Under the conditions of DA.1, the agent can absorb an additional per-step disturbance of:

% [Derived (DA.2, discrete adaptive reserve)]

$$\Delta\rho^*_{\text{step}} = (1 - \lvert\lambda\rvert) R - \rho_{\text{step}}$$

\textbf{Proof.} Identical to Prop A.2: the new $R^*_D = (\rho_{\text{step}} + \Delta\rho)/(1 - \lvert\lambda\rvert)$ must satisfy $R^*_D \leq R$. $\square$

The structure is identical to the continuous adaptive reserve $\Delta\rho^* = \alpha R - \rho$, with $\alpha$ replaced by the discrete contraction rate $(1 - \lvert\lambda\rvert)/\eta^*$ (which equals $\alpha$ in the fluid limit).


\medskip\noindent\textbf{Proposition DA.1S: Stochastic Bounded Mismatch (Discrete).}


\textbf{Statement.} Under DA2' with i.i.d. zero-mean disturbance $\mathbb{E}[w_k] = 0$, $\mathbb{E}[\lVert w_k\rVert^2] = \sigma^2_{\text{step}}$, the mismatch satisfies:

% [Derived (DA.1S, discrete stochastic bounded mismatch)]

$$\mathbb{E}[\lVert\delta_k\rVert^2] \leq \lambda^{2k}_{\text{eff}} \lVert\delta_0\rVert^2 + \frac{\sigma^2_{\text{step}}}{1 - \lambda^2_{\text{eff}}}$$

where $\lambda^2_{\text{eff}} = 1 - 2\eta^* c_{\min} + (\eta^*)^2 c^2_{\max}$.

\textbf{Proof.} Define $V_k = \lVert\delta_k\rVert^2$. Taking expectations of the squared update:

$$\mathbb{E}[V_{k+1} \mid \delta_k] = \lVert\delta_k - \eta^* F_d(\delta_k)\rVert^2 + \sigma^2_{\text{step}}$$

The first term satisfies:

$$\lVert\delta_k - \eta^* F_d(\delta_k)\rVert^2 = V_k - 2\eta^* \delta_k^T F_d(\delta_k) + (\eta^*)^2 \lVert F_d(\delta_k)\rVert^2$$

By DA2' (lower bound on $\delta^T F_d$ and upper bound implying $\lVert F_d\rVert^2 \leq c^2_{\max} V_k$):

$$\mathbb{E}[V_{k+1} \mid \delta_k] \leq \lambda^2_{\text{eff}} V_k + \sigma^2_{\text{step}}$$

This is a supermartingale (when $V_k$ is large enough). Iterating:

$$\mathbb{E}[V_k] \leq \lambda^{2k}_{\text{eff}} V_0 + \frac{\sigma^2_{\text{step}}}{1 - \lambda^2_{\text{eff}}}$$

The steady-state mean-square mismatch is $\sigma^2_{\text{step}} / (1 - \lambda^2_{\text{eff}})$. $\square$

\textbf{Recovery of continuous result.} In the fluid limit: $\sigma^2_{\text{step}} \to n\sigma^2_w / \nu$, $\lambda^2_{\text{eff}} \to 1 - 2\eta^* c_{\min}$, and $(1 - \lambda^2_{\text{eff}}) \to 2\eta^* c_{\min}$. The steady-state becomes $n\sigma^2_w / (2\nu \eta^* c_{\min}) = n\sigma^2_w / (2\alpha)$, recovering Prop A.1S exactly.

The discrete-to-continuous gap for Model S variance is $O((\eta^*)^2 c^2_{\max})$ — the $(\eta^*)^2 \lVert F_d\rVert^2$ term that vanishes in the fluid limit. This quantifies the error introduced by GA-5 and confirms it is small whenever $\eta^* c_{\max} \ll 1$.


\medskip\noindent\textbf{Fluid Limit Theorem.}


% [Derived (Conditional on Lipschitz regularity)]

\textbf{Statement.} Let $F_d$ be Lipschitz continuous with constant $L_F$ on $\mathcal B_R$. Let $\delta^{(\nu)}(t)$ denote the piecewise-constant interpolation of the discrete trajectory at event rate $\nu$, and $\delta(t)$ the solution of the continuous ODE $d\delta/dt = -F(\mathcal{T}, \delta) + w(t)$ with $F = \nu \eta^* F_d$. Then:

$$\sup_{t \in [0,T]} \lVert\delta^{(\nu)}(t) - \delta(t)\rVert \leq C \cdot \frac{\eta^* c_{\max}}{\nu^{1/2}}$$

for a constant $C$ depending on $T$, $L_F$, and $R$.

\textbf{Sketch.} This follows from the standard ODE-approximation theory for Euler schemes (Kushner \& Yin, 2003, Ch. 5). The discrete update $\delta_{k+1} = \delta_k - \eta^* F_d(\delta_k) + w_k$ is a forward Euler step for the ODE with step size $1/\nu$. The Lipschitz condition on $F_d$ ensures the local truncation error is $O(1/\nu^2)$. Summing over $O(\nu T)$ steps and applying the Gronwall inequality gives the $O(1/\nu^{1/2})$ bound.

For Model D (deterministic): the steady-state gap is exactly zero (both discrete and continuous converge to the same fixed point). The fluid-limit error affects only transients.

For Model S (stochastic): the steady-state variance gap is $O((\eta^*)^2 c^2_{\max})$, which equals $O(\eta^* c_{\max} / \nu)$ when expressed in terms of the event rate.


\medskip\noindent\textbf{Epistemic Status.}


\textbf{DA.1 and DA.2} are \textit{exact} — standard contraction-mapping results. The proofs require only the discrete sector condition DA2', which is itself derived from the gain principle + directional fidelity ( \S\,\ref{gain-sector-bridge}) for well-designed agents.

\textbf{DA.1S} is \textit{exact} under the stated i.i.d. assumption. The supermartingale argument is standard.

\textbf{Fluid limit theorem} is \textit{conditional} on Lipschitz regularity of $F_d$ — a standard regularity condition satisfied by all correction functions in the verified instances table ( \S\,\ref{gain-sector-bridge}). The convergence rate follows from classical ODE approximation theory (Kushner \& Yin, 2003); the application to the ACT mismatch dynamics is new but the mathematics is not.

\textbf{Max attainable:} exact for DA.1/DA.2/DA.1S; conditional for the fluid limit (Lipschitz cannot be removed).


\medskip\noindent\textbf{Discussion.}


\textbf{GA-5 is closed.} The fluid-limit bridging assumption is no longer required as an ungrounded assumption. For Model D, the discrete and continuous steady states are identical — the gap is zero. For Model S, the gap is $O(\eta^* c_{\max})$ in variance, quantitatively bounded and small in the regime where $\eta^* c_{\max} \ll 1$. The continuous-time results in \S\,\ref{sector-condition-derivation} are formally justified as the fluid limit of the discrete results here.

\textbf{The upper bound is the classical step-size condition.} The no-overshoot condition $c_{\max} < 2/\eta^*$ (equivalently $\eta^* < 2/c_{\max}$) is the same constraint as the step-size condition $\eta < 2/L$ for gradient descent, the stability condition for the Kalman gain, and the convergence condition for fixed-point iteration. For Bayesian updates, it is satisfied by construction. The discrete framework reveals this constraint, which is invisible in the continuous limit.

\textbf{No downstream results change qualitatively.} The persistence condition, adaptive reserve, and adversarial scaling laws derived in \S\,\ref{sector-condition-derivation} and \S\,\ref{adversarial-destabilization} hold as stated. The discrete framework provides sharper constants (replacing $\alpha$ with $(1 - \lvert\lambda\rvert)/\eta^*$) and a quantitative transient error bound, but the qualitative structure is unchanged.

\textbf{Section I's formal chain is now complete.} The prediction chain:

$$\text{gain principle} + \text{B1} \;\xrightarrow{\text{derived}}\; \text{sector condition} \;\xrightarrow{\text{Lyapunov/contraction}}\; \text{persistence, reserve, scaling}$$

holds in both continuous time (via \S\,\ref{sector-condition-derivation}) and discrete time (via this segment), with a quantitative bridge between them (the fluid limit theorem).

\end{derivation}

\begin{detail}[Simulation Results]
\label{simulation-results}
Six simulation variants validated, refined, and extended the analytical predictions from Section I, forcing regime splits and discovering scaling laws that the analytical derivations left ambiguous.



\medskip\noindent\textbf{Overview.}


The track-b simulation program tested Section I's analytical predictions about mismatch dynamics, adversarial coupling, observation noise, correction function robustness, and anisotropic persistence. The simulations modeled discrete-time mismatch dynamics as AR(1) processes with five correction functions (linear, saturating, threshold/dead-zone, sigmoid, structural breakdown), sweeping parameter spaces across gain, disturbance rate, coupling strength, observation noise, and dimensional structure.

The simulations were theory-shaping, not merely confirmatory. In particular:

\begin{itemize}
  \item The adversarial tempo exponent turned out to be regime-dependent (deterministic drift vs. stochastic noise, coupling-dominant vs. non-coupling-dominant), forcing a regime split that the mismatch ODE left ambiguous.
  \item Observation noise was found to gate the adversarial advantage, partially recoverable by optimal gain -- a finding that had no analytical treatment before the simulations.
  \item The scalar persistence condition was shown to overestimate adaptive capacity by 72\% in anisotropic systems, motivating the per-dimension persistence condition.
\end{itemize}



\medskip\noindent\textbf{Variant Summary.}


\begin{adjustbox}{max width=\textwidth}
\small
\begin{tabular}{llllll}
\toprule
Variant & Disturbance Model & Tested & Key Finding & Theory Impact & Promoted Segment \\
\midrule
A & \textbf{D} (deterministic drift) & Deterministic drift coupling; coupling-dominance sweep & Exponent $b \to 2.0$ in coupling-dominant limit (confirmed at 1.999) & Validates Model D derivation ($b = 2$) & \S\,\ref{adversarial-exponent-regimes} \\
B & \textbf{D+S} (interpolation) & Drift-noise interpolation across correction functions & Smooth transition between drift ($b = 2.0$) and noise ($b = 1.5$) regimes; coupling dominance is the key qualifier, not drift vs. noise per se & Confirmed that the coupling-dominant condition is quantitatively load-bearing & \S\,\ref{adversarial-exponent-regimes} \\
C & \textbf{S} (stochastic AR(1)) & Exponent vs. gain ($\eta$) in stochastic model & Exponent drops toward 0.5 as $\eta \to 0$ (away from coupling dominance at fixed $q_\text{base}$); discrete AR(1) exponent never exceeds 1.5 & Validates Model S derivation ($b = 3/2$) & \S\,\ref{adversarial-exponent-regimes} \\
D & \textbf{D vs. S} (separated) & Exponent vs. base noise ($q_\text{base}$) at fixed $\eta$ & Continuous ODE exponent $\to$ 2.0 (Model D), discrete AR(1) exponent $\to$ 1.5 (Model S), as $q_\text{base} \to 0$ & Definitively separated the two asymptotes; validates both derivations & \S\,\ref{adversarial-exponent-regimes} \\
E & \textbf{S} (stochastic AR(1)) & Observation noise; optimal gain validation & Observation noise collapses adversarial exponent from $\sim 1.0$ to $\sim 0.2$; Riccati-optimal gain restores it to $\sim 0.4$; 52\% mismatch reduction at moderate noise & Observation quality gates tempo advantage; optimal gain ( \S\,\ref{update-gain}) empirically validated & \S\,\ref{observation-gates-advantage} \\
F & \textbf{S} (stochastic AR(1)) & Multi-dimensional anisotropic correction; targeted adversarial attack & Per-dimension theory exact to 4 significant figures; scalar tempo overestimates by 72\%; targeted attack amplifies advantage by 17\% & Validates Model S per-dimension steady state & \S\,\ref{per-dimension-persistence} \\
Hafez bridge & \textbf{S} (stochastic AR(1)) & Bi-predictability $P$ vs. ACT mismatch in adversarial and non-adversarial settings & $P$ measures coupling architecture (scale-invariant); mismatch measures coupling performance; $P$ is blind to adversarial dynamics & $P$ and ACT mismatch are complementary diagnostics; $H_b$ (agent opacity) has no direct ACT analog -- potential gap for multi-agent work & -- \\
\bottomrule
\end{tabular}
\end{adjustbox}



\medskip\noindent\textbf{Methodology.}


\textbf{Discrete mismatch dynamics.} The environment follows a random walk $x_{t+1} = x_t + w_t$ with $w_t \sim N(0, q^2)$. The agent corrects via $\hat x_{t+1} = \hat x_t + \eta \cdot g(\delta_t)$, yielding the AR(1) mismatch process $\delta_{t+1} = (1 - \eta) \cdot \delta_t + w_t$ for linear $g$. This is the discrete-time analog of the mismatch ODE $d\Vert\delta\Vert/dt = -\mathcal{T} \cdot \Vert\delta\Vert + \rho$ from \S\,\ref{mismatch-dynamics}.

\textbf{Parameter sweeps.} Each variant swept its key parameter(s) across 7--20 values. Monte Carlo: 200 independent trials per parameter point, 10,000--20,000 timesteps per trial, with 2,000--5,000 step burn-in for steady-state convergence. Fixed random seeds for reproducibility.

\textbf{Exponent fitting.} Adversarial exponents were estimated by fitting $\log(\Vert\delta_B\Vert / \Vert\delta_A\Vert) = a + b \cdot \log(\mathcal T_A / \mathcal T_B)$ via weighted least squares, with weights inversely proportional to variance of each point's log-estimate. 95\% confidence intervals via bootstrap (1,000 samples).

\textbf{Correction functions.} Five functions $g: \mathbb{R} \to \mathbb{R}$ were tested, all satisfying $g(0) = 0$ and $g'(0) = 1$: linear ($g(\delta) = \delta$), saturating ($g(\delta) = \delta / (1 + \lvert\delta\rvert/R)$), threshold ($g(\delta) = \delta \cdot \mathbf{1}[\lvert\delta\rvert > \epsilon]$), sigmoid ($g(\delta) = R \cdot \tanh(\delta/R)$), and structural breakdown ($g(\delta) = \delta \cdot \mathbf{1}[\lvert\delta\rvert < R_\text{max}]$).

\textbf{Simulation code.} All code is in \texttt{../../msc/track-b-nonlinear-sims/}. Initial simulations: \texttt{sim1\_nonlinear\_mismatch.py} (single-agent), \texttt{sim2\_adversarial\_coupling.py} (two-agent). Variant extensions: \texttt{variants/variant\_ab\_drift.py}, \texttt{variants/variant\_cd\_regimes.py}, \texttt{variants/variant\_ef\_extensions.py}, \texttt{variants/variant\_hafez\_bridge.py}. Detailed result write-ups: \texttt{variants/variant\_ab\_results.md}, \texttt{variants/variant\_cd\_results.md}, \texttt{variants/variant\_ef\_results.md}, \texttt{variants/variant\_hafez\_results.md}.


\medskip\noindent\textbf{Key Findings.}



\medskip\noindent\textbf{The adversarial exponent is regime-dependent (now derived).}


The theory now distinguishes two disturbance models (Model D: bounded deterministic, GA-2; Model S: stochastic zero-mean, GA-2S), each producing a different steady-state scaling and therefore a different adversarial exponent:

\begin{itemize}
  \item \textbf{Model D, coupling-dominant:} $\lVert\delta\rVert_{ss} = \rho / \mathcal{T}$. Adversarial exponent $b = 2$ (derived from Prop A.1; Variant A confirms at 1.999).
  \item \textbf{Model S, coupling-dominant:} $\lVert\delta\rVert_{\text{rms}} = \sigma_w / \sqrt{2\mathcal{T}}$ (from Prop A.1S). Adversarial exponent $b = 3/2$ (derived; Variants C-D confirm at 1.481).
  \item \textbf{Non-coupling-dominant:} Exponent degrades smoothly toward 1.0 (Model D) or 0.5 (Model S) as base disturbance grows relative to adversarial coupling.
\end{itemize}


The original sim2 result ($b \approx 1.05$) was not a falsification of Corollary 11.2 but a measurement in the wrong regime -- stochastic noise coupling at moderate coupling dominance. Variants A--D systematically mapped the full regime space and validated the derived exponents. See \S\,\ref{adversarial-exponent-regimes} for the full treatment.


\medskip\noindent\textbf{Observation noise gates adversarial advantage.}


Adding observation noise $\sigma_\text{obs}$ to the mismatch signal collapsed the adversarial exponent from $\sim 1.0$ to $\sim 0.2$ at $\sigma_\text{obs} = 10 \times q_\text{env}$ (fixed gain). The Riccati-optimal gain ( \S\,\ref{update-gain}) partially restored the advantage to $\sim 0.4$, more than doubling it but not recovering the noise-free level. The optimal gain helps most in the moderate-noise regime, where it achieved a 52\% mismatch reduction over fixed gain. See \S\,\ref{observation-gates-advantage}.


\medskip\noindent\textbf{Per-dimension persistence is exact; scalar overestimates.}


In a 3-dimensional system with 5:1 gain variation across dimensions, the per-dimension AR(1) steady-state prediction matched simulation to 4 significant figures. The scalar persistence condition overestimated aggregate mismatch by 72\%, because it averages across dimensions while the weak dimension dominates the $L_2$ norm. Isotropic gain allocation reduced overall mismatch by 13\%; targeted adversarial attack on the weak dimension amplified the mismatch ratio by 17\%. See \S\,\ref{per-dimension-persistence}.


\medskip\noindent\textbf{Nonlinear correction creates thresholds, not lower exponents.}


Under deterministic drift, saturating, sigmoid, and breakdown correction functions did not simply reduce the adversarial exponent. Instead, they produced catastrophic divergence when disturbance exceeded the correction capacity ($\rho > \mathcal{T} \cdot R$). This is the persistence threshold failure from \S\,\ref{persistence-condition}, observed directly in simulation. The measured "exponents" above 3.0 for these functions were divergence artifacts, not meaningful scaling laws.


\medskip\noindent\textbf{Hafez bridge: architecture vs. performance.}


Bi-predictability $P$ (Hafez et al.) measures the informational architecture of agent-environment coupling; ACT mismatch measures the operational performance. In adversarial settings, $P$ remained nearly constant ($\sim 0.268$) across a 10:1 range of opponent tempo, while the mismatch ratio varied from 0.54 to 5.95. $P$ is scale-invariant after z-score normalization, making it insensitive to adversarial dynamics. The two metrics are complementary: $P$ diagnoses whether the coupling structure is sound; mismatch predicts how well it performs.


\medskip\noindent\textbf{Epistemic Status.}


\textit{Empirical.} Simulation results are reproducible (code in \texttt{../../msc/track-b-nonlinear-sims/}, fixed seeds, all results documented in variant write-ups). The key results -- regime-dependent exponents, observation noise gating, per-dimension exactness -- have been promoted to first-class segments ( \S\,\ref{adversarial-exponent-regimes}, \S\,\ref{observation-gates-advantage}, \S\,\ref{per-dimension-persistence}) with their own epistemic assessments. This appendix serves as reference for the simulation program as a whole.


\medskip\noindent\textbf{Discussion.}


\textbf{Internal validation, not external.} The simulations operate in the model's own terms: AR(1) processes with parameterized correction functions, Gaussian noise, and discrete-time dynamics. They validate the mathematical predictions within the model -- confirming that the analytical steady-state formulas, the exponent regimes, and the per-dimension theory are correct as statements about these stochastic processes. The gap between "the math predicts X within its own model" and "real agents exhibit X" remains an empirical question for each domain.

\textbf{What the simulations did not test.} The simulations did not test whether the AR(1) mismatch dynamics are a good model for any particular real-world adaptive system. They also did not test cross-dimensional coupling (off-diagonal correction), non-Gaussian disturbances, non-stationary parameters, or multi-agent composition dynamics. The interaction between observation noise and the deterministic-drift exponent regime ($b = 2.0$) was not tested -- Variant E used stochastic coupling only.

\textbf{Theory-shaping role.} The most important contribution of the simulations was not confirming predictions but forcing the theory to be more precise. The regime split in the adversarial exponent, the observation-noise gating mechanism, and the per-dimension bottleneck effect were all findings that the analytical derivations initially left ambiguous or unaddressed. The simulations forced the disturbance model split (Model D vs. Model S), which in turn enabled the analytical derivation of both exponents. The simulations now serve as validation of the derived results: Variant A validates $b = 2$ (Model D), Variants C-D validate $b = 3/2$ (Model S), Variant F validates the Model S per-dimension steady state.

\end{detail}

\section{Appendices: Operational Domains}

\begin{quote}\itshape
Operational-specific appendices and end-to-end domain instantiations validating the theory chain.
\end{quote}

\begin{detail}[Operationalization — Estimation Procedures]
\label{operationalization}
Estimation recipes for core ACT quantities, bridging the measurement gap between formal objects and practical deployment.



\medskip\noindent\textbf{Measurement Targets.}


\begin{adjustbox}{max width=\textwidth}
\small
\begin{tabular}{llll}
\toprule
Quantity & Role in ACT & Typical unit & Minimum data needed \\
\midrule
$U_M$ & Model uncertainty ( \S\,\ref{update-gain}) & domain-specific variance/entropy & Predictive posterior or ensemble spread \\
$U_o$ & Observation uncertainty ( \S\,\ref{update-gain}) & sensor variance/noise scale & Channel calibration or residual variance \\
$\rho_{\det}(t)$ & Mismatch injection rate bound, Model D ( \S\,\ref{mismatch-dynamics}) & surprise per time & Time-series of mismatch magnitudes; worst-case bounds \\
$\sigma_w$ & Disturbance noise scale, Model S ( \S\,\ref{mismatch-dynamics}) & surprise per $\sqrt{\text{time}}$ & Mismatch innovation variance (residual after correction) \\
$\rho_{\text{delib}}$ & Local mismatch drift during pauses ( \S\,\ref{deliberation-cost}) & surprise per time & Deliberation windows with no corrective action \\
$\alpha$ & Lower correction efficiency bound ( \S\,\ref{sector-condition-stability}) & inverse time & Vector mismatch trajectories + correction term \\
$R$ & Radius where local sector condition holds ( \S\,\ref{sector-condition-stability}) & surprise magnitude & Same as $\alpha$, plus breakdown detection \\
$\Vert\delta_{\text{critical}}\Vert$ & Functional adequacy threshold ( \S\,\ref{persistence-condition}) & surprise magnitude & Task-level performance curve vs mismatch \\
\bottomrule
\end{tabular}
\end{adjustbox}



\medskip\noindent\textbf{Estimator Cookbook.}



\medskip\noindent\textbf{Estimating $U_M$ and $U_o$.}


Use the most native uncertainty representation available in the domain:

\begin{adjustbox}{max width=\textwidth}
\small
\begin{tabular}{lll}
\toprule
Domain & $U_M$ estimator & $U_o$ estimator \\
\midrule
Kalman / linear Gaussian & Prior predictive variance $P_{t \vert t-1}$ & Measurement-noise variance $R_t$ (known or EM-estimated) \\
Conjugate Bayes & Posterior variance / inverse effective sample size & Likelihood variance (or precision inverse) \\
RL with ensembles & Across-head predictive variance $\operatorname{Var}_i[Q_i(s,a)]$ & TD-target noise variance over replay batches \\
Neural net regression & Ensemble or Laplace posterior variance & Aleatoric head output $\sigma^2(x)$ \\
PID / classical control & State-estimation covariance from observer & Sensor noise from calibration + residual PSD \\
\bottomrule
\end{tabular}
\end{adjustbox}


For the scalar gain heuristic ( \S\,\ref{update-gain}), normalize to common units and compute:

% [Operational Definition]

$$\hat{\eta}^*_t = \frac{\hat{U}_{M,t}}{\hat{U}_{M,t} + \hat{U}_{o,t}}$$


\medskip\noindent\textbf{Estimating Disturbance Parameters ($\rho_{\det}$, $\sigma_w$, $\rho_{\text{delib}}$).}


The disturbance model split (Model D vs. Model S; see \S\,\ref{mismatch-dynamics}, \S\,\ref{sector-condition-stability}) requires different estimation procedures depending on the disturbance character.

\textbf{Identifying the disturbance model.} The choice between Model D and Model S is domain-dependent:
\begin{itemize}
  \item \textbf{Model D} (bounded deterministic, GA-2): appropriate when the environment changes persistently and directionally. Estimate $\rho_{\det}$ from worst-case bounds on mismatch injection rate. Examples: adversary who maneuvers systematically, API that drifts, climate that shifts.
  \item \textbf{Model S} (stochastic zero-mean, GA-2S): appropriate when the environment fluctuates unpredictably around a stable mean. Estimate $\sigma_w$ from residual variance of mismatch innovations. Examples: market noise, sensor noise, random perturbations.
  \item \textbf{Mixed environments} use both: estimate $\rho_{\det}$ from the persistent component (trend) and $\sigma_w$ from the residual variance after detrending. The adversarial exponent interpolates smoothly between $b = 2$ (pure drift) and $b = 3/2$ (pure noise); see \S\,\ref{adversarial-exponent-regimes}.
\end{itemize}


Let $s_t = \lVert\delta_t\rVert$ in surprise units (e.g., negative log-likelihood residual scale).

**Model D: Global mismatch injection rate ($\hat{\rho}_{\det}$):**

% [Operational Definition]

$$\hat{\rho}_{\det}(t) = \left[\frac{s_{t+\Delta t} - s_t}{\Delta t} + \hat{\mathcal T}_t \, s_t\right]_+$$

where $[x]_+ = \max(x, 0)$ and $\hat{\mathcal T}$ evaluated at time $t$ is estimated adaptive tempo. Take the supremum (or a high quantile) over the estimation window.

**Model S: Noise scale ($\hat{\sigma}_w$):**

% [Operational Definition]

$$\hat{\sigma}_w^2 = \text{Var}\left[\frac{s_{t+\Delta t} - (1 - \hat{\eta}^*)s_t}{\sqrt{\Delta t}}\right]$$

i.e., the variance of the mismatch innovations (residuals after subtracting the expected correction). For AR(1) processes, this is the innovation variance $\hat{\rho}^2$ directly.

\textbf{Note on estimation sequencing.} Both estimators require $\hat{\mathcal T}_t$ (or $\hat{\eta}^*$), estimated from $\hat{\nu}$ and $\hat{\eta}^*$. Estimate the gain and event rate first (from the agent's internal statistics and observation timing), then use these to extract the disturbance parameters from the mismatch trajectory. This sequential structure avoids circularity but introduces sensitivity: errors in $\hat{\mathcal T}$ propagate into $\hat{\rho}_{\det}$ and $\hat{\sigma}_w$.

**Local pause-window drift for \S\,\ref{deliberation-cost} ($\hat{\rho}_{\text{delib}}$):**

% [Operational Definition]

$$\hat{\rho}_{\text{delib}} = \operatorname*{median}_{w \in \mathcal{W}_{\text{pause}}} \frac{s_{w,\text{end}} - s_{w,\text{start}}}{\Delta\tau_w}$$

using windows where corrective action is suspended or effectively delayed.


\medskip\noindent\textbf{Estimating $\alpha$ (sector lower bound).}


\S\,\ref{sector-condition-derivation} uses $\delta^T F(\mathcal{T}, \delta) \geq \alpha \Vert\delta\Vert^2$ for $\Vert\delta\Vert \leq R$. Operationally:

\begin{enumerate}
  \item Estimate $\dot{\delta}_t$ (finite differences or filtered derivative).
  \item Compute $\widehat F_t = -\dot{\delta}_t + w_t$ where disturbance proxy $w_t$ is estimated from exogenous perturbation channels or residual balancing.
  \item Form ratios $r_t = (\delta_t^T \widehat F_t) / \Vert\delta_t\Vert^2$ on bins of $\Vert\delta_t\Vert$.
  \item Set conservative lower bound $\hat{\alpha}$ as a low quantile (e.g., 10th percentile) of $r_t$ in the valid region.
\end{enumerate}



\medskip\noindent\textbf{Estimating $R$ (valid-region radius).}


Estimate $R$ as the largest radius for which sector inequality violations remain below tolerance:

% [Operational Criterion]

$$\hat{R} = \sup \left\{ r > 0 : \Pr\left(\delta^T \widehat{F} < \hat{\alpha}\Vert\delta\Vert^2 \,\middle\vert\, \Vert\delta\Vert \le r\right) \le \epsilon \right\}$$

with a chosen violation tolerance $\epsilon$ (e.g., 5\%).


\medskip\noindent\textbf{Estimating $\Vert\delta_{\text{critical}}\Vert$.}


Define a mission-level performance metric $J$ (reward rate, tracking error, service SLA, etc.). Set the critical mismatch threshold where performance crosses a minimum acceptable level $J_{\min}$:

% [Operational Definition]

$$\Vert\hat{\delta}_{\text{critical}}\Vert = \inf \left\{ d : \mathbb{E}[J \mid \Vert\delta\Vert = d] < J_{\min} \right\}$$

This anchors \S\,\ref{persistence-condition} to real task outcomes.


\medskip\noindent\textbf{Recommended Estimation Sequence.}


\begin{enumerate}
  \item Fix mismatch representation $\delta$ in one consistent unit system (prefer surprise-scale).
  \item Estimate $U_o$ from channel physics/calibration; estimate $U_M$ from model uncertainty.
  \item Validate gain behavior against \S\,\ref{update-gain} ($\hat{\eta}^*$ trend checks).
  \item Estimate $\rho_{\text{delib}}$ from pause windows ( \S\,\ref{deliberation-cost}) and $\rho(t)$ from full traces.
  \item Estimate $\alpha$ and $R$ from local correction dynamics ( \S\,\ref{sector-condition-derivation}).
  \item Estimate $\Vert\delta_{\text{critical}}\Vert$ from task-performance degradation.
  \item Compute derived diagnostics:
\end{enumerate}


\begin{itemize}
  \item Tempo margin (Model D): $\hat{\mathcal T} - \hat{\rho}_{\det}/\lVert\hat{\delta}_{\text{critical}}\rVert$
  \item Tempo margin (Model S): $\hat{\mathcal T} - n\hat{\sigma}_w^2/(2\lVert\hat{\delta}_{\text{critical}}\rVert^2)$
  \item Reserve: $\widehat{\Delta \rho^*} = \hat{\alpha}\hat{R} - \hat{\rho}_{\det}$
  \item Deliberation feasibility: $\Delta\eta^*(\Delta\tau)\Vert\delta_{\text{post}}\Vert - \hat{\rho}_{\text{delib}}\Delta\tau$
\end{itemize}



\medskip\noindent\textbf{Decision-Theoretic Procedures.}



\medskip\noindent\textbf{Estimating $\lambda$ (Exploration Price).}


The exploration weight $\lambda(M_t)$ in the unified policy objective ( \S\,\ref{ciy-unified-objective}) prices information in value-equivalent terms:

\begin{adjustbox}{max width=\textwidth}
\small
\begin{tabular}{lll}
\toprule
Context & $\lambda$ estimator & Source \\
\midrule
Finite bandits & Gittins index from dynamic programming & Exact (Gittins 1979) \\
Linear-Gaussian & Probing cost in quadratic objective & Exact (dual control) \\
Discrete MDP & $(\text{VoI})^2 / \text{info gain}$ & Information-directed sampling (Russo \& Van Roy) \\
General & $\hat{\lambda} = c \cdot \hat U_M / \hat U_o$ & Heuristic: scale CIY weight by relative uncertainty \\
\bottomrule
\end{tabular}
\end{adjustbox}


For the heuristic: when $U_M \gg U_o$ (highly uncertain model), exploration is cheap relative to exploitation risk, so $\lambda$ should be large. When $U_M \ll U_o$ (confident model, noisy observations), exploitation dominates. The constant $c$ is domain-specific.


\medskip\noindent\textbf{Deliberation Stopping Policy.}


From \S\,\ref{deliberation-cost}, deliberation of duration $\Delta\tau$ is warranted when $\Delta\eta^*(\Delta\tau) \cdot \Vert\delta_{\text{post}}\Vert > \rho_{\text{delib}} \cdot \Delta\tau$. Operationally:

\begin{enumerate}
  \item Estimate $\rho_{\text{delib}}$ from prior pause windows.
  \item Before each deliberation episode, estimate $\Vert\delta_{\text{post}}\Vert$ as current mismatch + $\rho_{\text{delib}} \cdot \Delta\tau_{\text{planned}}$.
  \item Estimate $\Delta\eta^*(\Delta\tau)$ from the diminishing-returns profile of past deliberation episodes.
  \item Stop deliberating when the marginal improvement rate $\partial \Delta\eta^* / \partial \Delta\tau$ drops below $\rho_{\text{delib}} / \Vert\delta_{\text{post}}\Vert$.
\end{enumerate}



\medskip\noindent\textbf{Structural-Switch Trigger.}


From \S\,\ref{structural-adaptation-necessity}, structural adaptation is indicated when parametric convergence leaves a mismatch floor. Operationally:

\begin{enumerate}
  \item Estimate the current mismatch floor $\Vert\delta\Vert_{\text{floor}}$ from converged residual statistics.
  \item Estimate post-switch expected mismatch as $\Vert\delta\Vert_{\text{new}} \approx \rho / \alpha'$ where $\alpha'$ is the sector bound under the candidate new model class.
  \item Estimate transition cost $C_{\text{switch}}$: knowledge loss, retraining time ($\Delta\tau_{\text{switch}}$), and accumulated mismatch during transition ($\rho \cdot \Delta\tau_{\text{switch}}$).
  \item Switch when: $(\Vert\delta\Vert_{\text{floor}} - \Vert\delta\Vert_{\text{new}}) \cdot T_{\text{horizon}} > C_{\text{switch}}$.
\end{enumerate}



\medskip\noindent\textbf{Estimator Uncertainty Guidance.}


**$\hat{\eta}^*$ (gain estimate).** Approximate variance via the delta method:

$$\text{Var}(\hat{\eta}^*) \approx \hat{\eta}^{*2}(1-\hat{\eta}^*)^2 \left[\frac{\text{Var}(\hat{U}_M)}{\hat{U}_M^2} + \frac{\text{Var}(\hat{U}_o)}{\hat{U}_o^2}\right]$$

Near 0 or 1, stable (dominated by one source). Near 0.5, most volatile.

**$\hat{\rho}$ (mismatch injection rate).** Finite-difference estimation amplifies noise. Recommend smoothing before differencing. Minimum ~20 observations for stable trend; ~50 for reliable variance. The $[\cdot]_+$ clipping introduces positive bias at small $\rho$.

**$\hat{\alpha}$ (sector lower bound).** The 10th-percentile approach gives approximately 90\% confidence under stationarity. Report quantile level, bin count, and bin width.

**$\hat{R}$ (sector-condition radius).** A sharp drop-off in sector-condition satisfaction is a strong signal; gradual decay suggests the sector condition may not hold cleanly.

\textbf{Nonstationarity caveat.} All estimators assume approximate stationarity over the estimation window. When environment regime changes are suspected ( \S\,\ref{structural-adaptation-necessity}), re-estimate from post-change data only.


\medskip\noindent\textbf{Epistemic Status.}


This is a \textit{procedures document} — estimation recipes, not theoretical claims. Each estimator inherits the epistemic status of the quantity it targets: the gain estimator's validity depends on \S\,\ref{update-gain}'s structural claim; the sector-bound estimator's validity depends on \S\,\ref{sector-condition-stability}'s assumptions. The estimation procedures themselves are \textit{conditional} on these theoretical foundations and on the stationarity and data-sufficiency assumptions noted above.

\end{detail}

\begin{remark}[Worked-example: Worked Example — 1D Active Tracking (Kalman Domain)]
\label{worked-example-kalman}
Every ACT quantity has an exact Kalman-filter counterpart. This is a \textit{validation} of the formal chain — all quantities are computable in closed form.



\medskip\noindent\textbf{System.}


The agent tracks scalar state $x_t$ and chooses sensor mode $a_t \in \{L, H\}$:

\begin{itemize}
  \item Dynamics: $x_{t+1} = x_t + v_t$, $v_t \sim \mathcal{N}(0, q)$, $q = 0.25$
  \item Observation: $y_t = x_t + n_t^{(a_t)}$
  \item Low mode noise: $r_L = 9$; High mode noise: $r_H = 1$
  \item Event rate: $\nu = 5 \text{ Hz}$
  \item High mode has higher energy cost
\end{itemize}



\medskip\noindent\textbf{Chain Instantiation.}



\medskip\noindent\textbf{Scope ( \S\,\ref{scope-condition}).}


\textit{Mapping: exact.}

$\Omega_t = x_t$ is partially observed via noisy channel. $\mathcal{A} = \{L, H\}$ is non-empty and causally affects observation quality. Residual uncertainty persists due to process and sensor noise.


\medskip\noindent\textbf{Causal Structure ( \S\,\ref{causal-structure}) + CIY ( \S\,\ref{causal-information-yield}).}


\textit{Mapping: exact.}

Action precedes observation and changes $P(y_t \mid do(a_t))$ through $r_{a_t}$. Using low mode as comparator action, the CIY is the interventional KL divergence:

% [Worked Quantity]

$$\text{CIY}(H) = D_{\mathrm{KL}}\!\big(P(y \mid do(H)) \,\Vert\, P(y \mid do(L))\big)$$

$$= \frac{1}{2}\left[\log\!\left(\frac{P^- + r_L}{P^- + r_H}\right) + \frac{P^- + r_H}{P^- + r_L} - 1\right]$$

With $P^- = 4.25$: $\text{CIY}(H) \approx 0.161 \text{ nats}$.


\medskip\noindent\textbf{Model ( \S\,\ref{agent-model}).}


\textit{Mapping: exact.}

Model state $M_t = (\hat x_{t\vert t}, P_{t\vert t})$ — a compression of interaction history with recursive update.


\medskip\noindent\textbf{Mismatch ( \S\,\ref{mismatch-signal}).}


\textit{Mapping: exact.}

$$\delta_t = y_t - \hat{x}_{t\vert t-1}$$


\medskip\noindent\textbf{Update Gain ( \S\,\ref{update-gain}).}


\textit{Mapping: exact.}

Scalar Kalman gain:

$$K_t = \frac{P^-_t}{P^-_t + r_{a_t}}$$

With $P^- = 4.25$: $K(H) \approx 0.810$, $K(L) \approx 0.321$.

The exact uncertainty ratio mapping: $U_M = P^-_t$, $U_o = r_{a_t}$, $\eta^* = K_t$.


\medskip\noindent\textbf{Exploration ( \S\,\ref{causal-information-yield}).}


% [Worked Objective]

$$a_t^* = \arg\max_a \left[\mathbb{E}[\text{value}(a)\mid M_t] + \lambda_t \, \mathbb{E}[\text{CIY}(a)\mid M_t]\right]$$

When uncertainty is high ($P^-$ large), CIY term favors high mode. As uncertainty falls, policy shifts toward low-cost $L$.


\medskip\noindent\textbf{Deliberation Threshold ( \S\,\ref{deliberation-cost}).}


Suppose a planning pause of $\Delta\tau = 0.5 \text{ s}$, with measured $\rho_{\text{delib}} = 0.40 \;\text{surprise/s}$.

Cost during pause: $0.20$ surprise units. If $\Vert\delta_{\text{post}}\Vert = 0.70$, deliberation is worthwhile when:

$$\Delta\eta^*(0.5)\cdot 0.70 > 0.20 \;\Longrightarrow\; \Delta\eta^*(0.5) > 0.286$$


\medskip\noindent\textbf{Structural Adaptation ( \S\,\ref{structural-adaptation-necessity}).}


Assume maneuvering regime change introduces sustained residual autocorrelation and mismatch floor. If estimated valid radius drops to $R = 0.08$ while $R^* = \rho/\alpha = 0.12$, parametric adaptation is no longer adequate ($R^* > R$), triggering model-class change (e.g., constant-velocity → constant-acceleration process model).


\medskip\noindent\textbf{Tempo + Persistence ( \S\,\ref{adaptive-tempo}, \S\,\ref{persistence-condition}).}


Using action mix $70\% H, 30\% L$:

$$\bar{\eta}^* = 0.7(0.810) + 0.3(0.321) = 0.663$$

$$\mathcal{T} = \nu \bar{\eta}^* = 5 \cdot 0.663 = 3.315 \;\text{s}^{-1}$$

With $\rho = 0.18 \text{ surprise/s}$ and $\Vert\delta_{\text{critical}}\Vert = 1$:

$$\mathcal{T} > \frac{\rho}{\Vert\delta_{\text{critical}}\Vert} \;\;\Rightarrow\;\; 3.315 > 0.18 \;\checkmark$$


\medskip\noindent\textbf{Lyapunov Bounds ( \S\,\ref{sector-condition-stability}, \S\,\ref{gain-sector-bridge}).}


\textit{Mapping: exact (derived, not estimated).}

By \S\,\ref{gain-sector-bridge}, the sector parameter for a scalar Kalman filter is $\alpha = \mathcal{T} = \nu \bar\eta^*$:

$$\alpha = \mathcal{T} = 3.315 \text{ s}^{-1}$$

This is derived from the Kalman gain, not estimated from data — the correction function $F(e) = Ke$ is exactly linear, so $\alpha = K$ per observation and $\alpha = \nu \bar K = \mathcal{T}$ in continuous time. With $R = 1.4$ (model class capacity, a domain parameter) and $\rho = 0.18$:

$$R^* = \frac{\rho}{\alpha} = \frac{0.18}{3.315} \approx 0.054 < R$$

$$\Delta\rho^* = \alpha R - \rho = 3.315(1.4) - 0.18 = 4.46$$

The agent is comfortably within its invariant region with substantial adaptive reserve.


\medskip\noindent\textbf{Mapping Quality Summary.}


\begin{adjustbox}{max width=\textwidth}
\small
\begin{tabular}{lll}
\toprule
ACT Concept & Kalman Mapping & Status \\
\midrule
Scope & Exact & Definitional \\
Causal structure + CIY & Exact & Closed-form KL \\
Model ($M_t$ as sufficient statistic) & Exact & Kalman state + covariance \\
Mismatch ($\delta_t$ = innovation) & Exact & Standard Kalman innovation \\
Gain ($\eta^* = K_t$) & Exact & Kalman gain IS uncertainty ratio \\
Tempo ($\mathcal{T} = \nu \bar{\eta}^*$) & Exact & Closed-form \\
Persistence condition & Exact & Linear ODE solution \\
Sector parameter ($\alpha = \mathcal{T}$) & Exact & Derived from Kalman gain via \S\,\ref{gain-sector-bridge} \\
Lyapunov bounds ($R^*$, $\Delta\rho^*$) & Exact & From derived sector parameter + domain $R$ \\
\bottomrule
\end{tabular}
\end{adjustbox}



\medskip\noindent\textbf{Epistemic Status.}


This is a \textit{worked instantiation}, not a theoretical claim. Every mapping is exact — the Kalman domain is the canonical case where ACT's formal chain has closed-form realizations. The example validates that the formal chain is internally consistent and instantiable.

\end{remark}

\begin{remark}[Worked-example: Worked Example — Nonstationary Bandit (RL Domain)]
\label{worked-example-bandit}
ACT's conceptual architecture applies beyond the conjugate-Bayesian regime. The mapping here is \textit{approximate} — it shows that ACT's concepts organize RL phenomena, but the quantitative relationships are structural analogies, not derivations. Where the mapping is exact versus approximate is marked explicitly.



\medskip\noindent\textbf{System.}


A $k$-armed bandit ($k = 4$) with slowly drifting reward means:


$$\mu_i(t+1) = \mu_i(t) + w_i(t), \quad w_i(t) \sim \mathcal{N}(0, q), \quad q = 0.01$$

\begin{itemize}
  \item Observation: Pulling arm $a_t = i$ yields reward $r_t = \mu_i(t) + \varepsilon_t$, $\varepsilon_t \sim \mathcal{N}(0, \sigma^2)$, $\sigma^2 = 1.0$.
  \item Event rate: $\nu = 1$ pull per step.
  \item Agent: Q-learning with fixed learning rate $\alpha$ and UCB action selection.
\end{itemize}



\medskip\noindent\textbf{Chain Instantiation.}



\medskip\noindent\textbf{Scope ( \S\,\ref{scope-condition}).}


\textit{Mapping: exact.}

$\Omega_t = (\mu_1(t), \ldots, \mu_4(t))$ — not directly observable. $\mathcal{A} = \{1, 2, 3, 4\}$. Residual uncertainty persists: means drift continuously, rewards are noisy, and only one arm is observed per step.


\medskip\noindent\textbf{Causal Structure ( \S\,\ref{causal-structure}) + CIY ( \S\,\ref{causal-information-yield}).}


\textit{Mapping: exact for causal ordering; approximate for CIY formula.}

The action (arm choice) temporally precedes the reward observation. Under the Gaussian reward model, CIY for arm $i$:

% [Discussion — Bandit CIY, Gaussian Case]

$$\text{CIY}(i;\, M_{t-1}) = \frac{1}{2\sigma^2} \mathbb{E}_{j \sim q}\!\left[(\hat{\mu}_i - \hat{\mu}_j)^2\right]$$

This measures how \textit{distinctive} arm $i$'s predicted reward is relative to alternatives — distinctiveness of predictions, not uncertainty about predictions. A proper uncertainty-aware CIY would require maintaining posterior variances per arm, which the basic Q-learning model does not track.


\medskip\noindent\textbf{Model ( \S\,\ref{agent-model}) and Sufficiency ( \S\,\ref{model-sufficiency}).}


\textit{Mapping: exact structural mapping; the model is deliberately impoverished.}

$$M_t = (\hat{\mu}_1, \ldots, \hat{\mu}_4,\; n_1, \ldots, n_4)$$

Model sufficiency $S(M_t) < 1$ for two reasons: (1) No drift tracking — the model treats reward means as stationary. (2) No uncertainty representation — unlike a Bayesian agent that would maintain posterior variances.


\medskip\noindent\textbf{Mismatch ( \S\,\ref{mismatch-signal}).}


\textit{Mapping: exact.}

$$\delta_t = r_t - \hat{\mu}_{a_t}$$

Decomposes exactly: $\mathbb{E}[\delta_t^2] = (\hat{\mu}_{a_t} - \mu_{a_t}(t))^2 + \sigma^2$.


\medskip\noindent\textbf{Update Gain ( \S\,\ref{update-gain}).}


\textit{Mapping: approximate. This is where the bandit agent is most deficient relative to ACT-optimal behavior.}

The Q-learning update

$$\hat{\mu}_{a_t} \leftarrow \hat{\mu}_{a_t} + \alpha \cdot \delta_t$$

uses a \textbf{degenerate constant gain}. It does not adapt to the agent's current uncertainty state.

\textbf{Optimal gain via effective window.} A fixed $\alpha$ is equivalent to exponential discounting with effective window $W = (1 - \alpha)/\alpha$. The per-arm model uncertainty balances estimation variance and drift-induced bias:

$$U_M \approx \frac{\sigma^2}{W} + q \cdot W$$

The optimal window $W^* = \sigma / \sqrt{q} = 10$ gives $\alpha^* \approx 0.091$, $U_M \approx 0.2$, and:

$$\eta^* = \frac{U_M}{U_M + U_o} = \frac{0.2}{1.2} \approx 0.167$$

The discrepancy between $\alpha^*$ and $\eta^*$ arises because $\alpha$ operates on raw prediction errors while $\eta^*$ accounts for the full uncertainty structure. The two converge when the model properly tracks its own uncertainty.


\medskip\noindent\textbf{Exploration ( \S\,\ref{causal-information-yield}).}


\textit{Mapping: approximate structural analogy.}

UCB: $a_t = \arg\max_i [\hat{\mu}_i + c\sqrt{\ln t / n_i}]$

\begin{adjustbox}{max width=\textwidth}
\small
\begin{tabular}{lll}
\toprule
ACT component & UCB analog & Quality \\
\midrule
$\mathbb{E}[\text{value}(a) \mid M_t]$ & $\hat{\mu}_i$ & Exact \\
$\lambda(M_t) \cdot \mathbb{E}[\text{CIY}(a)]$ & $c\sqrt{\ln t / n_i}$ & Approximate \\
\bottomrule
\end{tabular}
\end{adjustbox}


UCB captures the CIY structure (rarely-visited arms are more informative) but depends on visit count, not observation content.


\medskip\noindent\textbf{Tempo + Persistence ( \S\,\ref{adaptive-tempo}, \S\,\ref{persistence-condition}).}


\textit{Mapping: approximate.}

\textbf{Per-arm tempo.} With approximately uniform exploration, each arm sampled at rate $\nu/k = 1/4$:

$$\mathcal{T}_i = \frac{\nu}{k} \cdot \alpha = 0.25 \times 0.091 = 0.023 \;\text{step}^{-1}$$

\textbf{Per-arm drift rate:} $\rho_i = \sqrt{q} = 0.1 \;\text{reward units} \cdot \text{step}^{-1}$.

\textbf{Persistence check:} $\mathcal T_i = 0.023$ vs $\rho_i = 0.1$ — \textbf{fails}. Per-arm correction tempo is too low relative to drift. The agent cannot keep all arms' estimates current under uniform exploration.

\textbf{Interpretation.} This is expected and informative. ACT diagnoses exactly why: with 4 arms and one pull per step, each arm is visited too infrequently to track its drifting mean. Two remedies:

\begin{enumerate}
  \item **Increase $\eta^*$** (raise $\alpha$): Extract more per observation, but increase steady-state noise.
  \item **Concentrate $\nu$**: Abandon uniform exploration. Focus pulls on a subset, increasing per-arm $\nu/k_{\text{active}}$. This is exactly what a good UCB policy does.
\end{enumerate}


With focused exploration ($k_{\text{active}} = 2$, $\alpha = 0.2$): $\mathcal T_i = 0.5 \times 0.2 = 0.1$ — barely meets the threshold. The fundamental tension: with limited pulls and significant drift, the agent must accept either failing to track some arms or using high $\alpha$ that introduces noise.

\textbf{Aggregate failure.} $\mathcal{T} = \nu \cdot \alpha = 0.091$ vs $\rho = k \cdot \sqrt{q} = 0.4$. The agent's total adaptive capacity is outpaced by total environmental drift — a regime where model-based approaches (Bayesian bandits, Kalman bandit filters) with higher effective $\eta^*$ have a structural advantage.


\medskip\noindent\textbf{Mapping Quality Summary.}


\begin{adjustbox}{max width=\textwidth}
\small
\begin{tabular}{lll}
\toprule
ACT Concept & Bandit Mapping & Status \\
\midrule
Scope & Exact & Definitional \\
Causal structure & Exact & Structural \\
CIY & Approximate & Model lacks uncertainty representation \\
Model ($M_t$) & Exact & Definitional \\
Model sufficiency ($S < 1$) & Exact & Q-learner demonstrably insufficient \\
Mismatch ($\delta_t$) & Exact & Standard prediction error \\
Gain ($\alpha$ as degenerate $\eta^*$) & Approximate & Fixed, does not adapt \\
Exploration (UCB as approximate CIY) & Approximate & Structural analogy \\
Tempo ($\mathcal{T} = (\nu/k) \cdot \alpha$ per arm) & Approximate & Assumes uniform allocation \\
Persistence condition & Exact (qualitative) & Threshold structure is robust \\
\bottomrule
\end{tabular}
\end{adjustbox}



\medskip\noindent\textbf{Epistemic Status.}


This is a \textit{structural mapping}, not a derivation. The mapping is \textit{exact} for scope, causal ordering, model compression, and mismatch. It is \textit{approximate} for gain, exploration, and tempo — precisely because the Q-learner does not represent its own uncertainty, which is what \S\,\ref{update-gain} requires for optimal behavior. The gap between the Q-learner's fixed $\alpha$ and the ACT-optimal adaptive $\eta^*$ is precisely the gap between a fixed-gain PID controller and a Kalman filter.

The mapping status is \textit{conditional} because the quantitative relationships depend on the Gaussian reward model and specific parameterization. The qualitative conclusions (persistence failure under uniform exploration, the concentration-vs-noise tradeoff) should be robust.

\end{remark}

\begin{remark}[Nonstationary Bandit with Explicit Strategy (Section II Validation)]
\label{worked-example-strategy}
ACT's Section II machinery — objectives, strategy DAGs, the orient cascade, satisfaction gap, control regret, chain confidence decay, observability dominance, and edge update via gain — is exercised on a concrete system simple enough for closed-form analysis. The key payoff: Section II diagnostics tell the agent things that Section I quantities alone cannot.



\medskip\noindent\textbf{System.}


A 3-armed Bernoulli bandit with drifting success probabilities and an agent that maintains an explicit strategy DAG for arm selection.


\textbf{Environment.} Three arms with time-varying success probabilities:

$$\theta_k(t) \in [0, 1], \quad k \in \{1, 2, 3\}$$

At each step, pulling arm $k$ yields reward $r_t \in \{0, 1\}$ with $P(r_t = 1 \mid a_t = k) = \theta_k(t)$. The best arm changes slowly: at rate $\rho_\theta$, a random arm's success probability shifts by a random perturbation clipped to $[0.1, 0.9]$.

\textbf{Concrete parameterization.} At $t = 0$: $\theta_1(0) = 0.7$, $\theta_2(0) = 0.5$, $\theta_3(0) = 0.3$. Drift rate: $\rho_\theta = 0.05$ per step (one arm shifts by $\pm 0.05$ every 20 steps on average). Horizon: $N_h = 100$ steps.

\textbf{Agent.} Maintains Beta-Bernoulli beliefs per arm and an explicit strategy DAG encoding its plan for achieving a target reward rate.


\medskip\noindent\textbf{Section I Chain Instantiation (Summary).}


Section I quantities map as in \S\,\ref{worked-example-bandit}. The brief summary here establishes the epistemic baseline; the novel content is Section II.

\textbf{Scope} ( \S\,\ref{scope-condition}). $\Omega_t = (\theta_1(t), \theta_2(t), \theta_3(t))$; $\mathcal{A} = \{1, 2, 3\}$; residual uncertainty persists. \textit{Exact.}

\textbf{Model} ( \S\,\ref{agent-model}). $M_t = (\alpha_k, \beta_k)_{k=1}^{3}$ — Beta posteriors per arm, with exponential discounting (effective window $W = 20$):

$$\hat\theta_k = \frac{\alpha_k}{\alpha_k + \beta_k}$$

\textit{Exact structural mapping.}

\textbf{Mismatch} ( \S\,\ref{mismatch-signal}). $\delta_t = r_t - \hat\theta_{a_t}$. \textit{Exact.}

\textbf{Update gain} ( \S\,\ref{update-gain}). The Beta-Bernoulli conjugate update with discounted pseudo-counts gives:

$$\eta_k^* = \frac{1}{\alpha_k + \beta_k + 1}$$

For a well-observed arm with $\alpha_k + \beta_k \approx W = 20$: $\eta_k^* \approx 0.048$. \textit{Exact for the pulled arm; zero for unpulled arms.}

\textbf{Tempo} ( \S\,\ref{adaptive-tempo}). Per-arm tempo: $\mathcal T_k = \nu_k \cdot \eta_k^*$ where $\nu_k$ is the pull rate for arm $k$. Under uniform exploration: $\nu_k = 1/3$, so $\mathcal T_k \approx 0.016$ per step. \textit{Approximate.}


\medskip\noindent\textbf{Section II Chain Instantiation.}



\medskip\noindent\textbf{Objective ( \S\,\ref{objective-functional}).}


\textit{Mapping: exact.}

The agent's objective: achieve a target average reward rate over the horizon.


$$V_{O_t}(\tau_{0:N_h}) = \frac{1}{N_h}\sum_{t=0}^{N_h - 1} r_t$$

$$V_{O_t}^{\min} = 0.6$$

This is a utility-type objective with a satisfaction threshold: the agent counts the mission as successful if its average reward rate meets or exceeds $0.6$. The value functional evaluates trajectories on $[0, 1]$ (average reward).


\medskip\noindent\textbf{Value Object ( \S\,\ref{value-object}).}


\textit{Mapping: exact.}

Under the one-step improvement convention ($\pi_\text{cont} = \pi_\text{current}$):

$$V_O(M_t, \pi; N_h) = \mathbb{E}\!\left[\frac{1}{N_h}\sum_{s=t}^{t + N_h - 1} r_s \;\middle\vert\; M_t, \pi\right]$$

For a greedy policy that always pulls the estimated-best arm $k^* = \arg\max_k \hat\theta_k$:

$$V_O(M_t, \pi_\text{greedy}; N_h) = \hat\theta_{k^*}$$

This is tractable: the agent's expected average reward under greedy play equals its belief about the best arm's success probability. With $M_0$: $V_O = \hat\theta_1 = 0.7$.

\textbf{Action-value form:}

$$Q_O(M_t, k; \pi_\text{greedy}, N_h) \approx \frac{1}{N_h}\hat\theta_k + \frac{N_h - 1}{N_h}\hat\theta_{k^*}$$

The first pull contributes $1/N_h$ of the horizon; the rest follows the greedy policy.


\medskip\noindent\textbf{Objective Attainability ( \S\,\ref{satisfaction-gap}).}


\textit{Mapping: exact.}

$$A_O(M_t; \Pi, N_h) = \sup_{\pi \in \Pi} V_O(M_t, \pi; N_h)$$

For a bandit with Beta-Bernoulli beliefs, the optimal policy (Gittins index or Bayes-adaptive) achieves an attainability that is bounded:

$$\max_k \hat\theta_k \leq A_O \leq \max_k \hat\theta_k + \epsilon_\text{info}$$

where $\epsilon_\text{info}$ accounts for the information value of exploration (pulling a suboptimal arm now to identify the true best arm later). The upper bound is tight when exploration can reveal a better arm. For our initial beliefs: $A_O \approx 0.70 + \epsilon_\text{info}$.

\textbf{Satisfaction gap:}

$$\delta_\text{sat} = V_{O_t}^{\min} - A_O = 0.6 - 0.70 \approx -0.10$$

$\delta_\text{sat} < 0$: the objective is \textbf{attainable}. The agent believes it can achieve a $0.6$ average reward rate. The negative gap is margin — the agent has room to spare.


\medskip\noindent\textbf{Control Regret ( \S\,\ref{control-regret}).}


\textit{Mapping: exact.}

Under an explore-then-exploit policy that spends 30 steps exploring uniformly, then exploits the estimated-best arm:

$$V_O(M_t, \pi_\text{explore-exploit}; 100) = \frac{30}{100}\cdot\bar\theta + \frac{70}{100}\cdot\hat\theta_{k^*\!(30)}$$

where $\bar\theta = (\hat\theta_1 + \hat\theta_2 + \hat\theta_3)/3$ is the average expected reward during exploration and $\hat\theta_{k^*\!(30)}$ is the estimated-best arm after 30 observations. With initial beliefs: $V_O \approx 0.3 \cdot 0.5 + 0.7 \cdot 0.7 = 0.64$.

$$\delta_\text{regret} = A_O - V_O(\pi_\text{explore-exploit}) \approx 0.70 - 0.64 = 0.06$$

The agent is leaving $0.06$ reward-rate points on the table with its current explore-then-exploit strategy — it explores too much. This is the signal for strategy revision ( \S\,\ref{orient-cascade} step 3).

\textbf{The 2x2 diagnostic} ( \S\,\ref{control-regret} Discussion):

\begin{adjustbox}{max width=\textwidth}
\small
\begin{tabular}{lll}
\toprule
 & $\delta_\text{sat} \leq 0$ & $\delta_\text{sat} > 0$ \\
\midrule
$\delta_\text{regret} \approx 0$ & Success & Capability limit \\
$\delta_\text{regret} \gg 0$ & \textbf{Strategy problem} & Both \\
\bottomrule
\end{tabular}
\end{adjustbox}


Our agent is in the \textbf{Strategy problem} cell: the goal is achievable, but the current policy is suboptimal. Section I alone (which reports $\hat\theta_k$, $\eta^*_k$, and $\mathcal T_k$) cannot distinguish "good strategy, hard goal" from "bad strategy, easy goal." The 2x2 diagnostic is Section II's value-add.


\medskip\noindent\textbf{Strategy DAG ( \S\,\ref{strategy-dag}, \S\,\ref{and-or-scope}).}


\textit{Mapping: exact for structure and propagation; the DAG is simple enough for closed-form status.}

The agent's strategy for achieving $V_{O_t}^{\min} = 0.6$:

```
[root: "achieve ≥0.6 avg reward"] (AND)
/                    \
[Phase 1: "identify best arm"] (OR)    [Phase 2: "exploit best arm"]
/         |         \                      |
[arm 1       [arm 2       [arm 3         [pull identified
is best]     is best]     is best]        arm reliably]
```

\textbf{Nodes:}
\begin{itemize}
  \item $v_\text{root}$: AND-node. Both phases must succeed: identify the best arm AND exploit it. Terminal satisfaction condition: average reward $\geq 0.6$.
  \item $v_\text{id}$: OR-node (Phase 1). Identifying the best arm succeeds if ANY one arm is correctly identified as best.
  \item $v_1, v_2, v_3$: leaf condition nodes. "Arm $k$ is the best arm." Base credence: $p_k(M_t) = P(\theta_k > \theta_j \;\forall j \neq k \mid M_t)$.
  \item $v_\text{exploit}$: leaf action node (Phase 2). "Pulling the identified arm yields reward $\geq 0.6$." Base credence: $p_\text{exploit}(M_t) = P(\theta_{k^*} \geq 0.6 \mid M_t)$.
\end{itemize}


\textbf{Edge credences.} The edges from leaf nodes to their parents carry causal credences:
\begin{itemize}
  \item $p_{v_k, v_\text{id}}$: "correctly identifying arm $k$ as best enables Phase 1 success." These are structural (close to $1.0$) — if the agent correctly knows which arm is best, Phase 1 succeeds by definition. Set $p_{v_k, v_\text{id}} = 1.0$ for all $k$.
  \item $p_{v_\text{id}, v_\text{root}}$: "having identified the best arm advances the root goal." Set to $0.95$ (slight uncertainty about whether identification translates to exploitation success).
  \item $p_{v_\text{exploit}, v_\text{root}}$: "successful exploitation advances the root goal." Set to $0.90$ (the identified arm might drift below threshold during exploitation).
\end{itemize}


**Leaf base credences from $M_t$.** At $t = 0$ with $(\hat\theta_1, \hat\theta_2, \hat\theta_3) = (0.7, 0.5, 0.3)$ and moderate pseudo-counts $(\alpha_k + \beta_k \approx 5)$ from prior experience:

$$p_1(M_0) = P(\theta_1 > \theta_2, \theta_1 > \theta_3 \mid M_0) \approx 0.75$$

$$p_2(M_0) \approx 0.20, \quad p_3(M_0) \approx 0.05$$

$$p_\text{exploit}(M_0) = P(\theta_{k^*} \geq 0.6 \mid M_0) \approx 0.70$$

\textbf{Status propagation} ( \S\,\ref{strategy-dag}, \S\,\ref{and-or-scope}):

Phase 1 (OR-node):

$$s_{v_\text{id}} = 1 - \prod_{k=1}^{3}(1 - p_{v_k, v_\text{id}} \cdot p_k) = 1 - (1 - 0.75)(1 - 0.20)(1 - 0.05)$$

$$= 1 - (0.25)(0.80)(0.95) = 1 - 0.19 = 0.81$$

Phase 2 (leaf): $s_{v_\text{exploit}} = 0.70$.

Root (AND-node):

$$s_{v_\text{root}} = (p_{v_\text{id}, v_\text{root}} \cdot s_{v_\text{id}}) \cdot (p_{v_\text{exploit}, v_\text{root}} \cdot s_{v_\text{exploit}})$$

$$= (0.95 \times 0.81) \times (0.90 \times 0.70) = 0.770 \times 0.630 = 0.485$$

\textbf{Plan-confidence score:} $\hat P_\Sigma(M_0) = 0.485$. The strategy assigns about 49\% confidence to achieving the objective. This is distinct from $A_O$ (which optimizes over all policies) — $\hat P_\Sigma$ is the DAG's self-assessment of \textit{this specific plan}.


\medskip\noindent\textbf{Chain Confidence Decay ( \S\,\ref{chain-confidence-decay}).}


\textit{Mapping: exact (mathematical identity).}

The root node's confidence illustrates chain decay through an AND structure. The root requires two children to succeed, each with their own sub-chains:

$$\log \hat P_\Sigma = \log(p_{v_\text{id}, v_\text{root}} \cdot s_{v_\text{id}}) + \log(p_{v_\text{exploit}, v_\text{root}} \cdot s_{v_\text{exploit}})$$

$$= \log(0.770) + \log(0.630) = -0.261 + (-0.462) = -0.724$$

Each term is negative; the AND combination accumulates them additively in log-space. If the agent added a third AND-child (e.g., "maintain budget"), even with $s_\text{budget} = 0.95$, the plan confidence would drop to $0.485 \times 0.95 = 0.461$. The pressure toward shallow, narrow DAGs is structural.

\textbf{Contrast with the OR-node.} Phase 1 (OR-node) benefits from redundancy: three alternative arms provide resilience. The OR combination yields $s_{v_\text{id}} = 0.81$ even though no single arm has credence above $0.75$. The asymmetry between AND (multiplicative decay) and OR (complementary resilience) is the core structural result of \S\,\ref{chain-confidence-decay} combined with \S\,\ref{and-or-scope}.


\medskip\noindent\textbf{Orient Cascade ( \S\,\ref{orient-cascade}).}


\textit{Mapping: exact for ordering; the content of each step is tractable in this domain.}

Suppose at $t = 40$, the agent has been pulling arm 1 (the initially-best arm) and observes a run of failures: the last 8 of 10 pulls yielded $r_t = 0$. Meanwhile, arm 2's true probability has drifted to $\theta_2(40) = 0.8$ (unbeknownst to the agent) and arm 1 has drifted to $\theta_1(40) = 0.3$.

The orient cascade proceeds:

**Step 1: Reduce $\delta_\text{epistemic}$.** Update $M_t$ from the recent observations. The Beta posterior for arm 1 shifts:

$$\hat\theta_1: 0.70 \to 0.38 \quad (\text{after 8 failures in 10 trials with discounted prior})$$

Arms 2 and 3 are unobserved since the exploration phase ended. Their posteriors remain at approximately $\hat\theta_2 = 0.52$, $\hat\theta_3 = 0.28$ (slightly regressed toward the prior due to discounting). *The epistemic update proceeds without reference to $O_t$ or $\Sigma_t$* — directed separation ( \S\,\ref{directed-separation}) holds because the Beta update depends only on the observation and the prior, not on the agent's goals.

**Step 2: Evaluate $\delta_\text{sat}$.** With the updated $M_t$:

$$A_O(M_{40}) \approx \max_k \hat\theta_k + \epsilon_\text{info} = 0.52 + \epsilon_\text{info}$$

$$\delta_\text{sat} = 0.6 - (0.52 + \epsilon_\text{info})$$

If $\epsilon_\text{info} \approx 0.05$ (exploration can reveal a better arm): $\delta_\text{sat} \approx 0.6 - 0.57 = 0.03 > 0$. The objective is now \textbf{marginally unmet} — the agent believes it may not be able to achieve $0.6$ average reward. This triggers the disambiguation procedure ( \S\,\ref{satisfaction-gap}): is the goal truly infeasible, or is $M_t$ wrong about what's achievable?

**Step 3: Evaluate $\delta_\text{regret}$.** Under the current greedy policy (still pulling arm 1 despite updated beliefs):

$$V_O(M_{40}, \pi_\text{current}) = \hat\theta_1 = 0.38$$

$$\delta_\text{regret} = 0.57 - 0.38 = 0.19$$

Control regret is high. The agent is pulling arm 1 out of inertia, but $M_t$ now says arm 2 is better. \textit{This is the Section II payoff}: Section I would report the mismatch $\delta_t$ on each pull and the updated $\hat\theta_1$, but it cannot diagnose that the \textit{strategy} is the problem. The high $\delta_\text{regret}$ identifies the corrective action: revise $\Sigma_t$, not $O_t$.

**Step 4: Evaluate $\delta_\text{strategic}$.** The edge residual for the root AND-node's exploitation branch:

$$r_\text{exploit} = \mathbb{E}[\Delta V_O \mid \text{exploit arm 1}] - \Delta V_O^\text{observed}$$

The agent predicted reward $\approx 0.70$ per pull from arm 1; it observed $\approx 0.20$. The residual is large and positive, localizing the problem to the exploitation edge: the agent's causal belief that "pulling arm 1 yields high reward" is wrong.

Simultaneously, the leaf credences update: $p_1(M_{40}) \approx 0.15$ (arm 1 is probably not the best), $p_2(M_{40}) \approx 0.60$ (arm 2 is probably best, by default from the prior). The Phase 1 OR-node status shifts: the identification sub-goal may need to be revisited.

**Step 5: Revise $\Sigma_t$.** The agent restructures its strategy: switch exploitation target to arm 2 (the currently estimated-best arm). In DAG terms: $p_1(M_{40})$ has dropped to $0.15$ while $p_2(M_{40})$ has risen to $0.60$. The OR-node now routes through arm 2. The exploitation leaf updates: $p_\text{exploit}(M_{40}) = P(\theta_2 \geq 0.6 \mid M_{40}) \approx 0.40$.

After strategy revision: $\hat P_\Sigma(M_{40}) \approx (0.95 \times 0.73) \times (0.90 \times 0.40) = 0.693 \times 0.360 = 0.249$. Plan confidence has dropped from $0.485$ to $0.249$ — the agent correctly recognizes that its situation has deteriorated.

\textbf{What Section I alone cannot provide.} At $t = 40$, Section I reports: $\hat\theta_1 = 0.38$, $\hat\theta_2 = 0.52$, $\hat\theta_3 = 0.28$, and the per-arm mismatch signals. This tells the agent what it believes about each arm. It does NOT tell the agent:
\begin{itemize}
  \item Whether its \textit{goal} is still achievable ($\delta_\text{sat}$)
  \item Whether its \textit{strategy} is the problem or the goal is ($\delta_\text{regret}$)
  \item Which \textit{part} of the strategy is failing ($\delta_\text{strategic}$, edge residuals)
  \item What corrective action to take (the cascade's ordering)
\end{itemize}



\medskip\noindent\textbf{Observability Dominance ( \S\,\ref{observability-dominance}).}


\textit{Mapping: exact mechanism; approximate quantitative form.}

When the agent stops pulling arm 3 after the exploration phase (say at $t = 30$), the observation rate for arm 3 drops to $\nu_3 = 0$. By \S\,\ref{update-gain}:

$$\eta_3^* = \frac{U_{\text{edge},3}}{U_{\text{edge},3} + U_{\text{obs},3}}$$

With $U_{\text{obs},3} \to \infty$ (no observations): $\eta_3^* \to 0$. Arm 3's credence $p_3(M_t)$ is \textbf{frozen at its exploration-phase value}.

For the strategy DAG, the leaf node "arm 3 is best" has:
\begin{itemize}
  \item Nominal credence: $p_3(M_t) \approx 0.05$ (from $t = 30$ beliefs)
  \item Path observability: $\sigma_3 = 0$ (arm 3 is not being pulled)
  \item Observability-adjusted confidence: $\text{conf}_\text{obs}(\text{arm 3 path}) = 0.05 \times 0 = 0$
\end{itemize}


Arm 3's path through the strategy is \textbf{epistemically dead} — the agent cannot learn whether arm 3 has become the best arm. If $\theta_3$ drifts to $0.9$ (which it might, given nonstationarity), the agent will never discover this through the exploitation-only policy.

\textbf{The observability-dominance diagnostic.} Compare arm 2 and arm 3:

\begin{adjustbox}{max width=\textwidth}
\small
\begin{tabular}{lll}
\toprule
 & Arm 2 & Arm 3 \\
\midrule
Nominal credence $p_k$ & 0.52 & 0.28 (frozen) \\
Observability $\sigma_k$ & moderate (occasional pulls) & 0 (never pulled) \\
Adjusted confidence & $0.52 \times \sigma_2$ & 0 \\
Can be updated? & Yes & No \\
\bottomrule
\end{tabular}
\end{adjustbox}


Arm 2 is epistemically alive; arm 3 is epistemically dead. The agent should prefer exploring arm 2 (where it can learn) over arm 3 — unless it invests in making arm 3 observable again (pulling it despite low expected reward). This is the strategic case for exploration: not just expected value, but \textit{observability maintenance}.

\textbf{Section I vs Section II diagnosis.} Section I reports $\hat\theta_3 = 0.28$ with growing uncertainty (the posterior variance increases under discounting with no new data). Section II adds the diagnostic: this arm's strategy path is epistemically dead ($\sigma_3 = 0$), and restoring observability has value beyond the arm's nominal expected reward.


\medskip\noindent\textbf{Edge Update via Gain ( \S\,\ref{edge-update-via-gain}).}


\textit{Mapping: exact for the Beta-Bernoulli case.}

When the agent pulls arm 1 at $t = 35$ and observes failure ($r_t = 0$):

\textbf{Leaf credence update.} The Beta posterior for arm 1 updates:

$$\beta_1 \to \beta_1 + 1, \quad \hat\theta_1 = \frac{\alpha_1}{\alpha_1 + \beta_1 + 1}$$

This is the standard conjugate update — the gain principle reduces to Bayesian updating in the binary case ( \S\,\ref{edge-update-via-gain}).

\textbf{Strategy propagation.} The updated $\hat\theta_1$ changes $p_1(M_t)$ (the credence that arm 1 is best). This propagates through the OR-node to the root:

$$s_{v_\text{id}}^{\text{new}} = 1 - (1 - p_1^{\text{new}})(1 - p_2)(1 - p_3)$$

The OR-node's status is \textit{resilient} to single-arm credence drops: even if $p_1$ drops substantially, $p_2$ can compensate. This is the OR-node's structural advantage — redundancy absorbs single-edge failures.

\textbf{The exploitation edge.} The edge $p_{v_\text{exploit}, v_\text{root}} = 0.90$ is a causal credence: "pulling the identified-best arm reliably yields reward $\geq 0.6$." After observing repeated failures from arm 1:

$$\text{signal} = 0 \quad (\text{failure observed})$$

$$p_\text{exploit}^{\text{new}} = p_\text{exploit}^{\text{old}} + \eta_\text{edge} \cdot (0 - p_\text{exploit}^{\text{old}})$$

With $\eta_\text{edge} = U_\text{edge}/(U_\text{edge} + U_\text{obs})$, the exploitation edge's credence decreases. The rate of decrease is governed by the gain principle: well-established edges (many prior successes, $U_\text{edge}$ small) are slow to revise; newly formed edges are fast.


\medskip\noindent\textbf{Directed Separation ( \S\,\ref{directed-separation}).}


\textit{Mapping: exact. The agent is Class 1 (modular) by construction.}

The Beta-Bernoulli update for $M_t$ processes $(a_t, r_t)$ without reference to $(O_t, \Sigma_t)$:

$$(\alpha_{a_t}, \beta_{a_t}) \to (\alpha_{a_t} + r_t, \beta_{a_t} + 1 - r_t)$$

The update function $f_M$ has no $G_t$ argument. The strategy DAG then updates its leaf credences from the new $M_t$ — this is $f_G(G_t, M_{t+1}, e_t)$. The agent architecture is modular: a separate estimator (Beta posteriors) and planner (DAG evaluation), connected through the state-estimate interface.

$\kappa_\text{processing} = 0$ — no goal information flows into the epistemic update. Section II results apply exactly.


\medskip\noindent\textbf{Satisfaction Gap Dynamics.}


\textit{Mapping: exact definition; approximate computation.}

Tracking $\delta_\text{sat}$ over time reveals the agent's evolving assessment of goal feasibility:

\begin{adjustbox}{max width=\textwidth}
\small
\begin{tabular}{lllll}
\toprule
Time & $\hat\theta_{k^*}$ & $A_O$ (approx) & $\delta_\text{sat}$ & Interpretation \\
\midrule
$t = 0$ & 0.70 & 0.72 & $-0.12$ & Comfortably attainable \\
$t = 20$ & 0.68 & 0.70 & $-0.10$ & Still attainable \\
$t = 40$ & 0.52 & 0.57 & $+0.03$ & Marginally unmet \\
$t = 50$ (after exploring arm 2) & 0.75 & 0.77 & $-0.17$ & Attainable again — $M_t$ improved \\
$t = 70$ & 0.78 & 0.79 & $-0.19$ & Confidently attainable \\
\bottomrule
\end{tabular}
\end{adjustbox}


The key transition is $t = 40 \to 50$: $\delta_\text{sat}$ goes from positive (unmet) to negative (attainable) — not because the goal changed or the environment improved, but because $M_t$ improved. The agent explored arm 2, discovered $\hat\theta_2 \approx 0.75$, and its attainability assessment corrected upward. This matches the \S\,\ref{satisfaction-gap} disambiguation: the positive $\delta_\text{sat}$ at $t = 40$ was caused by a bad model, not an infeasible goal. The orient cascade's ordering (fix $M_t$ first, then reassess $\delta_\text{sat}$) produced the right diagnosis.


\medskip\noindent\textbf{The Section II Payoff.}


The purpose of this worked example is to demonstrate that Section II adds genuine diagnostic value beyond Section I. Here is the explicit comparison:

**Scenario at $t = 40$: arm 1 failing, arm 2 actually best, agent pulling arm 1.**

\textit{Section I diagnosis.} The mismatch signal $\delta_t$ is large. The posterior $\hat\theta_1$ has dropped. Per-arm tempo $\mathcal T_1 \approx 0.016$ is adequate for tracking arm 1's drift. The agent knows arm 1 looks bad. Section I prescribes: update $M_t$ faster (increase $\eta^*$) or observe more (increase $\nu$).

\textit{Section II diagnosis.} The orient cascade runs:
\begin{enumerate}
  \item $M_t$ updates correctly — arm 1's posterior drops to $0.38$. (Directed separation: this happens goal-blindly.)
  \item $\delta_\text{sat} = +0.03$ — the objective is marginally unmet. The agent's best-believed policy barely misses the target.
  \item $\delta_\text{regret} = 0.19$ — the current policy (pulling arm 1) is far from the best available (pulling arm 2).
  \item The 2x2 diagnostic says: "goal may be achievable, strategy is definitely suboptimal." Corrective action: revise $\Sigma_t$, then reassess $\delta_\text{sat}$.
  \item Edge residuals localize the problem: the exploitation branch's predicted-vs-observed value gap is large. The leaf credences identify arm 2 as the best candidate.
  \item Observability check: arm 3 is epistemically dead ($\sigma_3 = 0$). The agent should invest in observing arm 2 (or arm 3) before committing.
\end{enumerate}


Section II tells the agent: \textit{switch to arm 2, but first explore to confirm — the goal is probably still achievable, the current strategy is the problem, not the goal.} Section I tells the agent: \textit{arm 1 looks bad.} The difference is the difference between a diagnosis and a symptom.


\medskip\noindent\textbf{Mapping Quality Summary.}


\begin{adjustbox}{max width=\textwidth}
\small
\begin{tabular}{llll}
\toprule
ACT Concept & Bandit-Strategy Mapping & Status & Notes \\
\midrule
Scope ( \S\,\ref{scope-condition}) & Exact & Definitional & Same as \S\,\ref{worked-example-bandit} \\
Model $M_t$ ( \S\,\ref{agent-model}) & Exact & Structural & Beta posteriors per arm \\
Mismatch $\delta_t$ ( \S\,\ref{mismatch-signal}) & Exact & Standard & Prediction error \\
Update gain $\eta^*$ ( \S\,\ref{update-gain}) & Exact & Conjugate & Beta-Bernoulli gain \\
Tempo $\mathcal{T}$ ( \S\,\ref{adaptive-tempo}) & Approximate & Per-arm decomposition & Assumes known allocation \\
Directed separation ( \S\,\ref{directed-separation}) & Exact & Class 1 by construction & Beta update is goal-blind \\
Objective $O_t$ ( \S\,\ref{objective-functional}) & Exact & Utility with threshold & $V_{O_t}^{\min} = 0.6$ \\
Value object $V_O$ ( \S\,\ref{value-object}) & Exact & Closed-form (greedy) & $V_O = \hat\theta_{k^*}$ \\
Satisfaction gap $\delta_\text{sat}$ ( \S\,\ref{satisfaction-gap}) & Exact & Definition applied & $V_{O_t}^{\min} - A_O$ \\
Control regret $\delta_\text{regret}$ ( \S\,\ref{control-regret}) & Exact & Definition applied & $A_O - V_O(\pi_\text{current})$ \\
Strategy DAG $\Sigma_t$ ( \S\,\ref{strategy-dag}) & Exact & AND/OR with 7 nodes & Small enough for manual propagation \\
AND/OR scope ( \S\,\ref{and-or-scope}) & Exact & Natural fit & Phase 1 is genuinely OR; root is genuinely AND \\
Status propagation & Exact & Closed-form & Forward pass computable by hand \\
Plan confidence $\hat P_\Sigma$ ( \S\,\ref{strategy-dag}) & Exact (computation) & From propagation; systematically optimistic under correlated failure & $0.485$ at $t = 0$ \\
Chain confidence decay ( \S\,\ref{chain-confidence-decay}) & Exact & Mathematical identity & AND-node multiplicative; OR-node resilient \\
Orient cascade ( \S\,\ref{orient-cascade}) & Exact (ordering) & All 5 steps exercised & Content is tractable in this domain \\
Observability dominance ( \S\,\ref{observability-dominance}) & Exact (mechanism) & Unpulled arms frozen & $\sigma_3 = 0 \Rightarrow \eta_3^* = 0$ \\
Edge update via gain ( \S\,\ref{edge-update-via-gain}) & Exact (binary case) & Beta-Bernoulli conjugate & Gain principle = Bayesian update \\
Strategic calibration ( \S\,\ref{strategic-calibration}) & Approximate & Credit assignment tractable here & Single-arm attribution is clean \\
Satisfaction gap dynamics & Exact & Time series of $\delta_\text{sat}$ & Model improvement resolves positive gap \\
\bottomrule
\end{tabular}
\end{adjustbox}


\textbf{Quantities that map cleanly.} All Section II definitions (objective, value object, satisfaction gap, control regret, strategy DAG) have exact instantiations. Plan confidence $\hat P_\Sigma$ is exact as a computation (the AND/OR propagation is correct), but systematically overestimates actual success probability under correlated edge failures ( \S\,\ref{strategy-dag}). In this toy bandit, edge independence is a reasonable approximation; in complex domains the overestimate may be severe. The orient cascade ordering is exact and all five steps are exercised.

\textbf{Quantities that map approximately.} Strategic calibration is approximate because even in this simple domain, attributing value changes to specific DAG edges requires assumptions about the agent's execution fidelity. Observability dominance is exact in mechanism but the multiplicative form $\text{conf}_\text{obs} = \text{conf} \cdot \text{obs}$ is a first-order approximation.

\textbf{Quantities that don't map.} No Section II quantity fails to map in this domain. The system is simple enough that every quantity is computable (at least approximately). This is the advantage of the bandit setting: it is the simplest system that exercises all Section II machinery.


\medskip\noindent\textbf{Epistemic Status.}


This is a \textit{worked instantiation} — it demonstrates that the Section II formal chain is internally consistent and instantiable. The mapping quality is \textit{conditional}: the quantitative relationships depend on the Beta-Bernoulli reward model, the specific DAG structure, and the parameterization. The qualitative conclusions — especially the diagnostic value of the 2x2 table and the orient cascade's ordering — should be robust across reward models and DAG topologies.

The status is \textit{conditional} rather than \textit{exact} because the attainability $A_O$ is approximated (the exact Bayes-optimal policy for a nonstationary bandit is intractable), and the strategic calibration mapping inherits the discussion-grade status of \S\,\ref{strategic-calibration}. The core Section II definitions (satisfaction gap, control regret, plan confidence) are exact instantiations.

\end{remark}


\end{document}
